% Single-file manuscript. Compile twice with XeLaTeX.
% Requires standard TeX packages, kotex, and Noto Serif CJK KR or UnBatang.
% No external manuscript input, graphics, or BibTeX database is needed.
\documentclass[11pt,a4paper]{article}
\usepackage[a4paper,top=25mm,bottom=25mm,left=25mm,right=25mm,headheight=15pt]{geometry}
\usepackage{fontspec}
\usepackage{kotex}
\setmainfont{Latin Modern Roman}
\setsansfont{Latin Modern Sans}
\setmonofont{Latin Modern Mono}[Scale=MatchLowercase]
\IfFontExistsTF{Noto Serif CJK KR}{%
\setmainhangulfont{Noto Serif CJK KR}[BoldFont={Noto Serif CJK KR Bold},ItalicFont={Noto Serif CJK KR},BoldItalicFont={Noto Serif CJK KR Bold}]
\setsanshangulfont{Noto Serif CJK KR}[BoldFont={Noto Serif CJK KR Bold}]
}{\setmainhangulfont{UnBatang}}
\usepackage{amsmath,amssymb,amsthm,mathtools}
\usepackage{booktabs,tabularx,longtable,array,enumitem,xcolor,fancyhdr}
\usepackage{xurl}
\definecolor{inkblue}{HTML}{193F61}
\usepackage[unicode,colorlinks=true,linkcolor=inkblue,citecolor=inkblue,urlcolor=inkblue,
pdftitle={Arithmetic Alternatives for Euler's Constant: Integrated Research Manuscript},
pdfauthor={KEIO},pdfsubject={Arithmetic alternatives, functional transcendence and mixed specialization}]{hyperref}
\newcommand{\Q}{\mathbb Q}
\newcommand{\Z}{\mathbb Z}
\newcommand{\R}{\mathbb R}
\newcommand{\C}{\mathbb C}
\newcommand{\N}{\mathbb N}
\newcommand{\QB}{\overline{\mathbb Q}}
\newcommand{\Qbar}{\overline{\mathbb Q}}
\newcommand{\E}{\mathbf E}
\newcommand{\Ht}{\operatorname{H}}
\DeclareMathOperator{\Ein}{Ein}
\DeclareMathOperator{\Int}{Int}
\DeclareMathOperator{\trdeg}{trdeg}
\DeclareMathOperator{\rad}{rad}
\providecommand{\tightlist}{\setlength{\itemsep}{0pt}\setlength{\parskip}{0pt}}
\newtheorem{theorem}{Theorem}[section]
\newtheorem{lemma}[theorem]{Lemma}
\newtheorem{proposition}[theorem]{Proposition}
\newtheorem{corollary}[theorem]{Corollary}
\theoremstyle{definition}
\newtheorem{definition}[theorem]{Definition}
\newtheorem{example}[theorem]{Example}
\newtheorem{condition}[theorem]{Condition}
\theoremstyle{remark}
\newtheorem{remark}[theorem]{Remark}
\theoremstyle{plain}
\renewcommand{\proofname}{Proof}
\numberwithin{equation}{section}
\setlist{nosep,topsep=5pt,leftmargin=*}
\setlength{\emergencystretch}{3em}
\setlength{\parskip}{3pt}
\setlength{\parindent}{1.2em}
\linespread{1.055}
\allowdisplaybreaks
\setcounter{tocdepth}{1}
\renewcommand{\contentsname}{Contents}
\renewcommand{\refname}{References}
\clubpenalty=10000
\widowpenalty=10000
\pagestyle{fancy}\fancyhf{}
\fancyhead[L]{\small\sffamily KEIO \enspace | \enspace Arithmetic Alternatives}
\fancyhead[R]{\small\sffamily Integrated edition}
\fancyfoot[C]{\small\thepage}
\renewcommand{\headrulewidth}{0.3pt}
\pdfstringdefDisableCommands{%
\def\Ein{Ein}\def\gamma{gamma}\def\delta{delta}\def\eta{eta}%
\def\Qbar{Qbar}\def\QB{Qbar}\def\mathbb#1{#1}%
\def\operatorname#1{#1}\def\mathrm#1{#1}}

\providecommand{\rcEin}{\operatorname{Ein}}
\providecommand{\rcSpan}{\operatorname{span}}
\providecommand{\rcdiag}{\operatorname{diag}}
\providecommand{\rck}{\overline{\mathbb Q}}
\providecommand{\rcC}{\mathbb C}
\providecommand{\rcQ}{\mathbb Q}
\providecommand{\rcZ}{\mathbb Z}
\providecommand{\rcR}{\mathbb R}
\providecommand{\rcM}{\mathcal M}

\begin{document}

% BEGIN INTEGRATED SOURCE: frontmatter.tex
\begin{titlepage}
\centering
\vspace*{18mm}
{\small\sffamily KEIO RESEARCH PROGRAM\par}
\vspace{12mm}
{\Huge\bfseries Arithmetic Alternatives\\[5pt] for Euler's Constant\par}
\vspace{9mm}
{\Large Approximation, binomial arithmetic,\\[4pt]
functional transcendence and mixed specialization\par}
\vspace{12mm}
{\large An integrated research manuscript\par}
\vspace{3mm}
{\normalsize From the Two-World review to the rank-four\\
and cross-field investigations\par}
\vspace{14mm}
{\Large KEIO\par}
\vspace{4mm}
{23 September 2026\par}
\vfill
\begin{minipage}{0.87\textwidth}
\small
\textbf{Status.} This manuscript consolidates proof-bearing deductions,
conditional criteria, counterexamples and the limits of the investigated
methods. It does not establish irrationality or transcendence of Euler's
constant or Gompertz's constant, or linear independence of
$1,e,\delta,e\operatorname{Ein}(1)$ over $\overline{\mathbb Q}$.
Literature-wide novelty and journal-level significance are not asserted.
\end{minipage}
\vspace{10mm}
\end{titlepage}

\section*{Abstract}
\addcontentsline{toc}{section}{Abstract}
We integrate a sequence of investigations prompted by the evaluation of a
two-world approach to Euler's constant. Algebraicity and transcendence are
treated as alternative assumptions inside the same mathematics, while the
analytic definitions of the constants remain fixed. The quantitative
review gives the sharp fixed-framing exponent, a rational-assumption
floor, and a compressed degree-dependent norm bound. Binomial closures,
prime certificates, denominator growth, rational Borel kernels, Fourier
positivity and moment reconstruction provide exact auxiliary structures.
Independent Euler-integral values yield functional differential
transcendence, while a separate arithmetic-cost analysis records the
height and denominator losses in polynomial constructions. The designated
four-value problem is reduced to a precise rank-three alternative, and
its exceptional points are constrained by affine multiplicative
independence. A final set of imports from spectral theory, difference
Galois theory and $p$-adic differential equations supplies explicit
counterexamples and exact local diagnostics: a positive integer Jacobi
operator has a differentially transcendental Weyl function with a rational
value at the designated point, and the Euler initial-value family has an
exact $p$-adic radius dichotomy. Full proofs, corrected hypotheses and a
source-to-section record are retained. The missing arithmetic
specialization step is stated separately from what is proved.

\medskip
\noindent\textbf{Keywords.} Euler's constant; Gompertz's constant;
$E$-functions; arithmetic Gevrey series; binomial rings; differential
transcendence; Jacobi operators; $p$-adic differential equations.

\section*{한국어 요약과 읽기 안내}
\addcontentsline{toc}{section}{한국어 요약과 읽기 안내}
이 논문은 ``이거 평가해볼래? 혹시 이거 가지고 너가 더 추가로 뭔가
감마 닫을 수 있거나 초월성 증명할 수 있다면 도전''이라는 질문에 대한
검토 답변부터, 그 뒤의 A/B 세계 확장, Moneyball, Snowball, 핵심 재평가,
네 값 랭크 연구, 마지막 교차분야 시도까지를 통합한 연구 원고다.
기존 원고의 영문 정리와 증명을 본문으로 유지하되, 반복되는 입력을
공유하고 표기와 참고문헌을 통일했다. 대화의 모든 문장을 옮긴 기록이
아니라, 그 구간의 수학적 내용을 추적할 수 있게 정리한 논문이다.

\medskip
\noindent\textbf{현재의 수학적 결론.}
\[
 \gamma,\delta\text{의 개별 무리성ㆍ초월성 및}
 \quad\operatorname{rank}_{\overline{\mathbb Q}}
 \{1,e,\delta,e\operatorname{Ein}(1)\}=4
\]
는 이 논문에서 증명되지 않는다. 함수의 미분초월성, 보조값들의
대수적 독립성, 특정 값의 초월성을 구별한다. 완결된 보조 정리가 있다는
것과 감마 문제의 미증명 단계를 해결했다는 것은 별개의 주장이다.

\medskip
\noindent\begin{tabularx}{\textwidth}{@{}p{0.20\textwidth}X@{}}
\toprule
연구 단계 & 이 통합본에서 다루는 내용 \\
\midrule
초기 원고 검토 & 고정 구간의 정확한 지수 $1/2$, 유리수 가정 아래의 하한,
대수차수에 따른 노름 하한, 기존 표현의 수정. \\
A/B 세계 확장 & 이항계수 환, 소수별 역원 증명서, 분모와 밀도,
Borel 나눗셈, 감마 곱셈공식과 Fourier 양성 조건. \\
누적 구성과 전이 & 소수 분해, 양의 모멘트ㆍHankel 근사,
미분초월 함수, 역모멘트 문제, 유한 Fourier 재구성. \\
핵심 재평가 & 차수ㆍ높이ㆍ분모의 총비용, 고정 변형의 범위,
유한 실험과 무한족의 차이, 연구 가치에 대한 제한. \\
네 값 랭크 & 랭크 $3$ 세계의 정확한 분류,
예외점의 곱셈 관계 제한, 혼합 영점 제거의 미증명 단계. \\
교차분야 감사 & 정수 Jacobi 반례, 차분 갈루아 이론,
정확한 $p$진 반경, $p$곡률, 값 수송의 오류 점검. \\
\bottomrule
\end{tabularx}

\medskip
이 원고의 ``세계 A''는 실제 정의로 정해진 $\gamma$가 대수적이라는
가정이고, ``세계 B''는 초월적이라는 가정이다. 두 세계의 결과를
만든다는 것만으로 어느 쪽의 모순이 보장되지는 않는다. 소거할 수 있는
세계가 되려면 동일한 가정 아래에서 어떤 명제와 그 부정을 함께
증명해야 한다. 축소된 항등식 체계의 반모형을 실제 상수의 정의 전체에
대한 반모형으로 해석하지 않는다.

\medskip
원고 001ㆍ002와 그 이전 통합본은 연구의 배경이다. 이 논문에는
지정한 시작 질문 이후에 사용되는 입력과 증명을 재수록하며, 그 이전
문서들의 모든 절을 다시 싣지는 않는다. 마지막의 출처 대응표는
어느 연구 묶음이 어느 부분으로 들어왔는지를 명시한다.

\clearpage
\tableofcontents
\clearpage

\section{Scope, notation and the logic of arithmetic alternatives}
\label{int:scope}

This is a thematic, proof-bearing consolidation of the investigations
listed in the source record at the end. Repeated source front matter and
repeated statements of the research objective have been removed.
Established external theorems are cited; internal deductions are proved;
conditional implications keep their hypotheses. A computational example
is not used in place of an asymptotic quantifier.

The common normalization is
\begin{align}
 \gamma&=\lim_{n\to\infty}\left(\sum_{k=1}^{n}\frac1k-\log n\right),\\
 \Ein(z)&=\sum_{n\ge1}\frac{(-1)^{n-1}z^n}{n\,n!},
 &E_1(x)&=\int_x^\infty\frac{e^{-t}}{t}\,dt\quad(x>0),\\
 D(x)&=e^xE_1(x)=\int_0^\infty\frac{e^{-t}}{x+t}\,dt,
 &\delta&=D(1),\qquad \eta=\Ein(1).
\end{align}
Integration by parts also gives
$\delta=\int_0^\infty\log(1+t)e^{-t}\,dt$. The connection identities are
\begin{equation}\label{int:connection}
 \Ein(x)=\gamma+\log x+E_1(x),\qquad
 \delta=e\eta-e\gamma.
\end{equation}
All logarithms of positive real numbers are real logarithms.
Analytic continuation elsewhere uses the indicated branch.

Throughout, $\overline{\mathbb Q}$ denotes the field of algebraic numbers.
The synonymous macros $\QB$ and $\Qbar$ in the source notation have this
meaning. The letter $K$ is defined locally: it may denote a number field,
an exponential-value field, or a Jacobi operator, and is never silently
transported between these meanings. A denominator-clearing integer for
Taylor coefficients is not the same object as the reduced denominator
of a value at a point. Likewise, integer coefficients, a nonzero
coefficient vector and a nonzero evaluated form are distinct conditions.

\subsection{A, B and C are a proof organization, not separate axioms}
Let $C$ denote the ordinary background theory together with the analytic
definitions above, and write
\[
 A=C+\{\gamma\in\overline{\mathbb Q}\},\qquad
 B=C+\{\gamma\notin\overline{\mathbb Q}\}.
\]
These assumptions are exhaustive. To reject $A$ one needs a finite valid
derivation of a contradiction from $A$; to reject $B$ one needs such a
derivation from $B$. Producing many compatible consequences, even from
unrelated subjects, does not ensure that either derivation will be found.
An additional finite partition can refine the search: if
$A\Rightarrow C_1\lor\cdots\lor C_s$ and every $A\land C_i$ is
contradictory, then $A$ is contradictory. Both exhaustiveness and each
exclusion must be proved. For an infinite partition a uniform argument
or another theorem covering all cases is required.

The ``overlapping worlds'' viewpoint is therefore used as a search
discipline. It does not replace the proof of a separating arithmetic
statement. The rank-three classification and the spectral counterexample
later in the manuscript illustrate why strong-looking consequences can
remain compatible with the undesired arithmetic alternative.

\subsection{Four distinct levels of assertion}
\begin{enumerate}[leftmargin=*,itemsep=3pt]
\item \textbf{Established input:} an external theorem with its stated
coefficient field, singularity, regularity and integrality assumptions.
\item \textbf{Deduction proved here:} a statement with a complete argument
from those inputs. This label does not itself assert historical novelty.
\item \textbf{Conditional criterion:} an implication whose additional
hypothesis has not been established for the designated constants.
\item \textbf{Finite diagnostic:} an exact finite calculation or a numerical
observation. Neither proves the existence of a required infinite family.
\end{enumerate}

The familiar tests for small integer forms belong to a specific
Diophantine route: an explicit construction, a fixed output space,
adequate error relative to height, and nonvanishing of the evaluated
form. They are not necessary conditions for every possible proof of
transcendence. A relation-lifting theorem, a structural algebraization
theorem or an entirely different argument could in principle bypass that
route, provided its own hypotheses are proved.

\subsection{The designated target and its strength}
Define
\[
 \rho(1)=\operatorname{rank}_{\overline{\mathbb Q}}
 \{1,e,\delta,e\eta\}.
\]
The established independence of $e$ and $\eta$ gives $\rho(1)\ge3$.
If $\rho(1)=4$, then both $\gamma$ and $\delta$ are transcendental:
algebraicity of $\gamma$ would turn \eqref{int:connection} into a
nontrivial algebraic-coefficient relation, and algebraicity of $\delta$
would give a relation with $1$. The converse is not established by this
argument; indeed the rank-three classification below allows generic
coefficient cases in which $\gamma$ and $\delta$ are algebraically
independent. The target is therefore kept separate from the weaker
individual transcendence questions.

\subsection{Editorial and verification policy}
The publication-style exposition is assembled from the recovered source
manuscripts and proof notes. Mathematical assistance and internal
cross-checking used ChatGPT. Internal checking is not external peer
review or proof-assistant verification. Where an earlier wording required
correction, the corrected quantifier or scope is used here and recorded
in the integration notes. Source-reported computations that were not
rerun are not presented as new computations. All mathematical source
text and bibliography needed to typeset this edition are contained in
the accompanying single TeX file.

% END INTEGRATED SOURCE

\clearpage
\part{Foundations}

% BEGIN INTEGRATED SOURCE: snow_core/dependencies.tex
\section{Shared E-function foundations}
\label{sc:app:dependencies}\label{sc:sec:reused-inputs}

The following dependency proofs supply the inputs needed by the
subsequent deductions. They are reused results, not additional claims of
novelty. Only the order-one $\operatorname{Ein}$ theorem is used; no
higher-polyexponential assertion is required. 

\subsection{Order-one multipoint independence}
\label{sc:dep:ein}

We use
\[
 \operatorname{Ein}(z)=\sum_{n\ge1}\frac{(-1)^{n-1}z^n}{n\,n!},
 \qquad E_1(u)=\int_u^\infty\frac{e^{-t}}t\,dt\quad(u>0),
 \qquad K_{\exp}=\overline{\mathbb Q}
       (e^\alpha:\alpha\in\overline{\mathbb Q}).
\]
The connection formula and derivative are
\begin{equation}\label{sc:dep:connection}
 E_1(u)=\operatorname{Ein}(u)-\gamma-\log u,
 \qquad \operatorname{Ein}'(z)=\frac{1-e^{-z}}z.
\end{equation}
These normalizations agree with~\cite[\S6.2]{DLMF}.
The entire function $\operatorname{Ein}$ is an $E$-function: its
normalized Taylor coefficients are $(-1)^{n-1}/n$, whose common
denominators grow at most exponentially, and the displayed differential
identity places it in a finite linear differential system. In contrast,
$E_1$ is not an $E$-function at zero.

\begin{lemma}[Laurent derivative obstruction]\label{sc:dep:laurent}
Let $\nu_1,\ldots,\nu_d\in\overline{\mathbb Q}$ be linearly
independent over $\mathbb Q$. Endow
\[
 R=\overline{\mathbb Q}(z)[X_1^{\pm1},\ldots,X_d^{\pm1}],
 \qquad F=\operatorname{Frac}(R),\qquad DX_j=-\nu_jX_j.
\]
If $g\in F$ satisfies
$Dg=z^{-1}\sum_u c_uX^u$ with finitely supported
$c_u\in\overline{\mathbb Q}$, then all $c_u$ vanish.
The constant field of $F$ is $\overline{\mathbb Q}$.
\end{lemma}
\begin{proof}
An irreducible nonunit $P\in R$ cannot divide $DP$. Indeed, from
$DP=AP$, comparison of Newton polytopes gives
$A\in\overline{\mathbb Q}(z)$; the case $DP=0$ is included in the
coefficient argument. For any two nonzero Laurent coefficients
$a_u,a_v$ of $P$ one obtains
\[
 (a_u/a_v)'/(a_u/a_v)=(u-v)\cdot\nu.
\]
A rational function has no nonzero constant logarithmic derivative.
Independence of the $\nu_j$ therefore forces $u=v$, making $P$ a
Laurent monomial times a unit, a contradiction.

Any nonunit denominator factor of a reduced $g$ would consequently
have its pole order increased by differentiation. The proposed right
side has no such pole, so $g\in R$. Writing $g=\sum a_u(z)X^u$ gives
\[
 a_u'-(u\cdot\nu)a_u=c_u/z.
\]
For $u=0$, rational derivatives have zero residue. For $u\ne0$, a pole
of $a_u$ at zero produces a higher-order pole, while regularity of
$a_u$ at zero makes the left side regular there. Thus $c_u=0$ in all
cases. With zero right side, $u\ne0$ admits no nonzero rational
solution, and $a_0$ is constant. This proves the final assertion.
\end{proof}

\begin{theorem}[Multipoint $\operatorname{Ein}$ input]
\label{sc:dep:multipoint}
For distinct nonzero algebraic $\lambda_1,\ldots,\lambda_s$, put
\[d=\dim_{\mathbb Q}\operatorname{span}_{\mathbb Q}
\{\lambda_1,\ldots,\lambda_s\}.\]
Then
\[
 \operatorname{trdeg}_{\overline{\mathbb Q}}
 \overline{\mathbb Q}(e^{\lambda_i},\operatorname{Ein}(\lambda_i):
                      1\le i\le s)=d+s.
\]
In particular all the $\operatorname{Ein}(\lambda_i)$ are algebraically
independent over $K_{\exp}$, in the finite-subset sense.
\end{theorem}
\begin{proof}
Choose a $\mathbb Z$-basis $\nu_1,\ldots,\nu_d$ of the additive group
generated by the $\lambda_i$, with $\lambda_i=m_i\cdot\nu$.
Distinct exponential characters are linearly independent over
$\overline{\mathbb Q}(z)$, so $X_j=e^{-\nu_jz}$ realizes the field in
Lemma~\ref{sc:dep:laurent}. Set $Y_i(z)=\operatorname{Ein}(\lambda_i z)$;
then $DY_i=(1-X^{m_i})/z$. The generated field embeds in
$\overline{\mathbb Q}((z))$, and has no new differential constants.
The primitive Kolchin--Ostrowski theorem~\cite{Kolchin} turns any
algebraic dependence of the $Y_i$ over $F$ into
\[
 Dg=z^{-1}\left(\sum_i c_i-\sum_i c_iX^{m_i}\right),
 \qquad g\in F,\quad c_i\in\overline{\mathbb Q},
\]
with some $c_i\ne0$. This contradicts Lemma~\ref{sc:dep:laurent}, since
the characters $m_i$ are distinct and nonzero.

Extending constants from $\overline{\mathbb Q}$ to $\mathbb C$ does
not change this functional transcendence degree: any proposed relation
uses finitely many monomials with algebraic Taylor coefficients, whose
coefficient-matrix rank is unchanged by scalar extension. Include $1$,
the required exponential characters, and the $Y_i$ in a finite linear
system of $E$-functions. Its coefficients have no finite pole except
possibly zero. Siegel--Shidlovskii specialization at $z=1$, in the refined
form of Beukers~\cite{Beukers}, gives the asserted transcendence
degree. The exponential coordinates have exactly their toric relations,
so the $s$ remaining coordinates are independent over their field.
Finally every coefficient in a hypothetical relation over $K_{\exp}$
uses only finitely many algebraic exponential arguments. Add those
arguments to the finite system; the same result excludes the relation.
\end{proof}

A full-row-rank matrix over $\overline{\mathbb Q}$ therefore sends a
finite independent $\operatorname{Ein}$ tuple to an independent tuple
over $K_{\exp}$. This follows by extending its rows to an invertible
matrix and applying a linear change of indeterminates.

\subsection{The integer-point rank bound}
\label{sc:dep:integer-rank}

Write $U_n=\operatorname{Ein}(n)$, $y_n=E_1(n)$,
$I_n=e^ny_n$, and $\Omega(n)=\sum_pv_p(n)$.
For every composite integer $n$, the previously constructed residual is
\begin{align*}
 R_n&=y_n-\sum_{p\mid n}v_p(n)y_p+(\Omega(n)-1)y_1\\
 &=U_n-\sum_{p\mid n}v_p(n)U_p+(\Omega(n)-1)U_1.
\end{align*}
Both the logarithms and the common $\gamma$ term cancel by
\eqref{sc:dep:connection}. Each composite column $U_n$ is a unique pivot
in its own row, so the residuals are algebraically independent over
$K_{\exp}$. There are $N-\pi(N)-1$ such columns through $N$.
Consequently
\begin{equation}\label{sc:dep:I-rank}
 \operatorname{trdeg}_{K_{\exp}}K_{\exp}(I_1,\ldots,I_N)
 \ge N-\pi(N)-1,
\end{equation}
since multiplication by $e^n\in K_{\exp}^{\times}$ does not alter
the generated field. These residuals are also positive: with
$h(s)=E_1(e^s)$ one has $h''(s)=e^se^{-e^s}>0$, whence
$h(a+b)+h(0)>h(a)+h(b)$ for $a,b>0$. Iterate over the logarithms
of the prime factors of $n$, counted with multiplicity.


% END INTEGRATED SOURCE

\clearpage
\part{Quantitative review and arithmetic alternatives}

% BEGIN INTEGRATED SOURCE: review_ab/body.tex
\section{Logical scope and the review of the Two-World argument}
The supplied review addendum audits the 52-page manuscript \cite{KEIO}. The main logical and quantitative corrections below have also been compared with the recovered source text. The larger lattice searches are not independently reproduced here; exact lattice reduction alone does not certify the real input constants. The degree-twelve example preserved in the review archive has been checked again by exact rational arithmetic.

The arithmetic alternatives concern the same fixed real number
\[
 \gamma=\lim_{n\to\infty}\left(\sum_{j=1}^n\frac1j-\log n\right)
 =-\Gamma'(1),\qquad 0<\gamma<1.
\]
Write $\mathrm A$ for the hypothesis $\gamma\in\Qbar$, $\mathrm B$ for
$\gamma\notin\Qbar$, and $\mathrm C$ for the usual mathematical
background including this definition. A contradiction to $\mathrm A$
requires both implications
\[
 \mathrm C+\mathrm A\Longrightarrow\mathcal O\text{ has property }P,
 \qquad
 \mathrm C\Longrightarrow\mathcal O\text{ cannot have property }P.
\]
Rejecting an additional auxiliary property does not reject either
arithmetic alternative unless that property has first been derived from it.
The constructions below are classifications and conditional criteria;
none constructs the required infinite approximation family.

Write
\[
 \Ein(z)=\sum_{n\ge1}\frac{(-1)^{n-1}z^n}{n\,n!},\qquad
 T=e\Ein(1),\qquad \delta=eE_1(1)=\int_0^\infty\frac{e^{-t}}{1+t}\,dt=T-e\gamma.
\]
The scalar used to vary the framing will be $x$, so a fixed integer form becomes
\begin{equation}\label{ra:eq:affine}
 F_{a,b,c}(x)=a+bx+c(T-ex)=(a+cT)+(b-ce)x.
\end{equation}
At $x=\gamma$ it is the actual form $a+b\gamma+c\delta$.

\paragraph{Restricted models are not models of the full definition.}
The two-model theorem applies to the restricted theory which omits the identifications of its new symbols with the defining limits, integrals and canonical sums. It does not prove independence of the arithmetic nature of the actual, analytically defined $\gamma$, or impossibility for every available mathematical method. The manuscript itself partly recognizes this distinction; its abstract and concluding generalizations should retain it throughout.

\paragraph{Uniformity is essential.}
Every true strict inequality $|F_{a,b,c}(\gamma)|<\eta$ remains true in some neighborhood of $\gamma$, by continuity. Hence, with the manuscript's definition allowing an arbitrary finite enclosure, a single finite exceptional cancellation is locally shear-invariant. What the tie argument restricts is an entire named family with \emph{one fixed common interval}. The statement in Remark 18.3 that a finite cancellation is itself not shear-invariant needs this qualification.

\paragraph{Numerical exponents are not asymptotic exponents.}
The finite statistic $300/\log_{10}H_d$ in the proposition entitled ``Approximation exponents of $\gamma$'' does not establish $w_d^*(\gamma)=d$. The latter requires asymptotic quantifiers, nontrivial approximants and the actual approximation errors. Report the finite ratios only. Exact verification of LLL reduction on rounded input does not by itself certify enclosures of the original constants or root locations. Such claims need rational or ball intervals in addition to exact lattice calculations.

\paragraph{Other specific repairs and scope qualifications.}
\begin{enumerate}[leftmargin=*,itemsep=3pt]
\item In the finite-observation theorem, $S=\varnothing$ gives a false assertion of a degree-one algebraic integer in every neighborhood of $\gamma$. There is no such integer in $(0,1)$. Allow degree-one algebraic numbers (rationals), or require degree at least two for the algebraic-integer assertion in this case. The neighborhood proof also requires the interval enclosures to contain $\gamma$ in their interiors.
\item The ``one-sidedness'' proposition cannot infer that the program can only terminate on the algebraic-assumption side from finitely many excluded relations. A relation of larger height or degree has not thereby been ruled out. Under transcendence, existence of a derivation taking $\gamma$ to $1$ is also a positive existential consequence, contrary to its universal-negative description of every consequence.
\item The last item of the linear-hierarchy theorem needs the declared subspace $V$ restricted, for example, to the listed countable subspaces of $\E$. For arbitrary $V\subset\C$ with $e\in V$, take $V=\C$: its claimed exclusion is identically false.
\item The collision identity's citation to Fischler--Rivoal uses $a=1,s=r$, not $s=r-1$. With their notation, $\omega_{1,r}=(-1)^{r-1}f_{1,r+1;0}(1)/r!$.
\item The shadow proposition's title is broader than its proved conclusion: it establishes an $E$-value realization for $\delta_c$, while its body already notes logarithmic terms for general $\lambda$; higher entries also contain zeta values. Coefficient conjugation of a power series should not be identified with arbitrary Galois conjugation of its value. The dilation identity for the connection polynomials $P_r$ does not rescale the entire $\omega$-tower.
\item The mixed-value proposition establishes $\mathbf S\mathbf E$ membership; it does not assert canonical Borel membership in $\mathbf D$, and that stronger conclusion must not be inferred. A calculation of poles in the parameter $a$ is not a proof about all Stokes data in the variable $z$.
\item The $E$-membership decision theorem already requires an explicit certified $E$-function representation in its algorithmic branch. Its broader claim of termination on deciding membership must retain that requirement in the affirmative direction: bare nonconstructive membership does not supply the input. The $m$-fold multiplication twist factor is $e^{(m-1)\tau/2}$; $e^{\tau/2}$ is the duplication case.
\end{enumerate}
The central connection identities and the fixed-interval affine tie estimate survive these corrections.

\section{The fixed-interval exponent is exactly one half}
We use Shidlovskii's rational-coefficient linear independence measure, as stated in \cite[Theorem A]{FRliouville}: an independent $N$-component rational $E$-system, evaluated at a nonsingular nonzero rational point, has a lower bound $C_\varepsilon H^{-N+1-\varepsilon}$ for its nonzero integer linear forms.

\begin{lemma}\label{ra:lem:engine}
For each $\varepsilon>0$ there exists $C_\varepsilon>0$ such that
\[
 |h_0+h_1e+h_2T|\ge C_\varepsilon\max(1,|h_0|,|h_1|,|h_2|)^{-2-\varepsilon}
\]
for every nonzero integer triple. Moreover $1,e,T$ are linearly independent over $\QB$.
\end{lemma}
\begin{proof}
The vector $(1,e^z,F(z))$, $F(z)=e^z\Ein(z)$, satisfies
\[
 (1)'=0,\quad (e^z)'=e^z,\quad F'=F+(e^z-1)/z.
\]
It consists of $E$-functions with rational coefficients, with only $0$ a possible finite singularity. Its functions are linearly independent over $\C(z)$: otherwise, after division by $e^z$, $\Ein(z)$ would be a rational function plus a rational multiple of $e^{-z}$. On the positive real axis this is incompatible with $\Ein(z)=\log z+\gamma+O(e^{-z}/z)$. The remaining two functions are plainly independent. Apply the rational measure at $z=1$, and Beukers' linear independence theorem \cite{Beukers} for the algebraic-coefficient assertion.
\end{proof}

\begin{lemma}[Elementary simultaneous transference]\label{ra:lem:transfer}
For every $\varepsilon>0$ there is $C_\varepsilon>0$ such that, for all $q\ge1$ and integers $p,b$,
\begin{equation}\label{ra:eq:sim}
 \max(|qe-b|,|qT-p|)\ge C_\varepsilon q^{-1/2-\varepsilon}.
\end{equation}
\end{lemma}
\begin{proof}
Put $E_q=\max(|qe-b|,|qT-p|)$. The case $E_q\ge1$ is immediate. Otherwise $b/q$ and $p/q$ are bounded by absolute constants. Put $M=\lfloor\sqrt q\rfloor$. The $(M+1)^2>q$ residues $ub+vp\pmod q$, $0\le u,v\le M$, yield a nonzero pair $(h_1,h_2)$ with $|h_i|\le M$ and $q\mid h_1b+h_2p$. Set $h_0=-(h_1b+h_2p)/q$. Then $\max|h_i|\ll\sqrt q$ and
\[
 0<|h_0+h_1e+h_2T|
 =\left|\frac{h_1(qe-b)+h_2(qT-p)}q\right|
 \le \frac{2E_q}{\sqrt q}.
\]
Lemma~\ref{ra:lem:engine}, with exponent parameter $2\varepsilon$, bounds the left side below by $Cq^{-1-\varepsilon}$. Rearrangement proves the claim.
\end{proof}

\begin{theorem}[Sharp fixed framing interval]\label{ra:thm:half}
Let $I=[u,v]$ be a fixed bounded interval with $u<v$. For all $\varepsilon>0$ there exists $C_{I,\varepsilon}>0$ such that every nonzero integer triple, with $H=\max(1,|a|,|b|,|c|)$, satisfies
\begin{equation}\label{ra:eq:uniform}
 \sup_{x\in I}|a+bx+c(T-ex)|\ge C_{I,\varepsilon}H^{-1/2-\varepsilon}.
\end{equation}
Conversely, infinitely many triples with unbounded $H$ satisfy an upper bound $C_IH^{-1/2}$. Thus the supremum of the uniform approximation exponents for this particular pair is exactly $1/2$.
\end{theorem}
\begin{proof}
Write the supremum as $\eta$. The endpoint differences give
\[
 |b-ce|\le\frac{2\eta}{v-u},\qquad
 |a+cT|\le\eta\left(1+\frac{2\max(|u|,|v|)}{v-u}\right).
\]
If $c\ne0$, apply Lemma~\ref{ra:lem:transfer} with $q=|c|$, $b'=\operatorname{sgn}(c)b$, $p'=-\operatorname{sgn}(c)a$. It follows that $\eta\ge Cq^{-1/2-\varepsilon}\ge CH^{-1/2-\varepsilon}$. If $c=0$, either the integer $b$ has modulus at least one, giving $\eta\ge(v-u)/2$, or $a$ is a nonzero integer, giving $\eta\ge1$.

For the upper bound, simultaneous Dirichlet approximation gives arbitrarily large $Q$ and $1\le q\le Q$ with both errors at most $C Q^{-1/2}$. Take $(a,b,c)=(-p,b,q)$. Then $H\ll q\le Q$ and the uniform error is $O_I(Q^{-1/2})=O_I(H^{-1/2})$. The denominators must become unbounded because $e$ is irrational.
\end{proof}

\begin{remark}
This sharpens the manuscript's ceiling $1$ and its merely metric observation about $1/2$. It restricts the proved uniform bound, not the true error at the single point $x=\gamma$. It gives no universal bound on all approximations to $\gamma$.
\end{remark}

\begin{corollary}[Moving intervals]\label{ra:cor:moving}
Suppose $I_n\subset[-M,M]$ has length $0<\ell_n\le L$, and $H_n$ is the height of a nonzero integer triple. For every $\varepsilon>0$,
\[
 \eta_n:=\sup_{x\in I_n}|F_{a_n,b_n,c_n}(x)|
 \ge C_{M,L,\varepsilon}\ell_n H_n^{-1/2-\varepsilon}.
\]
Consequently, if $H_n\to\infty$, $\ell_n\ge H_n^{-\sigma+o(1)}$ and $\eta_n\le H_n^{-\kappa+o(1)}$, then $\kappa\le\sigma+1/2$.
\end{corollary}
\begin{proof}
The endpoint argument gives both engine errors at most $C_{M,L}\eta_n/\ell_n$. Use \eqref{ra:eq:sim}; the $c_n=0$ case follows directly from integrality as above. Take logarithms and then let $\varepsilon\downarrow0$.
\end{proof}
This is a genuine allowed way around the fixed-interval theorem. It spends increasing information about the position of $\gamma$, so it is not by itself an arithmetic construction or a proof of nonvanishing.

\section{A lower irrationality threshold without growth regularity}
\begin{theorem}[Rational-assumption floor]\label{ra:thm:rational}
Assume $\gamma\in\Q$. For every $\varepsilon>0$ there exists $C_\varepsilon>0$ such that every integer triple with $L=a+b\gamma+c\delta\ne0$ satisfies
\begin{equation}\label{ra:eq:rationalfloor}
 |L|\ge C_\varepsilon\max(1,|a|,|b|,|c|)^{-1-\varepsilon}.
\end{equation}
This uses the irrationality-exponent theorem for rational-coefficient $E$-functions in \cite[Theorem 1]{FRrational} (arXiv version 2, 2025).
\end{theorem}
\begin{proof}
Write $\gamma=r/s$ with $s>0$. The function $e^z(\Ein(z)-r/s)$ is an $E$-function with rational coefficients, and its value at $1$ is $\delta$. This value is irrational: otherwise $T-(r/s)e-\delta=0$ contradicts Lemma~\ref{ra:lem:engine}. Therefore \cite{FRrational} gives $\mu(\delta)=2$.

If $c\ne0$, put $m=sa+rb$ and $k=sc$. Then $sL=m+k\delta$. The assertion $\mu(\delta)=2$ implies $|m+k\delta|\ge C_\varepsilon|k|^{-1-\varepsilon}$ for all integer $m$ and nonzero integer $k$ (absorb the finitely many exceptional denominators into the positive constant). Since $|k|\le sH$, \eqref{ra:eq:rationalfloor} follows. If $c=0$, nonzero $L$ is a nonzero rational with denominator dividing $s$, so $|L|\ge1/s$.
\end{proof}

\begin{corollary}[Irregular-sequence irrationality test]\label{ra:cor:irrational}
If for some fixed $\kappa>1$ there are integer triples with $H_n\to\infty$ and
\[
 0<|a_n+b_n\gamma+c_n\delta|\le H_n^{-\kappa},
\]
then $\gamma$ is irrational. No regular growth of $H_n$ and no matching lower asymptotic for the forms are required.
\end{corollary}
\begin{proof}
Under rationality choose $0<\varepsilon<\kappa-1$ in Theorem~\ref{ra:thm:rational}; its lower bound contradicts the displayed upper bound for large $n$.
\end{proof}
\begin{remark}
The conclusion is irrationality, not transcendence: algebraic irrational values of $\gamma$ have not been excluded by this threshold. The nonzero condition cannot be dropped. Under rationality, infinitely many coefficient triples have exactly zero value with $c=0$.
\end{remark}

\section{Compressing the degree-dependent algebraic floor}
\begin{theorem}[Norm and symmetric powers]\label{ra:thm:norm}
Assume $\gamma$ is algebraic of degree $d\ge1$. For every $\varepsilon>0$, all integer triples with $L=a+b\gamma+c\delta\ne0$ satisfy
\begin{equation}\label{ra:eq:normfloor}
 |L|\ge C_{\gamma,\varepsilon}H^{-K_d-\varepsilon},\qquad
 K_d=d\binom{d+2}{2}-1.
\end{equation}
For $d=1$, Theorem~\ref{ra:thm:rational} improves $K_1=2$ to $1$.
\end{theorem}
\begin{proof}
Let $f(X)=A\prod_{j=1}^d(X-\gamma_j)\in\Z[X]$ be the primitive minimal polynomial, with $\gamma_1=\gamma$. Define $U=a+cY$, $V=b-cX$ and
\[
 P(X,Y)=(-1)^d V^d f(-U/V)
       =A\prod_{j=1}^d\bigl(a+b\gamma_j+cY-c\gamma_jX\bigr).
\]
The first expression is understood as its polynomial expansion even at $V=0$. Its coefficients are integers, its total degree is at most $d$, and $\Ht(P)\ll_\gamma H^d$.

At $(X,Y)=(e,T)$ the first factor is $L$. If $c\ne0$, each factor is nonzero by the algebraic linear independence in Lemma~\ref{ra:lem:engine}. If $c=0$, a zero factor would imply that a conjugate $\gamma_j$ is rational (unless $a=b=0$); hence $d=1$ and $L=0$, contrary to the hypothesis. Thus $P(e,T)\ne0$. Each of the other $d-1$ factors has modulus $O_\gamma(H)$.

Set $F(z)=e^z\Ein(z)$ and consider
\[
 F_{ij}(z)=e^{iz}F(z)^j,\qquad i,j\ge0,\quad i+j\le d.
\]
These $m_d=\binom{d+2}{2}$ rational-coefficient $E$-functions form a closed differential system:
\[
 F_{ij}'=(i+j)F_{ij}+\frac jz F_{i+1,j-1}-\frac jz F_{i,j-1},
\]
with the last two terms omitted for $j=0$. Use \cite[Theorem 1]{FRliouville} with coefficient field $\Q$ and $N=m_d$. This version only needs the evaluated form to be nonzero, already checked; no unproved high-order independence statement is needed. It gives
\[
 |P(e,T)|\ge C\Ht(P)^{-m_d+1-\varepsilon/d}
 \ge C' H^{-d(m_d-1)-\varepsilon}.
\]
Divide by the other factors, bounded by $O_\gamma(H^{d-1})$, and the fixed $A$. The resulting exponent is $dm_d-1+\varepsilon$.
\end{proof}

\begin{center}
\begin{tabular}{rrrrrr}
\toprule
Degree $d$&1&2&3&4&10\\
\midrule
Manuscript $d3^d-1$&2&17&80&323&590489\\
Norm bound $d\binom{d+2}{2}-1$&2&11&29&59&659\\
Best rational-assumption bound&1&--&--&--&--\\
\bottomrule
\end{tabular}
\end{center}
The constants depend on the alleged algebraic value and on $\varepsilon$. These are asymptotic bounds, not numerical certificates obtained by inserting a large height. This is a quantitative strengthening relative to the reviewed manuscript, by a standard norm construction.

\begin{corollary}[What an unbounded exponent would imply]
Suppose for every $\kappa>0$ there is an integer triple of height $H\ge2$ with
$0<|a+b\gamma+c\delta|<H^{-\kappa}$. Then $\gamma$ is transcendental.
\end{corollary}
\begin{proof}
If $\gamma$ had degree $d$, \eqref{ra:eq:normfloor} would bound all sufficiently large-height exponents by $K_d+1$. The finitely many remaining triples of bounded height with nonzero values cannot realize unbounded exponents either.
\end{proof}
No such family is constructed here. The improved floor does not convert the manuscript's finite or approximately $1/2$-exponent forms into a transcendence proof.

\section{Coefficient rank and conditional contradiction criteria}
\label{ra:sec:coefficient-rank}

Put $T=e\Ein(1)$ and $\delta=T-e\gamma$.  For an integer row
$v=(a,b,c)$ write
\[
 L(v)=a+b\gamma+c\delta,\qquad
 H(v)=\max(1,|a|,|b|,|c|).
\]
The results below replace some value-nonvanishing requirements by exact
coefficient-rank tests.  They do not construct the required infinite families.

\paragraph{Arithmetic inputs.}
Lemma~\ref{ra:lem:engine} supplies both algebraic linear independence of
$1,e,T$ and its integer-form lower bound. Theorem~\ref{ra:thm:rational}
supplies the rational-assumption floor. In the notation of this section,
these previously proved bounds read
\begin{equation}\label{ra:eq:rank-engine}
 |h_0+h_1e+h_2T|\ge C_\epsilon
 \max(1,|h_0|,|h_1|,|h_2|)^{-2-\epsilon}
 \quad(0\ne h\in\Z^3),
\end{equation}
and, under $\gamma\in\Q$,
\begin{equation}\label{ra:eq:rank-rational-floor}
 L(v)\ne0\Longrightarrow
 |L(v)|\ge C_{\gamma,\epsilon}H(v)^{-1-\epsilon}.
\end{equation}
Finally, if $\gamma$ is algebraic irrational, Roth's
theorem \cite{Roth} gives
\begin{equation}\label{ra:eq:rank-roth}
 |A+B\gamma|\ge C_{\gamma,\epsilon}
       \max(1,|A|,|B|)^{-1-\epsilon}
 \quad ((A,B)\in\Z^2\setminus\{(0,0)\}).
\end{equation}
All constants may depend on the alleged algebraic value of $\gamma$;
these bounds are asymptotic statements, not effective numerical certificates.

\begin{theorem}[Rank replaces individual nonvanishing]
\label{ra:thm:rank-irrational}
Suppose $H_n\to\infty$ and, for a fixed $\rho>0$, there are
$\Q$-linearly independent integer rows $v_n,w_n$ of heights at most
$H_n$ such that
\[
 |L(v_n)|,\ |L(w_n)|\le H_n^{-1-\rho}.
\]
Then $\gamma$ is irrational.  Neither value is separately assumed nonzero.
\end{theorem}
\begin{proof}
Assume $\gamma=r/s\in\Q$ in lowest terms.  Choose
$0<\epsilon<\rho$.  By \eqref{ra:eq:rank-rational-floor}, both displayed
values must be exactly zero for sufficiently large $n$.  Since $\delta$
is irrational under this assumption, $L(a,b,c)=0$ implies $c=0$ and
$sa+rb=0$.  Thus both rows lie on the single rational line
$\Q(-r,s,0)$, contradicting their independence.
\end{proof}

\begin{theorem}[A shared auxiliary coefficient]
\label{ra:thm:rank-shared-c}
Suppose $H_n\to\infty$, $c_n\in\Z\setminus\{0\}$, and for a fixed
$\rho>0$ there are distinct pairs
$(a_n,b_n),(a'_n,b'_n)\in\Z^2$ such that
\[
 H(a_n,b_n,c_n),\ H(a'_n,b'_n,c_n)\le H_n
\]
and
\[
 |a_n+b_n\gamma+c_n\delta|,
 \ |a'_n+b'_n\gamma+c_n\delta|\le H_n^{-1-\rho}.
\]
Then $\gamma$ is transcendental.  No regularity of the growth of $H_n$
or separate certification of value nonvanishing is required.
\end{theorem}
\begin{proof}
The rows are independent: proportional rows with the same nonzero last
coordinate would be identical.  Theorem~\ref{ra:thm:rank-irrational}
excludes rational $\gamma$.  If $\gamma$ were algebraic irrational,
subtracting the two forms would give
\[
 0<|A_n+B_n\gamma|\le2H_n^{-1-\rho},\qquad
 \max(|A_n|,|B_n|)\le2H_n,
\]
where $A_n=a_n-a'_n$, $B_n=b_n-b'_n$.  The pair is nonzero, and
irrationality makes its value nonzero.  Equation~\eqref{ra:eq:rank-roth},
with $0<\epsilon<\rho$, gives the incompatible lower bound
$C(2H_n)^{-1-\epsilon}$.
\end{proof}

\begin{corollary}[Conditional occupancy of a fixed coefficient]
If $\gamma$ is algebraic, then for each fixed $\rho>0$ and all
sufficiently large $H$, every fixed nonzero integer $c$ admits at most
one pair $(a,b)$ with $H(a,b,c)\le H$ and
$|a+b\gamma+c\delta|\le H^{-1-\rho}$.  Under rationality of $\gamma$
there are no such pairs.
\end{corollary}
\begin{proof}
In the rational case $c\ne0$ makes the value nonzero, so use
\eqref{ra:eq:rank-rational-floor}.  In the algebraic irrational case,
subtract two candidates and use \eqref{ra:eq:rank-roth} as above.  The
constants are independent of $c$.
\end{proof}

\begin{theorem}[Primitive content in auxiliary elimination]
\label{ra:thm:rank-content}
Suppose $v_n=(a_1,b_1,c_1)$ and $w_n=(a_2,b_2,c_2)$ are independent
integer rows of heights at most $H_n\to\infty$, and
\[
 |L(v_n)|,\ |L(w_n)|\le H_n^{-\kappa}
 \qquad\text{for a fixed }\kappa>1.
\]
Set
\[
 A_n=c_2a_1-c_1a_2,\qquad B_n=c_2b_1-c_1b_2.
\]
Suppose $(A_n,B_n)\ne(0,0)$ and, for some fixed $0\le\sigma<2$,
\[
 g_n:=\gcd(|A_n|,|B_n|)\ge H_n^\sigma.
\]
If $\kappa+2\sigma>3$, then $\gamma$ is transcendental.
\end{theorem}
\begin{proof}
Theorem~\ref{ra:thm:rank-irrational} excludes rationality.  Under algebraic
irrationality, divide both coefficients by $g_n$.  The resulting primitive
pair has nonzero value, height $h_n\le2H_n^{2-\sigma}$, and
\[
 0<\left|\frac{A_n+B_n\gamma}{g_n}\right|
 =\frac{|c_2L(v_n)-c_1L(w_n)|}{g_n}
 \le2H_n^{1-\kappa-\sigma}.
\]
Choose $\epsilon>0$ with
$\kappa+\sigma-1>(2-\sigma)(1+\epsilon)$.
Roth gives a lower bound $C h_n^{-1-\epsilon}$, which is at least
$C' H_n^{-(2-\sigma)(1+\epsilon)}$ and contradicts this upper bound.
\end{proof}

For $\sigma=0$ the sufficient threshold is $\kappa>3$; for
$\sigma=1$ it is $\kappa>1$.  The actual gcd is essential.  A large
raw determinant alone supplies no primitive gain.

\begin{proposition}[A power-gap consequence in the algebraic world]
\label{ra:prop:rank-gap}
Assume $\gamma$ is algebraic irrational.  Let $v,w$ be independent
integer rows with nonzero last coordinates, heights $H\le K$, and
$|L(v)|\le H^{-\kappa}$, $|L(w)|\le K^{-\kappa}$ for $\kappa>0$.
For every $\epsilon>0$ one has
\[
 K^{2+\epsilon}\ge C_{\gamma,\epsilon}
                      H^{\kappa-1-\epsilon}.
\]
\end{proposition}
\begin{proof}
With $A=c(w)a(v)-c(v)a(w)$ and
$B=c(w)b(v)-c(v)b(w)$, the pair $(A,B)$ is nonzero: otherwise the
nonzero last coordinates would make $v,w$ proportional.  Its value is
nonzero by irrationality, its height is at most $2HK$, and
\[
 |A+B\gamma|\le K H^{-\kappa}+H K^{-\kappa}
                      \le2K H^{-\kappa}.
\]
Combine this with \eqref{ra:eq:rank-roth} and rearrange.
\end{proof}
Consequently, an infinite family of distinct primitive rows with
nonzero last coordinates, exponent $\kappa>3$, and increasing heights
satisfying $\log H_{n+1}/\log H_n\to1$ would exclude algebraic
irrationality.  Such rows are pairwise independent after identifying
opposite signs, and \eqref{ra:eq:rank-rational-floor} excludes rationality
when their values are nonzero.  In the rational case the nonzero last
coordinate already guarantees this.  This is another sufficient
criterion, not a family produced here.

\begin{proposition}[An unconditional restriction on paired forms]
\label{ra:prop:rank-engine-pair}
Suppose independent rows of heights at most $H_n\to\infty$ have
$|L(v_n)|,|L(w_n)|\le H_n^{-\kappa}$.  Put
$C_n=a_1b_2-a_2b_1$, retain $A_n,B_n$ above, and let
$G_n=\gcd(|C_n|,|A_n|,|B_n|)$.
If $G_n\ge H_n^\tau$ for a fixed $0\le\tau<2$, then
$\kappa+3\tau\le5$.
\end{proposition}
\begin{proof}
The nonzero integer vector $(C_n,-A_n,-B_n)/G_n$ has height at most
$2H_n^2/G_n$.  Direct expansion gives
\[
 C_n-A_ne-B_nT
 =L(v_n)(b_2-c_2e)-L(w_n)(b_1-c_1e).
\]
Its value is nonzero by the independence of $1,e,T$.
Equation~\eqref{ra:eq:rank-engine} therefore implies
\[
 C_0 H_n^{1-\kappa}/G_n
 \ge C_\epsilon(H_n^2/G_n)^{-2-\epsilon},
 \qquad
 H_n^{5+2\epsilon-\kappa}\ge C'_{\epsilon}G_n^{3+\epsilon}.
\]
Insert $G_n\ge H_n^\tau$, let $n\to\infty$, and then
$\epsilon\downarrow0$.
\end{proof}
Here $G_n$ is the content of all three minors, while $g_n$ in
Theorem~\ref{ra:thm:rank-content} uses only two.  They must not be conflated.

\begin{remark}[What these criteria do not establish]
No required infinite family has been constructed.  In particular,
Theorems~\ref{ra:thm:rank-shared-c} and \ref{ra:thm:rank-content} are
\emph{sufficient}, not equivalent, transcendence criteria.  Their
eliminated primitive pairs give rational approximations to $\gamma$
with a fixed margin beyond exponent $2$, and therefore force
$\mu(\gamma)>2$.  Indeed their nonzero values tend to zero, so the
primitive pair heights are unbounded; small values make the heights
comparable with the reduced denominators.  The stated strict exponent
inequalities then give the asserted margin.  A transcendental $\gamma$
with $\mu(\gamma)=2$ would satisfy neither existence target.  Thus these
criteria rigorously bypass individual nonvanishing tests, but are not
evidence that a suitable family exists, and do not supersede the
logically different mixed $E$--$D$ separation target.
\end{remark}

\section{Exact information carried by the finite collision tower}
Let $\E$ denote the ring of $E$-values. To make the finite-ring statement
self-contained, we recall the connection identity and its proof from
\cite{KEIO}, with the corrected indexing. For $r\ge1$ set
\[
 \Ein_r(z)=\sum_{n\ge1}\frac{(-1)^{n-1}z^n}{n^r n!},\qquad
 P_r(\lambda)=[t^r]\Gamma(1-t)\lambda^t,
\]
and, for positive algebraic $\lambda$, define
\[
 \omega_{\lambda,r}=
 \int_0^\infty e^{-\xi}
 \frac{(-1)^{r-1}\log^{r-1}(1+\xi/\lambda)}
 {\lambda(r-1)!(1+\xi/\lambda)}\,d\xi.
\]
\begin{lemma}[Connection identity]\label{ra:lem:collision}
For all the indicated $\lambda,r$,
\[
 \omega_{\lambda,r}=e^\lambda(\Ein_r(\lambda)-P_r(\lambda)).
\]
\end{lemma}
\begin{proof}
Substitute $u=1+\xi/\lambda$. For $|t|$ small, absolute convergence
permits summation under the integral and gives
\[
 \sum_{r\ge1}\omega_{\lambda,r}t^r
 =e^\lambda t\int_1^\infty e^{-\lambda u}u^{-t-1}\,du
 =e^\lambda t\lambda^t\Gamma(-t,\lambda),
\]
where $\Gamma(\cdot,\lambda)$ is the upper incomplete Gamma function.
Its lower-series expansion gives
\[
 \Gamma(-t,\lambda)=\Gamma(-t)
 -\sum_{n\ge0}\frac{(-1)^n\lambda^{n-t}}{n!(n-t)}.
\]
Now $t\Gamma(-t)=-\Gamma(1-t)$; the $n=0$ term contributes
$\lambda^{-t}$; and $-t/(n-t)=-\sum_{k\ge1}t^k/n^k$ for $n\ge1$.
Therefore
\[
 t\Gamma(-t,\lambda)=-\Gamma(1-t)
 +\lambda^{-t}\left(1+\sum_{k\ge1}\Ein_k(\lambda)t^k\right).
\]
After multiplication by $e^\lambda\lambda^t$, comparison of the
coefficients of $t^r$ proves the identity.
\end{proof}

\begin{proposition}\label{ra:prop:ring}
If $R\ge1$ and $\Lambda\subset\QB_{>0}$ is finite with $1\in\Lambda$, then
\begin{align*}
 &\E[\omega_{\lambda,r}:\lambda\in\Lambda,\ 1\le r\le R]\\
 &\hspace{1cm}=\E[\gamma,\ \log\lambda\ (\lambda\in\Lambda),\ \zeta(2),\ldots,\zeta(R)].
\end{align*}
For $R=1$ the zeta list is empty. Therefore all the indicated collisions lie in $\E$ if and only if every displayed generator on the right lies in $\E$.
\end{proposition}
\begin{proof}
Put $g_0=1$ and $g_r=\Ein_r(1)-e^{-1}\omega_{1,r}$. The logarithmic generating function gives
\[
 \gamma=g_1,\qquad
 \zeta(k)=kg_k-\gamma g_{k-1}-\sum_{j=2}^{k-1}\zeta(j)g_{k-j}.
\]
This recursively recovers all the zeta values up to $R$ using ring operations only. The first-level relation also gives
\[
 \log\lambda=\Ein(\lambda)-e^{-\lambda}\omega_{\lambda,1}-\gamma.
\]
Conversely each $P_r(\lambda)$ is a rational polynomial in these generators, and $e^{\pm\lambda}$ and $\Ein_r(\lambda)$ are $E$-values. Both ring inclusions follow.
\end{proof}
This explains precisely what raising the tower order adds. Algebraicity of $\gamma$ does not supply membership of the other generators in $\E$. The proposition is a triangular reformulation, not a proof of the missing collisions.

\section{Binomial closure and exact arithmetic alternatives}
\label{ra:bin:section}

The polynomial ring generated by a number records only polynomial
relations with integer coefficients.  Adjoining generalized binomial
coefficients produces a larger ring whose finite-character theory is
more rigid.  This section applies the classical theory of
integer-valued polynomials; no literature-wide novelty claim is made.
The relevant framework is discussed in \cite{Elliott,Chabert}.
All statements hold for arbitrary real numbers before specialization
to Euler's constant.

\subsection{The enlarged ring and its reductions}

Set
\[
 \mathcal I=\{P\in\Q[X]:P(\Z)\subseteq\Z\},\qquad
 b_n(X)=\binom Xn,\qquad
 \mathcal B_x=\Z[b_n(x):n\ge0].
\]
The ring $\mathcal B_x$ contains $\Z[x]$.  Newton expansion gives
$\mathcal I=\bigoplus_{n\ge0}\Z b_n$, so evaluation induces a
surjection $\mathcal I\to\mathcal B_x$.  If $x$ is transcendental,
this map is injective.  If $x$ is algebraic with primitive minimal
polynomial $f\in\Z[X]$, its kernel is
\[
 J_f=f\Q[X]\cap\mathcal I.
\]
For a prime $p$ let $Z_p(f)=\{\beta\in\Z_p:f(\beta)=0\}$.
Write $C(\Z_p,\Z/p^k\Z)$ for the ring of continuous functions to
the discrete finite ring, with pointwise operations.  Such functions
are locally constant.

\begin{theorem}[Exact reduction algebras]\label{ra:bin:thm:reduction}
For every prime $p$ and integer $k\ge1$,
\[
 \mathcal B_x/p^k\mathcal B_x\simeq
 \begin{cases}
 C(\Z_p,\Z/p^k\Z),&x\notin\Qbar,\\[2mm]
 \displaystyle\prod_{\beta\in Z_p(f)}\Z/p^k\Z,&x\in\Qbar.
 \end{cases}
\]
The product over the empty set denotes the zero ring.  The unital
$\Z$-algebra homomorphisms from $\mathcal B_x$ to $\Z/p^k\Z$
are consequently indexed by $\Z_p$ in the transcendental case and
by $Z_p(f)$ in the algebraic case.  Reduction between successive
values of $k$ is a bijection on these homomorphism sets.
\end{theorem}

\begin{proof}
Every $P\in\mathcal I$ maps $\Z_p$ continuously into $\Z_p$:
apply continuity of its rational polynomial expression and density
of the integers.  Reduction of evaluation modulo $p^k$ has kernel
$p^k\mathcal I$.  Indeed, divisibility by $p^k$ at all integers
is precisely the assertion that $P/p^k$ is integer-valued.

We first prove that every locally constant function is represented.
It depends on a residue class modulo some $p^m$.  For
$a=\sum_{j<m}a_jp^j$ define
\[
 D_a(X)=\prod_{j=0}^{m-1}
 \left[1-\left(\binom X{p^j}-a_j\right)^{p-1}\right]\in\mathcal I.
\]
The digit form of Lucas's identity,
$\binom{t}{p^j}\equiv t_j\pmod p$ for
$t=\sum_{j\ge0}t_jp^j\in\Z_p$, shows that $D_a$ is $1$
modulo $p$ on $a+p^m\Z_p$ and $0$ elsewhere.  The identity first
holds for nonnegative integers, by comparing coefficients in
$(1+T)^t$ modulo $p$, and extends to $\Z_p$ by continuity.
The polynomial
\[
 E_a(X)=D_a(X)^{p^{k-1}}
\]
is the indicator of that ball modulo $p^k$.  To check this, use
$(1+pu)^{p^{k-1}}\equiv1\pmod{p^k}$ and
$p^{k-1}\ge k$; the case $k=1$ is immediate.
Integer linear combinations of the $E_a$ represent all locally
constant functions.  We have proved
\[
 \mathcal I/p^k\mathcal I\simeq C(\Z_p,\Z/p^k\Z).
\]
This proves the transcendental assertion.

In the algebraic case let $J$ be the image of $J_f$ in this function
ring.  Each function in $J$ vanishes on $Z_p(f)$.  Conversely,
fix $\beta\in\Z_p$ with $f(\beta)\ne0$, and set
$t=v_p(f(\beta))\ge0$.  On a sufficiently small residue ball
$U$ about $\beta$, the valuation of $f$ is constantly $t$.
The preceding indicator construction gives $q\in\mathcal I$
congruent to $1_U$ modulo $p^{t+1}$.  Then
\[
 g=\frac{fq}{p^t}\in J_f,\qquad g(\beta)\in\Z_p^\times.
\]
For the first assertion, at integers in $U$ the factor $f$ is
divisible by $p^t$, and outside $U$ the factor $q$ is divisible
by $p^{t+1}$.  Thus $g$ is integer-valued; the introduced denominator
has no other prime factor.  Its rational polynomial expression
is a multiple of $f$.

Let $h$ be locally constant modulo $p^k$ and vanish on $Z_p(f)$.
Its support is compact and disjoint from that finite root set.
Cover its support by finitely many clopen sets on each of which
one of the preceding functions $g$ is a unit.  Refine them to a
disjoint clopen partition.  On each part multiply its $g$ by
the locally constant function $h/g$, extended by zero outside
that part.  The resulting function belongs to $J$, and summing
the finitely many parts gives $h$.  Therefore $J$ is exactly
the ideal of functions vanishing on $Z_p(f)$.
Restriction to this finite root set is onto, since disjoint
clopen neighborhoods separate its points.  This proves the
displayed product formula.

A unital homomorphism from a locally constant function ring to
the local ring $\Z/p^k\Z$ sends the indicators of a finite
clopen partition to idempotents, selecting precisely one part.
These selections are compatible under refinement and determine
a unique point of the compact space $\Z_p$.  Constants then
force the homomorphism to be evaluation at that point.  In a
finite product the same argument selects one factor.  These
descriptions are independent of $k$, proving unique lifting.
\end{proof}

Taking inverse limits gives a further explicit consequence:
\[
 \widehat{\mathcal B_x}^{(p)}\simeq
 \begin{cases}
 C(\Z_p,\Z_p),&x\notin\Qbar,\\
 \Z_p^{\,\#Z_p(f)},&x\in\Qbar.
 \end{cases}
\]
Thus the transcendental alternative forces continuum-many characters
already modulo $p$, while the algebraic alternative permits at most
$\deg f$.  This is a structural consequence of the alternative,
not a computation of the characters of the actual gamma ring.
For example, $\mathcal B_{\sqrt{17}}$ has exactly two characters
to $\Z/2^k\Z$ for every $k$, with no transient congruence-root
branches.  In $\mathcal B_{\sqrt2}$ the number $2$ is a unit:
$\binom{\sqrt2}{2}=1-\sqrt2/2$ yields $\sqrt2/2$ and then
$1/2$.  Hence its reduction modulo $2$ is the zero ring.

\subsection{An explicit identity forced in the algebraic alternative}

For $N\ge1$ define the subgroup of finite sums
\[
 T_N(x)=\sum_{n\ge N}\Z\binom xn.
\]
Before evaluation this is the ideal of $\mathcal I$ vanishing at
$0,1,\ldots,N-1$: Newton expansion identifies vanishing at these
points with vanishing of the first $N$ binomial coefficients.
It follows that $T_N(x)$ is an ideal of $\mathcal B_x$.

\begin{theorem}[Tail quotient and constructive unit]
\label{ra:bin:thm:tail}
If $x$ is transcendental, then
\[
 \mathcal B_x/T_N(x)\simeq\Z^N.
\]
If $x$ is algebraic and $x\notin\{0,\ldots,N-1\}$, then
$T_N(x)=\mathcal B_x$.  More explicitly, if $f\in\Z[X]$
is its minimal polynomial of degree $d$ and
$M=\max_{0\le m<N}|f(m)|$, there are computable integers
$c_N,\ldots,c_K$ such that
\[
 1=\sum_{n=N}^Kc_n\binom xn,
 \qquad K\le N-1+dM.
\]
If $x=m\in\{0,\ldots,N-1\}$, the quotient instead equals $\Z$.
\end{theorem}

\begin{proof}
In the transcendental case evaluation is injective.  Restriction
of $\mathcal I$ to the $N$ integer points is onto $\Z^N$,
because the binomial interpolation matrix is unit triangular.
Its kernel is the tail ideal, proving the first assertion.

In the algebraic case under consideration each $f(m)$ is nonzero.
Set
\[
 L_m(X)=\binom Xm\binom{N-1-X}{N-1-m},\qquad
 P_m(X)=\binom{\operatorname{sgn}(f(m))f(X)}{|f(m)|}.
\]
Both polynomials are integer-valued.  On $0\le j<N$ one has
$L_m(j)=\delta_{mj}$, and $P_m(m)=1$.
Furthermore $f$ divides $P_m$ in $\Q[X]$, since its binomial
expression has a factor $f(X)$.  Consequently
\[
 F_N(X)=\sum_{m=0}^{N-1}L_m(X)P_m(X),\qquad
 Q_N(X)=1-F_N(X)
\]
satisfy $Q_N(j)=0$ for $0\le j<N$, $Q_N(x)=1$, and
$\deg Q_N\le N-1+dM$.  Expand $Q_N$ in the binomial basis:
\[
 Q_N(X)=\sum_{n=0}^{\deg Q_N}\Delta^nQ_N(0)\binom Xn.
\]
The coefficients are integers, and the first $N$ are zero.
This yields the required identity and proves that the tail ideal
contains $1$.  Finally, for $x=m\in\{0,\ldots,N-1\}$ the
ring is $\Z$ and every tail value is zero.
\end{proof}

\begin{corollary}[A proper-ideal target at Euler's constant]
\label{ra:bin:cor:gamma}
For Euler's constant, which satisfies $0<\gamma<1$, the following
are equivalent:
\begin{enumerate}
\item $\gamma$ is transcendental;
\item $T_N(\gamma)$ is a proper ideal for some $N\ge1$;
\item $\mathcal B_\gamma/T_N(\gamma)\simeq\Z^N$ for every $N\ge1$.
\end{enumerate}
For any fixed prime $p$, these conditions are also equivalent to
the existence of a unital homomorphism
\[
 \phi:\mathcal B_\gamma\longrightarrow\Z/p\Z,
 \qquad \phi\!\left(\binom\gamma n\right)=0\quad(n\ge1).
\]
\end{corollary}

\begin{proof}
The first three equivalences follow from Theorem~\ref{ra:bin:thm:tail}.
Under transcendence, evaluation at $X=0$ on $\mathcal I$, followed
by reduction modulo $p$, gives the indicated homomorphism.
Under algebraicity the tail-unit identity with $N=1$ would map
$1$ to $0$, so no such homomorphism exists.
\end{proof}

The corollary states an exact remaining target.  It does not prove
that target for $\gamma$.  Prescribing images of the generators
does not construct a homomorphism: preservation of all their
integer relations is essential, and the algebraic tail identity
exhibits the precise obstruction.

\subsection{Coefficient costs and the topology obstruction}

For every real $0<x<1$, independently of its arithmetic nature,
\[
 \left|\binom xn\right|
 =\frac{x\Gamma(n-x)}{\Gamma(1-x)\Gamma(n+1)}\quad(n\ge1).
\]
These values decrease strictly, since consecutive magnitudes have
ratio $(n-x)/(n+1)$.  The exact tail sum is
\[
 S_N(x)=\sum_{n\ge N}\left|\binom xn\right|
 =\frac{\Gamma(N-x)}{\Gamma(1-x)\Gamma(N)}
 =\prod_{j=1}^{N-1}\left(1-\frac{x}{j}\right).
\]
Indeed the product tends to zero, starts at $S_1=1$, and its
successive difference is $xS_N/N=|\binom xN|$.

\begin{proposition}[Necessary and constructive coefficient budgets]
\label{ra:bin:prop:cost}
Every finite identity $1=\sum_{n=N}^Kc_n\binom xn$ with
$0<x<1$ satisfies
\[
 \sum_{n=N}^K|c_n|\ge\left|\binom xN\right|^{-1},\qquad
 \max_{N\le n\le K}|c_n|>
 \frac{\Gamma(1-x)\Gamma(N)}{\Gamma(N-x)}.
\]
The respective lower bounds are asymptotic to
$\Gamma(1-x)N^{1+x}/x$ and $\Gamma(1-x)N^x$.

For the construction of Theorem~\ref{ra:bin:thm:tail}, put
$K_0=N-1+dM$ and $T=\max_{0\le j\le K_0}|f(j)|$.
There are at most $dM$ coefficient positions, and
\[
 \max_n|c_n|\le
 2^{K_0}\bigl[1+N(K_0+N)^{N-1}(T+M)^M\bigr],
 \qquad
 T\le(d+1)H(f)\max(1,K_0)^d.
\]
In particular the construction permits
$\log\max(1,\max_n|c_n|)=O_f(N^d\log(N+1))$.
\end{proposition}

\begin{proof}
The triangle inequality and the largest tail magnitude prove
the first lower bound.  With $H=\max|c_n|$, positivity of the
infinitely many omitted magnitudes gives
\[
 1\le H\sum_{n=N}^K\left|\binom xn\right|<HS_N(x),
\]
which proves the strict second bound.  The stated asymptotics
follow from the usual ratio asymptotic for the Gamma function.
For the upper estimate, at $0\le j\le K_0$ bound
$|L_m(j)|$ by $(K_0+N)^{N-1}$ and $|P_m(j)|$ by
$(T+M)^M$.  Thus
$|Q_N(j)|\le1+N(K_0+N)^{N-1}(T+M)^M$.
Use $|\Delta^nQ_N(0)|\le2^n\max_{j\le n}|Q_N(j)|$.
The bound for $T$ follows by bounding each monomial of $f$;
for fixed $f$, $M=O_f(N^d)$ and $K_0=O_f(N^d)$.
\end{proof}

The budgets do not conflict.  In particular, tail-unit identities with a coefficient bound independent of $N$
are impossible as $N\to\infty$ for every $x\in(0,1)$,
including rational numbers.  This does not exclude the algebraic
alternative, whose construction permits increasing coefficients.
The elementary identity $-8\binom{1/2}{2}=1$ already disproves
the claim that small tail values cannot have integer span
containing $1$; Theorem~\ref{ra:bin:thm:tail} gives such identities
starting arbitrarily far out.

\begin{remark}[Finite ring operations versus real infinite sums]
The positive weights $w_n(x)=(-1)^{n-1}\binom xn$ satisfy
\[
 \sum_{n\ge1}w_n(x)=1,\qquad
 \sum_{n\ge1}w_n(x)z^n=1-(1-z)^x
 \quad(0<x<1).
\]
These identities hold for rational and transcendental parameters
alike.  A character of $\mathcal B_x$ acts on finite ring
expressions; it need not preserve a real-convergent infinite sum.
In the transcendental case the character at the $p$-adic point
$0$ sends every $b_n(x)$, $n\ge1$, to zero.  Applying that
character termwise to the displayed real sum would give a false
contradiction.  Continuity for the real topology has not been
proved and is not part of a ring homomorphism.  Similarly, small
singularity sets of these generating functions do not supply the
integer-coefficient or denominator hypotheses of arithmetic
rationality criteria.
\end{remark}

The recovered companion archive contains
\texttt{check\_binomial\_tail.py}. Its eleven checks for
$f=2X-1$, $2X^2-1$, and $X^2+X-1$ were rerun successfully: each checks
integer Newton coefficients, vanishing below $N$, and exact divisibility
of $Q_N-1$ by $f$. These finite examples supplement the quantified
polynomial proof above.


\section{Rational Borel kernels and arithmetic division}
\label{ra:sec:borel}

The earlier Euler-series calculation extends to every rational Borel
kernel.  The resulting classification is an application of exponential
partial-fraction reduction and the regular-singularity theorem for
$G$-operators; it is not a claim of literature-wide priority.  Crucially,
it classifies formal arithmetic division, not the arithmetic nature of an
integral merely because the corresponding analytic quotient is regular.

For a formal series $u(x)=\sum_{n\ge0}u_nx^n$, put
\[
 \mathcal B u(z)=\sum_{n\ge0}\frac{u_n}{n!}z^n.
\]
We use the convention that $u$ is arithmetic Gevrey of order~$1$ if
$\mathcal B u$ is a $G$-function; see \cite{FRmixed}.  Let
$\alpha\in\Qbar\setminus\{0\}$, let
$R(z)=\sum_{n\ge0}r_nz^n\in\Qbar(z)$ be regular at~$0$, and define
\begin{equation}\label{ra:bor:def}
 f_R(x)=\sum_{n\ge0}n!r_nx^n,\qquad
 g_a(x)=\frac{f_R(x)-a}{x-\alpha},\qquad
 G_a=\mathcal B g_a,\quad a\in\Qbar.
\end{equation}
The quotient in \eqref{ra:bor:def} is taken in $\Qbar[[x]]$.

\begin{theorem}[Complete rational-kernel classification]\label{ra:bor:class}
The series $g_a$ is arithmetic Gevrey of order~$1$ if and only if
there exists $T\in\Qbar(z)$, regular at~$0$, such that
\begin{equation}\label{ra:bor:exact}
 R=T-\alpha T',\qquad a=T(0).
\end{equation}
In this case $T$ is unique and $G_a=T'$ is rational.
\end{theorem}

\begin{proof}
Write $g_a(x)=\sum c_nx^n$ and $b_n=c_n/n!$.  Coefficient comparison
in $(x-\alpha)g_a=f_R-a$ gives
\[
 b_0=\frac{a-r_0}{\alpha},\qquad
 \alpha n b_n=b_{n-1}-nr_n\quad(n\ge1).
\]
Consequently the convergent germ $G_a$ satisfies
\begin{equation}\label{ra:bor:ode}
 \alpha G_a'-G_a=-R',\qquad
 G_a(0)=\frac{a-R(0)}{\alpha}.
\end{equation}

At each finite pole $\beta$ of $R$, write the principal part as
$\sum_{j=1}^{m_\beta}r_{\beta,j}(z-\beta)^{-j}$ and put
\begin{equation}\label{ra:bor:residue}
 c_\beta=
 \sum_{j=1}^{m_\beta}
 \frac{r_{\beta,j}(-1/\alpha)^{j-1}}{(j-1)!}.
\end{equation}
Successively cancelling the highest pole orders, and then the polynomial
part, gives a rational function $T$, regular at~$0$, with
\begin{equation}\label{ra:bor:reduction}
 R=T-\alpha T'+\sum_\beta\frac{c_\beta}{z-\beta}.
\end{equation}
For completeness, a pole of order $j\ge2$ is removed using a multiple
of $(z-\beta)^{-j+1}$ in $T$.  The map $T\mapsto T-\alpha T'$ is
triangular with diagonal~$1$ on polynomials.  The remaining simple-pole
coefficient is exactly \eqref{ra:bor:residue}, since
\[
 \operatorname{Res}_{z=\beta}\bigl(e^{-z/\alpha}R(z)\,dz\bigr)
   =e^{-\beta/\alpha}c_\beta,
 \qquad
 e^{-z/\alpha}(T-\alpha T')
   =-\alpha\frac{d}{dz}(e^{-z/\alpha}T).
\]
These residues also prove uniqueness of the coefficients $c_\beta$.
The remaining ambiguity in $T$ would solve $T-\alpha T'=0$; a nonzero
solution is a constant multiple of $e^{z/\alpha}$ and is not rational.

An integrating factor in \eqref{ra:bor:ode} yields
\[
 G_a(z)=e^{z/\alpha}\left(
 \frac{a-R(0)}{\alpha}
 -\frac1\alpha\int_0^z e^{-t/\alpha}R'(t)\,dt\right).
\]
Since the residue of a derivative is zero,
\[
 \operatorname{Res}_\beta(e^{-t/\alpha}R'(t)\,dt)
 =\alpha^{-1}e^{-\beta/\alpha}c_\beta.
\]
Analytic continuation counterclockwise around $\beta$ therefore gives
\begin{equation}\label{ra:bor:monodromy}
 \Delta_\beta G_a(z)
 =-\frac{2\pi i}{\alpha^2}e^{(z-\beta)/\alpha}c_\beta.
\end{equation}
If $G_a$ were a $G$-function and some $c_\beta$ were nonzero, two
continued solutions of its minimal homogeneous differential equation
would have a nonzero exponential difference.  The minimal operator of
a $G$-function is regular singular, including at infinity; this standard
$G$-operator input is recalled in \cite[\S1, p.~3]{FRoperators}.
Every solution of such an operator has power--logarithmic local behaviour
after finite ramification.  The exponential $e^{z/\alpha}$ has an
irregular singularity at infinity and cannot be such a solution.
Thus every $c_\beta$ is zero.

We now have $R=T-\alpha T'$.  Substituting in \eqref{ra:bor:ode}, including
its initial condition, gives
\begin{equation}\label{ra:bor:solution}
 G_a(z)=T'(z)+\frac{a-T(0)}{\alpha}e^{z/\alpha}.
\end{equation}
Rational functions regular at~$0$ are $G$-functions.  If $a\ne T(0)$,
closure under subtraction would make $e^{z/\alpha}$ a $G$-function,
contradicting the same regular-singularity property.  Hence $a=T(0)$.
Conversely, \eqref{ra:bor:exact} makes \eqref{ra:bor:solution} equal to $T'$,
which proves arithmetic Gevrey membership and completes the proof.
\end{proof}

\begin{corollary}[Denominator obstruction]\label{ra:bor:denominator}
Fix a number field $K$ containing $a,\alpha$ and the coefficients of $R$.
Let $D_N$ be the least positive rational integer such that
$D_Nb_0,\ldots,D_Nb_N$ belong to the ring of integers of $K$.
If \eqref{ra:bor:exact} does not hold, then
\[
 \limsup_{N\to\infty}\frac{\log D_N}{N}=+\infty.
\]
\end{corollary}
\begin{proof}
Equation~\eqref{ra:bor:ode} gives holonomicity.  For every embedding of $K$
into $\C$, its rational forcing term is analytic in a neighbourhood
of~$0$ and its leading coefficient is a nonzero constant.  The
corresponding solution is analytic there.  Since there are finitely
many embeddings, Cauchy's estimates give exponential bounds for all
conjugates of the $b_n$.  Thus an exponential bound on $D_N$ would supply
the remaining $G$-function condition, contrary to Theorem~\ref{ra:bor:class}.
Failure of every such bound is the stated infinite limsup.
\end{proof}
We do not assert a limit, an effective growth rate, or a universal
factorial asymptotic for these denominators.

\subsection{What analytic evaluation does not supply}
Suppose now that $\alpha>0$ and $R$ has no poles on $[0,\infty)$.
The positive-direction Borel sum at $\alpha$ is
\[
 \mathcal L_\alpha(R)=\int_0^\infty e^{-t}R(\alpha t)\,dt.
\]
If \eqref{ra:bor:exact} holds, integration by parts gives
\begin{equation}\label{ra:bor:integral-exact}
 \mathcal L_\alpha(R)
 =\int_0^\infty e^{-t}
   \left(T(\alpha t)-\frac{d}{dt}T(\alpha t)\right)dt
 =T(0)=a.
\end{equation}
The rational function $T$ has no additional finite poles, so the
integration is legitimate and its boundary term vanishes at infinity.
The converse assertion
\begin{equation}\label{ra:bor:missing}
 \mathcal L_\alpha(R)=a\in\Qbar
 \quad\Longrightarrow\quad
 R=T-\alpha T',\quad a=T(0)
\end{equation}
is \emph{not} proved here.  It is a zero-removal assertion for this
family, within the framework of \cite[Conjecture~2]{FRmixed}.
For $R(z)=1/(1+z)$ and $\alpha=1$, the obstruction is $c_{-1}=1$,
whereas $\mathcal L_1(R)=\delta$.  Thus \eqref{ra:bor:missing} for this
one kernel would already prove the transcendence of $\delta$.

Arithmetic division is not always obstructed: for
\[
 R(z)=\frac1{1+z}+\frac1{(1+z)^2},\quad
 T(z)=\frac1{1+z},\quad \alpha=a=1,
\]
we have $G_a=T'=-1/(1+z)^2$ and $\mathcal L_1(R)=1$.
These examples distinguish the proved formal criterion from its
unproved analytic converse.

\subsection{The whole one-pole hierarchy has two generators}
\begin{proposition}[Exact one-pole reduction]\label{ra:bor:onepole}
Suppose the only possible finite pole of $R\in\Qbar(z)$ is $-1$.
If its principal part is $\sum_{j=1}^m r_j(1+z)^{-j}$, then
\[
 R=T-T'+\frac{c}{1+z},\qquad
 c=\sum_{j=1}^m\frac{(-1)^{j-1}r_j}{(j-1)!},
 \qquad
 \mathcal L_1(R)=T(0)+c\delta.
\]
Consequently the value space of the whole family is exactly
$\Qbar+\Qbar\delta$.  Under $\mathrm A$, this space has dimension~$2$
over $\Qbar$ and is contained in the ring $\mathbf E$ of $E$-values.
\end{proposition}
\begin{proof}
The reduction and integral formula follow from
\eqref{ra:bor:reduction} and integration by parts.  The kernels $1$ and
$1/(1+z)$ supply the two generators, proving equality of the value spaces.
Under $\mathrm A$, the function
\[
 F_\gamma(z)=e^z(\Ein(z)-\gamma)
\]
is an $E$-function with $F_\gamma(1)=\delta$.  Lemma~\ref{ra:lem:engine} gives algebraic linear independence of $1,e,T$.
If $\delta$ were algebraic, the identity
$T-\gamma e-\delta=0$ would contradict that independence.  Thus $\delta$ is transcendental under $\mathrm A$, and
the dimension is exactly~$2$.  An explicit $E$-function witness for
$\mathcal L_1(R)$ is $T(0)+cF_\gamma(z)$.
\end{proof}
Unconditionally the dimension is \emph{at most}~$2$, since this argument
does not determine the arithmetic nature of $\delta$.
Under $\mathrm A$ we obtain the consistent classification
\[
 \mathcal L_1(R)\in\Qbar
 \quad\Longleftrightarrow\quad c=0
 \quad\Longleftrightarrow\quad R=T-T'.
\]
Increasing the pole order or polynomial numerator in this family creates
no further linearly independent collision.  In particular it does not
contradict Theorem~\ref{ra:bor:class}.

The supplied companion also invokes a mixed function from earlier
manuscripts to motivate analytic division at $1$. That construction is not
reproduced or needed in the rational-kernel proof above. The
analytic-to-arithmetic implication in \cite[Conjecture~3]{FRmixed}
remains an additional hypothesis; the classification does not establish it.


\section{A stronger deformation and the normalization it still loses}
Write $\gamma_{\rm low}(a,z)$ for the lower incomplete gamma function, to avoid confusion with Euler's constant.
\begin{proposition}[A periodic deformation]\label{ra:prop:periodic}
For real $t$ put
\[
 h_t(a)=e^{t\sin(2\pi a)},\quad G_t(a)=h_t(a)\Gamma(a),\quad
 U_t(a,z)=G_t(a)-\gamma_{\rm low}(a,z),\qquad z>0.
\]
Then $G_t(a+1)=aG_t(a)$, $G_t(n+1)=n!$, and
\[
 G_t(a)G_t(1-a)=\frac\pi{\sin\pi a}.
\]
It has the same poles and residues as $\Gamma$, is positive on $a>0$, and $U_t(\cdot,z)$ is entire. Moreover
\[
 U_t(a+1,z)=aU_t(a,z)+z^ae^{-z},\qquad
 \partial_zU_t(a,z)=-z^{a-1}e^{-z},\qquad
 -G_t'(1)=\gamma-2\pi t.
\]
For $t\ne0$, it does not preserve the canonical upper-tail limit at infinity or logarithmic convexity.
\end{proposition}
\begin{proof}
The function $h_t$ is entire, zero-free, $1$-periodic, equals $1$ at integers, and satisfies $h_t(a)h_t(1-a)=1$. These facts give recurrence, integer values, reflection and unchanged residues. Subtracting the meromorphic lower incomplete function cancels every pole in $a$. Its recurrence and $z$-derivative give the two formulas for $U_t$. Since $h_t'(1)=2\pi t$, differentiation gives the shifted slope.

For $a>0$,
\[
 \lim_{z\to\infty}U_t(a,z)=(h_t(a)-1)\Gamma(a).
\]
If this is zero for every positive $a$, then $h_t\equiv1$, forcing $t=0$. Also
$(\log G_t)''(a)=\psi'(a)-4\pi^2t\sin(2\pi a)$; since $\psi'(a)\to0$ on the positive real axis, it becomes negative along an arithmetic progression when $t\ne0$.
\end{proof}
This strengthens the manuscript's deformation example while retaining more classical identities. It does not claim preservation of all Stokes data.

\begin{lemma}[All upper-tail positivity fixes the slope]
Let $G$ be differentiable on a real neighborhood of $1$, with $G(1)=1$. If
$G(a)-\gamma_{\rm low}(a,z)\ge0$ for every $z>0$ and every $a$ in that neighborhood, then $G'(1)=-\gamma$.
\end{lemma}
\begin{proof}
Let $z\to\infty$. Then $G(a)\ge\Gamma(a)$ and equality holds at $a=1$. The differentiable difference has a local minimum there, so its derivative is zero.
\end{proof}
Thus stronger analytic normalization does select the actual constant. The unresolved step is an arithmetic conclusion about that selected derivative. Merely identifying the correct member of a deformed family does not establish transcendence.

\section{Gamma multiplication, dynamics, and a finite contradiction}
\label{ra:sec:normalization}

This section imports Fourier analysis and invariant measures into the
normalization problem. All its conclusions are unconditional statements in
world C. They do not assume either arithmetic alternative for $\gamma$.

For a real $1$-periodic function $p$ and an integer $m\ge2$, write
\begin{equation}\label{ra:eq:Mm}
 (M_m)\qquad \sum_{j=0}^{m-1}p(x+j/m)=p(mx).
\end{equation}
Then $F(x)=\Gamma(x)e^{p(x)}$ retains the Gamma recurrence and its exact
Gauss $m$-multiplication formula. The additional conditions $p(0)=0$ and
$p(-x)=-p(x)$ give normalization at $1$ and the reflection formula,
respectively. We use the usual multiplication formula from \cite{DLMF}.

\subsection{A regularity boundary}

\begin{theorem}[Absolute-continuity rigidity]\label{ra:thm:AC}
An absolutely continuous real periodic solution of one equation $(M_m)$,
with $m\ge2$, is identically zero.
\end{theorem}
\begin{proof}
Integration gives $(m-1)\int_0^1p=0$. Let
$a_n=\int_0^1p(x)e^{-2\pi inx}\,dx$. Comparing the coefficient at
frequency $mn$ gives $m a_{mn}=a_n$. For $f=p'\in L^1$ we therefore
have $\widehat f(mn)=\widehat f(n)$. The Riemann--Lebesgue lemma along
$m^k n\to\infty$ makes every nonzero Fourier coefficient of $f$ vanish;
the constant coefficient also vanishes by periodicity. Fourier uniqueness
and absolute continuity imply that $p$ is constant, and its mean is zero.
\end{proof}

In particular, a positive $C^1$ function on $(0,\infty)$ satisfying the
Gamma recurrence and one exact Gauss multiplication formula must be
$\Gamma$: apply the theorem to $\log(F/\Gamma)$ and extend periodically.
Analytic Gamma uniqueness has classical antecedents; see
\cite[Section 2.11]{PerezMarco}. No originality is claimed for this
uniqueness principle.

\begin{theorem}[Continuous solutions retaining any finite list]
\label{ra:thm:smoothseries}
Let $S$ be a finite nonempty set of primes, and let
\[
 \mathcal N_S=\left\{\prod_{q\in S}q^{v_q}:v_q\in\Z_{\ge0}\right\},
 \qquad
 p_S(x)=\sum_{n\in\mathcal N_S}\frac{\sin(2\pi nx)}n.
\]
Then $p_S$ is a nonzero continuous, odd, $1$-periodic function with
$p_S(0)=0$. It satisfies $(M_m)$ for every $m\in\mathcal N_S$, $m\ge2$,
and is H\"older continuous of every exponent $0<\alpha<1$.
It is not absolutely continuous.
\end{theorem}
\begin{proof}
The series converges absolutely and uniformly because
$\sum_{n\in\mathcal N_S}n^{-1}=\prod_{q\in S}(1-q^{-1})^{-1}<\infty$.
Its coefficient at frequency $1$ is $1/(2i)$, so it is nonzero.
Summing a frequency over $x+j/m$ annihilates it unless $m\mid n$.
For the surviving frequencies put $n=m\ell$; division by an
$S$-smooth integer preserves membership in $\mathcal N_S$. This proves
$(M_m)$ term by term. Moreover,
\[
 |p_S(x+t)-p_S(x)|\le
 C_\alpha |t|^\alpha
 \prod_{q\in S}(1-q^{\alpha-1})^{-1}.
\]
This follows from $|\sin(u+v)-\sin u|\le C_\alpha |v|^\alpha$.
Absolute continuity would contradict Theorem~\ref{ra:thm:AC}.
\end{proof}

Given any finite list of integer multipliers, choose $S$ to contain their
prime factors. Thus $\Gamma(x)\exp(c p_S(x))$, for $c\ne0$, retains that
entire list as well as recurrence, normalization, positivity, and
reflection. These are real continuous functions, not meromorphic Gamma
functions with all their defining analytic properties. In particular,
no derivative at $1$ is being assigned to them.

\subsection{An exact dictionary with an external open problem}

Let $T_m(x)=mx\pmod1$ and $R(x)=-x\pmod1$ on the circle, and let
$\lambda$ be normalized Lebesgue measure.

\begin{theorem}[Invariant-measure dictionary]\label{ra:thm:dictionary}
The following classes are in bijection:
\begin{enumerate}
\item atomless Borel probability measures $\mu$ invariant under
$T_2,T_3$, and $R$;
\item continuous real $1$-periodic functions $p$ with $p(0)=0$,
$p(-x)=-p(x)$, $(M_2)$, $(M_3)$, and $x+p(x)$ nondecreasing on $[0,1]$.
\end{enumerate}
The maps are
\[
 p(x)=\mu([0,x])-x\quad(0\le x\le1),\qquad \mu=dx+dp.
\]
Lebesgue measure corresponds exactly to $p=0$.
\end{theorem}
\begin{proof}
For an atomless probability, its distribution function $F$ is continuous
with endpoints $0,1$. Reflection symmetry gives $F(1-x)=1-F(x)$,
so $p=F-x$ extends periodically to an odd function of mean zero.
For $n\ne0$, distributional differentiation gives
\begin{equation}\label{ra:eq:measureFourier}
 \widehat\mu(n)=2\pi in\widehat p(n).
\end{equation}
Invariance under $T_m$ is equivalent to
$\widehat\mu(mn)=\widehat\mu(n)$. Consequently
$m\widehat p(mn)=\widehat p(n)$. These are exactly the nonconstant
Fourier identities for $(M_m)$; its constant coefficient vanishes
because $p$ has mean zero. Fourier uniqueness and continuity prove
the identities pointwise.

Conversely, $F=x+p$ is a continuous distribution function. Its
Stieltjes measure is an atomless probability. Equation
\eqref{ra:eq:measureFourier} and $(M_m)$ prove invariance, and oddness
proves reflection symmetry. The constructions are inverse.
\end{proof}

The assertion that class (2) contains only $p=0$ is equivalent to the
atomless form of Furstenberg's $\times2,\times3$ measure-rigidity
conjecture, stated in \cite[Conjecture 1.3]{BadeaGrivaux}.
Indeed, given a non-Lebesgue invariant atomless measure, symmetrize it
as $\nu=(\mu+R_*\mu)/2$. If $\nu=\lambda$, then $\mu\le2\lambda$.
The Fourier coefficients of this absolutely continuous measure tend
to zero while being constant along powers of $2$; hence $\mu=\lambda$,
a contradiction. Thus symmetrization preserves nontriviality.
This is an equivalent reformulation, not a solution of that conjecture.

\subsection{A completed finite contradiction}

\begin{theorem}[Positive-measure obstruction]\label{ra:thm:gram}
For every finite nonempty $S$ and every real $c\ne0$, the function
$x+c p_S(x)$ is not nondecreasing on $[0,1]$.
\end{theorem}
\begin{proof}
Suppose it were nondecreasing. Then $\mu=dx+c\,dp_S$ would be a
probability measure, with Fourier coefficients
\[
 \widehat\mu(0)=1,\qquad
 \widehat\mu(n)=
 \begin{cases}
 a:=\pi c,& |n|\in\mathcal N_S,\\
 0,& |n|\notin\mathcal N_S,
 \end{cases}\quad n\ne0.
\]
These coefficients follow by distributional differentiation of the
uniformly convergent series; we have not assumed beforehand that
$dp_S$ is a measure.

Choose a prime $r\notin S$ and $q\in S$. With $L=r-1$, set
$n_j=q^{Lj}$ for $j\ge1$. Each $n_j$ is $S$-smooth, while each
nonzero difference $n_j-n_k$ is divisible by $r$ and is not
$S$-smooth. The Gram matrix of
$1,e^{2\pi in_1x},\ldots,e^{2\pi in_Kx}$ in $L^2(\mu)$ is therefore
\begin{equation}\label{ra:eq:Gram}
 M_K=\begin{pmatrix}1&a\boldsymbol1^T\\
 a\boldsymbol1&I_K\end{pmatrix}.
\end{equation}
Its Schur complement must satisfy $1-Ka^2\ge0$. Choosing $K>1/a^2$
contradicts this. Explicitly, the nonnegative trigonometric polynomial
\[
 U_K(x)=\left|1-a\sum_{j=1}^K e^{2\pi in_jx}\right|^2
\]
would obey $0\le\int U_K\,d\mu=1-Ka^2<0$.
\end{proof}

For $S=\{2,3\}$ take $r=7$, $q=2$, and $L=3$, so $n_j=8^j$.
At $c=1/\pi$ already $K=2$ gives the exact certificate
\[
 M_2=\begin{pmatrix}1&1&1\\1&1&0\\1&0&1\end{pmatrix},
 \qquad\det M_2=-1.
\]
This is a genuine contradiction to the added positivity assertion.
Neither A nor B has been proved to force that assertion, so it does
not reject either arithmetic world. Abstractly, the proof is Bessel's
inequality: infinitely many orthogonal unit vectors cannot all have
the same nonzero inner product with the vector $1$.

\begin{proposition}[Failure of bounded variation]\label{ra:prop:BV}
The function $p_S$ is not of bounded variation.
\end{proposition}
\begin{proof}
If it were, $\nu=dp_S/\pi$ would be a finite signed measure with
$\widehat\nu=1_{E_S}$, where $E_S=\mathcal N_S\cup(-\mathcal N_S)$.
Fourier uniqueness would give $\nu*\nu=\nu$. The classical idempotent
theorem, in the form stated in \cite[Theorem 1.1]{GreenSanders}, implies
that $E_S$ is in the coset ring of $\Z$. Such a set is periodic outside
a finite set: its finite Boolean description uses only residue
classes and singleton cosets. Every infinite such set has positive
symmetric density. But $E_S$ is infinite and
\[
 \#(E_S\cap[-X,X])\le
 2\prod_{q\in S}\left(1+\frac{\log X}{\log q}\right)=o(X).
\]
This is a contradiction.
\end{proof}

The deeper external theorem is used only in this strengthening.
The finite contradiction in Theorem~\ref{ra:thm:gram} is elementary.

\subsection{Why extremely small numerical defects do not settle the issue}

\begin{proposition}[Real approximation with any prescribed derivative]
\label{ra:prop:instability}
Fix $m\ge2$ and $a\in\mathbb R$. For $N\ge0$ let
\[
 q_N(z)=\sum_{j=0}^{N}m^{-j}\sin(2\pi m^jz),\qquad
 c_N=\frac{\gamma-a}{2\pi(N+1)},\qquad
 G_N(z)=\Gamma(z)e^{c_Nq_N(z)}.
\]
Then $G_N$ retains recurrence, reflection, normalization, positivity
on $(0,\infty)$, and the Gamma divisor and residues. Also
$-G_N'(1)=a$. Its real logarithmic multiplication defect is
\[
 R_N(x)=-c_Nm^{-N}\sin(2\pi m^{N+1}x),
 \qquad \|R_N\|_{\infty,\mathbb R}
 \le\frac{|\gamma-a|m^{-N}}{2\pi(N+1)}.
\]
The quotient $G_N/\Gamma$ converges uniformly to $1$ on the real axis.
\end{proposition}
\begin{proof}
The function $q_N$ is entire, odd, periodic, and vanishes at all
integers; $q_N'(1)=2\pi(N+1)$. This proves the exact properties.
Root-of-unity averaging gives the telescoping identity
\[
 \sum_{r=0}^{m-1}q_N(x+r/m)-q_N(mx)
 =-m^{-N}\sin(2\pi m^{N+1}x).
\]
It is precisely the log defect of the Gauss formula after multiplication
by $c_N$. Finally
$\|c_Nq_N\|_{\infty,\mathbb R}
\le |\gamma-a|/[2\pi(N+1)(1-m^{-1})]\to0$.
\end{proof}

When $a\ne\gamma$ these functions do not converge uniformly on a fixed
complex neighborhood of $1$; their fixed derivative discrepancy already
precludes that by Cauchy's estimates. They are not arithmetic models of
the actual Gamma function: the parameter $c_N$ uses the actual $\gamma$.
They prove precisely that tiny real defects cannot be promoted to
exact analytic normalization or to an arithmetic conclusion.

The missing implication for an A-world contradiction is therefore
substantive: A would have to force a nonzero absolutely continuous
solution of $(M_m)$, or some other object actually excluded by a theorem.
No such implication has been obtained.

\section{A finite algebraic camouflage certificate}
The review archive supplies an exact rational-arithmetic certificate for a
single degree-twelve example. Its polynomial is unambiguously specified by
\begin{align*}
 A_0&=577215664901532860606512090082,\\
 M_0&=6220800\cdot10^{45},\\
 f(X)&=f_0(X)+M_0(10^{30}X-A_0).
\end{align*}
where the coefficients of $f_0$ in ascending order are
\[
 \begin{split}
 (&3732480,1140480,1104480,371880,2140,38190,897935,\\
  &846090,606255,1211490,1925,746430,1).
 \end{split}
\]
The script verifies Eisenstein's criterion at $5$, coefficientwise
congruence to $X(X-1)\cdots(X-11)$ modulo $2^{10}3^5$, and the
minimal possible monic degree $12$ for that all-residue congruence.
All positive-degree coefficients of $f$ are positive. Opposite signs at
\[
 \frac{A_0}{10^{30}}-10^{-75},\qquad \frac{A_0}{10^{30}}
\]
therefore isolate its unique positive real root $\xi$. Defining
$N_{p,k}(\xi)=\#\{u\bmod p^k:f(u)\equiv0\pmod{p^k}\}$, exhaustive
integer enumeration verifies
\[
 N_{2,10}(\xi)=1024,\qquad N_{3,5}(\xi)=243,\qquad
 N_{2,11}(\xi)=512.
\]
For the real comparison, the certificate uses the positive-real digamma
remainder bound in \cite{DLMF}, with harmonic index $100$ and $20$
Bernoulli terms, together with rational logarithm bounds from $150$
terms of the expansion
\[
 \log x=2\sum_{j=0}^{L-1}\frac{z^{2j+1}}{2j+1}+R_L,
 \quad z=\frac{x-1}{x+1},\quad
 0<R_L<\frac{2z^{2L+1}}{(2L+1)(1-z^2)}\quad(x>1).
\]
The decomposition $\log100=6\log2+\log(25/16)$ then yields exact
rational endpoints for $\gamma$ and certifies
\[
 4.02\cdot10^{-31}<\gamma-\xi<4.03\cdot10^{-31}.
\]
These assertions were rerun successfully from the supplied
\texttt{verify\_camouflage.py}. They establish this finite example;
other LLL computations and asymptotic approximation exponents do not
follow from it.


% END INTEGRATED SOURCE

\clearpage
\part{Prime, moment and cross-domain constructions}

% BEGIN INTEGRATED SOURCE: extension_money/body.tex
% Audited Extension-II / Moneyball integration.
% Thematic ordering: binomial arithmetic; Fourier positivity; period ranks;
% Hankel extraction; analytic limits; differential obstruction; categories.

\section{Prime units, minimum certificate degree, and finite camouflage}
\label{em:pu:section}

Write $\mathcal I=\operatorname{Int}(\Z)$,
$\mathcal B_x=\Z[\binom xn:n\ge0]$, and
$U(x)=\{p\text{ prime}:1/p\in\mathcal B_x\}$.
Kalinov--Mandelshtam already characterize the algebraic binomial ring
by integral $p$-adic embeddings; their Corollary~6.3 and following
paragraph establish prime inverses and their positive density
\cite{KM}.  We do not claim these conclusions as new.  Here we give
finite certificates, prove a minimum-degree formula, and quantify
the limitations of finite prime observations.

\subsection{An explicit finite inverse certificate}

\begin{proposition}\label{em:pu:prop:certificate}
Let $f\in\Z[X]$ have degree $d\ge1$ and no root modulo $q=p^m$.
Then
\[
 W(X)=\frac{(-1)^{q-1}}p\binom{f(X)-1}{q-1}\in\mathcal I,
 \qquad \deg W=d(q-1).
\]
At every root $x$ of $f$, this gives $W(x)=1/p$, hence an explicit
finite integer-binomial expansion of $1/p$ in $\mathcal B_x$.
\end{proposition}
\begin{proof}
Put $J_q(Y)=(-1)^{q-1}\binom{Y-1}{q-1}$, so $J_q(0)=1$.
Lucas's digit identity gives, for all integers $y$,
\[
 J_q(y)\equiv\begin{cases}1,&q\mid y,\\0,&q\nmid y\end{cases}
 \pmod p.
\]
For negative integers the same conclusion follows by periodicity:
the coefficient of $T^{q-1}$ in $(1+T)^{y-1}$ modulo $p$ is
unchanged by replacing $y$ by $y+q$, because
$(1+T)^q\equiv1+T^q$.  The sign in $J_q$ is $1$ modulo $p$.
Thus $J_q(f(a))$ is divisible by $p$ for every integer $a$,
proving integer-valuedness of $W$.  Its value at a root is $1/p$.
Newton expansion supplies integer coefficients
$c_n=\Delta^nW(0)$ with $W(X)=\sum c_n\binom Xn$.
\end{proof}

For an algebraic $x$ with primitive minimal polynomial $f$, it follows
that
\begin{equation}\label{em:pu:eq:criterion}
 p\in U(x)\quad\Longleftrightarrow\quad
 f\text{ has no root in }\Z_p.
\end{equation}
Indeed, a root $\beta\in\Z_p$ gives a well-defined homomorphism
$\mathcal B_x\to\Z_p$, $P(x)\mapsto P(\beta)$, excluding $1/p$.
If no root exists, compactness of $\Z_p$ implies that some finite
depth $m$ has no congruence root, so the proposition applies.
There is an effective safe depth: choose $A,B\in\Z[X]$ and
$0\ne R\in\Z$ with $Af+Bf'=R$, and put $s=v_p(R)$.
A root modulo $p^{2s+1}$ would have $t=v_p(f'(a))\le s$;
strong Hensel lifting, since $v_p(f(a))>2t$, would give a root
in $\Z_p$.  Thus $m\le2s+1$ suffices in the rootless case.
One may use $R=\operatorname{Res}(f,f')$.

The certificate also has a crude explicit coefficient budget.  If
$K=d(q-1)$ and $T=\max_{0\le j\le K}|f(j)|$, then
\[
 \max_n|c_n|\le\frac{2^K}{p}(T+q)^{q-1},\qquad
 T\le(d+1)H(f)\max(1,K)^d.
\]
These follow by bounding $J_q(f(j))$ and using
$|\Delta^nW(0)|\le2^n\max_{j\le n}|W(j)|$.
The degree in Proposition~\ref{em:pu:prop:certificate} is an upper
bound, not generally an optimum.

\subsection{The exact minimum degree at large rootless primes}

For $p\in U(x)$ define
\[
 \kappa_p(x)=\min\{\deg P:P\in\mathcal I,\ P(x)=1/p\}.
\]
Equivalently, this is the least maximal binomial index of an inverse
certificate.

\begin{theorem}\label{em:pu:thm:minimum}
Let $x$ have primitive minimal polynomial $f$ of degree $d$.
Suppose $p>d$, $p$ does not divide the leading coefficient of $f$,
and $f$ has no root modulo $p$.  Let $A\in\mathbb F_p[X]$ be
the unique representative of degree less than $p$ of
$\bar f^{-1}$ modulo $X^p-X$, and set $t=\deg A$.  Then
\[
 \kappa_p(x)=d+t,\qquad p\le\kappa_p(x)\le p+d-1.
\]
If $d=2$, this specializes to $\kappa_p(x)=p+1$.
\end{theorem}
\begin{proof}
The no-root assumption makes $\bar f$ invertible modulo $X^p-X$.
Lift $A$ to $\widetilde A\in\Z[X]$ of the same degree.  Then
\[
 P(X)=\frac{1-f(X)\widetilde A(X)}p
\]
is integer-valued and has value $1/p$ at $x$, proving the upper
bound $\kappa_p(x)\le d+t$.

Suppose a certificate of smaller degree $K<d+t<2p$ existed.
Because $v_p(n!)\le1$ for $n\le K$, its integer binomial expansion
implies $pP\in\Z_{(p)}[X]$.  The relation at $x$ implies
\[
 Q(X)=\frac{pP(X)-1}{f(X)}\in\Q[X].
\]
It has coefficients in $\Z_{(p)}$: the polynomial $f$ has
$p$-adic Gauss norm $1$ by primitivity, and multiplicativity
of that norm applied to $fQ=pP-1$ bounds the norm of $Q$ by $1$.
Reduction modulo $p$ and evaluation at every $a\in\mathbb F_p$
give $\bar f(a)\bar Q(a)=-1$, since $P(a)\in\Z$.
Hence $\bar Q\equiv-A\pmod{X^p-X}$.  But
$\deg Q\le K-d<t<p$, a contradiction.  This proves minimality.
Finally, the nonzero polynomial $\bar f A-1$ is divisible by
$X^p-X$, so $d+t\ge p$, while $t\le p-1$.

For $d=2$, complete the square over $\mathbb F_p$:
after translation $\bar f$ is a nonzero constant times $Y^2-D$
with $D$ a nonsquare.  Its inverse function on $\mathbb F_p$
is even, and its unique representative of degree less than $p$
is therefore even.  Thus $t$ is even.  Since
$p-2\le t\le p-1$ and $p$ is odd, $t=p-1$.
\end{proof}

For example, at a root $x$ of $X^2-2$ the identity
\[
 \frac13=1-4\binom x2-12\binom x3-8\binom x4
\]
has the smallest possible maximal binomial index, namely $4$.
For any fixed quadratic algebraic parameter, the required index
at larger rootless odd primes grows exactly as $p+1$.
If the coefficients of $\widetilde A$ are chosen in
$\{0,\ldots,p-1\}$ and $K=d+t$, the same finite-difference
argument gives
\[
 \max_n|\Delta^nP(0)|\le\frac{2^K}{p}
 \left[1+(d+1)(t+1)(p-1)H(f)\max(1,K)^K\right].
\]
For fixed $f$ its logarithm is $O_f(p\log p)$.

\subsection{Frobenius density and the global alternatives}

Let $x$ be algebraic of degree $d>1$, and let $G$ be the Galois
group of the splitting field of $f$, acting transitively on its
roots.  Outside finitely many primes, an integral $p$-adic root
exists exactly when Frobenius fixes a root.  Chebotarev and
\eqref{em:pu:eq:criterion} give the natural density
\[
 \operatorname{dens}U(x)=
 \delta(G):=\frac{\#\{g\in G:g\text{ fixes no root}\}}{|G|},
 \qquad \frac1d\le\delta(G)\le1-\frac1d.
\]
The Frobenius correspondence and density theorem are standard
\cite{Milne}; their application to prime units is already in
\cite{KM}.  For completeness, the bounds follow from Burnside's
counting lemma.  If $F$ is the number of fixed points of a uniform
element of $G$, then $\mathbb EF=1$ and
$\mathbb E[F(F-1)]=r\ge1$, the number of orbits on ordered
distinct pairs.  Moreover
\[
 \mathbb E[(F-1)1_{F\ge1}]=\delta(G),\qquad
 F(F-1)\le d(F-1)\quad(F\ge1).
\]
Thus $r\le d\delta(G)$; and
$1=\mathbb EF\le d\Pr(F\ge1)$ gives the other bound.

For a reduced rational $x=a/b$, $b>0$, one has
$\mathcal B_x=\Z[1/b]$.  Indeed Bezout gives $1/b\in\Z[x]$,
while every binomial value at $a/b$ is $p$-integral when $p\nmid b$,
so its denominator uses only primes dividing $b$.
For transcendental $x$, the ring is isomorphic to $\mathcal I$,
whose evaluation at $0$ excludes every prime inverse.  Consequently
the complete classification is
\[
 \begin{array}{c|l}
 U(x)=\varnothing&x\in\Z\ \text{or}\ x\notin\Qbar,\\
 0<\#U(x)<\infty&x\in\Q\setminus\Z,\\
 \#U(x)=\infty&x\in\Qbar\setminus\Q.
 \end{array}
\]
In the last row the density is positive.  Since $0<\gamma<1$,
\[
 \gamma\notin\Qbar\ \Longleftrightarrow\ U(\gamma)=\varnothing,
 \qquad
 \gamma\in\Qbar\ \Longleftrightarrow\ U(\gamma)\ne\varnothing.
\]
An exact finite inverse certificate at $\gamma$ would therefore
prove algebraicity.  Excluding all such certificates would prove
transcendence.  Neither conclusion has been obtained.  Density
zero alone would not settle the issue, since the rational case
also has density zero.

\subsection{Every finite prime pattern has a quadratic realization}

\begin{proposition}\label{em:pu:prop:camouflage}
Let $S$ be a finite set of primes, $V\subseteq S$, and $J$ a
nonempty open real interval.  There is a quadratic real number
$x\in J$ with $U(x)\cap S=V$.
\end{proposition}
\begin{proof}
For each odd $p\in S$, prescribe a nonzero square residue for $D$
modulo $p$ when $p\notin V$, and a nonsquare residue when
$p\in V$.  At $2$, if present, prescribe $D\equiv1\pmod8$
or $D\equiv5\pmod8$, respectively.  Choose an auxiliary prime
$\ell\notin S$ and additionally require $D\equiv\ell\pmod{\ell^2}$.
The Chinese remainder theorem gives a positive such $D$.
Its valuation at $\ell$ is one, so it is not a rational square.

Take a bounded subinterval $(u,v)\subseteq J$ and a sufficiently
large positive integer $L$ coprime to every prime in $S$.
The interval $(Lu-\sqrt D,Lv-\sqrt D)$ has length exceeding one,
so contains an integer $a$.  Then
$x=(a+\sqrt D)/L\in J$ is quadratic.
At every $p\in S$, both $L$ and $D$ are units, so an integral
$p$-adic root of its minimal polynomial exists precisely when
$D$ is a square in $\Q_p$.  The prescribed congruences give this
condition exactly at $p\notin V$.  Apply \eqref{em:pu:eq:criterion}.
\end{proof}

Thus even an arbitrarily long finite list of nonunit primes is
compatible with a quadratic parameter in an arbitrarily narrow
interval about $\gamma$.  Unlike camouflage of full residue
counts, this binary statistic needs no increase in degree.
This does not exclude an all-primes argument or an argument with
independent height control.

The supplied exact-arithmetic program was rerun for the present integration.
All eleven checks passed: three finite-depth certificates, five
minimum-degree certificates, and three quadratic realizations of prescribed
patterns on $\{2,3,5,7\}$. It checks integer Newton coefficients,
exact polynomial divisibility, local congruences, and rational real
enclosures. The minimality assertions follow from
Theorem~\ref{em:pu:thm:minimum}, not from a bounded numerical search.
A second supplied program also passed nine certificate checks, ten
common-denominator snapshots, two hundred rational-denominator formula
checks, and the exact Fourier rejection gap $33/200$.

\section{Binomial denominators and a Galois-density invariant}
\label{em:sec:binomial-denominators}

For $x\in\Qbar$, put $K=\Q(x)$, let $L/\Q$ be its normal closure,
and define
\[
 D_N(x)=\min\left\{m\in\Z_{>0}:
       m\binom{x}{n}\in\mathcal O_K\quad(0\le n\le N)\right\}.
\]
This is a common denominator into the ring of integers, independent of
the choice of an integral basis.  It is different from the denominator
into $\Z[x]$ introduced below.  The results here are deductions from
local valuations and Chebotarev; no literature-wide priority is claimed.
Related binomial-value valuation descriptions occur in
\cite[Sections 5--6]{KM}; split-prime integrality and exceptional-prime
bounds occur explicitly in \cite[Lemma 4.4]{FP}.

\begin{theorem}[Exact good-prime denominators]
\label{em:thm:binomial-local}
For every $x\in\Qbar$ there is a finite set $B_x$ of rational primes
such that, for $p\notin B_x$ and all $N\ge0$,
\[
 v_p(D_N(x))=
 \begin{cases}
 0,&p\text{ splits completely in }L,\\
 v_p(N!),&\text{otherwise}.
 \end{cases}
\]
Moreover,
\[
 \sum_{p\in B_x}v_p(D_N(x))\log p=O_x(N).
\]
\end{theorem}
\begin{proof}
Choose $m\ge1$ with $\theta=mx\in\mathcal O_K$.  Include in $B_x$
all primes dividing $m$, the index $[\mathcal O_K:\Z[\theta]]$, or
the discriminant of $L$.  For $p\notin B_x$, every prime $\mathfrak p$
of $K$ above $p$ is unramified, $x$ is integral there, and its reduction
generates $\mathcal O_K/\mathfrak p$ over $\mathbb F_p$.

If some residue degree $f_{\mathfrak p}>1$, the reduction of $x$ is
not in $\mathbb F_p$.  Thus every $x-j$, $j\in\Z$, is a
$\mathfrak p$-unit, and, normalizing $v_{\mathfrak p}(p)=1$,
\[
 v_{\mathfrak p}\left(\binom{x}{n}\right)=-v_p(n!).
\]
This forces $v_p(D_N)\ge v_p(N!)$.  At every prime above $p$ the
numerator of the binomial coefficient is integral, so $v_p(N!)$ also
suffices to clear all the local denominators.  Equality follows.

If all residue degrees are one, every completion $K_{\mathfrak p}$ is
$\Q_p$ and $x\in\Z_p$.  For fixed $n$, the polynomial
$\binom{X}{n}$ maps $\Z_p$ into $\Z_p$: approximate its argument
by nonnegative ordinary integers and use continuity.  Thus
$v_p(D_N)=0$ in this case.

The residue degrees are the cycle lengths of Frobenius acting on the
conjugates of $x$.  This action of $\operatorname{Gal}(L/\Q)$ is
faithful, since those conjugates generate $L$.  All degrees are one
exactly when Frobenius is the identity, equivalently when $p$ splits
completely in $L$.  This proves the good-prime formula, including
mixed splitting types.

For the exceptional primes use the normalization
$v_{\mathfrak p}(p)=e_{\mathfrak p}$.  If
$v_{\mathfrak p}(x)\ge0$, then
$v_{\mathfrak p}\binom{x}{n}\ge-e_{\mathfrak p}v_p(n!)$.
If $v_{\mathfrak p}(x)<0$, ultrametricity gives the exact equality
\[
 v_{\mathfrak p}\left(\binom{x}{n}\right)
   =n v_{\mathfrak p}(x)-e_{\mathfrak p}v_p(n!).
\]
Consequently, with
$C_p=\max_{\mathfrak p\mid p}\max(0,-v_{\mathfrak p}(x)/e_{\mathfrak p})$,
\[
 0\le v_p(D_N)\le\lceil C_pN\rceil+v_p(N!)=O_{x,p}(N).
\]
Summing over the fixed finite set proves the last assertion.  In
particular, nonintegral algebraic arguments are covered.
\end{proof}

\begin{theorem}[Normal-closure asymptotic]
\label{em:thm:binomial-density}
For $g=[L:\Q]$ one has
\[
 \log D_N(x)=\left(1-\frac1g\right)N\log N+O_x(N),
 \qquad
 \lim_{N\to\infty}\frac{\log D_N(x)}{N\log N}=1-\frac1g.
\]
\end{theorem}
\begin{proof}
Let $S$ be the unramified completely split primes of $L$.  Their
Chebotarev density is $1/g$.  The fixed-field effective Chebotarev
estimate \cite[Theorem 1.1]{TZ}, followed by partial summation, gives
\[
 \sum_{\substack{p\le X\\p\in S}}\frac{\log p}{p}
      =\frac1g\log X+O_L(1).
\]
Its decaying error is integrable in this partial summation; a possible
fixed exceptional-zero contribution has exponent $\beta<1$ and is
integrable as well.  For $L=\Q$ use the usual prime estimate.
Legendre's formula now yields
\begin{align*}
 \sum_{\substack{p\le N\\p\in S}}v_p(N!)\log p
 &=N\sum_{\substack{p\le N\\p\in S}}\frac{\log p}{p}+O(N)\\
 &=\frac1gN\log N+O_L(N).
\end{align*}
Indeed the errors from $\lfloor N/p\rfloor-N/p$ total
$O(\sum_{p\le N}\log p)=O(N)$, and the higher prime powers contribute
at most $N\sum_p\log p/[p(p-1)]=O(N)$.  By
Theorem~\ref{em:thm:binomial-local},
\[
 \log D_N(x)=\log(N!)
       -\sum_{\substack{p\le N\\p\in S}}v_p(N!)\log p+O_x(N).
\]
Apply Stirling's formula.  Ordinary asymptotic Chebotarev, without an
effective error term, already proves the stated limit with an
$o_x(N\log N)$ remainder.
\end{proof}

\begin{remark}[Mixed splitting and lack of uniformity]
The constant is not the density of derangements.  For a non-Galois
cubic field with normal closure group $S_3$, the splitting types
$(1,1,1)$, $(1,2)$, and $(3)$ have densities $1/6$, $1/2$, and $1/3$.
Both latter types force $v_p(D_N)=v_p(N!)$, so the constant is $5/6$,
not $1/3$.  A single nonlinear completion suffices.  All error constants
depend on the fixed argument and field; no uniform effective bound in
an unknown degree or height of $\gamma$ is supplied here.
\end{remark}

\begin{proposition}[Exact rational denominator]
\label{em:prop:binomial-rational}
If $x=a/b\in\Q$, $(a,b)=1$ and $b\ge1$, then
\[
 D_N(a/b)=b^N\prod_{p\mid b}p^{v_p(N!)},\qquad
 \lim_{N\to\infty}\frac{\log D_N(a/b)}N
   =\log b+\sum_{p\mid b}\frac{\log p}{p-1}.
\]
These assertions hold for positive, zero, and negative $a$.
\end{proposition}
\begin{proof}
For $p\nmid b$, the argument belongs to $\Z_p$, so all binomial
values are $p$-integral.  If $p\mid b$, every numerator factor
$a-jb$ is a $p$-unit, whence
\[
 v_p\left(\binom{a/b}{n}\right)=-n v_p(b)-v_p(n!).
\]
The largest denominator exponent occurs at $n=N$.  This proves the
formula, including the empty product for $b=1$.  The limit follows
from $v_p(N!)/N\to1/(p-1)$.
\end{proof}

\begin{corollary}[The classical order-zero dichotomy]
For any $x\in\C$, the germ $(1+z)^x$ normalized to have value $1$ at
$0$ is a $G$-function if and only if $x\in\Q$.
\end{corollary}
\begin{proof}
A $G$-function must have algebraic coefficients, so its coefficient
of $z$ forces $x\in\Qbar$.  For algebraic $x$, the first-order
equation $(1+z)F'=xF$ gives holonomicity, and Cauchy's estimates on
a small disk give exponential bounds for all conjugate Taylor
coefficients.  The denominator condition holds for rational $x$ by
Proposition~\ref{em:prop:binomial-rational}, and fails for algebraic
irrational $x$ by Theorem~\ref{em:thm:binomial-density}, since $g\ge2$.
\end{proof}
Here $G$ means arithmetic Gevrey of order zero; ordinary analytic
Gevrey regularity is a different condition.  The classical dichotomy
is not claimed as a new classification theorem.

\subsection{A denominator defined in both arithmetic worlds}

For arbitrary complex $x$, define instead
\[
 J_N(x)=\min\left\{m\in\Z_{>0}:
        m\binom{x}{n}\in\Z[x]\quad(0\le n\le N)\right\}.
\]
This integer exists and divides $N!$.  Unlike $D_N$, it is defined
even when $x$ is transcendental.

\begin{theorem}[Quantized alternatives]
\label{em:thm:binomial-universal}
If $x$ is transcendental, $J_N(x)=N!$ exactly.  If $x$ is rational,
$J_N(x)=1$.  For algebraic $x$ with normal closure $L$,
\[
 \log J_N(x)=\left(1-\frac1{[L:\Q]}\right)N\log N+O_x(N).
\]
In particular the following limit exists for every complex $x$:
\[
 \lambda(x):=\lim_{N\to\infty}\frac{\log J_N(x)}{N\log N}
 =\begin{cases}
 0,&x\in\Q,\\
 1-1/[L:\Q],&x\in\Qbar\setminus\Q,\\
 1,&x\notin\Qbar.
 \end{cases}
\]
\end{theorem}
\begin{proof}
If $x$ is transcendental and $m\binom{x}{N}=P(x)$ with
$P\in\Z[X]$, injectivity of evaluation gives
$m\binom{X}{N}=P(X)$.  Its leading coefficient $m/N!$ must be an
integer.  Thus $N!\mid m$, while $N!$ always clears every coefficient
under consideration.  This proves exact equality.

Now let $x$ be algebraic and use $B_x$ from
Theorem~\ref{em:thm:binomial-local}.  For every $p\notin B_x$,
\[
 \Z_{(p)}[x]=\mathcal O_K\otimes_\Z\Z_{(p)}.
\]
Indeed $m$ and the index of $\Z[mx]$ in $\mathcal O_K$ are
$p$-local units.  To justify passage from this equality to denominators,
consider the homomorphism of additive $\Z$-modules
\[
 \phi_N:\Z\longrightarrow
       \bigoplus_{n=0}^N K/\Z[x],\qquad
 u\longmapsto\left(u\binom{x}{n}+\Z[x]\right)_{n=0}^N.
\]
Its kernel is $J_N(x)\Z$.  Replacing $\Z[x]$ by $\mathcal O_K$
defines a homomorphism with kernel $D_N(x)\Z$.  Localization is exact
and commutes with this finite direct sum, so the displayed equality
of localized rings makes their localized kernels equal.  Hence
\[
 v_p(J_N(x))=v_p(D_N(x))\qquad(p\notin B_x).
\]
No finite-generation assumption on $\Z[x]$ as a $\Z$-module is
needed; in particular this covers nonintegral $x$.  At exceptional
primes $0\le v_p(J_N)\le v_p(N!)$, because $J_N\mid N!$.
Their contribution is $O_x(N)$, and the same global calculation
as before proves the asserted asymptotic.

Finally, for $x=a/b$ in lowest terms, Bezout gives
$\Z[x]=\Z[1/b]$.  Proposition~\ref{em:prop:binomial-rational} shows
that every binomial value lies in this ring, so $J_N=1$.  The limit
in the transcendental case follows from Stirling's formula.
\end{proof}

\begin{remark}[Exact reformulation, not an arithmetic oracle]
Theorem~\ref{em:thm:binomial-universal} gives an equivalence:
$\gamma$ is transcendental if and only if $\lambda(\gamma)=1$.
However, computing $J_N(\gamma)$ exactly entails membership questions
in $\Z[\gamma]$, hence unknown polynomial relations.  Treating $\gamma$
as a formal indeterminate and assigning $J_N=N!$ would assume the
injectivity that must be proved.  No independent estimate of
$J_N(\gamma)$ is obtained here.  Likewise, proving $(1+z)^\gamma$
to be a $G$-function would prove rationality, whereas proving that it
is not a $G$-function would exclude rationality only.  The established
denominator alternatives therefore produce no contradiction in either
world and settle no designated constant problem.
\end{remark}

\section{Exact denominator defects and optimal joint camouflage}
\label{em:sec:denominator-bridge}

For $x\in\mathbb C$ define
\[
J_N(x)=\min\{m\ge1:m\tbinom xn\in\mathbb Z[x]\ (0\le n\le N)\}.
\]
It exists, divides $N!$, and $J_0=J_1=1$.  Let $x$ now be an algebraic
integer, with monic minimal polynomial $f$ of degree $d$.  Write
$F_n=X(X-1)\cdots(X-n+1)$ and $R_n=F_n\bmod f$, with $\deg R_n<d$.
The integral basis of the order $\mathbb Z[x]$, without any maximal-order
assumption, gives
\begin{equation}\label{em:eq:exact-J}
J_N(x)=\mathop{\rm lcm}_{0\le n\le N}
\frac{n!}{\gcd(n!,\operatorname{coeffs}(R_n))}.
\end{equation}
The convention for a zero coefficient vector is $\gcd(n!,0)=n!$.
All quantities are computed recursively by $R_{n+1}=(X-n)R_n\bmod f$.
The basic binomial-ring description below is prior work of
\href{https://arxiv.org/abs/2010.14494}{Kalinov--Mandelshtam (2021)};
we claim no literature-wide priority for the deductions proved here.

\begin{theorem}[Exact support of the factorial defect]\label{em:thm:first-defect}
For each prime $p$, put $e_p=\infty$ if $\bar f\in\mathbb F_p[X]$
has a nonlinear irreducible factor.  Otherwise let $e_p$ be its largest
linear-root multiplicity.  For every $N\ge0$,
\[
 p\mid N!/J_N(x)\quad\Longleftrightarrow\quad p e_p\le N.
\]
Consequently
\[
\operatorname{rad}(N!/J_N(x))=\prod_{p e_p\le N}p,
\qquad
\min\{N\ge2:J_N(x)<N!\}=\min_p p e_p.
\]
In particular $J_N(x)=N!$ whenever $N<d$.
\end{theorem}
\begin{proof}
Let $k=\lfloor N/p\rfloor$ and $n_0=kp$.  For $n_0\le n\le N$,
$v_p(n!)=v_p(N!)$; earlier indices have smaller valuation.  Reduction
modulo $p$ gives
\[
 \overline{F_{kp}}=(X^p-X)^k.
\]
If $\bar f$ does not divide this polynomial, the remainder $R_{kp}$
has a coefficient which is a $p$-unit.  Equation~\eqref{em:eq:exact-J}
forces $v_p(J_N)=v_p(N!)$.  If it does divide, every $R_n$ with
$n_0\le n\le N$ is coefficientwise divisible by $p$, since $F_{n_0}$
divides $F_n$.  Earlier factorials already have smaller valuation.
Thus all the denominators in \eqref{em:eq:exact-J} have exponent at most
$v_p(N!)-1$.  Finally, $X^p-X$ is the product of all distinct monic
linear polynomials over $\mathbb F_p$, so divisibility holds exactly
when $e_p\le k$.  This includes every ramified and nonmaximal-order
prime.  The minimum is finite because there are completely split good
primes, by Chebotarev.  For $N<d$, no reduction of $F_N$ by $f$ occurs,
and its leading coefficient is $1$.
\end{proof}

\begin{center}
\begin{tabular}{c|c|c}
$f$ & First defect index & Responsible $(p,e_p)$\\\hline
$X^2-2$ & $4$ & $(2,2)$\\
$X^2-X-1$ & $10$ & $(5,2)$\\
$X^3-X-1$ & $46$ & $(23,2)$\\
$X^3-2$ & $6$ & $(2,3)$
\end{tabular}
\end{center}
These examples were checked by exact remainder arithmetic through
$N=60$.  For $X^3-X-1$, the first completely split unramified prime
is $59$; keeping the ramified prime $23$ is essential for the exact answer.

\begin{remark}[An explicit differential interpretation]\label{em:defect-curvature}
Let $C$ be the integral companion matrix of $f$.  For the connection
$\nabla_{\partial_z}=\partial_z-C/(1+z)$, reduction modulo $p$ gives
\[
 \Psi_p=\frac{C-C^p}{(1+z)^p}.
\]
Indeed, its $p$th iterate on constant vectors is
$(-1)^p C(C+I)\cdots(C+(p-1)I)/(1+z)^p$, while
$\partial_z^p=0$ on $\mathbb F_p(z)$.  Cyclicity of $C$ shows that
$\Psi_p$ is nilpotent exactly when $e_p<\infty$, with nilpotency
index $e_p$.  Thus the first $p$-defect is $p$ times that index.
This is a concrete logarithmic-connection calculation, not a solution
of the general $p$-curvature conjecture.  A companion matrix for
$\gamma$ cannot be supplied without a polynomial relation for it.
\end{remark}

\begin{proposition}[Quadratic denominator camouflage]\label{em:prop:quadratic-J}
Given $N\ge2$ and a nonempty open real interval $I$, there is a
quadratic algebraic integer $x\in I$ satisfying $J_n(x)=n!$ for all
$n\le N$.  Its monic minimal polynomial can have coefficient height
at most $C_I\prod_{p\le N}p$.
\end{proposition}
\begin{proof}
Choose a bounded $(u,v)\subset I$ and set $M=\prod_{p\le N}p$.
Seek $f=X^2-AX+B$.  At each odd $p\le N$ prescribe $A=0$ and
$-B$ a nonzero nonsquare; at $2$ prescribe $A=B=1$.  Every reduction
is irreducible.  The Chinese remainder theorem gives residue classes
for $A,B$ modulo $M$.  Choose $A$ large positive in its class so that
\[
 A(v-u)-(v^2-u^2)>M,\qquad A>2\max(|u|,|v|)+2.
\]
There is an allowed $B$ strictly between $Au-u^2$ and $Av-v^2$.
Then $f(u)>0>f(v)$, giving a real root in $I$.  Irreducibility modulo
$2$ implies irreducibility over $\mathbb Q$, and
Theorem~\ref{em:thm:first-defect} proves the entire denominator prefix.
The choices require only $A,B=O_I(M)$.
\end{proof}

Define $\mathcal B_x=\mathbb Z[\binom xn:n\ge0]$ and
$U(x)=\{p:1/p\in\mathcal B_x\}$.  The binomial-ring criterion gives,
for algebraic $x$ with primitive minimal polynomial $f$,
\begin{equation}\label{em:eq:local-unit-criterion}
p\notin U(x)\quad\Longleftrightarrow\quad
f\text{ has a root in }\mathbb Z_p.
\end{equation}
For a transcendental parameter $t$, $U(t)=\varnothing$ and $J_n(t)=n!$.
The quadratic construction above does not imitate the first statistic:
it has prime inverses at all tested primes.  The following result
handles both statistics and is sharp in its degree.

\begin{theorem}[Optimal joint camouflage]\label{em:thm:joint-camouflage}
For every $N\ge2$ and nonempty open real interval $I$, there is a
cubic algebraic integer $x\in I$ with
\[
 U(x)\cap\{p:p\le N\}=\varnothing,
 \qquad J_n(x)=n!\quad(0\le n\le N).
\]
Its monic minimal polynomial can have height at most $C_I M\ell^2$,
where $M=\prod_{p\le N}p$ and $N<\ell\le2N$ is prime.
For $N\ge4$, degree three is minimal among algebraic integers.
\end{theorem}
\begin{proof}
Seek $f=X^3-AX+B$.  At $2$ prescribe $A=0,B=1$, giving
$(X+1)(X^2+X+1)$.  At every odd $p\le N$ prescribe $A$ to be a
nonzero nonsquare and $B=0$, giving $X(X^2-A)$.  In every case
there is a simple linear factor and an irreducible quadratic factor.
The linear root lifts to $\mathbb Z_p$ by Hensel, proving the
nonunit assertion by \eqref{em:eq:local-unit-criterion}; the nonlinear
factor proves the denominator assertion by Theorem~\ref{em:thm:first-defect}.

At an auxiliary prime $\ell>N$ require $A=0\pmod\ell$ and
$B=\ell\pmod{\ell^2}$.  Eisenstein ensures irreducibility.
These compatible prescriptions give a class for $A$ modulo $M\ell$
and one for $B$ modulo $M\ell^2$.  In a bounded $(u,v)\subset I$,
choose $A$ in its class such that
\[
 A(v-u)-(v^3-u^3)>M\ell^2,
 \qquad A>3\max(u^2,v^2)+2.
\]
Choose $B$ in its class strictly between $Au-u^3$ and $Av-v^3$.
Then $f(u)>0>f(v)$ and $f$ is strictly decreasing there, so it has a
unique root in this interval.  Also $A,B=O_I(M\ell^2)$.

For minimality, a degree-one algebraic integer has $J_2=1$.
For a quadratic algebraic integer, $2\notin U(x)$ gives a
$\mathbb Z_2$ root of its monic polynomial.  The other root is
also in $\mathbb Z_2$, since their sum is an integer.  Therefore
$\bar f$ splits fully modulo $2$ with $e_2\le2$, and
Theorem~\ref{em:thm:first-defect} gives $2\mid4!/J_4(x)$.
This contradicts the required prefix whenever $N\ge4$.
\end{proof}

\begin{example}[An exact twenty-step joint witness]\label{em:joint-witness}
The polynomial
\[
 f(X)=X^3-513113715440042X+296176197922605
\]
is Eisenstein at $23$.  It has one root in
$(57721/100000,28861/50000)$, because its endpoint values are
\[
 \frac{1830233458549489854616361}{10^{15}}>0,
 \qquad
 -\frac{412612961981365018758619}{125000000000000}<0,
\]
and its derivative is negative on the interval.  At every prime
$p\le20$ its reduction has splitting type $(1,2)$.  Thus this
specified cubic algebraic integer has no tested prime inverses and
$J_n=n!$ for every $n\le20$, exactly as a transcendental parameter.
The accompanying program checks the modular factors, endpoint signs,
Eisenstein conditions, and all denominators using only integer and
rational arithmetic.
\end{example}

\begin{remark}[Consequence for the research strategy]\label{em:finite-camouflage-scope}
The joint construction has
$\log H(f)\le\sum_{p\le N}\log p+2\log(2N)+O_I(1)$.
It rules out a finite-data inference from a small real interval,
absence of all tested prime inverses, and maximal tested denominators
to transcendence: a cubic algebraic integer reproduces every such
finite observation set.  It does not exclude an all-primes argument,
independent height control, or additional information specific to
$\gamma$.  Computing \eqref{em:eq:exact-J} requires a supplied polynomial
$f$; it supplies no arithmetic oracle for $\gamma$.
\end{remark}

\section{Finite Fourier ranges and robust positivity obstructions}

Let $S$ be a finite nonempty prime set, let $\mathcal N_S$ be the positive
$S$-smooth integers, and let $\mathcal R_S$ be the positive integers
coprime to every prime in $S$. Each positive integer is uniquely $ra$,
with $r\in\mathcal R_S$ and $a\in\mathcal N_S$.
For a continuous real $1$-periodic function $p$, consider
\[
 (M_q)\qquad \sum_{j=0}^{q-1}p(x+j/q)=p(qx),\quad q\in S.
\]
We use $\widehat\mu(n)=\int e^{-2\pi inx}\,d\mu(x)$ on $\mathbb T=\mathbb R/\Z$.

\begin{proposition}[Orbit description]
The equations $(M_q)$ hold precisely when $\widehat p(0)=0$ and there
are complex numbers $\beta_r$, $r\in\mathcal R_S$, such that
\[
 2\pi i(ra)\widehat p(ra)=\beta_r,
 \qquad
 2\pi i(-ra)\widehat p(-ra)=\overline{\beta_r}
 \quad(a\in\mathcal N_S).
\]
\end{proposition}
\begin{proof}
Fourier comparison gives $q\widehat p(qn)=\widehat p(n)$ and zero
mean. Equivalently the derivative coefficients are constant on the
multiplicative orbits. Unique factorization gives the assertion.
Conversely these equations annihilate every Fourier coefficient of
each continuous multiplication defect, which is therefore zero.
\end{proof}

This is a criterion for an existing continuous function, not a claim
that arbitrary orbit data define one. A useful sufficient realization is
\[
 U_S(x)=\sum_{a\in\mathcal N_S}\frac{\sin(2\pi ax)}a,
 \qquad V_S(x)=\sum_{a\in\mathcal N_S}\frac{\cos(2\pi ax)}a,
\]
\[
 p_\beta(x)=\sum_{r\in\mathcal R_S}
 \frac{\Re\beta_r\,U_S(rx)+\Im\beta_r\,V_S(rx)}{\pi r}.
\]
If $\sum_r|\beta_r|/r<\infty$, these series converge absolutely and
uniformly and give the specified coefficients. Indeed
$\sum_{a\in\mathcal N_S}a^{-1}=\prod_{q\in S}(1-q^{-1})^{-1}$,
and direct expansion verifies the coefficients. Real $\beta_r$ give
odd functions vanishing at integers. Prime equations suffice for all
$S$-smooth multipliers; a lone composite equation need not imply each
prime equation separately.

\begin{theorem}[Finite-range rigidity]\label{em:thm:finite-range}
Let $\sigma$ be a finite complex atomless measure on $\mathbb T$.
Suppose its Fourier coefficients take only finitely many values and
$(T_m)_*\sigma=\sigma$ for one integer $m\ge2$, where $T_m(x)=mx\pmod1$.
Then $\sigma=\sigma(\mathbb T)\lambda$, with $\lambda$ normalized
Lebesgue measure.
\end{theorem}
\begin{proof}
Let $V$ be the finite Fourier range. For each $v\in V$ the polynomial
\[
 P_v(t)=\prod_{w\in V\setminus\{v\}}\frac{t-w}{v-w}
\]
defines a convolution polynomial $P_v(\sigma)$; a constant term means
the corresponding multiple of $\delta_0$. Its Fourier transform is
the indicator of $E_v=\{n:\widehat\sigma(n)=v\}$. Fourier uniqueness
makes this an idempotent measure. Cohen's idempotent theorem, as
stated in \cite[Theorem 1.1]{GreenSanders}, puts $E_v$ in the coset ring of $\Z$.

Every coset-ring set of $\Z$ is periodic outside a finite set: its
finite description uses residue classes and singleton cosets. Taking
a common period for the finitely many $E_v$ shows that
$\widehat\sigma(n)$ equals a periodic sequence $b_n$ except at
finitely many integers. If its period is $Q$, the sequence $b_n$ is
the Fourier transform of the finite atomic measure
\[
 \alpha=\sum_{j=0}^{Q-1}w_j\delta_{j/Q},
 \qquad w_j=\frac1Q\sum_{k=0}^{Q-1}b_k e^{2\pi i kj/Q}.
\]
The remaining finitely supported Fourier sequence is the transform
of a trigonometric-polynomial density $f\lambda$. Thus
$\sigma=\alpha+f\lambda$. Atomlessness forces $\alpha=0$.
Finally invariance gives $\widehat\sigma(m^jn)=\widehat\sigma(n)$;
for $n\ne0$ the left-hand indices leave the finite Fourier support,
so all nonzero coefficients vanish.
\end{proof}

This is an application of a classical structural theorem with an
additional finite-range hypothesis. It does not settle the unrestricted
$\times2,\times3$ problem. The range hypothesis can also be replaced
by $\operatorname{dist}(\widehat\sigma(n),V)\to0$ for a fixed finite
$V$: recurrence along $m^jn$ then puts each nonzero-frequency
coefficient exactly in $V$ (with the zero-frequency coefficient added
as one more value if necessary). A probability measure whose coefficients
belong to a fixed uniformly discrete subset of $\C$ also has finite
Fourier range: all its coefficients lie in the closed unit disk, and a
uniformly discrete set has only finitely many points in a bounded disk.

\begin{corollary}[All nonzero finite-orbit constructions fail BV]
If only finitely many $\beta_r$ of a continuous solution are nonzero,
then bounded variation implies $p=0$.
\end{corollary}
\begin{proof}
The measure $dp$ would be atomless, have finite Fourier range, and be
invariant under any $T_q$, $q\in S$. Its total mass is zero.
Theorem~\ref{em:thm:finite-range} gives $dp=0$; the zero-mean condition
then gives $p=0$.
\end{proof}
The same argument excludes every nonzero continuous BV solution whose
derivative coefficients have finite range, even if infinitely many orbit
labels are present.

In particular, no finite selection of extra amplitudes or phases in
the previous smooth-frequency construction repairs its lack of
bounded variation. Positivity can be rejected without Cohen's theorem.

\begin{proposition}[An elementary finite certificate]
Let $p\ne0$ be a real finite-orbit solution. The distribution
$\lambda+dp$ is not a positive probability measure. If $p(0)=0$,
equivalently $x+p(x)$ is not nondecreasing on $[0,1]$.
\end{proposition}
\begin{proof}
Choose an orbit $r_0$ with $\beta_{r_0}=a\ne0$, and choose a prime
$\ell$ dividing neither the primes in $S$ nor any supported orbit
label. Fix $q\in S$ and set $n_j=r_0q^{j(\ell-1)}$.
A proposed measure would satisfy $\widehat\mu(n_j)=a$.
Every difference $n_i-n_j$, $i\ne j$, is divisible by $\ell$, so
it belongs to no supported signed orbit and has Fourier coefficient
zero. The characters $e_j(x)=e^{-2\pi in_jx}$ would be orthonormal.
For every $K$ this would imply
\[
 0\le\int\left|1-\overline a\sum_{j=1}^K e_j(x)\right|^2d\mu(x)
 =1-K|a|^2.
\]
Choose $K>|a|^{-2}$ for a contradiction.
\end{proof}
For example, with $S=\{2,3\}$ the two-orbit function
$p(x)=U_S(x)-U_S(5x)/5$ has derivative amplitudes $\pi$ and $-\pi$.
The prime $7$ and frequencies $n_j=8^j$ make every nonzero difference
divisible by $7$, outside both supported signed orbits. Thus the negative
second amplitude cannot restore positivity.

\begin{theorem}[A robust positivity test]\label{em:thm:robust}
Let $\mu$ be a probability measure and let $n_1,\ldots,n_K$ be
distinct integers. If
\[
 |\widehat\mu(n_j)-a|\le\varepsilon,
 \qquad |\widehat\mu(n_i-n_j)|\le\eta\quad(i\ne j),
\]
where $\varepsilon,\eta\ge0$, then
\begin{equation}\label{em:eq:robust}
 K(|a|-\varepsilon)_+^2\le1+(K-1)\eta.
\end{equation}
\end{theorem}
\begin{proof}
For $U=\sum_j e^{-2\pi in_jx}$, the hypotheses give
$|\int U\,d\mu|\ge K(|a|-\varepsilon)_+$ and
$\int|U|^2d\mu\le K+K(K-1)\eta$.
Cauchy--Schwarz on the probability space proves the result.
\end{proof}
When $0\le\eta<1$ and $(|a|-\varepsilon)_+^2>\eta$, rearrangement gives
\[
 K\le\frac{1-\eta}{(|a|-\varepsilon)_+^2-\eta}.
\]
A convenient sufficient condition for rejection is
$\varepsilon\le |a|/2$, $\eta\le |a|^2/8$, and
$K>8/|a|^2$: then the difference between the two sides of the necessary
inequality is at least $K|a|^2/8-1>0$.

For example, $a=1/2$, $\varepsilon=1/20$, $\eta=1/100$, and
$K=6$ give a left side $243/200$ and a right side $21/20$,
with positive gap $33/200$. Such certified Fourier data cannot
come from any probability measure. This test applies directly to
truncated moment and spectral reconstruction problems. It is an
elementary Gram inequality, not a newly claimed general positivity
theorem.

\begin{corollary}[An arithmetic sufficient condition]
An atomless probability measure invariant under one $T_m$ whose
Fourier coefficients have uniformly bounded algebraic degree and
absolute multiplicative Weil height is Lebesgue measure.
\end{corollary}
\begin{proof}
Degree at most $D$ and height at most $B$ bound the primitive
minimal-polynomial coefficients by $2^DB^D$, using elementary
symmetric functions and the product definition of height. Only
finitely many integer polynomials and roots remain. Apply
Theorem~\ref{em:thm:finite-range}.
\end{proof}
Algebraicity alone, or membership in one number field, does not give
this height bound. Neither arithmetic alternative for $\gamma$ has
been shown to force a nonzero measure in the excluded class.

\section{A complete transfer to finite Fourier moment reconstruction}
\label{em:sec:transfer-moments}

The following is a concrete use of classical Gram and circulant matrix
methods. It closes a finite moment problem and strengthens our earlier
necessary test; it is not presented as a new general moment theorem.
Write $c_n=\int_{\mathbb T}e^{-2\pi i n x}\,d\mu(x)$ for a positive
probability measure.

\begin{theorem}[Exact parity-moment feasibility]\label{em:thm:parity-moment}
Fix $K\ge2$ and $a,b\in\mathbb R$. The data
\[
c_{2j+1}=a\ (0\le j<K),\qquad c_{2j}=b\ (1\le j<K)
\]
have a representing probability measure if and only if
\begin{equation}\label{em:eq:parity-triangle}
b\le1,\qquad K|a|\le1+(K-1)b.
\end{equation}
One representing measure on $x=j/(2K)$ has weights
\[
w_0=\frac{1+(K-1)b+Ka}{2K},\qquad
w_K=\frac{1+(K-1)b-Ka}{2K},\qquad
w_j=\frac{1-b}{2K}\quad(j\notin\{0,K\}).
\]
\end{theorem}
\begin{proof}
The moment matrix $T=(c_{r-s})_{0\le r,s<2K}$ is positive semidefinite.
Reordering the monomials by parity gives
\[
T\sim\begin{pmatrix}
(1-b)I_K+bJ_K&aJ_K\\
aJ_K&(1-b)I_K+bJ_K
\end{pmatrix}.
\]
Its eigenvalues are $1-b$ with multiplicity $2K-2$ and
$1+(K-1)b\pm Ka$. This proves necessity. The displayed weights are
nonnegative under \eqref{em:eq:parity-triangle} and sum to one.
Finite Fourier summation, comparing with equal weights
$(1-b)/(2K)$ at all grid points, gives
\[
c_n=(b+a)/2+(-1)^n(b-a)/2\qquad(0<|n|<2K).
\]
This proves sufficiency without a general extension theorem.
\end{proof}

The feasible set is the triangle with vertices $(1,1)$, $(-1,1)$,
and $(0,-1/(K-1))$. For $b=0$ the exact bound is $|a|\le1/K$.
The earlier star Gram test, using only $1$ and the odd characters,
gave $|a|\le1/\sqrt K$. Thus using the full matrix extracts a
strictly stronger conclusion from the same moment data.

\begin{proposition}[Boundary uniqueness and support size]\label{em:moment-boundary}
Every boundary point of this triangle has a unique representing
measure, given by the weights above. Every interior point has
multiple representing measures. The smallest atomic support is
$\operatorname{rank}T$: it is $2K$ in the interior, $2K-1$ on either
open sloping edge, $2K-2$ at the bottom vertex, $2$ on the open top
edge, and $1$ at either top vertex.
\end{proposition}
\begin{proof}
Put $z=e^{-2\pi i x}$,
$E(z)=\sum_{j=0}^{K-1}z^{2j}$, and $O(z)=zE(z)$.
On the sloping edge with $a\ge0$,
\[
\int|E-O|^2d\mu=2K\{1+(K-1)b-Ka\}=0.
\]
The measure is supported on the zeros of $(1-z)E(z)$: the $2K$-th
roots except $-1$. Finite Fourier inversion fixes the weights.
The other edge uses $E+O$ and excludes $+1$. At the bottom vertex
both exclusions hold. On the top edge, $c_2=1$ implies support on
$\{1,-1\}$, and $c_1$ fixes the weights.

In the interior choose
$0<\theta<\min\{1,1-b,1+(K-1)b-K|a|\}$.
The parameters $(a/(1-\theta),b/(1-\theta))$ are feasible; let
$\nu$ represent them. Both
$\theta\lambda+(1-\theta)\nu$ and
$\theta\rho_M+(1-\theta)\nu$ represent the original data, where $\lambda$ is normalized Haar measure on the circle and
$\rho_M$ is uniform on $M$-th roots and $M>2K-1$. They are distinct.
Finally an $r$-atomic representation makes $T$ a sum of $r$ rank-one
matrices, so $r\ge\operatorname{rank}T$. The explicit measure attains
the bound in every listed case.
\end{proof}

\begin{theorem}[Exact feasibility with errors]\label{em:thm:parity-noise}
Fix $K\ge2$. Let $a,b\in\mathbb R$ and $\varepsilon,\eta\ge0$. The constraints
\[
|c_{2j+1}-a|\le\varepsilon\ (0\le j<K),\qquad
|c_{2j}-b|\le\eta\ (1\le j<K)
\]
admit a probability measure if and only if
\begin{equation}\label{em:eq:parity-robust}
b-\eta\le1,\qquad
K(|a|-\varepsilon)_+\le1+(K-1)\min\{1,b+\eta\}.
\end{equation}
The unknown moments may be complex.
\end{theorem}
\begin{proof}
Put $A=(|a|-\varepsilon)_+$ and $d=\min\{1,b+\eta\}$.
Necessarily $b-\eta\le1$. With $E,O$ as above,
\[
\int|E|^2d\mu
=K+2\sum_{j=1}^{K-1}(K-j)\Re c_{2j}
\le K+K(K-1)d.
\]
When $A>0$, the $K^2$ terms of $\int E\overline O\,d\mu$ all
have real parts of the sign of $a$ and magnitude at least $A$.
Therefore its modulus is at least $K^2A$. Cauchy--Schwarz and
$|E|=|O|$ bound this modulus by $\int|E|^2d\mu$.
If $A=0$, positivity of the latter integral suffices instead.
This proves the second necessary condition.

Conversely choose $a'=\operatorname{sgn}(a)A$, with
$\operatorname{sgn}(0)=0$, and $b'=d$.
Then $|a'-a|\le\varepsilon$ and $|b'-b|\le\eta$; the latter uses
$b-\eta\le1$. Theorem~\ref{em:thm:parity-moment} applies and explicitly
constructs a measure satisfying all the error bounds.
\end{proof}

For $a=1/2$, $b=0$, $\varepsilon=1/20$, and $\eta=1/100$,
$K=3$ already fails \eqref{em:eq:parity-robust} by $33/100$.
The previous star test does not reject until $K=6$ for this pattern.
At $K=2$ feasibility holds, with weights
$(191,99,11,99)/400$ on the fourth roots of unity.
This is an exact change in the number of moments needed to reject
the data, not an asymptotic or numerical optimization claim.
No consequence of either arithmetic alternative for $\gamma$ has
been shown to force these forbidden data.

\section{Several poles: cancellation and logarithmic rank}
\label{em:sec:multipole}

Let $\Lambda$ be a finite set of distinct positive algebraic numbers and put
\[
 I_\lambda=\int_0^\infty\frac{e^{-t}}{t+\lambda}\,dt
          =e^\lambda E_1(\lambda),\qquad I_1=\delta.
\]
For $R\in\Qbar(z)$ with finite poles only at $-\lambda$, $\lambda\in\Lambda$,
write its principal part there as
$\sum_{j=1}^{m_\lambda}r_{\lambda,j}(z+\lambda)^{-j}$ and set
\[
 c_\lambda=\sum_{j=1}^{m_\lambda}
             \frac{(-1)^{j-1}r_{\lambda,j}}{(j-1)!}.
\]
The rational-Borel reduction proved in \cite{Previous} gives a unique
rational function $T$, with no other finite poles, such that
\begin{equation}\label{em:mul:reduction}
 R=T-T'+\sum_{\lambda\in\Lambda}\frac{c_\lambda}{z+\lambda},
 \qquad
 \mathcal L(R):=\int_0^\infty e^{-t}R(t)\,dt
       =T(0)+\sum_{\lambda\in\Lambda}c_\lambda I_\lambda.
\end{equation}
All residues $c_\lambda$ are effectively computable algebraic numbers.
The integral identity follows by integration by parts; the exponential
kills the rational boundary term at infinity.

\subsection{The canonical gamma coefficient}
The connection formula gives
\begin{equation}\label{em:mul:connection}
 \mathcal L(R)=E_R-A_R\gamma-B_R,
\end{equation}
where
\[
 E_R=T(0)+\sum_\lambda c_\lambda e^\lambda\Ein(\lambda),\qquad
 A_R=\sum_\lambda c_\lambda e^\lambda,\qquad
 B_R=\sum_\lambda c_\lambda e^\lambda\log\lambda.
\]
Here $E_R$ belongs to the ring $\mathbf E$ of $E$-values.

\begin{theorem}[Exact canonical cancellation criterion]\label{em:mul:cancellation}
The coefficient $A_R$ vanishes if and only if every $c_\lambda$ vanishes,
equivalently if and only if $R=T-T'$ for a rational $T$.  In this case
$\mathcal L(R)=T(0)\in\Qbar$.
Equivalently, if $R(z)=\sum_{n\ge0}r_nz^n$ at zero and
$\mathfrak f_R(x)=\sum_{n\ge0}n!r_nx^n$, then there exists
$a\in\Qbar$ for which
$(\mathfrak f_R(x)-a)/(x-1)$ is arithmetic Gevrey of order one.
The only possible $a$ is $T(0)$.
\end{theorem}
\begin{proof}
The numbers $e^\lambda$ for distinct algebraic $\lambda$ are linearly
independent over $\Qbar$ by Lindemann--Weierstrass.  Thus
$\sum c_\lambda e^\lambda=0$ forces every algebraic coefficient to be zero.
Use the unique reduction \eqref{em:mul:reduction} for rational exactness
and the value formula. The arithmetic-division equivalence, including
uniqueness of $a=T(0)$, is the $\alpha=1$ case of the full rational-Borel
classification in Theorem~\ref{ra:bor:class}.
\end{proof}
This is a criterion for the coefficient in the particular identity
\eqref{em:mul:connection}, not uniqueness of every possible representation
of its value.  It does not prove that every nonexact kernel has a
transcendental integral: that statement includes the unknown arithmetic
nature of $\delta$.  The analytic-to-arithmetic implication remains the
zero-removal problem described in \cite{FRmixed}.

\begin{proposition}[Collision locations under $\mathrm A$]\label{em:mul:membership}
Assume $\gamma\in\Qbar$.  For every positive algebraic $\lambda$,
\[
 I_\lambda\in\mathbf E\quad\Longleftrightarrow\quad
 \log\lambda\in\mathbf E.
\]
Consequently the set of such locations is a multiplicative subgroup of
$\Qbar_{>0}$ closed under positive rational powers.  For fixed algebraic
$b>0$, $b\ne1$, either all $I_{b^j}$, $j\in\Z\setminus\{0\}$, belong
to $\mathbf E$, or none do, according as $\log b$ does or does not belong
to $\mathbf E$.
\end{proposition}
\begin{proof}
Use $e^{-\lambda}I_\lambda=\Ein(\lambda)-\gamma-\log\lambda$ and the
facts that $e^{\pm\lambda}$ and algebraic numbers belong to $\mathbf E$.
The subgroup assertion follows from the real logarithm laws and closure
of $\mathbf E$ under rational linear combinations.  Finally
$\log(b^j)=j\log b$, with $j\ne0$.
\end{proof}
When $1\in\Lambda$, there is also the unconditional ring identity
\[
 \mathbf E[I_\lambda:\lambda\in\Lambda]
 =\mathbf E[\gamma,\log\lambda:\lambda\in\Lambda].
\]
Indeed $\gamma=\Ein(1)-e^{-1}I_1$, and then
$\log\lambda=\Ein(\lambda)-\gamma-e^{-\lambda}I_\lambda$; the reverse
inclusion follows from the connection formulas. This restates the
first-level collision-ring mechanism in the multipole notation.
Thus $\mathrm A$ does not force every new pole to supply an $E/D$
collision.  Nor may Baker's logarithm estimates be applied to $B_R$ as
if its coefficients $c_\lambda e^\lambda$ were algebraic.  For geometric
locations $\lambda=b^j$, one can at least write
\[
 B_R=(\log b)\sum_j jc_j e^{b^j};
\]
Lindemann--Weierstrass shows that this vanishes exactly when $c_j=0$
for all $j\ne0$.  This does not settle membership of $B_R$ in $\mathbf E$.

\subsection{Independent residuals after removing common logarithms}
Define the field
\[
 K_{\exp}=\Qbar(e^\alpha:\alpha\in\Qbar).
\]
This field is not being identified with the ring $\mathbf E$.
We use the previously proved \emph{multipoint $\Ein$ input}: for distinct
nonzero algebraic $\lambda_i$, the values $\Ein(\lambda_i)$ are
algebraically independent over $K_{\exp}$; see the foundations reproduced
with \cite{Previous}.  This is an explicit proof dependency.  Its route
is to place finitely many exponentials in a Laurent differential field,
exclude rational primitives of the relevant distinct characters divided
by $z$, apply primitive differential-field independence, and then apply
Siegel--Shidlovskii/Beukers specialization at the nonsingular point $1$
\cite{Beukers}.  Enlarging the finite exponential field covers every
hypothetical relation over $K_{\exp}$.  We do not reproclaim that foundation
as a new result here.

\begin{theorem}[Logarithmic rank bound]\label{em:mul:rank}
Let $\lambda_1,\ldots,\lambda_m$ be distinct positive algebraic numbers.
Choose a basis $\rho_1,\ldots,\rho_r$ for the multiplicative group they
generate and write
\[
 \lambda_i=\prod_{h=1}^r\rho_h^{n_{ih}},\qquad
 N=(n_{ih}),\qquad s=\operatorname{rank}_{\Q}[\boldsymbol1\mid N].
\]
Then $\operatorname{rank}_{\Q}N=r$, $s\le r+1$, and
\begin{equation}\label{em:mul:rank-unconditional}
 \operatorname{trdeg}_{K_{\exp}}
 K_{\exp}(I_{\lambda_1},\ldots,I_{\lambda_m})
 \ge m-s\ge m-r-1.
\end{equation}
Under $\mathrm A$ the lower bound improves to $m-r$.
\end{theorem}
\begin{proof}
The finitely generated subgroup of $\Qbar_{>0}$ is torsion-free, so it
has the indicated basis.  Put $y_i=e^{-\lambda_i}I_{\lambda_i}$,
$U_i=\Ein(\lambda_i)$ and $\ell_h=\log\rho_h$.  The vector connection
identity is
\[
 y=U-\gamma\boldsymbol1-N\ell.
\]
Choose a rational matrix $B$ of row rank $m-s$ whose rows annihilate
$[\boldsymbol1\mid N]$.  Then $By=BU$.  A full-row-rank linear image
of an algebraically independent tuple is algebraically independent:
extend its rows to an invertible matrix and apply the corresponding
linear change of indeterminates.  Thus the field generated by $y$
contains $m-s$ algebraically independent elements over $K_{\exp}$.
Individual multiplication by $e^{\lambda_i}\in K_{\exp}^{\times}$ does
not change the field generated.  This proves
\eqref{em:mul:rank-unconditional}.

Under $\mathrm A$, instead take $B$ to have row rank $m-r$ and annihilate
only $N$.  Now $By=BU-\gamma B\boldsymbol1$; the translation belongs to
$\Qbar^{m-r}$ and preserves algebraic independence.
\end{proof}

\begin{corollary}[Geometric dilations]\label{em:mul:geometric}
Fix algebraic $b>0$, $b\ne1$.  Among the sequence
$(I_{b^j})_{j\ge0}$, removal of at most two terms leaves an algebraically
independent family over $K_{\exp}$.  Under $\mathrm A$, removal of at most
one term suffices.  Moreover the explicit values
\begin{align*}
 D_j&=E_1(b^j)-2E_1(b^{j+1})+E_1(b^{j+2})\\
    &=\Ein(b^j)-2\Ein(b^{j+1})+\Ein(b^{j+2}),\qquad j\ge0,
\end{align*}
are all strictly positive and algebraically independent over $K_{\exp}$.
\end{corollary}
\begin{proof}
For every prefix of length $m\ge2$, the logarithmic exponent column is
$(0,1,\ldots,m-1)^t$, so $r=1$ and $s=2$.  The rank bounds are $m-2$
and, under $\mathrm A$, $m-1$.  Process the sequence in increasing order,
retaining a term exactly when it is transcendental over the field
generated by retained terms.  Every skipped term is algebraic over that
field, so the number retained in each prefix is its transcendence degree.
The bounds permit at most two skips in total, or at most one under
$\mathrm A$.  The retained family is algebraically independent.

The common $\gamma$ term and the affine term $j\log b$ cancel in each
$D_j$.  For every finite prefix the second-difference matrix has full
row rank, so the multipoint input proves independence of the $D_j$.
Finally $h(t)=E_1(e^t)$ satisfies
$h''(t)=e^t e^{-e^t}>0$.  Strict convexity with the nonzero step $\log b$
proves $D_j>0$.
\end{proof}
The exceptional indices are not identified; the corollary therefore does
not decide the arithmetic nature of $I_1=\delta$ or any other designated
member.  The $D_j$ are $E$-values.  Expressing them in terms of $I_{b^j}$
uses coefficients $e^{-b^j}$, which are not algebraic, so this does not
put them in $\mathbf D$ or contradict Theorem~\ref{em:mul:cancellation}.

Positivity by itself also cannot provide the missing separation.  For
positive algebraic weights $a_\lambda$, the strictly positive kernel
\[
 R(z)=\sum_\lambda a_\lambda
       \left(\frac1{z+\lambda}+\frac1{(z+\lambda)^2}\right)
\]
has algebraic integral $\sum_\lambda a_\lambda/\lambda$, since
$R=T-T'$ with $T=\sum_\lambda a_\lambda/(z+\lambda)$.
The closed statements in this section concern canonical coefficients
and family ranks; no implication here excludes either gamma world.

\section{A prime decomposition and density-one independence}
\label{em:sec:bootstrap}

Write
\[
 K_{\exp}=\overline{\mathbb Q}(e^\alpha:
       \alpha\in\overline{\mathbb Q}),\qquad
 U_n=\operatorname{Ein}(n),\quad y_n=E_1(n),\quad I_n=e^ny_n.
\]
Our explicit input is the order-one multipoint theorem proved in the
foundational part of this volume: the $U_n$ are algebraically independent over $K_{\exp}$ in
the finite-subset sense. Its differential-field proof and
Siegel--Shidlovskii/Beukers specialization step are required dependencies
of the following arithmetic deductions. No analogous specialization for $E_1$
or for $\gamma$ is assumed. The connection identity is
\begin{equation}\label{em:boot:connection}
 y_n=U_n-\gamma-\log n.
\end{equation}
The deductions below make no claim of priority.

\subsection{Explicit positive residuals for every composite integer}
For composite $n$, set $\Omega(n)=\sum_pv_p(n)$ and
\begin{equation}\label{em:boot:residual}
 R_n=y_n-\sum_{p\mid n}v_p(n)y_p+(\Omega(n)-1)y_1.
\end{equation}

\begin{theorem}\label{em:boot:composites}
Every $R_n$ is strictly positive. The family of all $R_n$, indexed by
composite positive integers, is algebraically independent over
$K_{\exp}$. Each $R_n$ is an $E$-value.
\end{theorem}
\begin{proof}
The identity $\log n=\sum_pv_p(n)\log p$ and
\eqref{em:boot:connection} give
\[
 R_n=U_n-\sum_{p\mid n}v_p(n)U_p+(\Omega(n)-1)U_1.
\]
In any finite collection, the composite coordinate $U_n$ has coefficient
one in its own row and zero in every other row. All remaining columns
are prime coordinates and $U_1$. Thus the rows have full rank, and a
full-row-rank linear image of an algebraically independent tuple is
algebraically independent. The expression is also a rational linear
combination of $E$-values.

For positivity put $h(s)=E_1(e^s)$. Direct differentiation gives
\[
 h'(s)=-e^{-e^s},\qquad h''(s)=e^se^{-e^s}>0.
\]
Consequently, for $a,b>0$,
$h(a+b)-h(a)>h(b)-h(0)$. Iteration gives
\[
 h(a_1+\cdots+a_k)-\sum_{j=1}^k h(a_j)+(k-1)h(0)>0
 \quad(k\ge2, a_j>0).
\]
Take $a_j$ to be the logarithms of the prime factors of $n$, counted
with multiplicity. The left side is $R_n$.
\end{proof}

For example,
\[
 R_6=E_1(6)-E_1(2)-E_1(3)+E_1(1),\qquad
 R_{12}=E_1(12)-2E_1(2)-E_1(3)+2E_1(1).
\]
Here $\pi(N)$ denotes the number of primes at most $N$.
This construction combines different geometric bases with a single
triangular independence proof. Indeed the rows
\[
 c_n=(1,(v_p(n))_{p\le N}),\qquad 1\le n\le N,
\]
have rank $\pi(N)+1$, already witnessed by $n=1$ and the primes.
The residual rows for the composite $n\le N$ form a basis of the left
kernel. Hence they describe the entire rational linear cancellation
space for the common constant and the prime logarithms on this range.

\begin{corollary}\label{em:boot:density}
For $N\ge1$,
\[
 \operatorname{trdeg}_{K_{\exp}}
 K_{\exp}(I_1,\ldots,I_N)\ge N-\pi(N)-1.
\]
There is a subset $S\subseteq\mathbb N$ of natural density one such that
$\{I_n:n\in S\}$ is algebraically independent over $K_{\exp}$.
At most $\pi(N)+1$ of $I_1,\ldots,I_N$ can be algebraic over
$K_{\exp}$.
\end{corollary}
\begin{proof}
The field generated by $y_1,\ldots,y_N$ contains the
$N-\pi(N)-1$ independent composite residuals. Multiplication by
$e^n\in K_{\exp}^{\times}$ does not change this field.
Process the $I_n$ in increasing order, retaining $I_n$ exactly when it
is transcendental over the field generated by the retained predecessors.
The retained terms are algebraically independent and their number in
each prefix is that prefix's transcendence degree. Thus at most
$\pi(N)+1=o(N)$ terms are omitted. Every term algebraic over
$K_{\exp}$ is omitted.
\end{proof}
The set $S$ is not identified. Density one does not decide the
arithmetic nature of the designated term $I_1=\delta$.

\subsection{A sharper consequence in the algebraic-gamma world}
\begin{proposition}\label{em:boot:A}
If $\gamma$ is algebraic over $K_{\exp}$, then
\[
 \operatorname{trdeg}_{K_{\exp}}
 K_{\exp}(I_1,\ldots,I_N)\ge N-\pi(N),
\]
and $\delta$ is transcendental over $K_{\exp}$. In particular these
conclusions hold under $\mathrm A$, the assumption
$\gamma\in\overline{\mathbb Q}$.
\end{proposition}
\begin{proof}
Over the algebraic extension $K_{\exp}(\gamma)$, annihilate only the
$\pi(N)$ logarithm columns in \eqref{em:boot:connection}. This leaves
$N-\pi(N)$ independent combinations, since algebraic translation
preserves algebraic independence. Algebraic base extension preserves
transcendence degree. Finally $U_1=\gamma+\delta/e$ is transcendental
over $K_{\exp}$; both summands cannot be algebraic over that field.
\end{proof}
This is a consequence of $\mathrm A$, not a contradiction to it.

\subsection{Why combining bases does not remove the exceptional term}
Let $\rho_h>0$ be algebraic, let $n_{ih}\in\mathbb Z$, and assume the
locations $\lambda_i=\prod_h\rho_h^{n_{ih}}$ are distinct. The connection
identity reads
\[
 y_i=U_i-\gamma-\sum_h n_{ih}\log\rho_h.
\]
If selected $y_i$ are all algebraic over $K_{\exp}$, their rows
$(1,n_{i1},\ldots,n_{ir})$ must be rationally independent. Otherwise a
nonzero rational relation would express a nonzero linear combination
of the independent $U_i$ as an algebraic number over $K_{\exp}$.
Under the hypothesis of Proposition~\ref{em:boot:A}, the exponent rows
$(n_{ih})_h$ alone must be independent.

These constraints and the preceding bounds are sharp for the stated
algebraic inputs, as the following two formal models demonstrate.
They are not replacements for the actual analytic Euler integrals.
Let $K$ be any characteristic-zero field, take independent
indeterminates $(U_n)_{n\ge1}$, and prescribe constants $a_p\in K$
for all primes.

In the first model set $g=a\in K$, and define
\[
 \ell_p=U_p-g-a_p,\quad
 \ell_n=\sum_pv_p(n)\ell_p,\quad y_n=U_n-g-\ell_n.
\]
Then every $y_p=a_p$, whereas $y_1$ and all composite $y_n$ are
algebraically independent: each has a unique corresponding variable
$U_n$. Every prefix has transcendence degree $N-\pi(N)$.

In the second model choose also $a_1\in K$, set $g=U_1-a_1$, and use
the same definitions. Then $g$ is transcendental, $y_1=a_1$, and
$y_p=a_p$ for every prime. The composite $y_n$ are algebraically
independent, and every prefix has transcendence degree
$N-\pi(N)-1$. Both models obey all multiplicative logarithm-addition
relations and composite cancellation identities simultaneously.
The symbols $\ell_p$ are formal logarithmic coordinates, not actual
real logarithms; no analytic exponential equation has been imposed.

The countermodel below gives a stronger analytic warning using actual
$E$-values: algebraically independent positive approximants can have an
algebraic limit with the same leading dyadic error rate. Hence the
missing input must bind the exact Euler normalization to arithmetic,
not only to independence and convergence. The formal models above
assert no logical independence from a foundational axiom system.

Finally, the expression of $R_n$ through $I_n$ uses the coefficients
$e^{-n}$. It does not show membership in the $D$-value ring, whose
closure under arbitrary $E$-value multiplication is unavailable.
Consequently the positive independent residuals do not by themselves
contradict an $E/D$ intersection conjecture.

\section{A Hankel bootstrap: gamma as an isolated moment atom}
\label{em:sec:borel-bootstrap}

This construction imports the classical Christoffel variational principle
into the several-pole Laplace family.  It produces a computable infinite
sequence with a positive, geometrically contracting error.  Its inputs are
$E$-values and logarithms, rather than rational numbers; no arithmetic
conclusion for gamma is asserted.  Christoffel minimization also appeared
in the earlier integrated manuscript for a different, rational Gregory
moment family.  The present result concerns the tensor determinants of
the exponential-integral periods.

Fix positive algebraic numbers $a,h$.  For $u>0$ put
\[
 I_u=\int_0^\infty\frac{e^{-t}}{t+u}\,dt=e^uE_1(u),\qquad
 F(u)=\operatorname{Ein}(u)-\log u=\gamma+E_1(u).
\]
For $n\ge1$, with row and column indices $0,\ldots,n-1$, define
\[
 H_n=(E_1(a+(i+j)h))_{i,j},\quad
 F_n=(F(a+(i+j)h))_{i,j},\quad
 v_n=(1,\ldots,1)^{\mathsf T}.
\]
Thus $F_n=H_n+\gamma v_nv_n^{\mathsf T}$.

\begin{theorem}[Positive affine collapse of the tensor tower]
\label{em:bor:positive-affine}
The matrices $H_n$ are positive definite.  Set
\[
 c_n=v_n^{\mathsf T}\operatorname{adj}(H_n)v_n.
\]
Then $c_n>0$, and
\begin{align}
 \det(F_n-sv_nv_n^{\mathsf T})
   &=\det H_n+(\gamma-s)c_n,\label{em:bor:affine}\\
 c_n&=v_n^{\mathsf T}\operatorname{adj}(F_n)v_n.\label{em:bor:c-free}
\end{align}
Consequently
\begin{equation}\label{em:bor:tensor-factor}
 \det(I_{a+(i+j)h})_{i,j=0}^{n-1}
   =e^{na+hn(n-1)}\det H_n>0.
\end{equation}
Under the canonical substitution
$I_u=e^u(\operatorname{Ein}(u)-\log u-\gamma)$, this determinant
has degree exactly one in the displayed gamma parameter, for every $n$.
\end{theorem}
\begin{proof}
Put $q=e^{-h}\in(0,1)$ and introduce the finite positive measure
\[
 d\mu_{a,h}(x)=\frac{x^{a/h-1}}{-\log x}\,dx,
 \qquad 0<x<q.
\]
The substitutions $E_1(u)=\int_1^\infty e^{-ut}\,dt/t$ and
$x=e^{-ht}$ give
\[
 E_1(a+kh)=\int_0^q x^k\,d\mu_{a,h}(x).
\]
The measure is positive on an interval, so the squared norm of any
nonzero polynomial is positive.  This proves positive definiteness.
Now $\operatorname{adj}(H_n)=\det(H_n)H_n^{-1}$, whence $c_n>0$.
The rank-one determinant identity gives \eqref{em:bor:affine}; applying it
once with base matrix $F_n$ and comparing the coefficient of $s$ proves
\eqref{em:bor:c-free}.  Finally factor $e^{a+ih}$ from row $i$ and
$e^{jh}$ from column $j$ to obtain \eqref{em:bor:tensor-factor}.
\end{proof}

\begin{theorem}[Gamma-free formula and certified error]
\label{em:bor:christoffel}
The quantities
\[
 A_n=\frac{\det F_n}{v_n^{\mathsf T}\operatorname{adj}(F_n)v_n}
\]
are defined by the entire series for $\operatorname{Ein}$ and ordinary
real logarithms, without taking gamma as input.  They satisfy
\begin{align}
 A_n-\gamma=\varepsilon_n
 &=\frac{1}{v_n^{\mathsf T}H_n^{-1}v_n}\notag\\
 &=\min_{\substack{\deg p<n\\p(1)=1}}
      \int_0^q|p(x)|^2\,d\mu_{a,h}(x)>0,\label{em:bor:variational}\\
 0<\varepsilon_{n+1}&<q^2\varepsilon_n,\label{em:bor:contraction}\\
 \varepsilon_n&\le
 \frac{E_1(a)}{T_{n-1}(2e^h-1)^2}.
 \label{em:bor:chebyshev}
\end{align}
Here $T_k$ is the Chebyshev polynomial.  In particular $A_n$ decreases
strictly to gamma.  The fully explicit weaker upper bound
\[
 \varepsilon_n\le
 \frac{e^{-a}/a}{T_{n-1}(2e^h-1)^2}
\]
provides certified intervals for gamma using no $E_1$ input.
\end{theorem}
\begin{proof}
Equation \eqref{em:bor:affine} at $s=0$, together with
\eqref{em:bor:c-free}, gives $A_n=\gamma+\det(H_n)/c_n$.
For a coefficient vector $b$ of $p$, the objective in
\eqref{em:bor:variational} is $b^*H_nb$ and the constraint is
$v_n^{\mathsf T}b=1$.  The unique minimizer is
$b=H_n^{-1}v_n/(v_n^{\mathsf T}H_n^{-1}v_n)$; this proves the
variational identity.  If $p_n$ is the minimizing polynomial, $xp_n$
is feasible at level $n+1$.  Since $0<x<q$ on the support,
\[
 \varepsilon_{n+1}\le\int x^2|p_n(x)|^2d\mu_{a,h}(x)
       <q^2\varepsilon_n.
\]
For the sharper bound, use the feasible polynomial
\[
 p(x)=\frac{T_{n-1}(2x/q-1)}{T_{n-1}(2/q-1)}.
\]
Its numerator has modulus at most one on $[0,q]$; the measure's total
mass is $E_1(a)$.  Finally $E_1(a)<e^{-a}/a$ follows immediately
from its integral over $[a,\infty)$.
\end{proof}

The rate in \eqref{em:bor:chebyshev} is exponential: if
$\xi=2e^h-1$ and $\rho=\xi+\sqrt{\xi^2-1}>1$, then
$T_k(\xi)=(\rho^k+\rho^{-k})/2$.
The moments $F(a+kh)$ are precisely those of
\[
 \nu_{a,h}=\mu_{a,h}+\gamma\delta_1.
\]
Thus the construction recovers the mass of an isolated atom at $1$.
This is an application of the classical Christoffel principle, not a
claim of a new principle in moment theory.  A primary reference is
B.~Simon, \emph{The Christoffel--Darboux Kernel}~\cite{Simon}.

\subsection{The smallest case and the precise arithmetic gap}
At $n=2$, put $u_j=F(a+jh)$ for $j=0,1,2$.  Then
\[
 c_2=u_0-2u_1+u_2
     =E_1(a)-2E_1(a+h)+E_1(a+2h)>0,
\]
and
\[
 A_2=\frac{u_0u_2-u_1^2}{u_0-2u_1+u_2},\qquad
 A_2-\gamma
 =\frac{E_1(a)E_1(a+2h)-E_1(a+h)^2}{c_2}>0.
\]
The positivity of the numerator is the strict Cauchy--Schwarz
inequality for the two monomials $1,x$ in $L^2(\mu_{a,h})$.
For $a=h=1$, this tests the algebraic-coefficient tensor
$I_1I_3-I_2^2$.  Its quadratic gamma coefficient cancels because
$1+3=2+2$, but its linear coefficient is $-e^4c_2\ne0$.
Thus the assertion that the highest gamma coefficient can never cancel
in tensor products is false; the full tower nevertheless retains gamma.

For $a=h=1$ and $n=8$, the numerical diagnostic recorded in the Korean
companion gives $\varepsilon_8\simeq1.94354\times10^{-14}$ and the
fully explicit upper bound $9.40656\times10^{-14}$. Independent
160-digit arithmetic reproduces these figures. These decimal values
are diagnostics; the rigorous upper bound is the formula in
Theorem~\ref{em:bor:christoffel}.

The input count is $2n-1$ values of $\operatorname{Ein}$ and their
associated logarithms.  This formula does not produce rational approximants or
algebraic coefficients with a height bound.  The accompanying
independence argument in fact proves transcendence for a density-one
subfamily when $a=h=1$.  Therefore the positive error and its exponential
decrease are not small integer polynomial values at gamma.  An arithmetic compression theorem producing such
polynomials with an adequate height budget would be an additional
theorem, not a consequence of \eqref{em:bor:variational}.  Under an
algebraicity assumption on gamma the measure $\nu_{a,h}$ remains
perfectly legitimate, so its atom does not itself yield a contradiction.

\subsection{What algebraic-coefficient formal tensor cancellation allows}
Let $\lambda_1,\ldots,\lambda_m$ be distinct algebraic numbers and
$v_i=e^{\lambda_i}$.  For an algebraic-coefficient polynomial
$P\in\overline{\mathbb Q}[X_1,\ldots,X_m]$,
\begin{equation}\label{em:bor:formal-invariants}
 P(X+sv)=P(X)\text{ in }\mathbb C[X,s]
 \quad\Longleftrightarrow\quad P\text{ is constant}.
\end{equation}
Indeed differentiation gives
$\sum_i e^{\lambda_i}\partial_iP=0$.
The coefficient of each monomial is an algebraic linear combination of
the distinct exponentials.  Lindemann--Weierstrass forces every one of
its summands to vanish; hence every $\partial_iP$ is zero.
The same argument works for rational functions: for coprime
$A,B\in\overline{\mathbb Q}[X]$ with $B\ne0$, differentiate $A/B$
and multiply by $B^2$. Lindemann--Weierstrass then gives
$(\partial_iA)B-A\partial_iB=0$ for every $i$, so $A/B$ is constant.

The coefficient-field hypothesis matters.  Over $K_{\exp}$, the
nonconstant polynomial $e^{-\lambda_1}X_1-e^{-\lambda_2}X_2$
is invariant.  Also \eqref{em:bor:formal-invariants} is a formal identity
in independent $X_i$.  It is not a theorem that specialization at the
actual period values creates no additional identities.  Such a claim
would require arithmetic information not supplied by this calculation.

\subsection{Synthesis: the Hankel approximants themselves have density-one independence}
\label{em:sec:hankel-independence}

The multipoint $\operatorname{Ein}$ input and the Christoffel construction
combine to give arithmetic information about the approximants themselves.
Specialize $a=h=1$, put
\[
 m_k=\operatorname{Ein}(k+1)-\log(k+1),\qquad
 F_n=(m_{i+j})_{0\le i,j<n},\qquad
 A_n=\frac{\det F_n}{v_n^{\mathsf T}\operatorname{adj}(F_n)v_n},
\]
and retain the field $K_{\exp}$ from Section~\ref{em:sec:bootstrap}.

\begin{theorem}\label{em:boot:hankel-independent}
For every $N\ge1$,
\begin{equation}\label{em:boot:hankel-rank}
 \operatorname{trdeg}_{K_{\exp}}
 K_{\exp}(A_1,\ldots,A_N)\ge N-\pi(2N-1).
\end{equation}
There is a subset $S\subseteq\mathbb N$ of natural density one for which
$\{A_n:n\in S\}$ is algebraically independent over $K_{\exp}$.
At most $\pi(2N-1)$ of the first $N$ approximants are algebraic over
$K_{\exp}$.
\end{theorem}
\begin{proof}
First we show that the $A_n$ recover every even-index moment once the
odd-index moments are supplied. Fix $n$, put $d=n-1$, and regard the
largest moment $t=m_{2d}$ as variable, keeping the preceding moments
fixed. It occurs only in the bottom-right matrix entry. Both the
determinant and the sum of cofactors are affine in $t$, so
\begin{equation}\label{em:boot:mobius}
 A_n=\frac{\alpha t+\beta}{\kappa t+\lambda},
 \qquad
 \alpha,\beta,\kappa,\lambda\in
 K_{\exp}(m_0,\ldots,m_{2d-1}).
\end{equation}
For $n=1$ this reads $A_1=m_0$ over the empty preceding moment field.

At the actual moments $F_n$ is positive definite, since $\gamma>0$. Set
\[
 b=\frac{F_n^{-1}v_n}{v_n^{\mathsf T}F_n^{-1}v_n},\qquad
 p(x)=\sum_{j=0}^d b_jx^j.
\]
This polynomial minimizes the squared norm for
$\nu=\mu_{1,1}+\gamma\delta_1$ subject to $p(1)=1$. The atom adds the
constant $\gamma$ to the objective, so $p$ also minimizes the norm for
$\mu_{1,1}$. It must have degree exactly $d$. Otherwise, for $d\ge1$,
$(x-1)p$ would be an admissible variation of degree at most $d$
vanishing at $1$. Stationarity would give
\[
 \int_0^{e^{-1}}(x-1)p(x)^2\,d\mu_{1,1}(x)=0,
\]
which is impossible: the integrand is strictly negative away from
finitely many zeros of the nonzero real polynomial $p$.
For $d=0$, $p=1$ directly. Thus $b_d\ne0$.

Since $A_n=(v_n^{\mathsf T}F_n^{-1}v_n)^{-1}$, differentiation with
respect to the bottom-right entry gives
\[
 \frac{\partial A_n}{\partial t}=b_d^2>0.
\]
Together with \eqref{em:boot:mobius}, this proves
$\alpha\lambda-\beta\kappa\ne0$ at the preceding moment values.
The exact inverse
\[
 t=\frac{\beta-A_n\lambda}{A_n\kappa-\alpha}
\]
has nonzero denominator there, since
$A_n\kappa-\alpha=(\kappa\beta-\alpha\lambda)/
(\kappa t+\lambda)$. Induction on $n$ therefore proves the field identity
\begin{equation}\label{em:boot:moment-recovery}
 K_{\exp}(m_0,\ldots,m_{2N-2})
 =K_{\exp}(A_1,\ldots,A_N,m_1,m_3,\ldots,m_{2N-3}).
\end{equation}

On adjoining the $\pi(2N-1)$ numbers $\log p$ for primes
$p\le2N-1$, the moment field contains all
$\operatorname{Ein}(1),\ldots,\operatorname{Ein}(2N-1)$, because
$\log j=\sum_pv_p(j)\log p$. Their independence implies
\[
 \operatorname{trdeg}_{K_{\exp}}
 K_{\exp}(m_0,\ldots,m_{2N-2})
 \ge2N-1-\pi(2N-1).
\]
In \eqref{em:boot:moment-recovery} there are only $N-1$ extra odd-index
moments. Subtracting that possible transcendence-degree contribution
proves \eqref{em:boot:hankel-rank}. Greedy retention of a transcendence
basis, as in Corollary~\ref{em:boot:density}, omits at most
$\pi(2N-1)=o(N)$ terms and proves the remaining assertions.
\end{proof}

This is a synthesis of the two constructions: the new Hankel
approximants cannot all be rational, or even algebraic over $K_{\exp}$.
They contain a density-one independent subfamily while decreasing to
$\gamma$. This does not force their limit to be transcendental.
It rules out treating the raw $A_n$ as a sequence of rational or
algebraic approximants with an ordinary Diophantine height. Any
subsequent rationalization would still require its own error and
height analysis.

\section{An explicit countermodel to independence-based bootstrapping}
\label{em:sec:audit-countermodel}

The following construction uses the actual $E$-values already available
in the program. It identifies a collection of hypotheses that does not
determine the arithmetic nature of a limit.
Write $U_N=\Ein(2^N)$, $u_N=E_1(2^N)$, and
\[
 D_N=U_N-2U_{N+1}+U_{N+2}
     =u_N-2u_{N+1}+u_{N+2}.
\]
Here $D_N$ denotes a difference, not a binomial denominator.

\begin{proposition}[Algebraic limits with the same leading error]\label{em:algebraic-limit-countermodel}
For any real algebraic number $a$, the sequence
\[
 Z_N=a+(N+1)D_N,\qquad N\ge1,
\]
consists of $E$-values algebraically independent over $K_{\exp}$.
It decreases strictly to $a$ from above, and
\[
 Z_N-a\sim\frac{N+1}{2^N}e^{-2^N}.
\]
Its arguments have logarithmic height $O(N)$ and its linear
coefficients have logarithmic height $O(\log(N+1))$.
\end{proposition}
\begin{proof}
Every finite second-difference matrix has full row rank, so the
multipoint input proves independence of its $D_N$ values. Algebraic
translations and nonzero rational rescalings preserve independence;
closure of $E$-values under algebraic linear combinations gives the
first assertion.

Put $u(t)=E_1(t)$ and $d(t)=u(t)-2u(2t)+u(4t)$.
Substitution gives $u(2t)\le e^{-t}u(t)$. For $t\ge2$ this implies
\[
 (1-2e^{-t})u(t)\le d(t)<u(t),\qquad
 \frac{d(2t)}{d(t)}\le\frac{e^{-t}}{1-2e^{-t}}.
\]
Consequently
\[
 \frac{Z_{N+1}-a}{Z_N-a}
 \le\frac32\frac{e^{-2}}{1-2e^{-2}}<1\qquad(N\ge1).
\]
All errors are positive. The standard integration-by-parts estimate
$E_1(t)\sim e^{-t}/t$ and the same ratio bound give the asymptotic.
The displayed arguments and coefficients prove the height assertions.
\end{proof}

This is the same leading error as
\[
 Y_N=(N+1)\Ein(2^N)-N\Ein(2^{N+1})
 =\gamma+(N+1)u_N-Nu_{N+1}.
\]
Even the polynomial slope detector survives: for every nonzero
$P\in\mathbb C[X]$,
\[
 \lim_{N\to\infty}\frac{-\log|P(Z_N)|}{2^N}
 =\operatorname{ord}_{X=a}P.
\]
Indeed, factor $P(X)=(X-a)^mQ(X)$ with $Q(a)\ne0$.
Thus independence, actual $E$-value membership, positivity, monotonicity,
small coefficient heights, and the known leading decay do not together
force a transcendental limit. The countermodel changes the exact balanced
formula: its $\Ein$ coefficient sum is zero, whereas the sum for $Y_N$
is one. It does not rule out a future argument using this normalization
and additional arithmetic information.

\begin{lemma}[Algebraic tail intersection]\label{em:tail-intersection}
Let $x_1,x_2,\ldots$ be algebraically independent over a field
$K\subset\mathbb C$, and put $F_M=K(x_n:n\ge M)$. Then
\[
 \bigcap_M F_M=K,
 \qquad
 \bigcap_M\{z\in\mathbb C:z\text{ is algebraic over }F_M\}
 =\{z\in\mathbb C:z\text{ is algebraic over }K\}.
\]
\end{lemma}
\begin{proof}
Every expression involves finitely many variables. Disjoint independent
finite tuples $A,B$ have $K(A)\cap K(B)=K$, by clearing denominators
and comparing the resulting polynomial identity. This proves the first
claim. If $z$ is algebraic over both $K(A)$ and $K(B)$ but transcendental
over $K$, transcendence-degree submodularity gives
\[
 |A|+|B|\ge
 \operatorname{trdeg}_K K(A,B,z)+\operatorname{trdeg}_K K(z)
 \ge |A|+|B|+1,
\]
a contradiction. Taking one finite early tuple and one sufficiently
late tuple proves the second claim; the converse is immediate.
\end{proof}

For the independent $Y_N$ over $\overline{\mathbb Q}$, a proof that
$\gamma$ is algebraic over every tail field would prove algebraicity,
not transcendence. Ordinary convergence establishes no such algebraic
membership. These are elementary deductions used as stopping rules;
no novelty is asserted for the field-intersection lemma, and no gamma
alternative is rejected here.

\section{Importing zero theorems: a fresh-prime obstruction}
\label{em:sec:imports}

Put $D=d/dz$ and, for $r\in\overline{\mathbb Q}$, define
\[
 F_r(z)=e^z(\operatorname{Ein}(z)-r)
       =\sum_{n\ge0}(H_n-r)\frac{z^n}{n!}.
\]
The identity follows from $F_r'=F_r+(e^z-1)/z$ and $F_r(0)=-r$.
In the algebraic gamma world one may take $r=\gamma$.
Throughout this section a rational differential operator is held fixed
when its input parameter $r$ is varied.

An $E$-function $G(z)=\sum g_nz^n/n!$ is \emph{globally bounded} in
the convention of \cite[Introduction]{FRzeros} if some positive
integer $M$ makes every $M^{n+1}g_n$ an algebraic integer. Thus global
boundedness is imposed on the \emph{normalized} coefficients $g_n$,
equivalently on the coefficients of the associated ordinary
$G$-series $\sum g_nz^n$. The following result uses the weaker
condition of integrality outside a fixed finite set of primes.

\begin{theorem}[Finite prime support erases the input parameter]
\label{em:imp:erasure}
Let $r\in\overline{\mathbb Q}$ and
$L\in\overline{\mathbb Q}(z)[D]$. Suppose $G=L(F_r)$ is regular
at $0$, and write $G=\sum_{n\ge0}g_nz^n/n!$.
Suppose there are a number field $K$ containing every $g_n$ and a
finite set $S$ of rational primes such that every $g_n$ is integral
at every prime of $K$ above every $p\notin S$. Then
\[
 L(e^z)=0.
\]
In particular, for this fixed $L$ the expression $L(F_s)$ is independent
of $s\in\mathbb C$. The conclusion applies whenever $L(F_r)$ is a
globally bounded $E$-function.
\end{theorem}

\begin{proof}
Successive differentiation, using the displayed equation for $F_r$,
gives
\[
 L(F_r)=a(z)F_r+b(z)e^z+c(z),\qquad
 a,b,c\in\overline{\mathbb Q}(z),\qquad
 a(z)e^z=L(e^z).
\]
Indeed differentiation sends $a$ to $a'+a$, which is exactly the rule
for the coefficient of $e^z$ when differentiating $ae^z$.
Enlarge $K$ to contain $r$ and all these rational functions' coefficients.
Integrality away from $S$ is preserved by extension of number fields.
Choose a nonzero $q\in K[z]$ that clears their denominators, giving
\[
 qG=A(z)F_r+B(z)e^z+C(z),\qquad A,B,C\in K[z].
\]
Multiplication by the fixed polynomial $q$ only enlarges $S$ by
finitely many primes: normalized coefficients transform into finite
sums of algebraic coefficients times falling factorials and previous
normalized coefficients.

Suppose $A=\sum_{k=0}^d a_kz^k\ne0$. There is an $m\in\{0,\ldots,d\}$
such that
\[
 A_m=\sum_{k=0}^{\min(d,m)}a_km^{\underline{k}}\ne0,
\]
because the matrix $(m^{\underline{k}})_{0\le m,k\le d}$ is triangular
with nonzero diagonal. Write $qG=\sum h_nz^n/n!$.
Take any sufficiently large rational prime $p$. Exclude the finitely
many primes needed for integrality of the fixed coefficients and $r$,
and the numerator and denominator primes of the norm of $A_m$.
For any prime $\mathfrak p$ of $K$ above $p$, $A_m$ is then a unit.
For $p>m+d$, harmonic numbers satisfy
\[
 pH_{p+m-k}\equiv
 \begin{cases}1,& k\le m,\\0,&k>m\end{cases}
 \pmod{\mathfrak p}.
\]
Only the denominator $p$ contributes to this congruence, since all
indices in question are positive and less than $2p$.
The normalized coefficient of $A F_r$ at $n=p+m$ is
\[
 \sum_{k=0}^d a_k(p+m)^{\underline{k}}(H_{p+m-k}-r).
\]
The coefficients supplied by $B e^z+C$ are integral at
$\mathfrak p$. Reducing after multiplication by $p$ therefore gives
\[
 p h_{p+m}\equiv A_m\not\equiv0\pmod{\mathfrak p}.
\]
This contradicts integrality of $h_{p+m}$. Thus $A=0$, so $a=0$ and
$L(e^z)=0$. Finally $F_s-F_r=(r-s)e^z$ proves independence of $s$.
\end{proof}

If $L(e^z)\ne0$, the proof gives more than failure of global
boundedness. After the fixed polynomial multiplication, every
sufficiently large prime $p$ occurs in the denominator of the single
coefficient $h_{p+m}$, for one fixed offset $m\le\deg A$.
At unramified primes above such a $p$ its valuation is exactly $-1$.
No claim of literature-wide priority is made for this deduction.

The condition is necessary, not sufficient. The operator
$L=D-1$ erases the parameter but produces $(e^z-1)/z$, whose normalized
coefficients $1/(n+1)$ have infinitely many denominator primes.
The operator $L=z(D-1)$ does produce the globally bounded function
$e^z-1$. For the intended gamma application, $L$ must be constructed
with coefficients independent of gamma: allowing $L=\gamma z(D-1)$
under the algebraic hypothesis simply reinserts gamma as a coefficient.

\subsection{The recent zero-factorization theorem: exact applicability}

Fischler--Rivoal's \cite[Theorem~1.5]{FRzeros} assumes a nonzero globally
bounded $E$-function over $\mathbb Q$, with every zero of one common
multiplicity $m\ge2$ and a rational differential operator of order at
most $m+1$ annihilating it. Its conclusion is an $m$th-power
factorization by a globally bounded $E$-function of order at most two.
An output merely defined over $\overline{\mathbb Q}$ would not by
itself meet the rational-coefficient hypothesis. Their
\cite[Theorem~1.2 and Corollary~1.3]{FRzeros} also establish a special
case of entire quotient rigidity when the numerator is obtained by a
rational differential operator applied to the denominator.
Theorem~\ref{em:imp:erasure} closes one proposed preprocessing step:
no fixed gamma-independent rational differential operator makes
$F_\gamma$ globally bounded while retaining its input gamma parameter.
This conclusion does not exclude nonlinear constructions, several
independent framings, or a different input family. Nor does it claim
that global boundedness is necessary for every zero theorem.

\subsection{The exact mixed specialization and its logical cost}

Define
\[
 \widehat d(z)=\sum_{n\ge0}(-1)^nn!z^{-n-1},\qquad
 \Psi_r(z)=F_r(z)-\widehat d(z).
\]
Positive-direction summation yields
\[
 d_0(z)=e^zE_1(z),\qquad
 \Psi_{r,0}(z)=e^z(\gamma+\log z-r),\qquad
 \Psi_{r,0}(1)=e(\gamma-r).
\]

\begin{proposition}[Formal nondivisibility]\label{em:mixed-nondivisibility}
For every algebraic $r$, there is no mixed arithmetic Gevrey function
$Q$ such that $(z-1)Q=\Psi_r$.
\end{proposition}
\begin{proof}
Write $Q=G(z)+h(1/z)$ with $G$ an $E$-function and $h(1/z)$ having
only negative powers. If $c\in\overline{\mathbb Q}$ is the coefficient
of $z^{-1}$ in $h$, comparison of nonnegative powers gives
\[
 F_r(z)=(z-1)G(z)+c.
\]
Consequently $F_r(1)=c$. But $e$ and $\operatorname{Ein}(1)$ are
algebraically independent by the established $E$-function input;
thus $e(\operatorname{Ein}(1)-r)$ is transcendental, a contradiction.
\end{proof}

It follows that the statement
\[
 \forall r\in\overline{\mathbb Q},\quad
 \Psi_{r,0}(1)=0\ \Longrightarrow\
 \Psi_r\text{ is divisible by }z-1\text{ in the mixed class}
\]
is \emph{equivalent} to the transcendence of gamma. One direction
uses the proposition; the reverse direction is a vacuous implication.
This is the exact family specialization of
\cite[Conjecture~3]{FRmixed}, not a proved weakening of that
conjecture. Pure $E$-function division is a theorem
\cite[Proposition~4.1]{Beukers}; replacing its input class by mixed
functions is the missing arithmetic step.

\subsection{Worked \texorpdfstring{$p$}{p}-curvature calculation}

Put $Y=(\Psi_{r,0},e^z)^t$. Then
\[
 Y'=A(z)Y,\qquad A(z)=I+N/z,\qquad
 N=\begin{pmatrix}0&1\\0&0\end{pmatrix},\quad N^2=0.
\]
Its fundamental matrix is
\[
 e^z\begin{pmatrix}1&\log z\\0&1\end{pmatrix}.
\]
Neither the coefficient matrix nor this chosen fundamental matrix
depends on $r$ or gamma. The constant vector selecting $Y$ is
$(\gamma-r,1)^t$.

Define $A_0=I$ and $A_{n+1}=A_n'+A_nA$, so that $Y^{(n)}=A_nY$.
Direct differentiation gives
\[
 A_n=I+N\sum_{k=1}^n\binom nk(-1)^{k-1}(k-1)!z^{-k}.
\]
Modulo every prime $p$, the interior binomial coefficients disappear
and Wilson's congruence gives
\[
 A_p\equiv I-Nz^{-p}.
\]
This matrix is nonzero and invertible. With the opposite connection
sign convention the corresponding curvature changes sign; its
nonvanishing and gamma-independence are unchanged. Tensoring out the
exponential rank-one factor leaves the logarithmic curvature
$-Nz^{-p}$, which is nonzero nilpotent.

The effective order-one theorem of F\"urnsinn--Pannier
\cite[Theorem~1.3]{FP} concerns algebraicity of functions
solving a specified rational differential equation. It supplies no
arithmetic evaluation rule for the constant vector selecting a
particular solution of the present rank-two system. Thus its use here
does not distinguish the two gamma worlds. A result controlling the
arithmetic of the connection framing is still required.

\section{An imported categorical alternative and a failed positivity shortcut}

Let $\mathcal D_t=\operatorname{Rep}(S_t)$ denote the complex-linear
additive Karoubi envelope of the partition category, with its standard
symmetric tensor structure. This is not the representation category
of a literal group with nonintegral size. Put
\[
 R_+(t)=\left\{\sum_{n=0}^N a_n\binom tn:
                a_n\in\Z_{\ge0},\ N<\infty\right\}.
\]
We import two established results. For parameters outside $\Z_{\ge0}$,
Harman--Kalinov characterize the existence of complex-linear strong
symmetric monoidal functors $\mathcal D_s\to\mathcal D_t$ by the
condition $s\in R_+(t)$ \cite[Corollary 3.2.4]{HK}. Kalinov--Mandelshtam prove that $R_+(t)$
is a ring precisely when $t$ is algebraic and not a nonnegative integer
\cite[Theorem 1.3]{KM}. Here and below the parameters in the functor
statement are outside $\Z_{\ge0}$.

\begin{corollary}[A categorical reformulation]
For the actual $0<\gamma<1$, the following are equivalent:
\[
 \gamma\in\Qbar
 \quad\Longleftrightarrow\quad -1\in R_+(\gamma)
 \quad\Longleftrightarrow\quad
 \exists\,\mathcal D_{-1}\longrightarrow\mathcal D_\gamma
\]
where the arrow is a complex-linear strong symmetric monoidal functor.
\end{corollary}
\begin{proof}
The semiring $R_+(t)$ contains $1$, and is a ring exactly when it
contains $-1$, since multiplication by $-1$ then supplies every
additive inverse. Apply the two imported results. The finite relation
$1+\sum a_n\binom\gamma n=0$ is a nonzero rational polynomial
relation, since its polynomial has value $1+a_0>0$ at $0$.
\end{proof}

This is a direct corollary of known work, not an independently solved
problem. Under such a functor the source fundamental object maps to a
commutative Frobenius algebra of categorical dimension $-1$ in
$\mathcal D_\gamma$; this is not ordinary vector-space dimension.
A tempting positivity contradiction fails for an exact reason.
If $X$ is the fundamental object of $\mathcal D_t$, symmetric categorical
trace gives
\[
 \dim(\Lambda^2X)
 =\tfrac12\left(\operatorname{Tr}(\mathrm{id}_{X\otimes X})
                 -\operatorname{Tr}(c_{X,X})\right)
 =\frac{t(t-1)}2<0\quad(0<t<1).
\]
Thus no realization preserving this trace can land in a category where
all such dimensions are nonnegative. This obstruction applies to
\emph{both} arithmetic alternatives, and is not an axiom of
$\mathcal D_t$ itself. At $t=1/2$ the exact identity
$8\binom{1/2}{2}=-1$ gives the stated functor, even though
$\dim(\Lambda^2X)=-1/8$.

The word ``unitary'' alone cannot supply the missing trace-positivity
assumption: intrinsic $C^*$ dimensions and dimensions from a chosen
symmetric trace need not agree, as super symmetry illustrates.
The remaining arithmetic task is to prove existence or nonexistence
of the specified functor for $\gamma$, not to impose a positivity
axiom that already excludes valid algebraic examples.

% END INTEGRATED SOURCE

\clearpage
\part{Functional transcendence and arithmetic costs}

% BEGIN INTEGRATED SOURCE: snow_core/body.tex
\section{From arithmetic rank to differential transcendence}
\label{sc:sec:independence-dynamics}

Set $K=\overline{\mathbb Q}(e^\alpha:\alpha\in\overline{\mathbb Q})$.
The multipoint and Hankel recovery arguments reproduced in
Section~\ref{sc:sec:reused-inputs} give, for $N\ge1$,
\begin{align}
 \operatorname{trdeg}_K K(I_1,\ldots,I_N)&\ge N-\pi(N)-1,
 &I_n&=e^nE_1(n),\label{sc:dyn:I-rank}\\
 \operatorname{trdeg}_K K(A_1,\ldots,A_N)&\ge N-\pi(2N-1).
 \label{sc:dyn:A-rank}
\end{align}
Here $A_n$ denotes the $a=h=1$ Hankel approximant: with
$m_k=\operatorname{Ein}(k+1)-\log(k+1)$,
$F_n=(m_{i+j})_{0\le i,j<n}$ and $v_n=(1,\ldots,1)^{\mathsf T}$,
\[
 A_n=\frac{\det F_n}{v_n^{\mathsf T}\operatorname{adj}(F_n)v_n}.
\]
We transfer these rank bounds to differential equations and nonlinear
recurrences. No novelty claim is made for the general principles or
for their applications below.

\subsection{A coefficient-field lemma, including singular equations}
\begin{lemma}\label{sc:dyn:finite-field}
Let $L\subseteq\mathbb C$ be a characteristic-zero field. If a formal
power series $f(z)=\sum_{n\ge0}a_nz^n$ satisfies a nonzero polynomial
differential equation over $L(z)$, then for some $M$ every $a_n$ belongs
to $L(a_0,\ldots,a_M)$.
\end{lemma}
\begin{proof}
Clear denominators and choose a nonzero annihilator
$F\in L[z,Y_0,\ldots,Y_r]$ of least total degree in the $Y$ variables.
This degree is positive. At least one series
\[
 S_j(z)=\partial_{Y_j}F(z,f,\ldots,f^{(r)})
\]
is nonzero: a nonzero partial polynomial has smaller total degree and
cannot itself annihilate $f$. Put
\[
 v=\min_{S_j\ne0}(\operatorname{ord}_zS_j-j),\qquad
 b_j=[z^{v+j}]S_j,\qquad B(T)=\sum_{j=0}^r b_j(T)_j.
\]
Negative coefficient indices mean zero. At least one $b_j$ is nonzero,
and the falling factorials $(T)_j$ have distinct degrees, so $B\ne0$.
Its coefficients depend on only finitely many initial $a_n$.

Write $f_{<n}=\sum_{k<n}a_kz^k$ and $h=f-f_{<n}$. For sufficiently
large $n$, polynomial Taylor expansion gives
\[
 0=F(z,f_{<n},\ldots,f_{<n}^{(r)})
   +\sum_{j=0}^r S_j(z)h^{(j)}(z)+O(z^{2(n-r)}).
\]
Evaluating $S_j$ at $f$ instead of $f_{<n}$ changes the linear term
only by $O(z^{2(n-r)})$; every higher Taylor term has at least that
order. If $n\ge r$ and $n>v+2r$, the coefficient of $z^{n+v}$ gives
\begin{equation}\label{sc:dyn:coefficient-recursion}
 B(n)a_n=-[z^{n+v}]F(z,f_{<n},\ldots,f_{<n}^{(r)}).
\end{equation}
Indeed terms of $h$ with index greater than $n$ have order greater
than $n+v$ after linearization. The nonzero polynomial $B$ has only
finitely many nonnegative integer roots. Beyond these roots and the
finite initial indices defining $B$, equation
\eqref{sc:dyn:coefficient-recursion} determines each coefficient rationally
from earlier ones. Induction proves the assertion.
\end{proof}
This is the classical nonlinear coefficient-recursion mechanism. For a
primary modern account of Mahler's recurrence, see
Krattenthaler--Rivoal, \emph{Arithmetic properties of the Taylor
coefficients of differentially algebraic power series},
\url{https://arxiv.org/html/2502.09259v1}, Section~1, equation~(1.4).
The proof above supplies the precise general coefficient-field version
needed here; no assumption that the $a_n$ are algebraic is used.

\begin{corollary}\label{sc:dyn:infinite-rank}
If the coefficient field of $f\in\mathbb C[[z]]$ has infinite
transcendence degree over a fixed characteristic-zero subfield
$K\subseteq\mathbb C$, then $f$ is differentially transcendental over
$\mathbb C(z)$.
\end{corollary}
\begin{proof}
A putative differential equation contains only finitely many complex
constants $c_1,\ldots,c_s$ in its coefficients. Apply
Lemma~\ref{sc:dyn:finite-field} over $L=K(c_1,\ldots,c_s)$. All the Taylor
coefficients would then lie in $L(a_0,\ldots,a_M)$, a field of finite
transcendence degree over $K$, a contradiction.
\end{proof}
Multiplication of the individual Taylor coefficients by prescribed
nonzero elements of $K$ leaves their generated field unchanged.
In particular ordinary and exponential generating series have the
same obstruction.

\subsection{Explicit differentially transcendental integral transforms}
\begin{theorem}\label{sc:dyn:integral-transforms}
The analytic function and the entire function
\begin{align}
 \mathcal I(z)&=\int_0^1\frac{dx}{(1-zx)(1-\log x)}
       &&(|z|<1),\label{sc:dyn:stieltjes}\\
 \widehat{\mathcal I}(z)&=\int_0^1\frac{e^{zx}}{1-\log x}\,dx
       &&(z\in\mathbb C)\label{sc:dyn:laplace}
\end{align}
are differentially transcendental over $\mathbb C(z)$.
The ordinary and exponential generating functions of $E_1(n)$,
$\operatorname{Ein}(n)$, and $A_n$ have the same property.
\end{theorem}
\begin{proof}
Tonelli's theorem and $1/(n+t)=\int_0^1x^{n+t-1}dx$ give
\[
 I_n=\int_0^\infty\frac{e^{-t}}{n+t}\,dt
     =\int_0^1\frac{x^{n-1}}{1-\log x}\,dx.
\]
Geometric and exponential summation, uniform on compact subsets of the
respective domains, identify \eqref{sc:dyn:stieltjes} with
$\sum_{n\ge1}I_nz^{n-1}$ and \eqref{sc:dyn:laplace} with
$\sum_{n\ge1}I_nz^{n-1}/(n-1)!$. Equation~\eqref{sc:dyn:I-rank} and
Corollary~\ref{sc:dyn:infinite-rank} apply. Multiplication by
$e^{-n}\in K^\times$ gives the $E_1(n)$ assertion. The
$\operatorname{Ein}(n)$ are already algebraically independent over $K$.
Use \eqref{sc:dyn:A-rank} for $A_n$. All these ordinary generating series
have positive radius of convergence, and their exponential versions
are entire: $I_n=O(1/n)$, $E_1(n)=O(e^{-n}/n)$,
$\operatorname{Ein}(n)=O(\log(n+1))$, and $A_n$ is bounded.
\end{proof}
These are generating transforms of sampled values. The functions
$E_1(z)$ and $\operatorname{Ein}(z)$ themselves satisfy linear
differential equations; no contrary claim is being made. Nor does
functional differential transcendence decide the arithmetic nature of
an individual value such as $I_1=\delta$.
In particular, if $H(z)=\widehat{\mathcal I}(z)$, then
\[
H(0)=\int_0^1\frac{dx}{1-\log x}=I_1=eE_1(1)=\delta.
\]
The theorem classifies $H$ as a function, including its analytic germ at
zero; it does not classify the number $H(0)$. Differential transcendence
alone imposes no arithmetic classification on a prescribed value.

\subsection{Removing the gamma pole does not remove the obstruction}
Put
\[
 \mathcal A(z)=\sum_{n\ge1}A_nz^{n-1},\qquad
 \widehat{\mathcal A}(z)=\sum_{n\ge1}\frac{A_nz^{n-1}}{(n-1)!},
 \quad \xi=2e-1,\quad \rho=\xi+\sqrt{\xi^2-1}.
\]
The proved Christoffel bound is
\[
 0<A_n-\gamma\le\frac{E_1(1)}{T_{n-1}(\xi)^2}
                  \le4E_1(1)\rho^{-2(n-1)}.
\]
Consequently
\[
 \mathcal E_A(z)=\mathcal A(z)-\frac{\gamma}{1-z}
                =\sum_{n\ge1}(A_n-\gamma)z^{n-1}
\]
is analytic for $|z|<\rho^2$, and the entire function
\[
 \widehat{\mathcal E}_A(z)=\widehat{\mathcal A}(z)-\gamma e^z
\]
satisfies $|\widehat{\mathcal E}_A(z)|\le
4E_1(1)\exp(|z|/\rho^2)$. Both functions remain differentially
transcendental over $\mathbb C(z)$. Indeed finite transcendence degree
for their coefficient field would, on adjoining the single constant
$\gamma$, contradict \eqref{sc:dyn:A-rank}. Apply
Corollary~\ref{sc:dyn:infinite-rank}. Thus the simple pole at $1$ does not
account for the differential transcendence. Its residue's arithmetic
nature is still not determined.

\section{Quantitative exclusion of nonlinear recurrence templates}
\label{sc:sec:nonlinear-recurrence}

The argument below does not assume that a recurrence can be solved for
its last variable. It therefore includes singular or degenerate
specializations.

\begin{theorem}\label{sc:dyn:window}
Suppose $(a_n)_{n\ge1}\subseteq\mathbb C$ satisfies
\[
 \operatorname{trdeg}_K K(a_1,\ldots,a_N)\ge N-D(N).
\]
Let $L=K(c_1,\ldots,c_s)$, $d=\operatorname{trdeg}_K L$, and let
$P(T,X_0,\ldots,X_r)\in L[T,X_0,\ldots,X_r]$ be nonzero, with
positive degree in at least one $X_j$. Let $e$ be the number of
positive integers $n$ for which $P(n,X)$ is the zero polynomial;
then $e\le\deg_TP$. The number $Z_P(N)$ of indices $1\le n\le N-r$
satisfying $P(n,a_n,\ldots,a_{n+r})=0$ obeys
\begin{equation}\label{sc:dyn:window-budget}
 Z_P(N)\le(r+1)(D(N)+d)+e.
\end{equation}
\end{theorem}
\begin{proof}
Omit the identically specialized equations. Greedily select disjoint
satisfying intervals $[n,n+r]$ in increasing order of their left
endpoint. Each selected interval excludes at most $r+1$ possible
starts. If $M$ intervals are selected, then
$Z_P(N)-e\le(r+1)M$.
Each selected block has a nonzero algebraic relation over $L$, so its
$r+1$ entries contribute transcendence degree at most $r$. Taking the
union of disjoint blocks and adjoining all uncovered entries yields
\[
 \operatorname{trdeg}_L L(a_1,\ldots,a_N)\le N-M.
\]
By the tower formula and monotonicity,
\[
 N-D(N)\le\operatorname{trdeg}_K K(a_1,\ldots,a_N)
       \le d+N-M.
\]
Thus $M\le D(N)+d$, proving \eqref{sc:dyn:window-budget}.
\end{proof}
Rational-function coefficients in $T$ are covered after clearing
denominators and excluding their finitely many poles. Arbitrary
complex constants are covered because a single polynomial uses only
finitely many of them.

For $I_n$ and $E_1(n)$ take $D(N)=\pi(N)+1$; for $A_n$ take
$D(N)=\pi(2N-1)$. In either case every fixed finite-order algebraic
recurrence over $\mathbb C(n)$ can hold only on a density-zero set of
windows, and cannot hold eventually. This excludes linear
P-recursiveness and nonlinear polynomial/rational-in-$n$ templates.
For the completely independent sequence $\operatorname{Ein}(n)$, or
for the positive composite residuals enumerated increasingly, take
$D(N)=0$: any fixed such relation can hold on only finitely many
windows. This is an exact algebraic exclusion, not a test of approximate
numerical fits.

If a polynomial relation fits every window through $N$, the budget gives
\[
 N-r\le(r+1)(D(N)+d)+e,
 \qquad r+1\ge\frac{N+1-e}{D(N)+d+1}.
\]
A genuine forward algebraic generation rule obeys a stronger constraint.
If its first $N$ outputs are algebraic over a field generated over $K$
by $d$ independent parameters and $r$ initial seeds, then necessarily
\[
 r+d\ge N-D(N).
\]
A step at which its specialized defining polynomial vanishes
identically introduces an additional free seed, which must be counted.
These results do not exclude growing-order recurrences, new special
values introduced at each step, nonalgebraic coefficient sequences, or
the integral representations already available. They exclude an entire
class of fixed algebraic compression methods, without deciding the
arithmetic nature of $\gamma$.


\section{Nonlinear fresh-prime descent}
\label{sc:sec:nonlinear-imports}

Set $E=e^z$ and, for $r\in\overline{\mathbb Q}$, set
\[
 F_r=e^z(\operatorname{Ein}(z)-r)
     =\sum_{n\ge0}(H_n-r)\frac{z^n}{n!}.
\]
Then $E'=E$ and $F_r'=F_r+(E-1)/z$. Consequently
\[
 \mathcal A_r=\overline{\mathbb Q}(z)[E,E^{-1},F_r]
\]
is a differential algebra. It includes fixed polynomial nonlinear
combinations of all derivatives of its generators.
For a germ $G=\sum g_nz^n/n!$, write $[G]_n=g_n$.
Integrality of normalized coefficients will always mean integrality
at every prime above a specified rational prime, in one number field.

\subsection{The fresh-prime derivation}

\begin{lemma}\label{sc:nl:fresh}
Let $P\in K[z,T,T^{-1},Y]$, where $K$ contains $r$, and put
$G=P(z,E,F_r)$. For every fixed $m\ge0$ and every sufficiently large
rational prime $p$, at every prime $\mathfrak p$ of $K$ above $p$,
\[
 p[G]_{p+m}\equiv[E P_Y(z,E,F_r)]_m\pmod{\mathfrak p}.
\]
The same formula holds for
$P\in K[z,T,T^{-1},Y,V]$ evaluated at $V=J(z)$ if the normalized
coefficients $[J]_0,\ldots,[J]_{p+m}$ are integral at $\mathfrak p$.
There is then no derivative term in $V$.
\end{lemma}
\begin{proof}
It suffices to check $z^aE^bF_r^cJ^d$. All fixed algebraic
coefficients are integral after excluding finitely many primes.
Products of exponential generating series use integral multinomial
coefficients. Since $p+m<2p$, at most one $F_r$ factor can have
index at least $p$. Terms without such a factor are integral and
disappear after multiplication by $p$. For its index $p+k$,
\[
 p(H_{p+k}-r)\equiv1\pmod{\mathfrak p}.
\]
Lucas's congruence reduces the remaining multinomial coefficient to
the one of total index $m$, with the large index replaced by $k$.
There are $c$ choices of the large factor, giving $E\cdot cF_r^{c-1}$.
The remaining $E^b$ and $J^d$ factors have only low indices in these
terms. Multiplication by $z^a$ introduces a falling factorial. For
$a\le m$ it reduces to $m^{\underline a}$; for $a>m$ both sides vanish.
This proves the assertion. In the case with $J$, no factor of $J$
can supply a denominator, by the stated integrality hypothesis.
\end{proof}

The polynomial and its degree are fixed independently of $p$.
Choosing a different polynomial such as $F_r^p$ for each tested prime
is not a uniform application of this lemma: its derivative coefficient
vanishes at that particular prime.

\begin{lemma}\label{sc:nl:independence}
The functions $E$ and $F_r$ are algebraically independent over
$\overline{\mathbb Q}(z)$.
\end{lemma}
\begin{proof}
The exponential $E$ is transcendental over $\overline{\mathbb Q}(z)$.
Put $K_0=\overline{\mathbb Q}(z,E)$, with $DE=E$.
If $\operatorname{Ein}(z)$ were algebraic over $K_0$, averaging its
trace over that finite extension would yield $v\in K_0$ with
\[
 Dv=(1-E^{-1})/z.
\]
In the Laurent ring $\overline{\mathbb Q}(z)[E,E^{-1}]$, a nonunit
irreducible factor $B$ cannot divide $DB$: comparing two of its
nonzero Laurent coefficients would give a rational function with
nonzero constant logarithmic derivative. Thus any nonunit denominator
factor of $v$ would acquire a higher pole under differentiation, which
the Laurent-polynomial right side does not have. Hence $v$ is a Laurent
polynomial in $E$. Its coefficient of $E^0$ would be a rational
primitive of $1/z$, impossible by the residue at zero.
Since $F_r=E(\operatorname{Ein}-r)$, the assertion follows.
\end{proof}

\begin{theorem}[Algebraic descent from finite prime support]
\label{sc:nl:descent}
Let $K_0=\overline{\mathbb Q}(z,e^z)$ and let
$J(z)=\sum_{n\ge0}j_nz^n/n!$ be a germ regular at zero, algebraic over
$K_0(F_r)$. Assume all $j_n$ belong to one number field and are
integral outside a fixed finite set of rational primes. Then $J$ is
algebraic over $K_0$. In particular, if $J\in K_0(F_r)$, then $J\in K_0$.

More generally, retain regularity, algebraicity over $K_0(F_r)$, and
the one-number-field coefficient hypothesis, but drop the finite-prime
support assumption. If $J$ is not algebraic over $K_0$, there is a fixed
integer $m\ge0$ such that for every sufficiently large rational prime $p$,
at least one of $j_0,\ldots,j_{p+m}$ is nonintegral at a prime above $p$.
\end{theorem}
\begin{proof}
By Lemma~\ref{sc:nl:independence}, evaluation
\[
 K_0[Y,V]\longrightarrow K_0(F_r,J),\qquad(Y,V)\longmapsto(F_r,J)
\]
has a height-one prime kernel generated by an irreducible polynomial
$P(Y,V)$. Clear its $K_0$ coefficient denominators; the multiplier is
independent of $Y,V$ and is a unit in $K_0$. Irreducibility and the
kernel-generator assertion are used over $K_0[Y,V]$, not over the
polynomial coefficient ring after clearing. All coefficients then lie in
$K[z,T,T^{-1},Y,V]$ for a number field $K$, where $T=E$.
The algebraic extension of the field containing the $j_n$ preserves
integrality outside the original finite set.

Apply Lemma~\ref{sc:nl:fresh} to the identically zero series
$P(F_r,J)$. For each fixed $m$ the algebraic number
$[E P_Y(F_r,J)]_m$ reduces to zero at every sufficiently large prime.
It is therefore zero, since a nonzero algebraic number has only
finitely many prime divisors. Thus $P_Y(F_r,J)=0$.
The kernel description implies that $P$ divides $P_Y$. If $P_Y\ne0$,
its degree in $Y$ is smaller than that of $P$, which is impossible.
Hence $P_Y=0$ and $P\in K_0[V]$, proving the first assertion.
The rational case uses the fact that $K_0$ is relatively algebraically
closed in the pure transcendental field $K_0(F_r)$.

For the last assertion, if $J$ is not algebraic over $K_0$, then
$E P_Y(F_r,J)$ is nonzero. Choose a nonzero coefficient at a fixed
index $m$. At every sufficiently large prime it is a unit. If all
$j_0,\ldots,j_{p+m}$ were integral there, the finite version of
Lemma~\ref{sc:nl:fresh} would reduce that unit to zero. This contradiction
proves the asserted denominator appearance.
\end{proof}

This is a statement about an algebraic direction in a differential
field, not the arithmetic nature of gamma. Under the hypothesis that
gamma is algebraic one may set $r=\gamma$, consistently with the theorem.
If a construction inserts gamma into its algebraic coefficients or
branch-selection rule, the theorem does not erase that external input.

\subsection{Complete classification in the polynomial differential algebra}

\begin{theorem}\label{sc:nl:classification}
Let $\mathcal E_{\rm gb}$ be the ring of globally bounded $E$-functions,
where global boundedness concerns normalized Taylor coefficients.
Then
\[
 \mathcal A_r\cap\mathcal E_{\rm gb}
 =\overline{\mathbb Q}[z,e^z,e^{-z}].
\]
More generally, for $G\in\mathcal A_r$ regular at zero, integrality of
normalized coefficients outside finitely many primes already forces
$G\in\overline{\mathbb Q}(z)[e^z,e^{-z}]$.
\end{theorem}
\begin{proof}
Clear rational $z$ denominators and write $qG=P(z,E,F_r)$.
Multiplication by the fixed polynomial $q$ adds only finitely many
denominator primes to normalized coefficients.
Lemma~\ref{sc:nl:fresh} implies $E P_Y(z,E,F_r)=0$, and
Lemma~\ref{sc:nl:independence} gives $P_Y=0$ formally. Hence the displayed
general assertion holds, with Laurent-polynomial dependence on $E$.

Now suppose $G$ is a globally bounded $E$-function, hence entire.
Write $G=\sum_j a_j(z)e^{jz}$ with rational functions $a_j$.
At a nonzero finite algebraic pole $\alpha$, cancellation of the
highest principal parts would give
$\sum_j c_j e^{j\alpha}=0$ with algebraic $c_j$.
Lindemann--Weierstrass forces all $c_j=0$. Repeating at each possible
pole shows that each $a_j$ is a Laurent polynomial in $z$.

Suppose a negative power remains, and choose the maximal $M\ge1$:
\[
 G=\sum_{a=-M}^D z^a Q_a(e^z),\qquad
 Q_{-M}(T)=\sum_j c_jT^j\ne0.
\]
For fixed $k\ge1$ and large prime $p$, put $n=kp-M$.
Only the highest negative power has a factor $p$ in its denominator
$(n+1)\cdots(n+M)$; all smaller negative powers have integral
coefficients there. Fermat's congruence, including negative integer
frequencies, gives
\[
 p[G]_{kp-M}\equiv
 \frac{\sum_j c_jj^k}{k(-1)^{M-1}(M-1)!}
 \pmod{\mathfrak p}.
\]
Thus $\sum_jc_jj^k=0$ for every $k\ge1$.
Vandermonde invertibility at the distinct nonzero integer frequencies
forces $c_j=0$ for $j\ne0$. The remaining term $c_0z^{-M}$ cannot be
canceled at zero by any lower negative power, contradicting entire
regularity. Hence no negative power remains.
The reverse inclusion follows directly: polynomial multiples of
$e^{jz}$, $j\in\mathbb Z$, have globally bounded normalized coefficients.
\end{proof}

For example, $(e^z-1)^2/z$ has normalized coefficients
$(2^{n+1}-2)/(n+1)$. The apparent cancellation at $n=p-1$ is misleading:
at $n=2p-1$ every odd prime gives $p[G]_{2p-1}\equiv1\pmod p$.
The second prime block detects the surviving denominator.

Finite prime support alone does not imply that a rational expression
in $E$ is an $E$-function. For example $1/(2-e^z)$ has integral
normalized coefficients $\sum_{k=0}^n k!S(n,k)$, obtained by expanding
the geometric series in $e^z-1$, but has poles at
$\log2+2\pi i\mathbb Z$. This distinguishes the weak arithmetic
hypothesis in Theorem~\ref{sc:nl:descent} from the entire-function
hypothesis used in Theorem~\ref{sc:nl:classification}.

\subsection{Import boundary and remaining scope}

The convention for global boundedness above is exactly that used in
the introduction of Fischler--Rivoal~\cite{FRzeros}.
Their Theorem~1.5 uses this condition in a zero-factorization result;
Theorem~1.2 and Corollary~1.3 concern a different entire-quotient theorem.
The present classification concerns our preprocessing algebra, and
does not claim to generalize those theorems or establish historical
priority.

The descent theorem goes beyond fixed linear operators to rational
and algebraic operations. It still does not cover arbitrary new
transcendental functions, varying-degree constructions with uniform
height control, several independently normalized connection constants,
or a direct arithmetic theorem for mixed Gevrey values. None of the
results here decides the algebraic or transcendental world for gamma.

Exact finite checks are supplied in \texttt{check\_nonlinear\_imports.py}:
288 nonlinear fresh-prime congruences and seven second-block Laurent
checks. These are checks of the formulas, not substitutes for the proofs.


\section{Sharp decay and a two-point inverse moment theorem}
\label{sc:sec:sharp-christoffel}

The previous Hankel construction gave positive upper approximants to
gamma.  Here its exponential rate is determined exactly, and a second
Christoffel trace is shown to recover all the moment information lost by
a single trace.  The arguments use classical polynomial inequalities,
orthogonal polynomials, and Jacobi matrices; no priority claim is made
for these methods or their consequences.

\subsection{The exact exponential rate}

Fix $a,h>0$, put $q=e^{-h}$, and let
\[
 d\mu(x)=\frac{x^{a/h-1}}{-\log x}\,dx\quad(0<x<q),\qquad
 \lambda_n^\mu(z)=\min_{\substack{\deg p<n\\p(z)=1}}
                         \int |p(x)|^2\,d\mu(x).
\]
For $z>q$ set
\[
 \xi_q(z)=\frac{2z}{q}-1,\qquad
 \rho_q(z)=\xi_q(z)+\sqrt{\xi_q(z)^2-1}.
\]

\begin{theorem}[Sharp root asymptotic]\label{sc:sc:root-rate}
For every fixed $z>q$,
\begin{equation}\label{sc:sc:root}
 \lim_{n\to\infty}\lambda_n^\mu(z)^{1/n}=\rho_q(z)^{-2}.
\end{equation}
In particular the earlier approximants $A_n=\gamma+\varepsilon_n$
satisfy
\[
 \lim_{n\to\infty}-\frac1n\log\varepsilon_n
 =2\log\rho_q(1)=4\operatorname{arcosh}(e^{h/2}).
\]
The exponent is independent of $a$.  This is a fixed-parameter limit;
it is not asserted uniformly when $a$ or $h$ varies with $n$.
\end{theorem}
\begin{proof}
The normalized Chebyshev polynomial on $[0,q]$ gives
\[
 \lambda_n^\mu(z)\le
 \frac{E_1(a)}{T_{n-1}(\xi_q(z))^2},
\]
so the upper root bound in \eqref{sc:sc:root} follows from
$T_m(\xi)=(\rho^m+\rho^{-m})/2$.

For the lower bound fix $0<\delta<q$.  The density has a positive
minimum $w_\delta$ on $[\delta,q]$, whose length is $L=q-\delta$.
Let $p$ be real, of degree at most $m=n-1$, with $p(z)=1$, and put
$M=\max_{[\delta,q]}|p|$.  Chebyshev's exterior extremal inequality gives
\[
 M\ge\frac1{T_m(\xi_\delta(z))},\qquad
 \xi_\delta(z)=\frac{2z-q-\delta}{q-\delta}.
\]
For $m\ge1$, Markov's inequality gives
$\max|p'|\le2m^2M/L$.  From a point where $|p|=M$ there is a one-sided
interval of length $L/(4m^2)$ contained in $[\delta,q]$ on which
$|p|\ge M/2$.  Therefore
\[
 \int |p|^2\,d\mu\ge
 \frac{w_\delta L}{16m^2T_m(\xi_\delta(z))^2}.
\]
Complex coefficients cannot improve the minimum: taking the real part
preserves the constraint and can only decrease the objective.  Taking
$n$th roots yields a lower limit at least
$(\xi_\delta(z)+\sqrt{\xi_\delta(z)^2-1})^{-2}$.
Let $\delta\downarrow0$ to finish the proof.
\end{proof}

\begin{corollary}[The second trace has its own sharp rate]
\label{sc:sc:second-rate}
Let $\nu=\mu+\gamma\delta_1$, and let $d>1$. Then
\[
 \lim_{n\to\infty}\lambda_n^\nu(d)^{1/n}=\rho_q(d)^{-2}.
\]
Moreover $0<\lambda_{n+1}^\nu(d)<d^{-2}\lambda_n^\nu(d)$.
For $n\ge2$ one has the explicit upper bound
\[
 \lambda_n^\nu(d)\le
 \frac{E_1(a)}{(d-1)^2T_{n-2}(\xi_q(d))^2}.
\]
\end{corollary}
\begin{proof}
Adding the positive atom gives $\lambda_n^\nu(d)\ge\lambda_n^\mu(d)$.
If $p$ is its minimizing polynomial, $xp/d$ is feasible at the next
level. Since $x^2\le1$ on the support, with strict inequality on the
continuous part, this proves the strict contraction.
For the reverse root bound use
\[
 p(x)=\frac{x-1}{d-1}
       \frac{T_{n-2}(2x/q-1)}{T_{n-2}(2d/q-1)}.
\]
It has degree at most $n-1$, equals one at $d$, and vanishes at the
atom. Since $|x-1|\le1$ on $[0,q]$, it gives the stated bound.
\end{proof}

Only the exponential rate has been proved here.  A full leading
constant or a ratio asymptotic is not inferred from a root limit.
The density is also in the classical Szeg\H{o} class after an affine
change to $[-1,1]$, but no unverified normalization of a stronger
Szeg\H{o} asymptotic is needed for these conclusions.

The root limit also proves the exact convergence radius
\[
R\!\left(\sum_{n\ge1}(A_n-\gamma)z^{n-1}\right)
=\rho_{e^{-1}}(1)^2
\]
for $a=h=1$, by Cauchy--Hadamard. The series remains differentially
transcendental by the coefficient-field argument above.

\subsection{Two exterior traces determine the measure}

Let $\sigma$ be a finite positive measure with infinite support contained
in a compact interval $[L,U]$.  Let $p_n$ be its orthonormal polynomials
with positive leading coefficient, and let
\[
 K_n(t,t)=\sum_{j=0}^{n-1}p_j(t)^2,
 \qquad \lambda_n(t)=K_n(t,t)^{-1},\qquad K_0(t,t)=0.
\]
The convention uses polynomials of degree strictly below $n$.

\begin{theorem}[Two-trace reconstruction]\label{sc:sc:two-trace}
Fix $U\le c<d$.  The two sequences
$(\lambda_n(c))_{n\ge1}$ and $(\lambda_n(d))_{n\ge1}$ uniquely determine
$\sigma$.  The first $N$ entries of both traces determine its moments
through degree $2N-2$ using arithmetic operations and positive square
roots.  No rational-field equality is asserted.
\end{theorem}
\begin{proof}
Every zero of every $p_n$ lies strictly between the extreme points of
the support. Thus $p_n(c),p_n(d)>0$, even when $c=U$.  The data recover
these values without a sign ambiguity:
\begin{equation}\label{sc:sc:square-root}
 p_n(t)=\sqrt{\lambda_{n+1}(t)^{-1}-\lambda_n(t)^{-1}},
 \quad t\in\{c,d\},\quad n\ge0,
\end{equation}
where $\lambda_0(t)^{-1}=0$.

Write the Jacobi recurrence as
\[
 a_{n+1}p_{n+1}(t)+b_np_n(t)
       =tp_n(t)-a_np_{n-1}(t),\qquad a_{n+1}>0,
\]
with $a_0p_{-1}=0$.  Suppose $a_n$ is already known. Evaluating at
$c,d$ gives a two-by-two linear system for $a_{n+1},b_n$.
Its determinant is
\[
 p_n(c)p_n(d)
 \left(\frac{p_{n+1}(c)}{p_n(c)}-
       \frac{p_{n+1}(d)}{p_n(d)}\right)<0.
\]
Indeed each ratio $r_n(t)=p_n(t)/p_{n-1}(t)$ is positive and strictly
increasing for $t\ge U$.  Induction follows from
\[
 a_nr_n(t)=t-b_{n-1}-\frac{a_{n-1}}{r_{n-1}(t)},\qquad
 a_nr_n'(t)=1+\frac{a_{n-1}r_{n-1}'(t)}{r_{n-1}(t)^2}>0,
\]
with $r_1'(t)=1/a_1$. Thus every reconstruction system is nonsingular.

The mass $M=\sigma(\mathbb R)=\lambda_1(c)=\lambda_1(d)$ is known.
The first $N$ trace entries give
$a_1,\ldots,a_{N-1}$ and $b_0,\ldots,b_{N-2}$.  Complete the
$N\times N$ truncated Jacobi matrix by assigning any value to the last
diagonal entry.  For $0\le k\le2N-2$ its moments are
\[
 \int x^k\,d\sigma(x)=M\,(J_N^k)_{00}.
\]
A walk starting and ending at index zero cannot use the last diagonal
entry in fewer than $2N-1$ steps, or reach the next omitted index in
fewer than $2N$ steps.  This proves the finite assertion.  All moments
determine a compactly supported measure uniquely by polynomial density
in continuous functions, proving uniqueness of $\sigma$.
\end{proof}

The endpoint allowance is essential for the gamma application:
$\nu=\mu+\gamma\delta_1$ has support bounded above by $1$ and
$\lambda_n^\nu(1)=A_n$.

\begin{proposition}[An admissibility criterion]\label{sc:sc:admissible}
Fix $L<U\le c<d$. Two positive proposed traces arise from a positive
infinite-support measure on $[L,U]$ if and only if all of the following
hold: their first entries agree; their reciprocal increments in
\eqref{sc:sc:square-root} are positive; every reconstruction system is
nonsingular and gives $a_n>0$, $b_n\in\mathbb R$; and the reconstructed
Jacobi matrix satisfies $LI\le J\le UI$ on finitely supported vectors.
The resulting measure is unique.
\end{proposition}
\begin{proof}
Necessity follows from the previous proof and the representation of
$J$ as multiplication by $x$ on the polynomial subspace of $L^2(\sigma)$.
Conversely the quadratic-form bounds produce a bounded selfadjoint
Jacobi operator with spectrum in $[L,U]$. Its spectral measure at the
first coordinate, multiplied by the common prescribed mass, has infinite
support because every $a_n$ is positive. Its orthonormal polynomials
satisfy the reconstructed recurrence. Induction makes their values at
$c,d$ equal the positive roots prescribed by \eqref{sc:sc:square-root}.
Summing their squares recovers the reciprocal traces by telescoping.
\end{proof}
This is a countable exact characterization, not a finite stopping
algorithm for arbitrary real input sequences.

\subsection{One trace is provably insufficient}

\begin{proposition}[An explicit continuum of identical traces]
\label{sc:sc:one-trace-counterexample}
There is a continuum of distinct probability measures, all with infinite
support contained in $[5/16,319/400]$, for which
\[
 \lambda_N(1)=\frac{15}{16^N-1}\qquad(N\ge1).
\]
\end{proposition}
\begin{proof}
For $9/100\le t\le11/100$, take the bounded Jacobi matrix with
\[
 a_n=t\ (n\ge1),\qquad
 b_0=1-4t,\qquad b_n=1-\frac{17t}{4}\ (n\ge1).
\]
The diagonal-minus/plus-row-radius bounds put its spectrum in
$[5/16,319/400]$.  Its first-coordinate spectral measure is a probability
measure with infinite support.  The recurrence, starting with $p_0=1$,
gives $p_n(1)=4^n$ for every $n$, because
$4t+(1-17t/4)+t/4=1$.  Hence the kernel is
$\sum_{j=0}^{N-1}16^j=(16^N-1)/15$.  The measures are distinct since
their means are $b_0=1-4t$.
\end{proof}

\subsection{The second trace preserves the full arithmetic rank}

Specialize $a=h=1$ and set
\[
 m_k=\operatorname{Ein}(k+1)-\log(k+1),\quad
 F_n=(m_{i+j})_{0\le i,j<n},\quad
 v_n(z)=(1,z,\ldots,z^{n-1})^{\mathsf T}.
\]
These are moments of $\nu=\mu_{1,1}+\gamma\delta_1$. Define
\[
 A_n=\frac1{v_n(1)^{\mathsf T}F_n^{-1}v_n(1)},\qquad
 B_n=\frac1{v_n(2)^{\mathsf T}F_n^{-1}v_n(2)}.
\]
The duplicate first value is $A_1=B_1=m_0$.
Let $K_{\exp}=\overline{\mathbb Q}(e^\alpha:\alpha\in
\overline{\mathbb Q})$ and retain the explicitly stated companion input
that distinct nonzero algebraic arguments give algebraically independent
$\operatorname{Ein}$ values over this field.

\begin{theorem}[Lossless rank transfer through two traces]
\label{sc:sc:arithmetic-rank}
For every $N\ge1$,
\begin{align*}
 &\operatorname{trdeg}_{K_{\exp}}
 K_{\exp}(A_1,\ldots,A_N,B_2,\ldots,B_N)\\
 &\qquad=\operatorname{trdeg}_{K_{\exp}}
 K_{\exp}(m_0,\ldots,m_{2N-2})
 \ge2N-1-\pi(2N-1).
\end{align*}
In the ordered list $A_1,A_2,B_2,A_3,B_3,\ldots$, there exists a set of
indices of natural density zero whose deletion leaves an algebraically
independent family over $K_{\exp}$.
\end{theorem}
\begin{proof}
The trace field is a subfield of the moment field because its entries
are rational expressions in the moments. By Theorem~\ref{sc:sc:two-trace}
with $c=1,d=2$, the moments are algebraic over the trace field: the
reconstruction uses positive square roots followed by rational
operations. Thus their transcendence degrees agree. Adjoining the
$\pi(2N-1)$ prime logarithms $\log p$, $p\le2N-1$, to the moment
field recovers the $2N-1$ independent values
$\operatorname{Ein}(1),\ldots,\operatorname{Ein}(2N-1)$.
This proves the bound. Greedy retention of an algebraically independent
subfamily skips at most $\pi(2N-1)=o(N)$ entries in each completed block;
an incomplete block differs by only one entry.
\end{proof}

This is stronger than the single-trace dimension bound: no $N-1$
odd-moment allowance is lost. The two limits nevertheless satisfy
$A_n\downarrow\gamma$ and $B_n\downarrow0$; independence of these finite
data does not establish the arithmetic nature of either limit.

\subsection{Parameter tuning and the arithmetic cost that is actually proved}

At fixed $h$, the gamma error has exponent
$c(h)=4\operatorname{arcosh}(e^{h/2})$. It obeys
\[
 c(h)\sim4\sqrt h\quad(h\downarrow0),\qquad
 c(h)=2h+4\log2+O(e^{-h})\quad(h\to\infty).
\]
Larger spacing therefore increases the limiting exponent per Hankel
dimension, but
also increases the largest input argument $a+(2n-2)h$. No uniform
fixed-parameter asymptotic is being applied to moving $h_n$.

There is a rigorous algebraic information cost: any collection of
constants from which all $2N-1$ nonduplicate trace entries can be
obtained by algebraic operations over $K_{\exp}$ must generate a field
of transcendence degree at least $2N-1-\pi(2N-1)$. A fixed finite list of
auxiliary constants cannot encode this entire tower algebraically.
This is an algebraic-rank statement, not a lower bound for the
Diophantine heights of subsequent rational approximations.

For comparison, if independently supplied rational moments have common
denominator $Q$ and absolute value at most $C\ge1$, their $A_n$, when its denominator is nonzero, is a
rational number with an unreduced representation whose numerator and
denominator have absolute values at most
\[
 n! (QC)^n,\qquad Q\,n^2(n-1)!(QC)^{n-1},
\]
respectively. This follows by applying the Leibniz bound to the integer
matrix $QF_n$ and its cofactors. It is an upper bound for one
rationalization procedure, with possible further gcd cancellation;
it neither proves a height obstruction nor supplies the rational
moments or their required error tolerance. Obtaining simultaneous
height and error bounds remains the missing gamma bridge.

\paragraph{References and scope.}
For the classical variational principle, Jacobi recurrence, and kernel
machinery see B.~Simon, \emph{The Christoffel--Darboux Kernel},
\url{https://arxiv.org/abs/0806.1528}.  The Szeg\H{o}-class comparison is
described in G.~Szeg\H{o}, \emph{Orthogonal Polynomials}, fourth edition,
Chapter~12, and in the NIST DLMF,
\url{https://dlmf.nist.gov/18.2#xi}.  The exact rate proof above is
self-contained apart from the classical Chebyshev and Markov
inequalities; the reconstruction and counterexample are proved in full.
No named external open problem, or the arithmetic nature of gamma, is
claimed to be settled.


\section{Global camouflage densities and complete local denominators}
\label{sc:sec:global-arithmetic}

Keep
$J_N(x)=\min\{m\ge1:m\binom xn\in\mathbb Z[x],\ 0\le n\le N\}$
and $U(x)=\{p:1/p\in\mathbb Z[\binom xn:n\ge0]\}$.
Define the simultaneous camouflage primes by
\[
\mathcal C(x)=\{p:p\notin U(x),\quad p\nmid p!/J_p(x)\}.
\]
For transcendental $x$ this is the set of all primes.  For rational
$x$ it is empty, because $J_p(x)=1$.  The prime-unit description
used below is prior work of
\href{https://arxiv.org/abs/2010.14494}{Kalinov--Mandelshtam}; the
Frobenius interpretation is the standard local factorization theorem.
No literature-wide priority is claimed for the following deductions.

\begin{theorem}[Sharp joint density bound]\label{sc:thm:joint-density}
Let $x$ be algebraic of degree $d\ge2$, and let $G$ be the group
of the normal closure, in its faithful transitive action on the
$d$ roots.  Put $g=|G|$ and let $\delta(G)$ be the proportion of
derangements.  Then $\mathcal C(x)$ has natural density
\[
 c(x)=1-\delta(G)-\frac1g\le\frac{d-2}{d-1}.
\]
Equality holds if and only if the action is sharply $2$-transitive.
\end{theorem}
\begin{proof}
Outside a finite prime set, absence of a prime inverse means that
Frobenius has a fixed point, while absence of a $p$-defect at index
$p$ means that Frobenius is not the identity.  Chebotarev gives the
stated density.

For the bound, let $F(s)$ be the number of fixed points of $s\in G$.
The case $d=2$ is immediate.  For $d\ge3$, Burnside's lemma gives
\[
 \sum_sF(s)=g,\qquad
 r:=g^{-1}\sum_sF(s)(F(s)-1)\ge1,
\]
where $r$ counts the orbits on ordered distinct pairs.  The identity
contributes $d$ to the first sum, so $c\le1-d/g$.  If
$g\le d(d-1)$ this is at most $(d-2)/(d-1)$, strictly so for
$g<d(d-1)$.

A nonidentity permutation has at most $d-2$ fixed points.  Hence
\[
 r\le(d-2)\delta(G)+\frac{2(d-1)}g
   =(d-2)(1-c)+\frac dg.
\]
For $g\ge d(d-1)$ this yields
\[
 c\le\frac{d-3+d/g}{d-2}\le\frac{d-2}{d-1},
\]
strictly if $g>d(d-1)$.  Equality therefore requires $g=d(d-1)$
and equality in the first fixed-point estimate.  Thus every
nonidentity with a fixed point fixes exactly one point.  Every
two-point stabilizer is trivial, and its orbit has $g=d(d-1)$
elements; the action is sharply $2$-transitive.  Conversely, a
sharply $2$-transitive action has $d(d-2)$ nonidentity elements
with a fixed point, giving the stated equality.
\end{proof}

This is a direct application of the classical fixed-point moment
method; compare the Cameron--Cohen bound discussed in
\href{https://arxiv.org/abs/1403.7666}{Burness--Tong-Viet}.
For every prime $q$, the splitting field of $X^q-2$ has the affine
group $\operatorname{AGL}_1(\mathbb F_q)$, so the bound is sharp
for actual algebraic numbers in infinitely many degrees.  Indeed,
for odd $q$ the polynomial remains Eisenstein at a prime above $2$
in $\mathbb Q(\zeta_q)$, which is unramified at $2$.  Its splitting
field therefore has degree $q(q-1)$ and the full affine action.

\begin{center}
\begin{tabular}{c|c|c|c|c}
$G$ & $d$ & $\delta(G)$ & $1/g$ & $c$\\\hline
$C_2$ & $2$ & $1/2$ & $1/2$ & $0$\\
$C_3$ & $3$ & $2/3$ & $1/3$ & $0$\\
$S_3$ & $3$ & $1/3$ & $1/6$ & $1/2$\\
$A_4$ & $4$ & $1/4$ & $1/12$ & $2/3$\\
$S_4$ & $4$ & $3/8$ & $1/24$ & $7/12$
\end{tabular}
\end{center}

The preceding cubic camouflage construction has splitting type
$(1,2)$ at a tested prime, so its irreducible cubic has group $S_3$.
It can imitate every prescribed finite initial observation set,
but its global camouflage density is exactly $1/2$.  There is no
contradiction: the first visible discrepancy depends on arithmetic
height.  The density $c(x)$ exists for every complex $x$, and
\[
 x\notin\overline{\mathbb Q}\quad\Longleftrightarrow\quad c(x)=1.
\]
This separates rational values too, unlike prime-unit density alone.
It does not compute $c(\gamma)$ from real approximations.

\subsection{Complete exponents at maximal-order primes}

Kalinov--Mandelshtam already use the local valuation method in
Section 5.3, Lemmas 5.12--5.14.  Here it gives an explicit common
denominator formula when the local order is maximal.

\begin{theorem}[Exact local denominator formula]\label{sc:thm:local-full}
Let $x$ be an algebraic integer with monic minimal polynomial $f$
and $K=\mathbb Q(x)$.  Suppose $p\nmid[\mathcal O_K:\mathbb Z[x]]$,
and factor
$\bar f=\prod_i\varphi_i^{e_i}$ over $\mathbb F_p$, with
$f_i=\deg\varphi_i$.  For all $N\ge0$,
\[
v_p(J_N(x))=
\begin{cases}
 v_p(N!),&\text{some }f_i>1,\\
 0,&\text{all }f_i=e_i=1,\\
 v_p(N!)-\lfloor N/(pe)\rfloor,
 &\text{all }f_i=1,\ e=\max_i e_i\ge2.
\end{cases}
\]
This includes wild ramification.  If $\mathbb Z[x]=\mathcal O_K$,
it gives the exact denominator at every prime.
\end{theorem}
\begin{proof}
The index assumption identifies $\mathbb Z_p[x]$ with the product
of the full local integer rings.  The numbers $f_i,e_i$ are their
residue degrees and ramification indices.  If some $f_i>1$, the
reduction of $x$ in that residue field is not in $\mathbb F_p$.
Every $x-j$ is a unit, so the corresponding valuation of
$\binom xn$, normalized by $v(p)=1$, is $-v_p(n!)$.  This gives
the first case.  A component with $f_i=e_i=1$ is $\mathbb Q_p$,
where $x\in\mathbb Z_p$ and all binomial values are integral.

Consider a component with $f_i=1$ and $e_i=e\ge2$.  Let $a$ be
an integer lifting its residue root.  Its integer ring is
$\mathbb Z_p[x_i]$, whose maximal ideal is $(p,x_i-a)$: the
quotient by those generators is $\mathbb F_p$.  In the discrete
valuation normalization $v_{\mathfrak p}(p)=e$, this implies
$v_{\mathfrak p}(x_i-a)=1$.  Put $\pi=x_i-a$.

The common denominator through index $N$ is unchanged by integer
translation of its argument.  The binomial Vandermonde formulas
for the shifts $a$ and $-a$ give inverse triangular integer
transformations.  We may thus use $\pi$ in this component.  Since
$v(\pi)=1/e<1$,
\[
 v(\pi-j)=\begin{cases}1/e,&p\mid j,\\0,&p\nmid j,\end{cases}
 \qquad
 -v\binom\pi n=v_p(n!)-\frac1e\lceil n/p\rceil.
\]
Within the block $pt\le n<p(t+1)$ the maximum is at $pt$.
The successive block maxima increase by $v_p(pt)-1/e>0$.
Writing $k=\lfloor N/p\rfloor$, the maximum through $N$ is
$v_p(N!)-k/e$.  Its ceiling is
$v_p(N!)-\lfloor N/(pe)\rfloor$.  Taking the maximum over local
components selects the largest ramification index $e$.
\end{proof}

In particular, for a root of an Eisenstein polynomial at $p$ of
degree $d\ge2$, its local order is maximal and
\[
 v_p(J_N)=v_p(N!)-\lfloor N/(pd)\rfloor.
\]
For $\sqrt2$ this gives the full exponent
$v_2(N!)-\lfloor N/4\rfloor$.  For a root of $X^3-X-1$, whose
order is maximal, the ramified contribution at $23$ is
$v_{23}(N!)-\lfloor N/46\rfloor$.

\begin{example}[Why the index hypothesis is necessary]
For $x=\sqrt{17}$, the order has index $2$ and both local roots
are in $\mathbb Z_2$.  Thus all binomial values are integral in
the maximal order, and the index gives $v_2(J_N)\le1$.
But $\binom x2=(17-x)/2\notin\mathbb Z_2[x]$, so
$v_2(J_N)=1$ for every $N\ge2$.  Applying the excluded formula
only to $\bar f=(X+1)^2$ would give a linearly growing answer.
In general the maximal-order denominator $D_N$ satisfies
\[
 v_p(D_N)\le v_p(J_N)
 \le v_p(D_N)+v_p([\mathcal O_K:\mathbb Z[x]]).
\]
\end{example}

\subsection{Finite detection with supplied degree and height}

Write $\nu(x)=\min\{N\ge2:J_N(x)<N!\}$.
For an algebraic integer of degree $d\ge2$ and normal-closure
absolute discriminant $\Delta_L$,
\begin{equation}\label{sc:eq:effective-defect}
 \nu(x)\le d\,\Delta_L^{310}.
\end{equation}
Indeed, \href{https://arxiv.org/pdf/2105.14181}{Kadiri--Wong,
Theorem 1} supplies an unramified completely split prime
$p\le\Delta_L^{310}$.  All conjugates of $x$ lie in $\mathbb Z_p$,
so $\bar f$ splits linearly, with maximal multiplicity at most $d$.
The exact first-defect theorem gives a defect by $N=pd$, including
possible nonmaximal-order primes.  The exponent $310$ is an explicit
literature bound, not a claim of optimality or practical efficiency.

For completeness a bound using just monic degree $d$ and coefficient
height $H\ge1$ is fully explicit.  Put $D=d!$ and
\[
 \mathfrak D(d,H)=
 [dD(D-1)(1+H)]^{D(D-1)}.
\]
Then $\Delta_L\le\mathfrak D(d,H)$, whence
$\nu(x)\le d\,\mathfrak D(d,H)^{310}$.
To see the discriminant bound, let $g=[L:\mathbb Q]\le D$ and
$M=g(g-1)/2$.  Among the vectors $c\in\{0,\ldots,M\}^d$ there
is one such that $\theta=\sum c_i\alpha_i$, for the roots
$\alpha_i$ of $f$, has $g$ distinct conjugates.  Each equality
between two embeddings is a nonzero linear equation excluding at
most $(M+1)^{d-1}$ vectors; the $M$ equations cannot cover the
grid.  This integral primitive element has conjugates bounded by
$dM(1+H)$, so its discriminant, divisible by $\Delta_L$, is at most
$[2dM(1+H)]^{g(g-1)}\le\mathfrak D(d,H)$.

For nonintegral algebraic irrational parameters one can use the
prime-unit statistic directly.  Jordan's theorem gives a derangement
conjugacy class.  Kadiri--Wong supplies an unramified prime in that
class with $p\le\Delta_L^{310}$.  There is no $\mathbb Q_p$ root,
hence no $\mathbb Z_p$ root, and therefore
\[
 \min U(x)\le\Delta_L^{310}.
\]
If the primitive minimal polynomial has degree $d$ and height $H$,
with leading coefficient $a$, the algebraic integer $ax$ has the
same normal closure and a monic polynomial of height at most $H^d$.
Consequently $\min U(x)\le\mathfrak D(d,H^d)^{310}$.
For a nonintegral rational of height at most $H$, a prime dividing
its denominator is at most $H$.  These are finite certificate
ranges when bounds and exact prime-unit observations are supplied.

No independent degree, height, or normal-closure discriminant bound
for an allegedly algebraic $\gamma$ has been supplied.  Such bounds
would be new arithmetic input; normalization of the existing
$\operatorname{Ein}$ differential system does not provide them.
Thus effective Chebotarev gives a finite certificate range for
bounded algebraic candidates, not a finite decision algorithm for
$\gamma$.

The absence of a height bound from the common differential system
can itself be demonstrated.  Every algebraic $r$ gives an $E$-function
$F_r=e^z(\operatorname{Ein}(z)-r)$ satisfying
\[
 D(D-1)\bigl[z(D-1)F_r\bigr]=0.
\]
These functions are uniformly bounded on each fixed complex compact
set when $r$ ranges in a bounded real interval.  Nevertheless,
$r_n=(\lfloor nc\rfloor+\sqrt2)/n$, for a rational interior
point $c$, gives distinct degree-two parameters in that interval
for all sufficiently large $n$, with unbounded minimal-polynomial
height.  Thus the common equation, the $E$-function property, and
single-embedding analytic bounds alone cannot provide the missing
height bound.  Extra uniform arithmetic data could change this
conclusion and are not excluded here.

The accompanying exact-arithmetic program checks the density bound
and equality condition on $29$ permutation groups, and $50{,}128$
local denominator equalities through $N=240$ across four maximal
orders, together with the excluded $\sqrt{17}$ example.  These
checks supplement the quantified proofs above.


\section{Residue-profile moments and joint transfer stability}
\label{sc:sec:generalized-transfer}

These are structured applications of classical finite Fourier and
positive-matrix methods. They generalize the earlier parity example;
no priority claim or implication for the arithmetic nature of $\gamma$
is made.

\subsection{Complete reconstruction for arbitrary cyclic profiles}
Fix $d,K\ge2$, put $M=dK$, and let
$a_0=b\in\mathbb R$, $a_{d-r}=\overline{a_r}$.
For $c_n=\widehat\mu(n)$ prescribe
\[
c_0=1,\qquad c_n=a_{n\bmod d}\quad(1\le n<M),
\qquad \lambda_l=\sum_{r=0}^{d-1}a_r e^{2\pi ilr/d}\in\mathbb R.
\]
The Hermitian restriction on the target profile is part of the model.

\begin{theorem}[Residue-profile simplex]\label{sc:thm:residue-simplex}
The data admit a positive probability measure if and only if
\begin{equation}\label{sc:eq:residue-feasible}
1-b\ge0,\qquad 1-b+K\lambda_l\ge0\quad(0\le l<d).
\end{equation}
An explicit representing measure at $x=j/M$ has weights
\[
w_{Kl}=\frac{1-b+K\lambda_l}{M},\qquad
w_j=\frac{1-b}{M}\quad(K\nmid j).
\]
The feasible profiles form a $d$-simplex. Every boundary profile has a
unique representing measure; every interior profile has more than one.
The minimum atomic support is
\[
d(K-1)\boldsymbol1_{b<1}
+\#\{l:1-b+K\lambda_l>0\}.
\]
\end{theorem}
\begin{proof}
Order the moment matrix by exponent residues modulo $d$. It becomes
\[
T=(1-b)I_M+A\otimes J_K,
\qquad A_{rs}=a_{r-s\bmod d}.
\]
Its eigenvalues are $1-b$, with multiplicity $d(K-1)$, and
$1-b+K\lambda_l$. Necessity follows from positivity. The weights
above are nonnegative, sum to one because $\sum_l\lambda_l=db$,
and yield the prescribed moments by Fourier inversion.

The simplex vertices are the profiles of the $d$ point masses at
$l/d$ and the profile $b=-1/(K-1)$, $a_r=0$ for $r\ne0$.
The latter is represented by the uniform measure on the $M$-th roots
other than the $d$-th roots. The barycentric weights are respectively
$w_{Kl}$ and $(K-1)(1-b)/K$.

Use $z=e^{-2\pi ix}$ on the circle. Put $\zeta_j=e^{-2\pi ij/M}$ and
$q_j(z)=\sum_{n=0}^{M-1}\zeta_j^{-n}z^n$.
Its prescribed squared integral is $M^2w_j$.
At a boundary point some $w_j=0$, forcing the support onto the
zeros of $q_j$, namely the other $M$-th roots. All zero weights can
be excluded in this way, and finite Fourier inversion fixes the
remaining weights. In the interior choose
\[
0<\theta<\min\{1,1-b,\min_l(1-b+K\lambda_l)\}.
\]
The rescaled profile $a/(1-\theta)$ is feasible. Mix its grid measure
with weight $\theta$ of either Lebesgue measure or the uniform measure
on $R$-th roots, $R>M-1$. The resulting distinct measures represent
the original data. Finally, any $s$-atomic representation has
$\operatorname{rank}T\le s$; the grid measure attains that rank.
\end{proof}

A failed inequality in \eqref{sc:eq:residue-feasible} supplies a $j$
with $w_j<0$. Then $|q_j|^2$ is an explicit nonnegative polynomial
whose proposed integral is negative. Thus the same calculation gives
both reconstruction and a finite dual rejection certificate.
Feasibility costs one length-$d$ Fourier transform and $d+1$
comparisons, rather than a general $M\times M$ positivity test.

\begin{theorem}[Exact reduction of noisy moments]\label{sc:thm:residue-noise}
Assume $\varepsilon_r\ge0$ and
$\varepsilon_{d-r}=\varepsilon_r$. There is a probability measure
such that
\[
|c_n-a_{n\bmod d}|\le\varepsilon_{n\bmod d}
\quad(1\le n<M)
\]
if and only if there is a Hermitian profile $a'$ satisfying
\eqref{sc:eq:residue-feasible} and
$|a'_r-a_r|\le\varepsilon_r$ for every $r$.
\end{theorem}
\begin{proof}
For necessity average the feasible moment matrix under arbitrary
permutations within its $d$ residue classes and under cyclic
permutation of the classes. The result is positive semidefinite,
has diagonal one, and has the form
$(1-b')I_M+A'\otimes J_K$ with a Hermitian residue profile.
Each entry is an average of entries in the same target disk; for
negative frequency differences the Hermitian target and symmetric
radius assumptions give the same conclusion. Convexity of the disks
preserves every error bound. Theorem~\ref{sc:thm:residue-simplex}
now applies. Sufficiency follows from its explicit grid measure.
\end{proof}

This is a convex feasibility problem in $d$ real dimensions: a
second-order cone problem for complex disks, or a linear program
for real symmetric targets after averaging with complex conjugation.
It is exact as a mathematical reduction; numerical boundary decisions
still require certification.
For completeness, let $C$ be the feasible simplex and use
$\langle u,a\rangle=\sum_r\Re(\overline{u_r}a_r)$.
Its support function is
\[
h_C(u)=\max\left\{-\frac{u_0}{K-1},\
\max_l\sum_r\Re(\overline{u_r}e^{-2\pi ilr/d})\right\}.
\]
Convex separation gives the equivalent dual criterion
\[
\langle u,a\rangle\le h_C(u)+\sum_r\varepsilon_r|u_r|
\quad\text{for every Hermitian }u.
\]
A violating $u$ certifies infeasibility; a feasible profile constructs
a measure. With real algebraic inputs, finite semialgebraic decision
methods can resolve even boundary cases.

\subsection{The exact joint stability constant for two multipliers}
For $N\ge1$ and distinct primes $p,q$ define
\[
(P_pf)(x)=\frac1p\sum_{j=0}^{p-1}f((x+j)/p),
\qquad \widehat{P_pf}(n)=\widehat f(pn).
\]
Let $C_N(p,q)$ be the best constant, over trigonometric polynomials
of degree at most $N$, in
\begin{equation}\label{sc:eq:joint-stability}
\|f-\widehat f(0)\|_2\le C_N(p,q)
\bigl(\|f-P_pf\|_2^2+\|f-P_qf\|_2^2\bigr)^{1/2}.
\end{equation}
Put $D_N=\{(a,b)\ge0:p^aq^b\le N\}$ and let $H_N$ be the matrix
of the energy
\[
\sum_{(a,b)\in D_N}
\bigl(|u_{a,b}-u_{a+1,b}|^2+|u_{a,b}-u_{a,b+1}|^2\bigr),
\]
where $u$ is extended by zero outside $D_N$.
Its diagonal entry at $(a,b)$ is $2+\boldsymbol1_{a>0}
+\boldsymbol1_{b>0}$, and neighboring entries are $-1$.

\begin{theorem}[Joint spectral reduction and optimal order]
\label{sc:thm:joint-stability}
One has $C_N(p,q)=\lambda_{\min}(H_N)^{-1/2}$, with
\[
\dim H_N=\frac{(\log N)^2}{2\log p\log q}+O_{p,q}(\log N)
\quad(N\longrightarrow\infty).
\]
Set $L_p=1+\lfloor\log_pN\rfloor$,
$L_q=1+\lfloor\log_qN\rfloor$,
and $L=1+\lfloor\log_{pq}N\rfloor$. Then
\[
\frac1{2\sqrt2\sin(\pi/(4L+2))}
\le C_N(p,q)\le
\frac1{2\sqrt{\sin^2(\pi/(4L_p+2))+
\sin^2(\pi/(4L_q+2))}}.
\]
In particular $C_N(p,q)=\Theta_{p,q}(\log N)$. The exact optimum
and the lower-bound examples can be realized by nonnegative real
probability densities.
\end{theorem}
\begin{proof}
Decompose every nonzero frequency uniquely as $rp^aq^b$, with
$p\nmid r$ and $q\nmid r$. Parseval splits the joint defect into
the energies on $D_{N/|r|}$. These are lower ideals; zero extension
from $D_{N/|r|}$ into $D_N$ preserves energy. Hence their least
Rayleigh quotients are no smaller than that on $D_N$. The orbit
$r=1$ attains the minimum. The matrix is positive definite, since
zero energy propagates a constant along the positive coordinate
edges to the zero exterior. This proves the spectral formula.
The dimension estimate is the elementary lattice-point count in
$a\log p+b\log q\le\log N$.

On a one-prime chain of length $s$, the difference matrix $I-S$
has least squared singular value $4\sin^2(\pi/(4s+2))$.
Indeed the vector $u_j=\cos((j+1/2)\pi/(2s+1))$ solves the
second-difference equation with first endpoint $u_{-1}=u_0$ and
last endpoint $u_s=0$. Applying this to each $p$-chain and each
$q$-chain and adding the resulting inequalities proves the upper
bound. For the lower bound, the square $0\le a,b<L$ lies inside
$D_N$. The tensor product of the two length-$L$ minimizing vectors,
extended by zero, has energy-to-norm ratio
$8\sin^2(\pi/(4L+2))$.

A real minimizing vector supplies a cosine polynomial $g$ on its
Fourier orbit. For sufficiently small $t\ne0$, $1+tg$ is a
nonnegative probability density, with the same homogeneous norm
ratio. This also applies to the square test vector.
\end{proof}

Thus two maps improve constants but cannot give a uniform or
$o(\log N)$ homogeneous Lipschitz estimate of the form
\eqref{sc:eq:joint-stability}, even for positive densities.
This is a bandlimited norm statement. It does not settle an
unrestricted invariant-measure conjecture or provide the missing
arithmetic bridge for $\gamma$.




% Core results: source proofs retained.
\section{Moving-degree fresh-prime removal: an explicit arithmetic
budget}\label{sc:md:moving-degree-fresh-prime-removal-an-explicit-arithmetic-budget}

\subsection{Scope}\label{sc:md:scope}

This note addresses polynomial preprocessing of \[
E=e^z,\qquad F_r=e^z(\operatorname{Ein}(z)-r)
=\sum_{n\ge0}(H_n-r)\frac{z^n}{n!},
\] where \(r\) is algebraic. Write \([G]_n=n![z^n]G\). The preprocessing
polynomial is allowed to change with the prefix length, and both its
degree and coefficient height may grow.

The conclusion concerns \textbf{removing prime denominators from
normalized coefficients}, not merely obtaining an E-function. In fact,
\(F_r\) is already an E-function for algebraic \(r\); its harmonic
denominators have exponential least-common-multiple growth. A lower
bound for the cost of making a finite prefix integral is an arithmetic
budget, not by itself an obstruction to every E-function argument or a
proof about gamma.

The functional algebraic independence of \(E,F_r\) over \(\mathbb C(z)\)
also has a short direct proof. As \(z\to+\infty\), \[
F_r(z)=e^z(\log z+\gamma-r)+O(1/z).
\] In a nonzero polynomial relation, collect the largest total
\((E,F_r)\)-degree \(s\). Its contribution is \(e^{sz}\) times a nonzero
polynomial in \(z\) and \(\log z+\gamma-r\); all lower-degree
contributions, and errors from the displayed expansion, are
exponentially smaller. A nonzero polynomial in \(z,\log z+c\) has a
leading power of \(z\) times a nonzero polynomial in \(\log z\), so it
cannot be exponentially small. This excludes the relation, after
clearing any rational \(z\)-denominators. It uses no arithmetic
assumption on \(\gamma\). The ingredients below are an explicit uniform
zero estimate, an integer jet certificate independent of \(r\), and a
degree-height bound valid for changing polynomials.

\subsection{The fresh-prime congruence is uniform in polynomial
degree}\label{sc:md:the-fresh-prime-congruence-is-uniform-in-polynomial-degree}

Let \(P\in\mathbb Z[z,T,Y]\), put \(G=P(z,E,F_r)\), and let
\(\mathfrak p\) lie over a rational prime \(p\) at which \(r\) is
integral. For \textbf{every} \(0\le m<p\), \[
\boxed{p[G]_{p+m}\equiv [E P_Y(z,E,F_r)]_m\pmod{\mathfrak p}.}\tag{MD.1}\label{sc:md:eq:1}
\] No bound on the degree of \(P\) is needed for this identity. Its
right side may vanish because the degree is large, as in \(P=Y^p\), but
that is a detected arithmetic possibility rather than a failure of the
identity.

\textbf{Proof.} It is enough to expand \(z^aE^bF_r^c\). Since
\(p+m<2p\), at most one factor \(F_r\) can receive an index at least
\(p\) in its exponential-generating-series convolution. All terms with
only smaller indices are \(\mathfrak p\)-integral and disappear after
multiplication by \(p\). For the unique large index \(p+k\), \[
p(H_{p+k}-r)\equiv1\pmod{\mathfrak p}.
\] Lucas reduction removes the leading base-\(p\) digit from the
multinomial coefficient. There are \(c\) choices of the large factor,
yielding \(E\,cF_r^{c-1}\). When the factor \(z^a\) is restored, its
falling factorial reduces to \(m^{\underline a}\) for \(a\le m\). If
\(a>m\), it either contributes a factor \(p\) or the remaining total
degree is below \(p\); both sides of \eqref{sc:md:eq:1} are zero modulo
\(\mathfrak p\). This also covers \(a\ge p\) and \(c\ge p\). \(\square\)

\subsection{An explicit zero estimate without a Padé-perfectness
assumption}\label{sc:md:an-explicit-zero-estimate-without-a-paduxe9-perfectness-assumption}

Let \(D\ge1,A\ge0\), and set \[
L=\frac{(D+1)(D+2)}2,\qquad
M_*=LA+\frac{L(L-1)}2.
\]

\begin{theorem}[Uniform polynomial zero estimate]\label{sc:md:uniform-zero}

For every nonzero polynomial \(Q(z,T,Y)\) of \(z\)-degree at most \(A\)
and total \((T,Y)\)-degree at most \(D\), \[
\boxed{\operatorname{ord}_{z=0}Q(z,E,F_r)\le M_*.}\tag{MD.2}\label{sc:md:eq:2}
\] The bound is independent of \(r\), of the coefficient height, and of
the chosen polynomial.

\end{theorem}
\textbf{Proof.} Let \(V\) be the column of the \(L\) monomials
\(E^bF_r^c\), \(b+c\le D\), and let \(\theta=z\,d/dz\). Their equations
are \[
\theta(E^bF_r^c)=z(b+c)E^bF_r^c
+c(E^{b+1}F_r^{c-1}-E^bF_r^{c-1}).
\] Thus \[
\theta V=(B_0+zB_1)V,
\] where \(B_0\) is nilpotent: each of its nonzero transitions strictly
lowers the exponent \(c\).

Write \(G=Q(z,E,F_r)=\ell_0(z)V\), with \(\deg\ell_0\le A\), and
recursively set \[
\ell_{j+1}=\theta\ell_j+\ell_j(B_0+zB_1).
\] Then \(\ell_jV=\theta^jG\) and \(\deg\ell_j\le A+j\). Let \(h\le L\)
be the first index for which \(\ell_0,\ldots,\ell_h\) are dependent over
\(\mathbb C(z)\). The first \(h\) rows are independent. Their rational
span \(W\) is connection-stable: the first dependence puts the
connection of \(\ell_{h-1}\) in that span, and the other rows map to
their successors; the Leibniz rule handles rational coefficients.

Over the discrete valuation ring \(R=\mathbb C[[z]]\), take the
saturation \[
\mathcal S=(W\otimes_{\mathbb C(z)}\mathbb C((z)))\cap R^L.
\] It is free of rank \(h\), is a direct summand of \(R^L\), and is
stable under the row connection. Reduction modulo \(z\) embeds it into
the ambient residue space; the induced residue is a restriction of the
nilpotent matrix \(B_0\), and is therefore nilpotent.

Choose a basis matrix \(S\) of \(\mathcal S\), and put \(y=SV\). It is a
formal power-series solution of \(\theta y=C(z)y\), with
\(C(z)\in\operatorname{Mat}_h(R)\) and \(C(0)\) nilpotent. It is nonzero
because \(G\ne0\). Here \(G\ne0\) follows from the previously
established algebraic independence of \(E,F_r\). A nonzero solution
\(y\) must have at least one unit coordinate: if its minimum order were
\(m>0\), comparison of leading terms would make \(m\) an eigenvalue of
the nilpotent matrix \(C(0)\).

Let \(s\) be the index valuation of the row module generated by
\(\ell_0,\ldots,\ell_{h-1}\) inside its saturation. By Smith normal
form, \[
s=\min\{\operatorname{ord}_0\Delta:
\Delta\text{ a nonzero maximal minor of }(\ell_j)_{j<h}\}.
\] The degree of every such minor is at most \[
\sum_{j=0}^{h-1}(A+j)=hA+h(h-1)/2,
\] so \(s\) is no larger. The adjugate formula also gives
\(z^s\mathcal S\) inside the original row module. If \(G\) has order
\(m\), every \(\theta^jG\) has order at least \(m\). Therefore every
coordinate of \(z^sy\) has order at least \(m\). Since one coordinate of
\(y\) is a unit, \(m\le s\). Finally \(h\le L\) proves \eqref{sc:md:eq:2}. \(\square\)

This argument uses the nilpotent regular-singular residue, not an
assertion that the first as-many-as-dimension jets are nonsingular. That
stronger perfectness assertion is neither needed nor claimed.

The same proof is reusable beyond these functions: if an analytic vector
\(V\) satisfies \(\theta V=(B_0+zB_1)V\) with nilpotent \(B_0\), then
every \textbf{nonzero evaluated} expression \(\ell_0(z)V\),
\(\deg\ell_0\le A\), has order at most \(LA+L(L-1)/2\). Algebraic
independence is used only to ensure nonvanishing for all nonzero
polynomial inputs in the present application. No priority claim relative
to general multiplicity-estimate literature is made. For established
effective zero estimates and their
Shidlovskii--Bertrand--Chirskii--Yebbou background, see
\href{https://arxiv.org/html/1906.05529v4}{Bostan--Rivoal--Salvy, §1,
``Minimal differential equations'', equation (4)}. This note gives a
self-contained special-case argument; it does not claim to have invented
the zero-estimate method or established priority for the combined
arithmetic budget.

\subsection{A finite integer certificate independent of the algebraic
shift}\label{sc:md:a-finite-integer-certificate-independent-of-the-algebraic-shift}

Use the full basis \[
\mathcal B_{A,D}=\{z^aT^bY^c:0\le a\le A,\ b+c\le D\},
\qquad B=(A+1)L.
\] Theorem~\ref{sc:md:uniform-zero} makes evaluation into jets of orders \(0,\ldots,M_*\)
injective. Select any \(B\) independent rows with indices \[
0\le m_1<\cdots<m_B\le M_*;
\qquad M=\max_i m_i.
\] For each basis monomial form its normalized coefficient at order
\(m_i\), multiply that row by \(m_i!\), and compute the determinant at
\(r=0\). Denote this nonzero integer by \(\Delta_{A,D}\).

Two facts make this a reusable certificate.

\textbf{Integrality.} The ordinary coefficients of \(F_0\) lie in
\((n!)^{-2}\mathbb Z\), because \(H_n\) has denominator dividing \(n!\).
The same is true of \(E\), and this coefficient class is closed under
products: the factorial ratios produce squares of integer multinomial
coefficients. Hence a normalized coefficient of any integral polynomial
in \(z,E,F_0\) has denominator dividing \(n!\). The row scaling
therefore gives an integer matrix.

\textbf{Shift invariance.} The substitution \(F_r=F_0-rE\) acts on the
complete monomial basis by a triangular matrix with determinant one,
separately at each \(z\)-degree. Thus exactly the same row determinant
is \(\Delta_{A,D}\) for every algebraic \(r\). When \(r\) is integral at
\(\mathfrak p\), this change of basis and its inverse are integral
there.

An explicit height bound follows from Hadamard. Put
\(H_M=\sum_{j=1}^M1/j\). For a basis monomial, its normalized
coefficient through order \(M\) has absolute value at most \[
M^A D^M H_M^D.
\] Indeed the harmonic factors in the exponential convolution are at
most \(H_M\), the sum of multinomial coefficients is at most \(D^M\), and
multiplication by \(z^a\) contributes at most \(M^A\). Therefore \[
\boxed{\log|\Delta_{A,D}|\le
B\left\{\tfrac12\log B+\log(M!)+A\log M
+M\log D+D\log H_M\right\}.}\tag{MD.3}\label{sc:md:eq:3}
\] Here the selected \(M\) can be used; replacing it by \(M_*\) gives a
closed-form bound. In particular, \[
\log|\Delta_{A,D}|=
O\bigl((A+1)D^4(A+D^2)\log(A+D+2)\bigr).\tag{MD.4}\label{sc:md:eq:4}
\] The implicit constant is absolute. Keeping the exact determinant can
be much stronger than this deliberately coarse estimate.

\subsection{A degree-height budget for changing
polynomials}\label{sc:md:a-degree-height-budget-for-changing-polynomials}

For \(P\in\mathbb Z[z,T,Y]\) retaining \(Y\), define \[
C_Y(P)=\gcd\{c\,p_{abc}:c\ge1\}>0,
\] where \(p_{abc}\) are its monomial coefficients. Let \(H(P)\) be
their maximum absolute value. Then \(C_Y(P)\le D H(P)\).

\begin{theorem}[Finite-prefix prime-removal budget]\label{sc:md:finite-budget}

Assume \(\deg_zP\le A\), \(\deg_{T,Y}P\le D\), and let \(G=P(z,E,F_r)\).
Let \(\mathcal P\) be any set of rational primes satisfying \[
M<p\le T-M,
\] such that for each \(p\in\mathcal P\) there is a prime
\(\mathfrak p\mid p\) at which \(r\) and every \([G]_n\),
\(0\le n\le T\), are integral. Then \[
\boxed{\prod_{p\in\mathcal P}p\ \mid\
\Delta_{A,D}\,C_Y(P).}\tag{MD.5}\label{sc:md:eq:5}
\] Consequently, \[
\boxed{\sum_{p\in\mathcal P}\log p
\le\log|\Delta_{A,D}|+\log D+\log H(P).}\tag{MD.6}\label{sc:md:eq:6}
\]

\end{theorem}
\textbf{Proof.} Put \(Q=TP_Y\). It has the same degree caps as \(P\);
its coefficient vector in \(\mathcal B_{A,D}\) is integral and has
content \(C_Y(P)\). For \(m=m_i\), congruence \eqref{sc:md:eq:1} and the integral
prefix give \([Q(z,E,F_r)]_{m_i}=0\) modulo \(\mathfrak p\), since
\(p+m_i\le T\) and \(m_i<p\). Multiplying by \(m_i!\), the
integer-scaled jet matrix at shift \(r\) kills the vector of \(Q\)
modulo \(\mathfrak p\). Its determinant is \(\Delta_{A,D}\). The
adjugate gives \[
\Delta_{A,D}\,u_j\equiv0\pmod{\mathfrak p}
\] for every integer coefficient \(u_j\) of \(Q\). These are rational
integers; hence \(p\) divides every \(\Delta_{A,D}u_j\), and so divides
\(\Delta_{A,D}C_Y(P)\). Multiply over the distinct primes. \(\square\)

Only prime factors of \(\Delta\) exceeding \(M\) can matter in \eqref{sc:md:eq:5};
stripping its small-prime factors sharpens the numerical budget. The
displayed height estimate simply uses the full determinant to avoid
factoring it.

\subsubsection{Combined denominator-height budget}\label{sc:md:corollary-combined-denominator-height-budget}

Let \(K=\mathbb Q(r)\), and let \(d_T\) be the least positive rational
integer such that \[
d_T[G]_n\in\mathcal O_K\qquad(0\le n\le T).
\] It exists because the prefix is finite and consists of algebraic
numbers. Let \(S_r\) be the finite set of rational primes at which \(r\)
is not integral at at least one prime of \(K\). Then \[
\boxed{\prod_{\substack{M<p\le T-M\\p\notin S_r}}p
\ \mid\ d_T\,\Delta_{A,D}\,C_Y(P).}\tag{MD.7}\label{sc:md:eq:7}
\] Indeed, if \(p\nmid d_T\), the full prefix is integral at every prime
above \(p\), so Theorem~\ref{sc:md:finite-budget} applies. If \(p\mid d_T\), the prime is
already accounted for. Consequently, with
\(\vartheta(x)=\sum_{p\le x}\log p\), \[
\boxed{\log(d_T H(P))\ge
\vartheta(T-M)-\vartheta(M)
-\sum_{p\in S_r}\log p
-\log|\Delta_{A,D}|-\log D.}\tag{MD.8}\label{sc:md:eq:8}
\] Thus the bound measures the combined cost of residual denominators
and preprocessing coefficients, rather than assuming that preprocessing
removes every denominator.

\subsubsection{\texorpdfstring{Corollary (slowly growing nonlinear
preprocessing still costs
\(e^{T-o(T)}\))}{Corollary (slowly growing nonlinear preprocessing still costs e\^{}\{T-o(T)\})}}\label{sc:md:corollary-slowly-growing-nonlinear-preprocessing-still-costs-et-ot}

Fix algebraic \(r\). Suppose \(P_T\) retains \(Y\) and has degree caps
\(A_T,D_T\). Define \(d_T\) for its normalized prefix as above. If \[
(A_T+1)D_T^4(A_T+D_T^2)\log(A_T+D_T+2)=o(T),\tag{MD.9}\label{sc:md:eq:9}
\] then \[
\boxed{\log\bigl(d_T H(P_T)\bigr)\ge T-o(T).}\tag{MD.10}\label{sc:md:eq:10}
\] In particular, if preprocessing makes the entire prefix integral,
then \(\log H(P_T)\ge T-o(T)\). For \(A_T=0\), the sufficient growth
condition simplifies to \[
D_T^6\log(D_T+2)=o(T).
\]

\textbf{Proof.} Condition \eqref{sc:md:eq:9} implies both \(M=o(T)\) and
\(\log|\Delta|=o(T)\). Apply \eqref{sc:md:eq:8} and use the prime number theorem in the
form \(\vartheta(T-M)-\vartheta(M)=T-o(T)\). Also \(\log D=o(T)\).
\(\square\)

The leading constant one is sharp in either allocation of cost. For
algebraic integral \(r\), \(P_T(Y)=Y\) has height one and
\(d_T=\operatorname{lcm}(1,\ldots,T)\). The latter equality follows by
subtracting successive normalized coefficients, whose difference is
\(1/n\). Alternatively, \[
P_T(Y)=\operatorname{lcm}(1,\ldots,T)Y
\] has integral normalized coefficients through order \(T\), and its
logarithmic height is \(\psi(T)=T+o(T)\). This is a sharp asymptotic
cost statement within the specified regime.

\subsubsection{Degree escapes are accounted for, not assumed
absent}\label{sc:md:degree-escapes-are-accounted-for-not-assumed-absent}

If \(P=Y^L\), then \(C_Y(P)=L\). Thus a prime dividing \(L\) can make
the first fresh-prime derivative vanish even at height one; removing
many such primes requires their product to divide \(L\), so it costs
degree. More generally, \(P_Y=0\) modulo \(p\) means \(P\bmod p\) is a
polynomial in \(Y^p\). Formula \eqref{sc:md:eq:5} captures this mechanism through
derivative content. The theorem does not pretend to exclude all
fast-degree constructions outside \eqref{sc:md:eq:9}.

\subsection{The sharper pure-polynomial
subfamily}\label{sc:md:the-sharper-pure-polynomial-subfamily}

For comparison, when \(P\in\mathbb Z[Y]\) has degree \(d\), a direct jet
argument is stronger. Exclude primes at which \(r\) is nonintegral or
\(1-r\) is not a unit; assume \(r\ne1\). If \(p>d\) and \(p+d-1\le T\),
an integral prefix gives \[
[E P'(F_r)]_m\equiv0\quad(0\le m<d).
\] The germ \(F_r+r\) has linear coefficient \(1-r\), a unit. Formal
inversion modulo \(p\), together with the unit factor \(E\), forces the
polynomial \(P'\) of degree at most \(d-1\) to be zero modulo \(p\).
Since \(p>d\), every nonconstant coefficient of \(P\) is divisible by
\(p\). Hence \[
\log H(P)\ge\vartheta(T-d+1)-\vartheta(d)-O_r(1).
\] For \(d=o(T)\), this again gives \(\log H\ge T-o(T)\), with a much
wider allowed degree regime. The multivariate theorem above pays for
separating the additional \(z,E\) directions and does not claim the same
optimal degree threshold.

The exceptional shift \(r=1\) can also be covered:
\(F_1+1=z^2/4+O(z^3)\). Away from \(p=2\), a nonzero polynomial
\(P'\bmod p\) of degree at most \(d-1\) has composition order at most
\(2(d-1)\). Test those jets instead, requiring \(p>\max\{d,2(d-1),2\}\)
and \(p+2(d-1)\le T\). This retains the same \(d=o(T)\) asymptotic
conclusion.

\section{A fixed-shift obstruction for the Laguerre approximants}
\label{sc:hp:sec:laguerre-shift-obstruction}

We use the Hessami Pilehrood $k!$ family, which is distinct from
Aptekarev's original $(n+k)!$ family:
\[
 q_n=\sum_{k=0}^n\binom nk^2k!=n!L_n(-1),\qquad
 p_n=\sum_{k=0}^n\binom nk^2k!(2H_{n-k}-H_k),\qquad
 E_n=p_n-\gamma q_n.
\]
The source is Kh.~Hessami Pilehrood and T.~Hessami Pilehrood,
\emph{On a continued fraction expansion for Euler's constant},
\href{https://arxiv.org/pdf/1010.1420}{arXiv:1010.1420},
equations (21)--(23), Lemmas 2--3, and Theorem 1.
Writing $S u_n=u_{n+1}$ and
\[
 R=S^2-2(n+2)S+(n+1)^2,
\]
their formulas give $Rq=0$, $Rp=RE=-n/(n+2)$, and
\begin{equation}\label{sc:hp:eq:hp-error-growth}
 E_n=n!n^{-1/4}e^{-2\sqrt n}
       \left(\sqrt{\pi/e}+O(n^{-1/2})\right).
\end{equation}
In particular $E_n$ is eventually positive and grows faster than any
power of $n$.

\begin{theorem}[Fixed rational-function shifts]
\label{sc:hp:thm:fixed-hp-shifts}
Let $T\in\mathbb Q(n)[S]$ be fixed.  If $Tq$ is not the zero
sequence, then for some $c\ne0$ and $\rho\in\frac12\mathbb Z$,
\[
 (TE)_n=c n^\rho E_n(1+o(1)).
\]
Consequently $|(TE)_n|\to\infty$.  Adding a fixed rational function
of $n$ does not change this conclusion.
\end{theorem}

\begin{proof}
We first prove
\begin{equation}\label{sc:hp:eq:hp-adjacent-ratio}
 \frac{E_{n+1}}{E_n}=n-\sqrt n+o(\sqrt n).
\end{equation}
The positive Whittaker integral in the cited paper has the form
\[
 I_n=\frac{n!}{e}\int_0^\infty
       \frac{e^{-t}}{1+t}\left(\frac t{1+t}\right)^n\,dt,
 \qquad I_n-E_n=O((n+1)^{-2}).
\]
The integral identity follows by substituting the standard integral for
$U(n+1,1,1)$ into its equation (28); the error bound is its Lemma 2.
Set $J_n=I_n/n!$.  The positive moment representation implies
$J_n^2\le J_{n-1}J_{n+1}$ by Cauchy--Schwarz.  Therefore
$f_n=\log J_n$ is convex.  Equation~\eqref{sc:hp:eq:hp-error-growth} gives
\[
 f_n=-2\sqrt n-\tfrac14\log n+\tfrac12\log(\pi/e)+O(n^{-1/2}).
\]
For $h=\lfloor n^{3/4}\rfloor$, the backward and forward $h$-step
secant slopes bound $f_{n+1}-f_n$.  Substitution yields
\[
 f_{n+1}-f_n=-n^{-1/2}+O(n^{-3/4}),
\]
and the negligible correction $I_n-E_n$ gives
\eqref{sc:hp:eq:hp-adjacent-ratio}.

Right division by the monic operator $R$ in the shift Ore ring gives
\[
 T=UR+A(n)S+B(n),\qquad A,B\in\mathbb Q(n).
\]
The hypothesis and $Rq=0$ imply $(A,B)\ne(0,0)$.  Since $RE$ is a
rational function,
\[
 (TE)_n=A(n)E_{n+1}+B(n)E_n+C(n),\qquad C\in\mathbb Q(n).
\]
If $A=0$, the conclusion follows from the leading term of $B$.
Otherwise let $a=\deg A$ and $B_0=B+nA$.  Equation
\eqref{sc:hp:eq:hp-adjacent-ratio} gives
\[
 A(n)\frac{E_{n+1}}{E_n}+B(n)
 =B_0(n)-A(n)\sqrt n\,(1+o(1)).
\]
The nonzero rational function $B_0$, if present, has integral degree;
$A\sqrt n$ has half-integral degree $a+1/2$.  Their leading terms
cannot cancel.  The rational function $C$ is negligible relative to
the factorial growth of $E_n$.  This proves the theorem.
\end{proof}

The scope is essential: the theorem does not bound the common divisor
after integerization, and it does not cover a shift order increasing
with $n$.  The exact identity $(n+2)RE=-n$ is not a small gamma form,
because simultaneously $(n+2)Rq=0$.

\paragraph{The primitive arithmetic budget.}
Write $p_n=A_n/b_n$ in lowest terms, $b_n>0$, and set
$g_n=\gcd(A_n,q_n)$.  Then
\[
 P_n=A_n/g_n,\qquad Q_n=b_nq_n/g_n
\]
are coprime integers and
\[
 P_n-\gamma Q_n=\frac{b_n}{g_n}E_n.
\]
Hence their nonzero forms tend to zero along an infinite subsequence
if and only if, along that subsequence,
\begin{equation}\label{sc:hp:eq:hp-primitive-budget}
 \log g_n-\log b_n-\log(n!)+2\sqrt n+\tfrac14\log n
 \longrightarrow+\infty.
\end{equation}
This identity locates the missing arithmetic saving; it does not
establish it.

\begin{proposition}[Elementary congruences]
For every $n\ge1$ and prime $\ell$,
\[
 q_n\equiv1\pmod n,\qquad
 q_n\equiv q_{n\bmod\ell}\pmod\ell.
\]
If $\ell\le n<2\ell$, then in $\mathbb Z_{(\ell)}$,
\[
 \ell p_n\equiv2q_{n-\ell}\equiv2q_n\pmod\ell.
\]
\end{proposition}

\begin{proof}
For $k\ge1$, write
$\binom nk^2k!=\binom nk n(n-1)\cdots(n-k+1)$ to obtain divisibility
by $n$.  For reduction modulo $\ell$, discard $k\ge\ell$ and apply
Lucas congruence to the remaining binomial coefficients.
For the last claim put $n=\ell+r$, $0\le r<\ell$.  In
$\ell(2H_{n-k}-H_k)$, the residue is $2$ for $k\le r$.
For $r<k<\ell$, both harmonic numbers are $\ell$-integral and the
binomial coefficient also vanishes modulo $\ell$.
For $k\ge\ell$, the factorial supplies a factor $\ell$ while the
scaled harmonic expression is $\ell$-integral.  Only $k\le r$
contributes, giving $2q_r$.
\end{proof}

For odd $\ell$ in this interval, $q_n\not\equiv0\pmod\ell$ forces
the reduced denominator $b_n$ to contain $\ell$ exactly once.
Uniform nonvanishing cannot simply be assumed: at $n=\ell-1$,
\[
 q_{\ell-1}\equiv\sum_{k=0}^{\ell-1}k!=!\ell\pmod\ell.
\]
Its nonvanishing for every odd prime is Kurepa's left-factorial
hypothesis; see R.~Me\v strovi\'c,
\emph{Variations of Kurepa's left factorial hypothesis},
\href{https://arxiv.org/abs/1312.7037}{arXiv:1312.7037}.
This is not an equivalence between that hypothesis and the gamma
problem.

Another tempting shortcut is actually false. Exact recurrence
arithmetic gives
\[
 \boxed{g_{224}=233>224.}
\]
For a compact certificate of the common prime,
\[
 q_{224}\equiv208\cdot233\pmod{233^2},\qquad
 A_{224}\equiv31\cdot233\pmod{233^2},\qquad
 b_{224}\equiv16\pmod{233}.
\]
The full integer gcd is computed by the accompanying exact script.
Thus neither ``all common primes are at most $n$'' nor $g_n\mid n!$
is valid.  The same script finds only 27 nontrivial $g_n$ through
$n=3000$, with maximum 233; these finite observations prove no
asymptotic gcd bound.

\paragraph{Moving order remains an arithmetic problem.}
Pascal's identity gives, for every $n,m\ge0$,
\[
 \Delta^m L_n(-1)=\sum_{j=0}^n\frac{\binom nj}{(j+m)!}
 =\frac{n!}{(n+m)!}L_n^{(m)}(-1),
\]
so $(n+m)!\Delta^m(q_n/n!)$ is integral.  This is an exact
identity, not a proof that a substantially smaller denominator
suffices.  No factorial-size cancellation for moving order, and no
complete primitive-gcd obstruction, is established here.

\section{Positive Binet kernels: a global floor and extremal
zeros}\label{sc:pk:positive-binet-kernels-a-global-floor-and-extremal-zeros}

\subsection{Definitions}\label{sc:pk:definitions}

Write

\[
Q_M(u)=\prod_{m=2}^M(u+m^2),\qquad
F(u)=\sum_{m=2}^M\frac{R_m}{u+m^2}=\frac{P(u)}{Q_M(u)}.
\]

The normalization \(\sum_m R_m=1\) makes \(P\) monic of degree
\(d=M-2\). Positivity of $F$ on $[0,\infty)$ is exactly positivity of $P$ there. Log balance imposes

\[
\sum_m R_m v_p(m)=0\quad\text{for every prime }p\le M.
\]

The Binet energy is

\[
E(F)=\int_0^\infty\frac{F(u)}{e^{2\pi\sqrt u}-1}\,du.
\]

For balanced kernels, the supplied Binet identity gives
\(E(F)=\gamma-\sum_m R_m h_m\), with \(h_m=H_{m-1}+1/(2m)\). The
feasible set is the nonempty compact set established by the preceding
positive-kernel construction for \(M\ge4\); its minimum is \(E_M^*>0\).

\begin{lemma}[Reused nonemptiness and compactness input]
The normalized log-balanced nonnegative kernel set is nonempty and compact
for every $M\ge4$.
\end{lemma}
\begin{proof}
Take $R_2=2$, $R_4=-1$ and all other residues zero. Then
$\sum R_m=1$, $\sum R_m\log m=0$, and
$F(u)=2/(u+4)-1/(u+16)>0$ for $u\ge0$.
The feasible residue set is a closed convex subset of a finite-dimensional
space. A nonzero recession direction would give a nonnegative kernel
$F_R$ with $\sum R_m=0$ and $\sum R_m\log m=0$.
It is integrable on the half-line because $F_R(u)=O(u^{-2})$, and
\[
\int_0^\infty F_R(u)\,du=-2\sum_mR_m\log m=0.
\]
Nonnegativity then forces $F_R=0$, hence every residue is zero, a
contradiction. A closed finite-dimensional convex set with no nonzero
recession direction is bounded. This proves compactness.
\end{proof}

\subsection{A full-cone monic polynomial norm
lemma}\label{sc:pk:a-full-cone-monic-polynomial-norm-lemma}

If $P$ is any monic real polynomial of degree $d$ nonnegative on $[0,\infty)$, then for every $T>0$,

\[
\int_0^T P(u)\,du
\ge
\frac{T^{d+1}}
{(d+1)\binom{d}{\lfloor d/2\rfloor}^2}
\ge\frac{T^{d+1}}{(d+1)4^d}.
\tag{PK.1}\label{sc:pk:eq:1}
\]

\textbf{Proof.} Every such $P$ admits a half-line sum-of-squares
representation $P=A^2+uB^2$. This can be proved directly by factoring: a
negative real factor $u+c$ is $(\sqrt c)^2+u$; a conjugate pair $(u-\alpha)^2+\beta^2$ equals $(u-r)^2+2(r-\alpha)u$ with $r=\sqrt{\alpha^2+\beta^2}$;
positive roots occur to even order; and representations are preserved by

\[
(A^2+uB^2)(C^2+uD^2)=(AC-uBD)^2+u(AD+BC)^2.
\]

If $d=2k$, $A$ has degree $k$ with leading coefficient $\pm1$, while $\deg B\le k-1$.
Therefore the integral is bounded below by the minimum squared norm of a
monic degree-k polynomial for Lebesgue measure on $[0,T]$. If $d=2k+1$, $B$ is monic up to sign of degree $k$ and the lower bound is the corresponding monic norm for weight $u$.

For $\epsilon=0$ or $1$, shifted Jacobi orthogonality gives the monic
degree-k norm

\[
\min_{A\text{ monic},\deg A=k}\int_0^T u^\epsilon A(u)^2du
=\frac{T^{2k+\epsilon+1}}
{(2k+\epsilon+1)\binom{2k+\epsilon}{k}^2}.
\]

For completeness the normalized orthogonal polynomial on $[0,1]$ is
supplied by

\[
R_k(x)=\frac{x^{-\epsilon}}{k!}
\frac{d^k}{dx^k}\{x^{k+\epsilon}(x-1)^k\}.
\]

Integration by parts proves orthogonality to every lower-degree
polynomial. Its leading coefficient is $\binom{2k+\epsilon}{k}$; its squared norm is $1/(2k+\epsilon+1)$, by one further beta-integral computation.
Scaling gives the displayed formula. This proves \eqref{sc:pk:eq:1}. The argument
allows arbitrary signs of the coefficients of $P$.

\subsection{The global exponential
floor}\label{sc:pk:the-global-exponential-floor}

For $y>0$ define

\[
J(y)=\int_0^1\log(y^2+t^2)dt
=\log(1+y^2)-2+2y\arctan(1/y),
\]

\[
C(y)=2\pi y+J(y)+\log4-2\log y.
\]

\textbf{Theorem.} For every $M\ge4$ and every normalized nonnegative kernel
above,

\[
E(F)\ge
\frac{16y^2}{(M-1)(1+y^2)}\exp(-C(y)M).
\tag{PK.2}\label{sc:pk:eq:2}
\]

This statement does not even require log balance, so it applies
uniformly to the full balanced cone.

\textbf{Proof.} Put $T=y^2M^2$. On $0\le u\le T$,

\[
\frac1{e^{2\pi\sqrt u}-1}>e^{-2\pi yM},\qquad Q_M(u)\le Q_M(T).
\]

Thus $E(F)\ge e^{-2\pi yM}Q_M(T)^{-1}\int_0^TP(u)\,du$. Apply \eqref{sc:pk:eq:1}.
To estimate the denominator, $f(t)=\log(y^2+t^2)$ is increasing, and a
right-Riemann-sum bound gives

\[
\sum_{m=2}^M f(m/M)
\le M\int_0^1f(t)dt+f(1)-2f(0).
\]

Consequently

\[
Q_M(T)\le M^{2(M-1)}e^{MJ(y)}\frac{1+y^2}{y^4}.
\]

Since $d=M-2$, substitution and cancellation of all powers of $M$ yields
exactly \eqref{sc:pk:eq:2}.

Differentiation gives

\[
C'(y)=2\pi+2\arctan(1/y)-2/y.
\]

The unique minimizer is $y=1/v$, where $v>0$ is the unique
solution of

\[
v-\arctan v=\pi.
\]

At that point $C_{\min}=\log(4(1+v^2))$, so

\[
\boxed{
E_M^*\ge
\frac{16}{(1+v^2)(M-1)}[4(1+v^2)]^{-M}.
}
\tag{PK.3}\label{sc:pk:eq:3}
\]

The numerical values, only explanatory evaluations of the exact
formulas, are

\begin{align*}
v&=4.493409457909064\ldots,\\
C_{\min}&=4.439858114896194\ldots,\\
\frac{16}{1+v^2}&=0.755047187612980\ldots.
\end{align*}

Thus \(\limsup_{M\to\infty}(-\log E_M^*)/M\le4.439858114896194\ldots\).
The bound is not asserted sharp. A stronger lower bound matching the lcm
exponent, or a matching asymptotic of $E_M^*$, has not been obtained.

\subsection{Extremal zero-count theorem for the full log-balanced
cone}\label{sc:pk:extremal-zero-count-theorem-for-the-full-log-balanced-cone}

\textbf{Theorem.} Every extreme point of the feasible cone slice has a
numerator P whose zeros on {[}0,infinity), counted with multiplicity,
total at least

\[
M-\pi(M)-2.
\tag{PK.4}\label{sc:pk:eq:4}
\]

There exists an energy minimizer satisfying \eqref{sc:pk:eq:4}.

Equivalently an extreme minimizing numerator can be written

\[
P(u)=u^\epsilon A(u)^2 W(u),\qquad \epsilon\in\{0,1\},
\quad \deg W\le\pi(M),\quad W(u)>0\ (u\ge0),
\]

where all zeros of $A$ lie on $[0,\infty)$. Thus the part not accounted
for by a square and a possible factor u has degree at most pi(M), rather
than M. This is a necessary structure at some global minimizer; it does
not determine the zeros or prove an algorithm finds that minimizer.

\textbf{Proof.} Partial fractions identify $R$ with all polynomials of degree at most $d=M-2$. The affine constraints have rank $r=\pi(M)+1$: if $a_0+\sum_pa_pv_p(m)=0$ for every $m$, inserting $m=p$ gives $a_p=-a_0$, while $m=4$ gives $a_0=0$. Thus every $a_p$ is zero.

Factor $P=GW$, where $G$ consists of all roots on $[0,\infty)$, including the full multiplicity at zero; its positive roots have even multiplicity. Put $s=\deg G$. Then $G\ge0$ on $[0,\infty)$, and $W$ is strictly positive there and has degree $d-s$. For every polynomial $H$ of degree at most $d-s$, $|H|/W$ is bounded on $[0,\infty)$. Hence for sufficiently small $\eta>0$, both $G(W+\eta H)$ and $G(W-\eta H)$ remain nonnegative.

The space of perturbations $GH$ has dimension $d-s+1$. If it exceeded $r$, some nonzero such perturbation would preserve all $r$ affine constraints, contradicting extremality. Thus $s\ge d+1-r=M-\pi(M)-2$. The objective is linear and the
feasible slice is compact, so a minimizing face contains an extreme
point of the slice.

This theorem does not say all roots are at zero, nor that the remaining
coefficients are positive. It therefore does not reduce to the old
coefficientwise cone. Optimizing a monomial-positive numerator misses
the possible movable positive double roots forced by this full-cone
geometry.

\subsection{What the arithmetic budget does and does not
imply}\label{sc:pk:what-the-arithmetic-budget-does-and-does-not-imply}

At the stated target denominator $D_M=2L_M$, $\log D_M=M+o(M)$. Equation
\eqref{sc:pk:eq:3} implies that the guaranteed positive error cannot have an unbounded
exponent measured against this \textbf{unreduced prescribed budget}:

\[
\limsup\frac{-\log E(F_M)}{\log(2L_M)}\le4.439858114896194\ldots
\]

for any feasible sequence $F_M$. Thus a superexponential gain in $M$ or an
unbounded-exponent argument using that budget is ruled out. This is not
a bound in terms of the reduced denominator of the rational target:
large gcd cancellations are not controlled here.

Crucially, the constant $4.4398\ldots$ exceeds $1$. The bound is
compatible with $E_M^*$ much smaller than $1/(2L_M)$, so it does not
close the natural irrationality window. Even an upper bound
$E_M^*<1/(2L_M)$ alone would not prove the target
intersection

\[
E_M^*\le\frac{\{2L_M\gamma\}}{2L_M}.
\]

The fractional part could vanish in the rational gamma world. No claim
that positivity, compactness, or small unconstrained energy forces the
needed rational target is valid.

\section{Rational-center Gamma identities and the free first
cumulant}\label{sc:me:rational-center-gamma-identities-and-the-free-first-cumulant}

The one-center translation and finite-jet invariant-ring argument is a
reused input. The extension below considers all rational-center
reflection, shift, and multiplication identities simultaneously. Its
conclusion applies to the listed identities, not to a second global
Gamma function.

\subsection{Simultaneous rational-center jet
deformation}\label{sc:me:simultaneous-rational-center-jet-deformation}

For each positive rational \(q\), let

\[
G_q(z)=\Gamma(q+z).
\]

These are separately labelled convergent power-series germs. For a real
or complex parameter \(c\), define

\[
G_q^{[c]}(z)=e^{cz}G_q(z).
\tag{ME.2.1}\label{sc:me:eq:2.1}
\]

\begin{proposition}[Simultaneous rational-center deformation]\label{sc:me:deformation} The deformation \eqref{sc:me:eq:2.1}, with the same \(c\) at
every center, preserves all the following identities and all their
finite derivative consequences:

\[
G_{q+1}(z)=(q+z)G_q(z),
\tag{ME.2.2}\label{sc:me:eq:2.2}
\]

\[
G_q(z)G_{1-q}(-z)=\frac{\pi}{\sin\pi(q+z)}\qquad(0<q<1),
\tag{ME.2.3}\label{sc:me:eq:2.3}
\]

\[
G_1(z)G_1(-z)=\frac{\pi z}{\sin\pi z},
\tag{ME.2.4}\label{sc:me:eq:2.4}
\]

and, for every positive integer \(m\),

\[
G_{mq}(mz)
=(2\pi)^{(1-m)/2}m^{mq+mz-1/2}
\prod_{k=0}^{m-1}G_{q+k/m}(z).
\tag{ME.2.5}\label{sc:me:eq:2.5}
\]

It fixes every value \(\Gamma(q)\) and every polygamma value
\(\psi^{(r)}(q)\), \(r\ge1\), while replacing \(\psi(q)\) by
\(\psi(q)+c\). Consequently the designated first cumulant
\(-G_1'(0)=\gamma\) becomes \(\gamma-c\).

\end{proposition}
\emph{Proof.} In \eqref{sc:me:eq:2.2} both sides acquire \(e^{cz}\). In \eqref{sc:me:eq:2.3} and \eqref{sc:me:eq:2.4}
the two factors cancel. In \eqref{sc:me:eq:2.5} both sides acquire \(e^{cmz}\). Finally

\[
\log G_q^{[c]}(z)=cz+\log G_q(z),
\]

which proves all coefficient claims. These are identities of germs, so
comparison of any finite number of Taylor coefficients is valid. \(\square\)

This includes the usual digamma distribution identities. Their constant
terms do not isolate the common shift: the sum of \(m\) digamma shifts
agrees with the shifted term \(m\psi(mq)\). Higher polygamma
distribution identities contain no shift at all.

\subsubsection{Lower incomplete-Gamma data can remain
fixed}\label{sc:me:lower-incomplete-gamma-data-can-remain-fixed}

For a positive cutoff \(\alpha\), set

\[
L_{q,\alpha}(z)=\int_0^\alpha e^{-t}t^{q+z-1}\,dt,
\qquad
U_{q,\alpha}(z)=G_q(z)-L_{q,\alpha}(z).
\]

Keep every \(L_{q,\alpha}\) fixed, and put

\[
U_{q,\alpha}^{[c]}(z)=G_q^{[c]}(z)-L_{q,\alpha}(z).
\tag{ME.2.6}\label{sc:me:eq:2.6}
\]

Then the split identity, cutoff derivative, and shape recurrence remain
true:

\[
G_q^{[c]}=L_{q,\alpha}+U_{q,\alpha}^{[c]},
\quad
\partial_\alpha U_{q,\alpha}^{[c]}
=-e^{-\alpha}\alpha^{q+z-1},
\]

\[
U_{q+1,\alpha}^{[c]}
=(q+z)U_{q,\alpha}^{[c]}+\alpha^{q+z}e^{-\alpha}.
\tag{ME.2.7}\label{sc:me:eq:2.7}
\]

Indeed the lower recurrence is
\(L_{q+1,\alpha}=(q+z)L_{q,\alpha}-\alpha^{q+z}e^{-\alpha}\), and
subtraction from \eqref{sc:me:eq:2.2} gives \eqref{sc:me:eq:2.7}. In particular, every actual
algebraic relation or independence statement involving \textbf{only the
unchanged lower-jet constants} is preserved.

\subsubsection{Exact elimination
consequence}\label{sc:me:exact-elimination-consequence}

Let \(B\) be the subfield of \(\mathbb C\) generated by the fixed
constants occurring above: zero-order Gamma values, higher polygamma
values, \(\psi(q)+\gamma\), and lower jets, together with the elementary
constants in the identities. For rational \(q\), the combinations
\(\psi(q)+\gamma\) have the classical expressions in \(\pi\) and
logarithms of algebraic numbers; see
\href{https://dlmf.nist.gov/5.4}{DLMF 5.4}.

An explicit one-parameter model over \(B[g]\) is

\[
\mathcal G_q^{(g)}(z)=\Gamma(q)
\exp\left((\psi(q)+\gamma-g)z+
\sum_{j\ge2}\frac{\psi^{(j-1)}(q)}{j!}z^j\right).
\tag{ME.2.8}\label{sc:me:eq:2.8}
\]

At \(q=1\), this is

\[
\mathcal G_1^{(g)}(z)
=\exp\left(-gz+\sum_{j\ge2}\frac{(-1)^j\zeta(j)}jz^j\right).
\]

All identities \eqref{sc:me:eq:2.2}--\eqref{sc:me:eq:2.7} hold coefficientwise in this model, with
lower jets unchanged. The construction may be specialized to
\textbf{any} real algebraic \(g=a\), as well as to any real
transcendental \(g\). Thus polynomial elimination using exactly this
listed identity system cannot produce a nonzero polynomial constraint on
the free parameter \(g\), nor exclude all algebraic choices for it.

This is a statement about the specified system of identities. It does
not assert that all unknown arithmetic relations among actual Gamma and
upper incomplete-Gamma values are preserved.

\subsubsection{Essential scope restriction: the model is not a second
global Gamma
function}\label{sc:me:essential-scope-restriction-the-model-is-not-a-second-global-gamma-function}

The deformed centered germs need not glue to a single function. Where
the original germs overlap, their actual consistency is
\(G_q(z+r)=G_{q+r}(z)\); the deformation instead gives

\[
G_q^{[c]}(z+r)=e^{cr}G_{q+r}^{[c]}(z).
\tag{ME.2.9}\label{sc:me:eq:2.9}
\]

For real \(c\), compatibility at any nonzero rational translation
already forces \(c=0\). For complex \(c\), compatibility at all
sufficiently small rational translations also forces \(c=0\). Likewise
the actual upper-tail condition

\[
U_{q,\alpha}(z)\longrightarrow0\quad(\alpha\to\infty)
\]

is lost: \eqref{sc:me:eq:2.6} tends to \((e^{cz}-1)G_q(z)\).

Accordingly, this is \textbf{not} a countermodel to the definition of
the actual Gamma function or actual gamma. It proves that the selected
local algebraic jet identities alone leave a free mean. Global analytic
compatibility removes that freedom, but says nothing by itself about the
arithmetic nature of the selected mean.

\subsection{A complete positivity characterization of the remaining
parameter}\label{sc:me:a-complete-positivity-characterization-of-the-remaining-parameter}

The following result identifies exactly where analytic normalization
enters the generalized-moment approach.

For \(n\ge0\), write

\[
D_n=\int_0^\infty e^{-t}\log^n(1+t)\,dt,
\quad
G_n=(-1)^n\Gamma^{(n)}(1),
\]

\[
\eta_n=n!\sum_{k\ge1}\frac{(-1)^{k-1}}{k!\,k^n}.
\]

Here \(D_0=G_0=1\), \(\eta_0=1-e^{-1}\), and \(D_1\) is the
Euler--Gompertz constant. Splitting the defining Gamma integral at 1
gives

\[
G_n+(-1)^{n+1}\frac{D_n}{e}=\eta_n.
\tag{ME.3.1}\label{sc:me:eq:3.1}
\]

For a real parameter \(g\), define polynomials \(G_n(g)\) by

\[
M_g(t):=\sum_{n\ge0}G_n(g)\frac{t^n}{n!}
=\exp\left(gt+\sum_{j\ge2}\frac{\zeta(j)}j t^j\right)
=e^{(g-\gamma)t}\Gamma(1-t),
\tag{ME.3.2}\label{sc:me:eq:3.2}
\]

and set

\[
D_n(g)=e(-1)^{n+1}(\eta_n-G_n(g)).
\tag{ME.3.3}\label{sc:me:eq:3.3}
\]

\begin{theorem}[Positivity selects the actual first cumulant]\label{sc:me:positivity} For real \(g\), the following conditions are
equivalent:

\begin{enumerate}
\def\labelenumi{\arabic{enumi}.}
\tightlist
\item
  \(g=\gamma\).
\item
  \(D_n(g)\ge0\) for every \(n\ge1\).
\item
  \(D_n(g)/n!\to0\).
\end{enumerate}

If \(g\ne\gamma\), one parity of \(D_n(g)\) is strictly negative for all
sufficiently large indices of that parity.

\end{theorem}
\emph{Proof.} The series definition immediately gives

\[
\frac{\eta_n}{n!}=1+O(2^{-n}).
\tag{ME.3.4}\label{sc:me:eq:3.4}
\]

The meromorphic function in \eqref{sc:me:eq:3.2} has its nearest pole at \(t=1\), with
principal part \(e^{g-\gamma}/(1-t)\). After subtracting this principal
part, it is holomorphic in \(|t|<2\). Therefore, for each fixed
\(1<R<2\), Cauchy's estimate gives

\[
\frac{G_n(g)}{n!}=e^{g-\gamma}+O_{g,R}(R^{-n}).
\tag{ME.3.5}\label{sc:me:eq:3.5}
\]

Combining \eqref{sc:me:eq:3.3}--\eqref{sc:me:eq:3.5},

\[
\frac{D_n(g)}{n!}
=e(-1)^{n+1}\left(1-e^{g-\gamma}+O_{g,R}(R^{-n})\right).
\tag{ME.3.6}\label{sc:me:eq:3.6}
\]

For real \(g\ne\gamma\), the leading constant in parentheses is nonzero.
The claimed eventual negative parity follows, and convergence to zero is
impossible. At \(g=\gamma\), \eqref{sc:me:eq:3.1} identifies \(D_n(g)\) with the
strictly positive defining integral; \eqref{sc:me:eq:3.6} also gives the normalized
limit zero. \(\square\)

The theorem can be phrased without gamma in the construction: use the
first expression in \eqref{sc:me:eq:3.2}, the fixed zeta values, and the explicit
E-values \(\eta_n\). Gamma enters only to prove which real parameter the
resulting positivity condition selects. No novelty claim is made for
this elementary analytic characterization.

\subsubsection{Why this does not settle the algebraic
world}\label{sc:me:why-this-does-not-settle-the-algebraic-world}

A proof that every real algebraic \(a\) violates condition 2 would prove
gamma transcendental. But Theorem~\ref{sc:me:positivity} alone does not show this violation
for any unspecified algebraic \(a\): its eventual-sign argument requires
first knowing \(a\ne\gamma\).

Moreover, \textbf{every finite collection of strict moment-positivity
tests admits algebraic choices of the parameter}. The actual sequence
\(D_n\) is the moment sequence of the positive measure

\[
d\nu(y)=e^{1-e^y+y}\,dy\qquad(y>0).
\tag{ME.3.7}\label{sc:me:eq:3.7}
\]

Its finite Hankel and shifted Hankel matrices are positive definite,
because the integral of any nonzero polynomial square, also after
multiplication by \(y\), is strictly positive. Their entries \(D_n(g)\)
are polynomials in \(g\), so any finitely many such matrices remain
positive definite on a neighborhood of \(\gamma\). That neighborhood
contains rational, hence algebraic, parameters. The same applies to any
finite list of the individual strict inequalities \(D_n(g)>0\).

Full-moment positivity is even less restrictive: \eqref{sc:me:eq:3.2} is the
moment-generating function of a standard Gumbel random variable
translated by \(g-\gamma\). Thus its full Hamburger moment matrices are
positive for every real \(g\); they cannot select the mean at all.

\subsection{Differential transcendence and dimension counts do not
repair the
elimination}\label{sc:me:differential-transcendence-and-dimension-counts-do-not-repair-the-elimination}

Hölder's differential-transcendence theorem concerns the Gamma function,
not the arithmetic of a selected derivative value. Each individual germ
\(e^{cz}\Gamma(q+z)\) remains differentially transcendental over
\(\mathbb C(z)\): otherwise multiplying it by the differentially
algebraic function \(e^{-cz}\) would make \(\Gamma(q+z)\) differentially
algebraic. Hence this property also survives the free-mean deformation.
See the primary discussion in
\href{https://arxiv.org/abs/2303.02767}{Wang--Yao, Difference
independence of Euler gamma function}.

Likewise reflection expresses the even zeta cumulants through \(\pi\),
but does not constrain the first cumulant. Under the hypothesis that
gamma is algebraic, the identity

\[
G_2=\gamma^2+\zeta(2)
\]

simply makes \(G_2\) transcendental. Higher transcendence can reside in
the zeta cumulants and the upper incomplete-Gamma moments without
producing a contradiction at the first cumulant.

\section{Published approximation benchmarks and the primitive budget}
\label{sc:literature:benchmarks}

This section records the source review dated 20 September 2026. The
published input statements below are source-reviewed benchmarks, not
newly proved approximation theorems or newly measured numerical data.
They are retained to distinguish real convergence, integerization, and
primitive cancellation. The review was selective; it establishes no
classification of all possible approaches.

\subsection{The Van Assche--Wolfs benchmark}

The source review uses version 3 of Van Assche--Wolfs,
arXiv:2404.09799 (15 April 2025)~\cite{sc:VW}, Theorem~1.2 and
Propositions~2.5--2.6. Its normalization and approximation exponent
must not be mixed with those of version 2. Write
\[
d_n=\operatorname{lcm}(1,\ldots,n),\qquad
A_n=F^{(I)}_{n;1}(1),\qquad B_n=F^{(I)}_{n;2}(1)>0.
\]
These letters refer only to the published family in this subsection.
The quoted estimates and integrality are
\begin{align*}
\left|\gamma-\frac{A_n}{B_n}\right|
 &=O\!\left(e^{-4n^{3/4}+n^{1/2}+(3/8)n^{1/4}}\right),\\
B_n&\asymp n^{-9/8}
 e^{4n^{3/4}-(1/2)n^{1/2}-(3/8)n^{1/4}},\\
P_n&=n!d_nA_n\in\mathbb Z,\qquad
Q_n=n!d_nB_n\in\mathbb Z.
\end{align*}
Multiplying the first estimate by the quoted denominator gives
\begin{equation}\label{sc:literature:vw-error}
|Q_n\gamma-P_n|
\le Cn!d_n n^{-9/8}e^{\sqrt n/2}.
\end{equation}
This guaranteed upper bound does not tend to zero. It does not prove
that every actual residual diverges: the residuals may oscillate, and
an upper bound alone excludes no exceptional small subsequence.
The leading guaranteed quotient-error exponent relative to the
unreduced $Q_n$ is
$4/(n^{1/4}\log n)\to0$; in the linear-form convention the
corresponding exponent is this quantity minus one.
The repeated-differentiation discussion in the cited paper's
Section~5 does not supply a uniform theorem for a parameter increasing
with $n$.

\begin{lemma}[Exact coefficient identity used in the primitive audit]
\label{sc:literature:vw-identity}
For integers $n\ge1$, $k\ge0$,
\[
\sum_{\ell=1}^{n}\frac{(-1)^\ell}{\ell}
\binom{n+k}{n-\ell}
=-\binom{n+k}{k}(H_{n+k}-H_k).
\]
Consequently the published integer pair can be written
\begin{align*}
T_{n,k}&=\binom nk^2\binom{n+k}{k}\frac{n!}{k!},\\
Q_n&=d_n\sum_{k=0}^{n}T_{n,k},\\
P_n&=d_n\sum_{k=0}^{n}T_{n,k}
 (4H_k-2H_{n-k}-H_{n+k}).
\end{align*}
The latter formulas rearrange the cited definitions; they do not
provide a new approximating family.
\end{lemma}
\begin{proof}
Compute $D^n(x^{n+k}\log x)/n!$ in two ways. Differentiating
$D^n x^{n+k+s}/n!=\binom{n+k+s}{n}x^{k+s}$ with respect to $s$
at zero gives
\[
\binom{n+k}{n}x^k
 \bigl(\log x+H_{n+k}-H_k\bigr).
\]
Leibniz's rule gives the logarithmic term plus
\[
x^k\sum_{\ell=1}^n
 \frac{(-1)^{\ell-1}}{\ell}\binom{n+k}{n-\ell}.
\]
Comparison proves the identity. Substitution of that identity in the
published coefficient sums yields the displayed $P_n,Q_n$ formulas.
\end{proof}

For $g_n=\gcd(P_n,Q_n)$, the source's exact-integer computation reports
\begin{center}
\begin{tabular}{r r r r}
\toprule
$n$&$g_n$&$\log(Q_n/g_n)$&$\log g_n/(n\log n)$\\
\midrule
10&192&33.048606&0.228330\\
25&90&112.826729&0.055918\\
50&96&257.116922&0.023335\\
100&1692&562.817931&0.016142\\
200&220&1259.532706&0.005090\\
400&28836&2724.602499&0.004285\\
\bottomrule
\end{tabular}
\end{center}
The integer pairs and gcds were computed exactly; the logarithms are
rounded diagnostics. These six values are not an asymptotic theorem.
Neither $\log g_n=o(n\log n)$ nor $\log g_n=O(n)$ is established.
Conversely, if an infinite subsequence has certified nonzero residuals
and, for some fixed $\varepsilon>0$,
\[
g_n\ge n!d_n e^{(1/2+\varepsilon)\sqrt n},
\]
then dividing \eqref{sc:literature:vw-error} by $g_n$ produces nonzero
integer linear forms tending to zero. This is a sufficient
factorial-scale saving, not a necessary condition and not an obtained
lower bound.

\subsection{The conditions on external value theorems remain necessary}

The source audit distinguishes the Hessami Pilehrood $k!$ sum used
above from Aptekarev's original
$\widetilde q_n=\sum_k\binom nk^2(n+k)!$ family.
The identity $q_n=n!L_n(-1)$ is already explicit in the cited Hessami
Pilehrood paper. It is not a new structural discovery.
The source-reviewed 2026 survey by Aptekarev and the geometric-series
paper by Burnol likewise supply no quoted irrationality theorem for
$\gamma$. In particular, Burnol's fixed-parameter geometric expansion
approximates a combination $\gamma+(\ell-1)\log2$; geometric
convergence alone supplies no nonzero integer form in $\gamma$ alone.

Fischler--Rivoal's \emph{Rational approximations to values of
$E$-functions}, Theorem~1 in the manuscript dated 3 July 2025~\cite{FRrational}, is quoted
for irrational values at rational points of $E$-functions with
\emph{rational} Taylor coefficients. Under the stronger hypothesis
$\gamma\in\mathbb Q$, one has
\[
F_\gamma(z)=\sum_{n\ge0}(H_n-\gamma)\frac{z^n}{n!}
\in\mathbb Q[[z]],\qquad F_\gamma(1)=\delta.
\]
The value is transcendental under that hypothesis by the
$E$-function specialization input, and the quoted approximation
result then gives $\mu(\delta)=2$. This is a conditional consequence,
not a contradiction. The rational-coefficient theorem cannot simply
be applied under only $\gamma\in\overline{\mathbb Q}\setminus\mathbb Q$.

The source-reviewed 8 September 2026 Fischler--Rivoal measure theorem~\cite{sc:FRBessel} for the logarithmic derivative of $J_0$ requires a permitted value
representation. None expressing $\gamma$ as such a value is given.
In particular,
$z(F_r'-F_r)=e^z-1$ does not have the polynomial right-hand side
required by the quoted first-order inhomogeneous alternative.
Similarly, the source-reviewed 17 September 2026 Rivoal arithmetic-difference-operator theorem~\cite{sc:RivoalSigma} requires its Mellin-associated
operator to be a $G$-operator. The theorem's input conditions have
not been verified for the gamma approximants at issue. Neither result
can be applied solely because a differential or difference equation
has been written down.

Finally, the source review of Powers, arXiv:2511.01849v7~\cite{sc:Powers}, reports
family-level transcendence results for Gamma derivatives, generalized
eta constants, and generalized Gompertz constants. Its stated
proportion bound for transcendental Gamma derivatives is of order
$N^{-1/2}$, not a positive-density assertion. Those quoted results
are not used in the elementary rational-center deformation or
positivity proofs above. For the generalized constants discussed in that source,
$\eta_0=1-e^{-1}$, so independence over algebraic numbers must not
be upgraded to independence of the entire indexed family over a field
already containing $e$.

\section{Arithmetic reassessment: primitive targets and limits}\label{sc:au:arithmetic-reassessment-primitive-targets-and-scoped-obstructions}

\subsection{A Dirichlet existence theorem is not an upper
ceiling}\label{sc:au:a-dirichlet-existence-theorem-is-not-an-upper-ceiling}

The claim that every triple of numbers satisfies an approximation
ceiling \(\tau\le2\) is false. Dirichlet's principle supplies good
approximations; it does not prohibit much better approximations.

An elementary explicit counterexample is enough. Let

\[
\ell=\sum_{k\ge1}10^{-k!},\qquad
(x_0,x_1,x_2)=(1,\sqrt2,\sqrt2+\ell).
\]

This triple is linearly independent over $\mathbb Q$, since a nontrivial relation
would make the Liouville number $\ell$ algebraic. Put \(q_m=10^{m!}\) and
let \(p_m/q_m\) be its m-th partial sum. Then

\[
0<|-p_m-q_m\sqrt2+q_m(\sqrt2+\ell)|
=|q_m\ell-p_m|<2q_m^{-m}.
\]

All three integer coefficients are nonzero. Thus no universal exponent-2
ceiling exists even with linear independence and all coordinates active.
This does not establish any superior exponent for the actual
gamma/Gompertz triple; it invalidates an asserted universal ban.

This correction was already acknowledged in the older integrated
archive; it is retained here to prevent reuse of the invalid universal
ceiling, without a novelty claim.

\subsection{Primitive elimination has two distinct height
effects}\label{sc:au:primitive-elimination-has-two-distinct-height-effects}

Suppose

\[
R_i=a_i+b_i\gamma+c_i\delta,\quad
|a_i|,|b_i|,|c_i|\le H,\quad |R_i|\le H^{-\tau}.
\]

Eliminating delta gives

\[
c_2R_1-c_1R_2=A+B\gamma,
\quad |A|,|B|\le2H^2,
\quad |A+B\gamma|\le2H^{1-\tau}.
\]

If $(A,B)$ is nonzero and $g=\gcd(A,B)$, the primitive output satisfies

\[
h\le2H^2/g,\qquad
|A/g+(B/g)\gamma|\le2H^{1-\tau}/g.
\]

Consequently the frequently quoted exponent \((\tau-1)/2\) describes
only a raw generic budget. If \(g=H^{\sigma+o(1)}\), with \(\sigma<2\)
and \(\tau+\sigma-1>0\), the displayed estimates certify, up to
constants and o(1), an exponent at least

\[
\frac{\tau+\sigma-1}{2-\sigma}
\]

against the primitive height. Actual height drops or better error
cancellation can improve this further. Neither such gcd growth nor
nonvanishing is automatic.

This also does not defeat the separate fact that a restricted collection
of E-function expressions might involve gamma and delta only through one
fixed combination: eliminating that combination could eliminate gamma at
the same time. These are different obstructions and require different
proofs.

The archived Sondow calculation retains both circuit saturation and
final target gcd. Its exact quotient norm is invariant under positive
rational rescaling of columns. Therefore a large unnormalized quotient
index, a large common denominator, or a huge raw coefficient vector is
not by itself a primitive-height impossibility theorem.

The source audit reports one archived finite, exactly certified
inequality

\[
0<p-q\gamma<q^{-101/100}
\]

at a 37-digit denominator. This is compatible with gamma being
algebraic. It neither produces an infinite family nor proves eventual
denominator growth.

\subsection{Gevrey and window-width estimates need their
hypotheses}\label{sc:au:gevrey-and-window-width-estimates-need-their-hypotheses}

Subexponential convergence for fixed-pole approximants, or factorial
costs for one moving-pole construction, does not exclude all
varying-degree, mixed, nonlinear, or differently normalized
constructions. A rigorous obstruction must specify its input ring,
coefficient normalization, degree growth, denominator support, and the
permitted transformations.

The earlier common-lcm window budget giving \(w<0.386\) is presentation
dependent. In the integerized Sondow family one may instead use each
column's own multiplier \(d_{2n}\). On $[N,2N]$, the old counting
argument already yields exponentially small upper bounds with exponent
\(\log16-2>0\). The difficulty is not the absolute impossibility of such
a wide window. It is that all the short vectors supplied by counting can
map to exact zero.

Likewise, a statement that only two families in all the literature
qualify is not established by inspecting two families. It should be
replaced by an explicitly delimited list of families actually analyzed.

\subsection{The exact Sondow quotient, including both algebraic
worlds}\label{sc:au:the-exact-sondow-quotient-including-both-algebraic-worlds}

Use the established identity

\[
e_n=d_{2n}I_n=a_n\gamma-b_n+\sum_p\beta_{p,n}\log p,
\qquad a_n,b_n,\beta_{p,n}\in\mathbb Z,
\]

where $I_n$ is Sondow's positive double integral and
$d_m=\operatorname{lcm}(1,\ldots,m)$. On a finite block define

\[
\Lambda=\ker_{\mathbb Z}B,\quad
\Phi(u)=(a\cdot u,b\cdot u),\quad
K_0=\ker_{\mathbb Z}[B;a;b],\quad
\Omega=\Phi(\Lambda)\subseteq\mathbb Z^2.
\]

The first isomorphism theorem gives, unconditionally,

\[
\Lambda/K_0\simeq\Omega,\qquad
\operatorname{rank}\Omega\le2.
\]

This quotient is torsion free. The evaluation of a target $(Q,P)$ is $Q\gamma-P$.

\textbf{If gamma is irrational}, whether algebraic or transcendental,
its evaluation kernel on $\Lambda$ is exactly $K_0$. Indeed, $Q\gamma=P$
with integers $Q,P$ forces $Q=P=0$.

\textbf{If $\gamma=a/b$ is rational in lowest terms}, the additional
evaluation kernel is

\[
\ker(u\mapsto\textstyle\sum u_ne_n)/K_0
\simeq\Omega\cap\mathbb Z(b,a).
\]

If $\operatorname{rank}\Omega=2$ this is a rank-one lattice. For every vector whose
certified error is less than $1/b$, the image must lie on that rational
line. It need not be the zero pair.

Thus the statement ``algebraic gamma forces all sufficiently small
linear forms to vanish'' is wrong without a quantitative,
degree-dependent height threshold. Rationality has a fixed spacing $1/b$.
Algebraic irrationality has nonzero linear forms and height-dependent
lower bounds. The previously established polynomial-collapse result
concerns a different family $P_D$ with very strong degree/height/error
inequalities; it must not be transferred silently to these linear forms.

\subsubsection{Why the existing counting theorem does not survive the
quotient}\label{sc:au:why-the-existing-counting-theorem-does-not-survive-the-quotient}

On $[N,2N]$ there are $N+1$ columns and $O(N/\log N)$ prime rows, all entries $\exp(O(N))$. The elementary counting proof of Siegel's lemma
yields nonzero $u\in\Lambda$ with

\[
\log\|u\|_\infty=O(N/\log N),\qquad
\left|\sum u_ne_n\right|
\le\exp[-(\log16-2)N+o(N)].
\]

Appending the two rows $a$ and $b$ changes neither asymptotic. The same
theorem therefore yields nonzero $u\in K_0$ at exactly this scale. This
is a proven mechanism for producing short useless vectors, not merely a
hypothetical failure.

One can record the rank of the span of the images of short vectors; it
is 0, 1, or 2. Naming that statistic is not an advance. Proving an
all-level lower bound on it with the required representative error would
be an advance. In particular, proving rank 2 for two errors tending to
zero would prove irrationality. A rank-1 family can also suffice if its
primitive denominators provably tend to infinity while errors tend to
zero.

The stronger assertion that short first-coordinate-dominant vectors
exist is a sufficient condition. The converse ``gamma irrational implies
this particular dominant family exists'' has not been proved, so the
condition should not be called equivalent to irrationality.

\subsubsection{Rank two by itself is inexpensive and
insufficient}\label{sc:au:rank-two-by-itself-is-inexpensive-and-insufficient}

The source audit reports that the archived exact \(n=1,\ldots,8\)
calculation gives two rational-kernel vectors with target determinant

\[
\frac{258848000709925137756895171}{3003}\ne0
\]

before column integerization. Clearing denominators transfers them to
the integerized lattice without destroying independence. Padding these
vectors by zeros proves $\operatorname{rank}\Omega=2$ for every prefix $1,\ldots,M$ with $M\ge8$. This elementary observation is not a tail estimate:
both witnesses can remain supported in the first eight columns, so their
errors need not decrease.

It does not prove rank two with small representatives in moving windows
$[N,2N]$. Nor does a finite moving-window rank computation prove an
asymptotic estimate.

\subsection{A successful super-Roth sequence must allow sparse
scales}\label{sc:au:a-successful-super-roth-sequence-must-allow-sparse-scales}

For any real x, let p/q and p'/q' be distinct reduced fractions, with
\(q'\ge q>0\), both satisfying

\[
|qx-p|\le q^{-1-\eta},\qquad
|q'x-p'|\le(q')^{-1-\eta},\qquad\eta>0.
\]

Their nonzero integer determinant gives

\[
1\le |qp'-q'p|
\le q' q^{-1-\eta}+q(q')^{-1-\eta}
\le2q'q^{-1-\eta}.
\]

Hence

\[
q'\ge\tfrac12q^{1+\eta}.
\]

This is an unconditional restriction, including for transcendental x. A
proposed construction cannot produce distinct such approximants at every
consecutive level with \(\log q_N=cN+o(N)\), \(c>0\). It must allow
repeated outputs or sufficiently sparse new levels. Requiring a uniform
``new super-Roth target at every nearby height'' would wrongly reject
every real number, not just an algebraic gamma world.

For two independent target pairs with both denominators \textless=H and
both errors \textless=E, the simpler restriction is \(1\le2HE\). Thus no
proof should demand two independent super-Roth approximants at one
common comparable height.

\paragraph{The resulting status of the central arithmetic budget.}
For fixed algebraic $r$, the moving-degree theorem proves
$\log(d_NH(P_N))\ge N-o(N)$ when
$D_N^6\log(D_N+2)=o(N)$ in the $z$-independent case. Its leading
constant one is attained by $P_N=Y$ for algebraic-integral $r$.
This is a coefficient denominator-and-height theorem. It is compatible
with $F_r$ being an $E$-function, permits no automatic inference about
cancellation after evaluation at a point, and leaves fast-degree,
rational-function, and different-representation constructions outside
its stated scope. Substitution $r=\gamma$ under an algebraicity
hypothesis supplies no contradiction.

The other completed branches constrain fixed shifts, full positive
Binet kernels, and a specified local Gamma identity system. None
constructs an infinite family with both certified nonzero arithmetic
output and sufficiently small error relative to primitive height.
Historical priority is unestablished; internal audits and finite exact
checks do not constitute external peer review. These results do not
prove irrationality or transcendence of $\gamma$ and do not establish
the later rank-four target theorem.

% END INTEGRATED SOURCE

\clearpage
\part{The four-value problem and cross-field tests}

% BEGIN INTEGRATED SOURCE: rank_cross/body.tex
\section{The four-value rank and its exceptional points}\label{rc:sec:rank}

This continuation separates unconditional deductions from the remaining
specialization problem.  Throughout this part, put $k=\rck$ and use the
principal logarithm on $\rcC\setminus(-\infty,0]$.  Our normalization is
\begin{align}
 \rcEin(z)&=\sum_{n\ge1}\frac{(-1)^{n-1}z^n}{n\,n!},&
 E(z)&=e^z,&F(z)&=e^z\rcEin(z),\label{rc:eq:normalization}\\
 E_1(x)&=\int_x^\infty\frac{e^{-t}}t\,dt,&
 D(x)&=e^xE_1(x)=\int_0^\infty\frac{e^{-t}}{x+t}\,dt
 \quad(x>0).\nonumber
\end{align}
The connection formula and the distinguished values are
\begin{equation}\label{rc:eq:connection}
 \rcEin(x)=\gamma+\log x+E_1(x),\qquad
 \eta=\rcEin(1),\qquad \delta=D(1)=e(\eta-\gamma).
\end{equation}
For positive algebraic $x$, define
\begin{equation}\label{rc:eq:rho}
 \rho(x)=\dim_k\rcSpan_k\{1,e^x,D(x),F(x)\}.
\end{equation}
The designated assertion $\rho(1)=4$ is not proved here.  It would imply
the transcendence of both $\gamma$ and $\delta$: algebraicity of the former
would make $\delta-e\eta+\gamma e=0$ a nonzero algebraic-coefficient
relation, and algebraicity of the latter would already relate $1$ and
$\delta$.  The converse implication is unavailable, as the exact
classification below makes particularly clear.

\subsection{A finite E-function input with two functional proofs}

\begin{lemma}\label{rc:lem:evalues}
If $x_1,\ldots,x_m$ are distinct positive algebraic numbers, then
\[
 1,\ e^{-x_1},\ldots,e^{-x_m},\
 \rcEin(x_1),\ldots,\rcEin(x_m)
\]
are linearly independent over $k$.
\end{lemma}

\begin{proof}
Set $X_i(z)=e^{-x_i z}$ and $U_i(z)=\rcEin(x_i z)$.  These are
E-functions and satisfy the rational differential system
\begin{equation}\label{rc:eq:finite-esystem}
 X_i'=-x_iX_i,\qquad U_i'=\frac{1-X_i}{z},\qquad 1'=0.
\end{equation}
For clarity, the normalized Taylor coefficients of $U_i$ are
$(-1)^{n-1}x_i^n/n$.  In a fixed number field all their conjugates have
at most exponential growth, and clearing the coefficients through $n$
requires, in addition to an exponential denominator for powers of $x_i$,
only $\operatorname{lcm}(1,\ldots,n)$, which is exponentially bounded.
The differential equations give holonomicity.  Thus the arithmetic
E-function hypotheses hold, rather than merely entire-function growth.

First $1,X_1,\ldots,X_m$ are independent over $\rcC(z)$.  Clear rational
denominators in a relation, order the positive $x_i$, and let $z$ tend to
negative real infinity.  The largest exponential with a nonzero polynomial
coefficient dominates all smaller ones, a contradiction; after removing
all exponentials the polynomial constant term also vanishes.

Suppose a rational relation involving $U_i$ exists, and choose one with
the smallest positive number of nonzero $U_i$ coefficients.  Normalize
one of them to $1$ and write
\[
 \sum_{i\in I}R_iU_i=R_0+\sum_{j=1}^m S_jX_j,
 \qquad R_i,R_0,S_j\in\rcC(z).
\]
Differentiation, followed by \eqref{rc:eq:finite-esystem}, produces a
relation with fewer $U_i$ terms unless every $R_i'$ vanishes.  Minimality
forces $R_i=c_i\in\rcC$.  Independence of the exponentials now gives
\[
 S_i'-x_iS_i=-c_i/z.
\]
For $c_i\ne0$ this has no rational solution: a pole of $S_i$ at zero
would yield a pole of strictly higher order in $S_i'$, uncancelled by
$-x_iS_i$, whereas if $S_i$ is regular there the left side cannot have a
simple pole.  Thus every $c_i$ is zero, contradicting the normalization.
The $2m+1$ functions are consequently independent over $\rcC(z)$.

The system has no finite singularity outside zero.  Apply the linear
ordinary-point case of Beukers' refinement of the Siegel--Shidlovskii
theorem at $z=1$ \cite[Theorem 1.3]{Beukers}.  It gives the asserted
independence of values.  This step uses an established theorem for
E-functions; it does not use a mixed Gevrey specialization conjecture.
\end{proof}

The following second functional proof records the coefficient audit of
the same input.  Suppose, after clearing denominators, that
\begin{gather*}
 P_0(z)+\sum_iP_i(z)e^{-x_i z}+\sum_iQ_i(z)\rcEin(x_i z)=0,\\
 P_i=\sum_jp_{ij}z^j,\qquad Q_i=\sum_jq_{ij}z^j.
\end{gather*}
For $n$ larger than all degrees, the normalized Taylor coefficients give
\[
 \sum_i(-x_i)^nR_i(n)=0.
\]
Here, with $(T)_j=T(T-1)\cdots(T-j+1)$ and $(T)_0=1$, we put
\begin{equation}\label{rc:eq:coefficient-audit}
 R_i(T)=\sum_jp_{ij}(-x_i)^{-j}(T)_j
       -\sum_jq_{ij}(-x_i)^{-j}\frac{(T)_j}{T-j}.
\end{equation}
If a nonzero $R_i$ with largest $x_i$ were present, division by its
exponential term would leave a nonzero rational function, of polynomial
asymptotic size, equal to exponentially small terms.  Hence every $R_i$
vanishes.  Its residue at $T=j$ is exactly
$-q_{ij}(-x_i)^{-j}j!$, because all other summands are regular there.
Thus all $q_{ij}$ vanish; independence of the falling-factorial basis
then gives all $p_{ij}=0$, followed by $P_0=0$.  Beukers' theorem completes
the specialization as before.

\subsection{The unique real triple and the complete rank-three alternatives}

\begin{proposition}\label{rc:prop:bad-representation}
For positive algebraic $x$ one has $\rho(x)\in\{3,4\}$.  Moreover,
\begin{equation}\label{rc:eq:bad-representation}
 \rho(x)=3\quad\Longleftrightarrow\quad
 \gamma+\log x=a_x+b_xe^{-x}+c_x\rcEin(x)
 \quad(a_x,b_x,c_x\in k).
\end{equation}
When it exists, this triple is unique and belongs to $(k\cap\rcR)^3$.
\end{proposition}

\begin{proof}
Multiplication of the $m=1$ case of Lemma~\ref{rc:lem:evalues} by
$e^x$ proves independence of $1,e^x,F(x)$, so the rank is at least three.
Any relation on the four values must therefore have a nonzero coefficient
of $D(x)$.  Solving for it, dividing by $e^x$, and using
\eqref{rc:eq:connection} gives \eqref{rc:eq:bad-representation}, and the
steps are reversible.  Subtraction of two triples gives a relation on
$1,e^{-x},\rcEin(x)$, proving uniqueness.  Complex conjugation fixes
all three values and the left side, so uniqueness also forces every
coefficient to be real.
\end{proof}

We use, for the next stronger conclusion, the established algebraic
independence of $e$ and $\eta$ over $k$.  One primary source is
Rivoal~\cite[\S5(ii)]{rc:rivoal}: the series denoted there by $E(-1)$
is $-\rcEin(1)$, and replacing $e^{-1}$ by its inverse preserves algebraic
independence.  This input is stronger than Lemma~\ref{rc:lem:evalues}.

\begin{theorem}[Complete classification at rank three]\label{rc:thm:rank3}
Assume $\rho(1)=3$, and let $(a,b,c)\in(k\cap\rcR)^3$ be its unique
triple.  Then
\begin{equation}\label{rc:eq:rank3-pair}
 \gamma=a+b/e+c\eta,\qquad
 \delta=(1-c)e\eta-ae-b.
\end{equation}
Exactly the following disjoint alternatives occur:
\begin{enumerate}
 \item If $b=c=0$, then $\gamma=a\in k$ and
 $k(e,\delta)=k(e,\eta)$; hence $e,\delta$ are algebraically independent.
 \item If $a=0$ and $c=1$, then $\delta=-b\in k$ and
 $k(e,\gamma)=k(e,\eta)$; hence $e,\gamma$ are algebraically independent.
 \item In all other cases,
 $k(\gamma,\delta)=k(e,\eta)$, and $\gamma,\delta$ are algebraically
 independent.
\end{enumerate}
Furthermore, under the rank-three assumption, $e,\gamma$ are algebraically
independent exactly when $c\ne0$, and $e,\delta$ are algebraically
independent exactly when $c\ne1$.
\end{theorem}

\begin{proof}
The pair of equations is Proposition~\ref{rc:prop:bad-representation}
at $x=1$ and \eqref{rc:eq:connection}.  If $\gamma\in k$, subtract it
from the first equation and use the independence of $1,e^{-1},\eta$;
this forces $b=c=0$.  The converse is immediate.  In this case
$\delta=e(\eta-a)$, proving the first field equality.  Similarly,
the independence of $1,e,e\eta$ applied to the second equation gives
$\delta\in k$ if and only if $a=0,c=1$.  Then
$\gamma=\eta+b/e$, proving the second field equality.  The two cases
are disjoint.

For the remaining cases eliminate $\eta$ using
$\eta=\gamma+\delta/e$ to obtain
\begin{equation}\label{rc:eq:birational}
 e\bigl((1-c)\gamma-a\bigr)=c\delta+b.
\end{equation}
The denominator $(1-c)\gamma-a$ can vanish only in the two cases already
listed.  Indeed, if $c=1$ its vanishing means $a=0$; if $c\ne1$ it
makes $\gamma=a/(1-c)$ algebraic, and the first case applies.  Outside
these cases the exact inverse transformation is
\[
 e=\frac{c\delta+b}{(1-c)\gamma-a},\qquad
 \eta=\gamma+\delta/e.
\]
Together with \eqref{rc:eq:rank3-pair}, it proves equality of the two
fields.  Their transcendence degree is two, proving the third alternative.
Finally, the first equation of \eqref{rc:eq:rank3-pair} solves for $\eta$
from $(e,\gamma)$ when $c\ne0$, and puts $\gamma$ in $k(e)$ when $c=0$.
The second equation has the analogous properties for $(e,\delta)$
according as $c\ne1$ or $c=1$.
\end{proof}

No triple has been excluded for the actual constants by this
classification.  In particular, obtaining transcendence or even algebraic
independence of $\gamma,\delta$ under a generic rank-three hypothesis is
not a contradiction.  To isolate the logical issue, the artificial
numbers $g=\eta/2$ and $d=e\eta/2$ satisfy $\eta=g+d/e$ and are
algebraically independent, but $\{1,e,d,e\eta\}$ has rank three.
They are not the Euler and Gompertz constants.  The classification also
implies the conditional statement that if the actual $\gamma,\delta$
were both transcendental and algebraically dependent, then $\rho(1)=4$;
no such dependence is established here.

\subsection{Affine multiplicative restrictions on all deficient points}

\begin{theorem}\label{rc:thm:badset}
Let $B=\{x\in k\cap\rcR_{>0}:\rho(x)=3\}$.  For distinct
$x_1,\ldots,x_s\in B$, no nonzero integer vector $(m_1,\ldots,m_s)$
satisfies both
\begin{equation}\label{rc:eq:circuit}
 \sum_i m_i=0,\qquad \prod_i x_i^{m_i}=1.
\end{equation}
Consequently, if the $\rcQ$-span of $\{\log x:x\in G\}$ has finite
dimension $d$ for a multiplicative group $G\subset k\cap\rcR_{>0}$,
then $\#(B\cap G)\le d+1$.
\end{theorem}

\begin{proof}
Multiply the representation \eqref{rc:eq:bad-representation} for $x_i$
by $m_i$ and add.  Since all logarithms are real,
\eqref{rc:eq:circuit} cancels both $\gamma$ and the logarithms, leaving
\[
 0=\sum_i m_i a_i+\sum_i m_i b_i e^{-x_i}
                         +\sum_i m_i c_i\rcEin(x_i).
\]
Lemma~\ref{rc:lem:evalues} gives $b_i=c_i=0$ whenever $m_i\ne0$.
Thus $\gamma+\log x_i=a_i\in k$ on the support.  The support contains
two distinct points, since a single nonzero coefficient cannot have sum
zero.  Subtracting their equations makes $\log(x_i/x_j)$ a nonzero
algebraic number.  The Hermite--Lindemann theorem would then make
$x_i/x_j$ transcendental, contrary to its algebraicity.

For the cardinality bound, $d+2$ vectors $(1,\log x_i)$ in a rational
space of dimension $d+1$ have a nontrivial rational linear relation.
Clear denominators and exponentiate its logarithmic component to obtain
\eqref{rc:eq:circuit}.  Finite generation of $G$ is not required.
\end{proof}

\begin{corollary}\label{rc:cor:geometric}
For every positive algebraic $r\ne1$,
\[
 \#\{n\in\rcZ:\rho(r^n)=3\}\le2.
\]
In particular $\max\{\rho(1),\rho(2),\rho(4)\}=4$, and among
$1,\ldots,N$ at most $\pi(N)+1$ points have rank three.  The positive
integer points of rank four therefore have natural density one.
\end{corollary}

\begin{proof}
The logarithms of $r^n$ span a one-dimensional rational space.  More
explicitly, for distinct $n_1,n_2,n_3$ the nonzero coefficients
$(n_2-n_3,n_3-n_1,n_1-n_2)$ annihilate both $1$ and $n$, so three such
points contradict Theorem~\ref{rc:thm:badset}.  For $1,2,4$ the coefficients
are $(1,-2,1)$.  The integers up to $N$ lie in the positive rational
group generated by primes up to $N$, whose rank is $\pi(N)$.
The final conclusion uses the classical fact $\pi(N)=o(N)$.
\end{proof}

At each full-rank point both $\gamma+\log x$ and $D(x)$ are
transcendental, since algebraicity of either would give a dependence in
\eqref{rc:eq:rho}.  None of these bounds singles out $x=1$.  They also
do not constitute an independent new tail-transcendence mechanism:
the earlier finite-rank tail-defect arguments already yield related
cofinite conclusions for $E_1$ values.  The precise object here is the
set of failures of the four-value rank, with its affine multiplicative
restriction.  No priority claim is made for this short consequence of
established E-function theory.

\section{Functional rank, mixed division, and the exact missing implication}
\label{rc:sec:mixed}

\subsection{A differential system does not determine value rank}

\begin{proposition}\label{rc:prop:functionalrank}
The functions $1,E,F,D$ have rank four over $\rcC(z)$ and satisfy
\begin{equation}\label{rc:eq:four-system}
 1'=0,\qquad E'=E,\qquad
 F'=F+(E-1)/z,\qquad D'=D-1/z.
\end{equation}
The only finite singularity of this displayed system is zero.
\end{proposition}

\begin{proof}
The equations follow by differentiation.  If
$A+BE+CF+JD=0$ with rational functions $A,B,C,J$, continue once
counterclockwise around zero.  Only $D$ changes, by $-2\pi iE$.
Subtracting the two relations gives $J=0$.  If $C\ne0$, divide by $CE$
to write $\rcEin=R+Se^{-z}$ with rational $R,S$.  Differentiation and
independence of $1,e^{-z}$ imply $R'=1/z$ and $S'-S=-1/z$.
A rational derivative has zero residue at every pole, so $R'=1/z$ is
impossible.  Hence $C=0$, and $A+BE=0$ forces $A=B=0$.
\end{proof}

For any $s\in\rcC$, replacement of $D$ by $D_s=D+sE$ preserves the
system, the functional rank, and the monodromy.  Taking $s=\gamma$ gives
\begin{equation}\label{rc:eq:shift-countermodel}
 D_\gamma=F-E\log z,\qquad D_\gamma(1)=F(1),
\end{equation}
so the values have rank exactly three although $1$ is still an ordinary
point.  This is a countermodel to an inference from these differential
facts alone.  It does not retain the prescribed positive-ray asymptotic
normalization: $D(z)\sim1/z$, whereas $D_s(z)$ has the additional term
$se^z$ for $s\ne0$.  Nor is the arithmetic admissibility of the shift
by $\gamma$ known.  The missing information concerns the arithmetic of
the canonical solution, not regularity of the evaluation point.

\subsection{Formal arithmetic nondivisibility}

Use the formal series
\[
 \widehat D(z)=\sum_{n\ge0}(-1)^n n!z^{-n-1}.
\]
Its positive-direction Borel sum is $D$.  We employ the normalized mixed
arithmetic space $\rcM$ consisting of expressions
$H(z)+h(1/z)$, where $H$ is an E-function, $h(w)=\sum_{n\ge1}h_nw^n$
has algebraic coefficients, and its coefficientwise Borel transform
$\mathcal B h(t)=\sum_{n\ge1}h_nt^n/n!$ is a G-function.  Thus the
negative part has no constant term.  Arithmetic Gevrey order one here
includes holonomicity, exponential conjugate bounds and exponential
common-denominator bounds after Borel transformation; analytic factorial
growth alone is insufficient.  This is the mixed setting of
\cite[Definition 1 and \S2]{FRmixed}.

\begin{theorem}\label{rc:thm:nondivision}
For $a,b,c,d\in k$, put
$\Psi=a+bE+cF+d\widehat D$.  Then
\begin{equation}\label{rc:eq:nondivision}
 \Psi\in(z-1)\rcM\quad\Longrightarrow\quad a=b=c=d=0.
\end{equation}
\end{theorem}

\begin{proof}
Suppose $\Psi=(z-1)(H+\sum_{n\ge1}h_nz^{-n})$ in the stated arithmetic
class.  Comparison of nonnegative powers gives
\[
 a+bE+cF=(z-1)H+h_1.
\]
Since $H$ is entire, substitution of $z=1$ gives
$a+be+ce\eta=h_1\in k$.  The independence of $1,e,e\eta$ gives
$b=c=0$ and $h_1=a$, and then $H=0$.  Comparison of the negative powers
now yields
\[
 h_{n+1}-h_n=d(-1)^{n-1}(n-1)!\quad(n\ge1).
\]
With $S_n=\sum_{j=0}^n(-1)^jj!$ this means
$h_{n+2}=a+dS_n$.  If $\mathcal Bh$ is a G-function, so is its second
derivative
\begin{equation}\label{rc:eq:B-ad}
 B_{a,d}(t)=\sum_{n\ge0}\frac{a+dS_n}{n!}t^n,
 \qquad B_{a,d}'-B_{a,d}=-\frac d{(1+t)^2},
 \qquad B_{a,d}(0)=a+d.
\end{equation}
The coefficient recurrence verifies the equation directly.  If $d\ne0$,
the integrating-factor equation
$(e^{-t}B_{a,d})'=-de^{-t}/(1+t)^2$ has residue $de$ at $t=-1$.
Continuation once around $-1$ consequently changes $B_{a,d}$ by
$2\pi i d e^{t+1}$.

A G-function has a minimal homogeneous differential equation regular
singular at every point of the projective line, including infinity
\cite[\S1]{FRoperators}.  Analytically continued
branches solve the same equation, as does their difference.  A nonzero
multiple of $e^t$ cannot be such a solution: solutions of a regular
singular equation at infinity have power-logarithmic, rather than
exponential, growth on fixed sectors.  This contradiction excludes
$d\ne0$.  If $d=0$, equation \eqref{rc:eq:B-ad} gives $B_{a,0}=ae^t$,
which violates the same condition unless $a=0$.  All coefficients vanish.
\end{proof}

One useful special case makes the index conventions explicit.  For every
algebraic $\alpha$,
\begin{align}
 \frac{\widehat D(z)-\alpha}{z-1}
 &=\sum_{m\ge1}b_mz^{-m},&
 b_m&=-\alpha+\sum_{j=0}^{m-2}(-1)^jj!,\label{rc:eq:pure-quotient}\\
 B(t)&=\sum_{m\ge1}\frac{b_m}{m!}t^m,&
 B'-B&=-\alpha+\log(1+t),\qquad B(0)=0.\nonumber
\end{align}
The empty sum is zero.  For $Y=B'$ one obtains
$Y'-Y=1/(1+t)$ and $Y(0)=-\alpha$.  The residue of $e^{-t}/(1+t)$
at $-1$ is $e$, so continuation adds $2\pi i e^{t+1}$ to $Y$.
The same regular-singularity argument proves that the quotient in
\eqref{rc:eq:pure-quotient} is never arithmetic Gevrey one.  The statement
holds independently of whether its Borel-summed numerator vanishes at one.

\begin{proposition}[The restricted zero-division gate]\label{rc:prop:gate}
For the four-dimensional algebraic-coefficient family above, the assertion
\begin{equation}\label{rc:eq:Z4}
 (Z_4):\qquad \Psi_0(1)=0\Longrightarrow\Psi\in(z-1)\rcM
\end{equation}
is equivalent to $\rho(1)=4$.  The subscript $0$ denotes summation in
the positive real direction.
\end{proposition}

\begin{proof}
If $(Z_4)$ holds, Theorem~\ref{rc:thm:nondivision} makes every
algebraic-coefficient relation
$a+be+ce\eta+d\delta=0$ trivial.  Conversely, full value rank makes
the hypothesis in \eqref{rc:eq:Z4} possible only for the zero expression,
whose quotient belongs to $\rcM$.
\end{proof}

Fischler--Rivoal's general mixed zero-removal statement is Conjecture~3,
and their Theorem~4 assumes that conjecture before deducing ordinary-point
linear specialization \cite{FRmixed}.  The present system
meets its differential and summation conditions at $1$ in direction zero,
so the conjecture would imply the target.  Its conditional hypothesis
cannot be dropped.  Formula \eqref{rc:eq:nondivision} is an unconditional
formal obstruction; \eqref{rc:eq:Z4} is the unresolved passage from an
analytic zero to membership in the arithmetic quotient class.  In
particular, formal independence in $\rcM/(z-1)\rcM$ does not establish
injectivity of analytic evaluation on these classes.

\subsection{Why analytic division does not supply algebraic coefficients}

Let $A(z)=a+bE(z)+cF(z)$ with $a,b,c\in k$.  Entire analytic division
does give $H_A(z)=(A(z)-A(1))/(z-1)$.  If
$H_n=\sum_{j=1}^n1/j$ and $H_0=0$, then
$F(z)=\sum_{n\ge0}H_nz^n/n!$, and coefficient comparison gives
\begin{equation}\label{rc:eq:analytic-tails}
 [z^n]H_A=
 b\left(e-\sum_{j=0}^n\frac1{j!}\right)
 +c\left(e\eta-\sum_{j=0}^n\frac{H_j}{j!}\right).
\end{equation}
Indeed the coefficient of a quotient $(A-A(1))/(z-1)$ is the tail of the
Taylor series of $A$ beyond degree $n$, which proves the formula.  If
$(b,c)\ne(0,0)$, each displayed coefficient is transcendental: its
algebraicity would give a forbidden nontrivial relation on $1,e,e\eta$.
Analytic tail estimates therefore cannot be substituted for E-function
arithmetic conditions.  Under a hypothetical relation $A(1)+d\delta=0$,
replacing $A(1)$ by $-d\delta$ only moves the same unknown connection
constant; it does not place these coefficients in a number field.
When $b=c=0$ the coefficients in the negative quotient are algebraic,
but Theorem~\ref{rc:thm:nondivision} still excludes the required
G-function denominator bounds.

This also prevents an application of a fresh-prime theorem whose parameter
is required to be algebraic to the expression $a+b/e+c\eta$ solely
because it is hypothesized to equal $\gamma$.  That expression need not
be algebraic outside the exceptional strata of Theorem~\ref{rc:thm:rank3}.
The differential equation
\[
 \Psi_0'-\Psi_0=-a+\frac{cE-c-d}{z}
\]
is itself compatible with a zero at the ordinary point one; it merely
specifies an initial condition.  Finally, the target
$\delta\notin\rcSpan_k\{1,e,e\eta\}$ has two separate arithmetic
aspects: excluding algebraic $\delta$, and excluding a transcendental
$\delta$ in that E-value span.  An intersection assertion that common
E- and divergent values are algebraic would address the second aspect;
pure divergent zero removal would address the first.  Neither is supplied
here by the differential system alone.

\section{Spectral reformulation and its precise limitations}\label{rc:sec:spectral}

For a nonnegative self-adjoint Jacobi operator $J$ on
$\ell^2(\rcZ_{\ge0})$, use the convention
\[
 m_J(z)=\langle e_0,(J+zI)^{-1}e_0\rangle\qquad(z>0).
\]
Here $a_n$ is the entry between coordinates $n$ and $n+1$, and $b_n$
is the diagonal entry at $n$.  This fixes the resolvent sign throughout.

\subsection{Laguerre determinants and an exact raw integer error}

For the probability measure $e^{-t}\,dt$ on $[0,\infty)$, the
orthonormal polynomials $p_n(t)=(-1)^nL_n(t)$ satisfy
\begin{equation}\label{rc:eq:laguerre-recurrence}
 tp_n=(n+1)p_{n+1}+(2n+1)p_n+np_{n-1}.
\end{equation}
These standard Laguerre formulas can be found in
\cite[\S\S18.2,18.9]{DLMF}.  Thus the Jacobi coefficients are
$a_n=n+1$, $b_n=2n+1$.  The Jacobi Carleman criterion
$\sum_na_n^{-1}=\infty$ makes the operator on finite sequences
essentially self-adjoint; its spectral representation identifies its
closure with multiplication by $t$ for the displayed measure
\cite[Chapter 2]{rc:teschl}.  In particular,
\begin{equation}\label{rc:eq:laguerre-weyl}
 m_J(z)=D(z),\qquad m_J(1)=\delta.
\end{equation}
Schur complementation gives the classical continued fraction
\[
 \delta=\cfrac1{2-\cfrac{1^2}{4-\cfrac{2^2}{6-
                    \cfrac{3^2}{8-\cdots}}}}.
\]
The minus signs agree with the $J+zI$ convention; compare
\cite[\S6.9]{DLMF}.

\begin{proposition}\label{rc:prop:laguerre-error}
Let $J_N$ be the leading $N\times N$ truncation, and set
\[
 Q_N=\det(J_N+I)=N!L_N(-1),\qquad
 \frac{P_N}{Q_N}=\langle e_0,(J_N+I)^{-1}e_0\rangle.
\]
For $N=0$ use $Q_0=1,P_0=0$.  Then $P_N,Q_N$ are integers,
\begin{equation}\label{rc:eq:laguerre-integer-recurrence}
 X_{N+1}=2(N+1)X_N-N^2X_{N-1},\quad
 (Q_0,Q_1)=(1,2),\quad(P_0,P_1)=(0,1),
\end{equation}
and
\begin{equation}\label{rc:eq:laguerre-exact-error}
 Q_N\delta-P_N
 =N!\int_0^\infty\frac{e^{-s}s^N}{(1+s)^{N+1}}\,ds>0.
\end{equation}
For $N\ge1$,
\begin{equation}\label{rc:eq:laguerre-lower}
 Q_N\delta-P_N\ge
 \frac{N!}{\sqrt N+2}e^{-2\sqrt N-1}\longrightarrow\infty.
\end{equation}
\end{proposition}

\begin{proof}
The monic characteristic polynomial is $(-1)^NN!L_N(t)$, by
\eqref{rc:eq:laguerre-recurrence}.  Substitution of $t=-1$ gives $Q_N$;
the cofactor formula gives integer $P_N$.  Determinant or continuant
expansion gives \eqref{rc:eq:laguerre-integer-recurrence} and its initial
values.  For a direct second-kind remainder calculation define
\[
 I_N=N!\int_0^\infty\frac{e^{-t}L_N(t)}{1+t}\,dt,
 \qquad \widetilde P_N=Q_N\delta-I_N.
\]
The Laguerre recurrence, and $\int_0^\infty e^{-t}L_N(t)\,dt=0$
for $N\ge1$, show that $I_N$ obeys the same recurrence for $N\ge1$:
indeed, writing $2N+1-t=2N+2-(1+t)$ produces exactly the vanishing
orthogonality integral.  The initial values are $I_0=\delta$ and
$I_1=2\delta-1$.  Hence $\widetilde P_0=0$, $\widetilde P_1=1$, and
uniqueness for the recurrence identifies $\widetilde P_N=P_N$.

Insert $(1+t)^{-1}=\int_0^\infty e^{-(1+t)s}\,ds$.  Interchanging the
integrals is justified by absolute integrability, since $L_N$ is a
polynomial.  Its Laplace transform is
\[
 \int_0^\infty e^{-(1+s)t}L_N(t)\,dt
       =\frac{s^N}{(1+s)^{N+1}}.
\]
For completeness, expansion
$L_N(t)=\sum_{j=0}^N\binom Nj(-t)^j/j!$ makes the integral equal to
$(1+s)^{-1}\sum_j\binom Nj[-(1+s)^{-1}]^j$, proving the formula.
It follows that $I_N$ is the right side of
\eqref{rc:eq:laguerre-exact-error}, which is positive.

Restrict that integral to $s\in[\sqrt N,\sqrt N+1]$.  On this interval
$e^{-s}\ge e^{-\sqrt N-1}$, $(1+s)^{-1}\ge(\sqrt N+2)^{-1}$, and
\[
 \left(\frac{s}{1+s}\right)^N
 \ge\exp\{-N\log(1+1/\sqrt N)\}\ge e^{-\sqrt N}.
\]
Integration over an interval of length one proves
\eqref{rc:eq:laguerre-lower}; its divergence follows from factorial
growth, for example from Stirling's formula.
\end{proof}

The initial convergents are $1/2$, $4/7$, and $20/34$.  This is the
classical Laguerre denominator with its homogeneous Gompertz companion,
also related to the constructions in \cite{sc:HP}.
Positivity certifies nonvanishing but does not make these raw integer
linear forms small.  Dividing $P_N,Q_N$ by their common divisor changes
the arithmetic estimate; no uniform gain from that primitive
normalization is established by the lower bound above.

\subsection{An integer Jacobi family with rational designated values}

\begin{theorem}\label{rc:thm:integer-jacobi}
For each integer $B\ge2$, let $K_B$ be the self-adjoint closure of the
Jacobi matrix
\begin{equation}\label{rc:eq:KB}
 a_n=2(n+1),\qquad b_0=B,\qquad b_n=5n\quad(n\ge1).
\end{equation}
It is nonnegative, has compact resolvent and infinite spectrum, and all
its moments $\langle e_0,K_B^ne_0\rangle$ are integers.  Its canonical
Weyl function is meromorphic on $\rcC$, has infinitely many poles, and
satisfies
\begin{equation}\label{rc:eq:KB-value}
 m_{K_B}(1)=1/B,\qquad \lim_{z\to+\infty}zm_{K_B}(z)=1.
\end{equation}
In particular it is neither algebraic over $\rcC(z)$ nor D-finite.
\end{theorem}

\begin{proof}
Since $\sum_n1/[2(n+1)]=\infty$, the Jacobi Carleman criterion gives
essential self-adjointness.  It suffices first to prove positivity for
$K=K_2$, since $K_B=K+(B-2)P_0$, where
$P_0v=\langle e_0,v\rangle e_0$, is a nonnegative bounded rank-one
perturbation.  Conjugate $K$ by the diagonal unitary $Ue_n=(-1)^ne_n$
to make its off-diagonal entries negative.  Choose
\[
 h_0=h_1=1,\qquad h_n=(3/4)^{n-1}\quad(n\ge1).
\]
Expansion of the squares gives, for finitely supported $v$,
\begin{align*}
 \langle v,UKU^{-1}v\rangle
 &=\sum_{n\ge0}a_nh_nh_{n+1}
       \left|\frac{v_n}{h_n}-\frac{v_{n+1}}{h_{n+1}}\right|^2
   +\sum_{n\ge0}R_n|v_n|^2,\\
 R_n&=b_n-a_{n-1}\frac{h_{n-1}}{h_n}
                -a_n\frac{h_{n+1}}{h_n},
\end{align*}
where the first subtracted term is omitted for $n=0$.  Direct calculation
gives $R_0=R_1=0$ and
$R_n=5n/6-3/2>0$ for $n\ge2$.  This proves nonnegativity on a core,
and hence for the closure.

Set $w_n=(-1/2)^n$.  For $n\ge1$, division of $((K_B+I)w)_n$ by
$w_n$ gives $-4n+(5n+1)-(n+1)=0$.  Its zeroth component is
$(B+1)-1=B$.  Thus
\[
 (K_B+I)w=Be_0.
\]
The vector $w$ is square summable, and its formal image under $K_B$ is
$Be_0-w\in\ell^2$.  It belongs to the maximal Jacobi domain, which
essential self-adjointness identifies with the self-adjoint domain.
Since $K_B+I$ is invertible, its $(0,0)$ resolvent entry is $w_0/B=1/B$.

The elementary inequality $2\lvert u\overline v\rvert\le|u|^2+|v|^2$
applied to adjacent coordinates gives
\[
 K_B\ge\rcdiag(n-2)_{n\ge0}
\]
as a quadratic-form inequality.  Consequently a bounded subset of the
form domain, with norm coming from $K_B+3I$, has uniformly small
$\ell^2$ tails: its tail from $N$ onward is bounded by a constant
divided by $N+1$.  Finite-coordinate projections then prove compactness
of the form-domain embedding, and hence compactness of the resolvent.
Its spectrum is an infinite discrete sequence tending to infinity.

All $a_n$ are nonzero.  The recurrence shows inductively that every
$e_n$ belongs to the span of $e_0,K_Be_0,\ldots,K_B^ne_0$, so $e_0$
is cyclic.  Every eigenvalue therefore has a nonzero spectral weight
in its probability spectral measure.  The spectral representation of
$m_{K_B}$ gives a meromorphic function with a pole at the negative of
each eigenvalue, and thus infinitely many poles.  An algebraic function
over $\rcC(z)$ has only finitely many possible finite singularities;
a solution of a nonzero linear differential equation with polynomial
coefficients likewise has singularities only among the finitely many
zeros of its leading coefficient.  The infinite-pole property excludes
both algebraicity and D-finiteness.

Finally $z/(z+t)\to1$, bounded between zero and one for $z>0,t\ge0$,
so dominated convergence for the probability spectral measure proves
the normalization in \eqref{rc:eq:KB-value}.  Since $K_B^ne_0$ has
finite support and integer coordinates, its zeroth coordinate, the
$n$th moment, is an integer.
\end{proof}

This counterexample keeps positivity, integer coefficients, essential
self-adjointness, unbounded spectrum, integer moments, canonical
normalization and functional transcendence.  These general properties
therefore do not imply irrationality of a designated resolvent value.
The example is not the Laguerre operator, and it makes no contrary claim
about the actual arithmetic nature of $\delta$.

\subsection{Meixner identification and full differential transcendence}

Non-D-finiteness by itself does not rule out nonlinear algebraic
differential equations.  We now retain the full argument establishing
the stronger property for $K=K_2$.

\begin{lemma}[Direct Meixner spectral representation]\label{rc:lem:meixner}
The canonical Weyl function of $K$ in Theorem~\ref{rc:thm:integer-jacobi}
is
\begin{equation}\label{rc:eq:meixner-R}
 m_K(z)=\frac{R(z)}{1+2R(z)},\qquad
 R(z)=\frac34\sum_{j\ge0}\frac{4^{-j}}{3j+z-1}.
\end{equation}
\end{lemma}

\begin{proof}
Fix $0<c<1$ and put $\mu_j=(1-c)c^j$.  Define polynomials $M_n(j)$
by the convergent generating function near $u=0$,
\[
 G_j(u)=\sum_{n\ge0}M_n(j)u^n
       =\frac1{1-u}\left(\frac{1-u/c}{1-u}\right)^j.
\]
A geometric sum, absolutely convergent for sufficiently small $u,v$,
gives
\begin{align*}
 \sum_{j\ge0}\mu_jG_j(u)G_j(v)
 &=\frac{1-c}{(1-u)(1-v)-c(1-u/c)(1-v/c)}\\
 &=\frac1{1-uv/c}.
\end{align*}
Coefficient comparison proves
$\sum_j\mu_jM_n(j)M_m(j)=c^{-n}\delta_{nm}$.  Thus
$p_n(j)=(-\sqrt c)^nM_n(j)$ are orthonormal and have positive leading
coefficients.  Differentiation of $G_j$ gives
\[
 (1-c)jG_j=(c-u)G_j-(c-u)(1-u)G_j'.
\]
Comparing coefficients, with $M_{-1}=p_{-1}=0$, yields
\begin{align*}
 jM_n&=-\frac{c(n+1)}{1-c}M_{n+1}
       +\frac{(1+c)n+c}{1-c}M_n-\frac n{1-c}M_{n-1},\\
 jp_n&=a_np_{n+1}+b_np_n+a_{n-1}p_{n-1},\\
 a_n&=\frac{\sqrt c(n+1)}{1-c},\qquad
 b_n=\frac{(1+c)n+c}{1-c}.
\end{align*}

We also need completeness, so that this is an identification of the
actual spectral measure.  If $f\in L^2(\mu)$ is orthogonal to every
polynomial, the series
$H(s)=\sum_{j\ge0}\mu_j f(j)s^j$ is analytic for
$|s|<c^{-1/2}$ by Cauchy--Schwarz.  Its derivatives at $s=1$ vanish by
orthogonality to the falling factorials; termwise differentiation is
justified on compact subdisks.  Since $1$ is interior to the disk,
$H\equiv0$, and hence $f=0$.  Complex conjugation, if required by the
inner-product convention, does not change the argument.

The unitary map $e_n\mapsto p_n$ consequently identifies the closure of
the Jacobi matrix with multiplication by $j$ on $L^2(\mu)$.
Indeed multiplication extends the displayed recurrence, and Carleman's
criterion ensures that this self-adjoint extension is unique.  At
$c=1/4$ the coefficients are $a_n=2(n+1)/3$ and $b_n=(5n+1)/3$.
Therefore $L=3J-I$ has off-diagonal entries $2(n+1)$, diagonal entries
$5n$, and spectral measure $\mu_j=(3/4)4^{-j}$ at $3j-1$.  Its Weyl
function is precisely $R$ in \eqref{rc:eq:meixner-R}.

Since $K=L+2P_0$, the rank-one resolvent identity gives, wherever the
resolvents exist,
\[
 (K+zI)^{-1}=(L+zI)^{-1}
 -\frac{2(L+zI)^{-1}P_0(L+zI)^{-1}}{1+2R(z)}.
\]
Taking the $(0,0)$ entry yields $m_K=R-2R^2/(1+2R)=R/(1+2R)$.
Meromorphic continuation proves the identity throughout its domain.
\end{proof}

\begin{lemma}\label{rc:lem:no-shift-rational}
For any field $C_0$ of characteristic zero, the rational difference
equation
\begin{equation}\label{rc:eq:shift-rational}
 f(t+1)-4f(t)=-1/t
\end{equation}
has no solution in $C_0(t)$, where the shift fixes $C_0$.
\end{lemma}

\begin{proof}
Work over an algebraic closure of $C_0$.  If $f$ has a finite pole,
partition its finite pole set into orbits under translation by integers.
Within any occupied orbit choose a representative $\alpha$ and the
smallest and largest integer indices $p\le q$ for which
$\alpha+p,\alpha+q$ are poles.  The difference in
\eqref{rc:eq:shift-rational} has a noncancelling pole at
$\alpha+p-1$, contributed only by $f(t+1)$, and another at
$\alpha+q$, contributed only by $-4f(t)$.  They are distinct, and poles
from other orbits cannot cancel them.  The right side has exactly one
pole, a contradiction.  If $f$ has no finite poles it is polynomial,
which also cannot produce the pole on the right side.
\end{proof}

\begin{theorem}\label{rc:thm:differential-transcendence}
The canonical Weyl function $m_K$ is differentially transcendental over
$\rcC(z)$: it satisfies no nonzero polynomial differential equation
with coefficients in $\rcC(z)$.  Nevertheless $m_K(1)=1/2$.
\end{theorem}

\begin{proof}
The series defining $R$ converges normally on compact sets avoiding its
simple poles, so $R$ is meromorphic on $\rcC$.  Shifting the summation
index gives
\[
 R(z+3)=4R(z)-\frac3{z-1}.
\]
Thus $y(t)=R(3t+1)$ is meromorphic and satisfies
$y(t+1)=4y(t)-1/t$.

Let $\mathcal P$ be the field of one-periodic meromorphic functions of
$t$.  Hardouin--Singer's Proposition~3.9(2) states, in the form needed
here, that a meromorphic solution $y\notin\mathcal P(t)$ of
$y(t+1)=ay(t)+b(t)$, with constant $a\ne0$ and $b\in\rcC(t)$,
is differentially algebraic over $\mathcal P(t)$ if and only if
$b=f(t+1)-af(t)$ for some $f\in\rcC(t)$
\cite[Proposition 3.9(2)]{rc:hardouin-singer}.  A nonzero constant $a$
is in the standard form required by that proposition.

We check the nonmembership hypothesis explicitly.  The coordinate $t$
is transcendental over $\mathcal P$: a nonzero polynomial relation with
periodic coefficients, after every integer translation, would make the
same polynomial have infinitely many distinct roots $t+n$ in the field
of meromorphic functions.  Thus $\mathcal P(t)$ is an abstract rational
function field on which translation fixes the coefficient field.  If
$y$ belonged to it, Lemma~\ref{rc:lem:no-shift-rational} with
$C_0=\mathcal P$ would be contradicted.  The same lemma with
$C_0=\rcC$ excludes the rational solution required by
Hardouin--Singer.  The proposition therefore proves differential
transcendence of $y$ over $\mathcal P(t)$, and hence over $\rcC(t)$.
An affine change of the independent variable preserves differential
algebraicity, so $R$ is differentially transcendental over $\rcC(z)$.

Finally $R=m_K/(1-2m_K)$, and the denominator is not identically zero
since $m_K(z)\to0$ on the positive real ray.  Differentially algebraic
functions form a field, so differential algebraicity of $m_K$ would
imply that of $R$, a contradiction.  At $z=1$, the $j=0$ summand of
$R$ has a simple pole of residue $3/4$ and the remaining summands are
regular.  The quotient $R/(1+2R)$ has a removable singularity with
value $1/2$, consistently with \eqref{rc:eq:KB-value}.
\end{proof}

This conclusion uses the cited nonlinear differential-algebraic theorem,
not an inference from infinitely many poles.  It is an application of
existing theory to an explicit comparison operator, and supplies no new
arithmetic statement about the Euler resolvent.

\subsection{A nearby arithmetic Gevrey example and a degenerate change of basis}

The probability density $(1+t)e^{-t}/2$ gives another useful comparison:
\begin{equation}\label{rc:eq:nearby-D}
 \widetilde D(z)=\int_0^\infty\frac{(1+t)e^{-t}}{2(z+t)}\,dt
 =\frac{1+(1-z)D(z)}2,\qquad\widetilde D(1)=\frac12.
\end{equation}
The identity follows from $1+t=(z+t)+(1-z)$.  Its moments are
$(n+2)n!/2$, and the shifted Borel kernel of its inverse-power expansion
is
\[
 \sum_{n\ge0}\frac{(-1)^n}{n!}\frac{(n+2)n!}{2}u^n
 =\frac{u+2}{2(1+u)^2}.
\]
Thus it retains a probability measure with unbounded support, the
normalization $1/z$, and an arithmetic Gevrey-one expansion.  It is not
rational, since \eqref{rc:eq:nearby-D} would then make $D$ rational;
the nonzero logarithmic monodromy of $D$ excludes that possibility.
However, direct differentiation gives
\[
 \widetilde D'=\frac z{z-1}\widetilde D+
                         \frac{1-2z}{2z(z-1)},\qquad
 \frac{\widetilde D-1/2}{z-1}=-D/2.
\]
The corresponding formal quotient is already arithmetic Gevrey one.
The change of basis from $(1,D)$ to $(1,\widetilde D)$ degenerates at
$z=1$, where the displayed scalar equation is singular.  This example
therefore does not contradict an ordinary-point mixed lifting statement
for the original, nondegenerate four-function system.

\section{Exact local radii, denominator growth, and the foliation test}
\label{rc:sec:adelic}

\subsection{A number-field initial-value family}

For $\alpha\in k$ consider the solution at zero of
\begin{equation}\label{rc:eq:initial-ode}
 Y_\alpha'=Y_\alpha-\frac1{1+t},\qquad Y_\alpha(0)=\alpha.
\end{equation}
Write $S_n=\sum_{j=0}^n(-1)^jj!$, with $S_{-1}=0$.  Coefficient
comparison, or induction in the differential equation, gives
\begin{equation}\label{rc:eq:initial-series}
 Y_\alpha(t)=\sum_{n\ge0}\frac{\alpha-S_{n-1}}{n!}t^n.
\end{equation}
For example the recurrence for $c_n=[t^n]Y_\alpha$ is
$(n+1)c_{n+1}=c_n-(-1)^n$, and multiplication by $n!$ telescopes
to the formula.  In the complex realization uniqueness gives
\begin{equation}\label{rc:eq:initial-real}
 Y_\alpha(t)=D(1+t)+(\alpha-\delta)e^t.
\end{equation}
The local complex radius is exactly one: the equation gives analyticity
on $|t|<1$, and a radius greater than one would make its left side
regular at $t=-1$, contradicting the inhomogeneous pole there.  If
$\delta$ were algebraic, $\alpha=\delta$ would be the initial value
selected by the real positive-ray normalization.  This statement alone
does not impose any nonarchimedean initial-value matching.

\begin{theorem}[Exact nonarchimedean radius]\label{rc:thm:radius}
Fix a prime $p$, a finite extension $L/\rcQ_p$, and an image of
$\alpha$ in $L$.  Normalize the absolute value by $|p|_p=p^{-1}$, and
put
\[
 \delta_p=\sum_{j\ge0}(-1)^jj!\in\rcZ_p.
\]
The radius of convergence of \eqref{rc:eq:initial-series} is
\begin{equation}\label{rc:eq:radius}
 R_p(\alpha)=
 \begin{cases}
  1,&\alpha=\delta_p,\\
  p^{-1/(p-1)},&\alpha\ne\delta_p.
 \end{cases}
\end{equation}
\end{theorem}

\begin{proof}
The factorial series converges because $v_p(j!)\to\infty$.  If
$\alpha\ne\delta_p$, its coefficient numerators
$\alpha-S_{n-1}$ converge to a nonzero element.  Their valuations are
eventually constant, since an ultrametric-convergent sequence with
nonzero limit eventually has the absolute value of that limit.
Legendre's formula gives $v_p(n!)/n\to1/(p-1)$.  Cauchy--Hadamard
then gives the second radius in \eqref{rc:eq:radius}.

If $\alpha=\delta_p$, then
$\alpha-S_{n-1}=\sum_{j\ge n}(-1)^jj!$ belongs to $n!\rcZ_p$,
so all Taylor coefficients are integral and the radius is at least one.
If it were greater than one, the series and its derivative would be
analytic at $t=-1$.  Multiplying \eqref{rc:eq:initial-ode} by $1+t$
and substituting $t=-1$ would give $0=-1$, a contradiction.  The radius
is therefore exactly one.
\end{proof}

Equality modulo $p$ or any fixed power of $p$ does not imply the first
case: the improved radius requires equality with the entire $p$-adic
sum.  This distinction will also appear in the curvature calculation.

\begin{theorem}[Zero finite-place radius product]\label{rc:thm:radius-product}
Let $K=\rcQ(\alpha)$, $d=[K:\rcQ]$, and for every finite place
$v\mid p$ set $n_v=[K_v:\rcQ_p]/d$.  Interpret the radius in
Theorem~\ref{rc:thm:radius} in $K_v$.  Then
\begin{equation}\label{rc:eq:radius-product}
 \prod_{v\text{ finite}}R_v(\alpha)^{n_v}=0.
\end{equation}
If $d>1$, the following stronger estimate holds prime by prime:
\begin{equation}\label{rc:eq:radius-per-prime}
 \prod_{v\mid p}R_v(\alpha)^{n_v}
 \le p^{-(1-1/d)/(p-1)}.
\end{equation}
If $\alpha\in\rcQ$, the mismatching primes instead satisfy
\begin{equation}\label{rc:eq:mismatch-divergence}
 \sum_{p:\alpha\ne\delta_p}\frac{\log p}{p-1}=\infty.
\end{equation}
\end{theorem}

\begin{proof}
Suppose first that $d>1$.  At a matching place $\alpha=\delta_p$
lies in $\rcQ_p$.  Since $K=\rcQ(\alpha)$, its completion $K_v$ is
generated topologically over $\rcQ_p$ by the image of $\alpha$;
therefore $K_v=\rcQ_p$ and $n_v=1/d$.  There is at most one matching
place over a given $p$: two embeddings of $\rcQ(\alpha)$ into
$\rcQ_p$ sending its generator to the same prescribed element
$\delta_p$ coincide and induce the same place.  Thus the total
matching weight is at most $1/d$.  The weights above $p$ sum to one,
and Theorem~\ref{rc:thm:radius} gives
\eqref{rc:eq:radius-per-prime}.  The classical divergence of
$\sum_p\log p/(p-1)$ proves the zero product.

Now write $\alpha=A/B\in\rcQ$ in lowest terms with $B>0$, and let
$M=\{p:\alpha\ne\delta_p\}$.  Suppose contrary to
\eqref{rc:eq:mismatch-divergence} that
$C=\sum_{p\in M}\log p/(p-1)<\infty$.  At matching primes all
coefficients of \eqref{rc:eq:initial-series} are integral.  At the other
primes the factorial denominators show that a common denominator through
degree $N$ is bounded by
\[
 B\prod_{p\in M}p^{v_p(N!)}\le B\exp(CN),
\]
where only finitely many factors are nontrivial and
$v_p(N!)\le N/(p-1)$.  The complex analytic radius is positive, so
the rational coefficients have exponential absolute-value bounds.
They are holonomic: applying $(1+t)\partial_t+1$ to
$(\partial_t-1)Y_\alpha=-1/(1+t)$ gives a homogeneous linear
differential equation with rational coefficients.  These are precisely
the G-function requirements, so $Y_\alpha$ would be a G-function.

But $(e^{-t}Y_\alpha)'=-e^{-t}/(1+t)$ has residue $-e$ at $-1$.
Continuation around $-1$ changes $Y_\alpha$ by $-2\pi i e^{t+1}$.
As in Theorem~\ref{rc:thm:nondivision}, this is incompatible with the
global regular-singularity theorem for G-functions.  Thus
\eqref{rc:eq:mismatch-divergence} holds, and the radius dichotomy proves
\eqref{rc:eq:radius-product} also for rational $\alpha$.
\end{proof}

The rational case is the denominator obstruction expressed in local
radius language, rather than a new distribution theorem for factorial
residues.  For primary background on real and $p$-adic factorial sums
and their approximation, see \cite{rc:ernvall-hytonen}.  The proof above
specifies its own use of the global G-function theorem.

\begin{proposition}[Common denominator consequences]\label{rc:prop:denominators}
Let $q_N$ be the least positive rational integer for which all
coefficients of \eqref{rc:eq:initial-series} through degree $N$ become
algebraic integers after multiplication by $q_N$.  If $\alpha\in\rcQ$,
then
\[
 \frac{\log q_N}{N}\longrightarrow\infty.
\]
If $\alpha$ is algebraic irrational, then there are constants
$c_\alpha,C_\alpha>0$ such that
\[
 c_\alpha N!\le q_N\le C_\alpha N!\qquad(N\ge1).
\]
\end{proposition}

\begin{proof}
For rational $\alpha$, at each fixed mismatching prime $p$ the valuations
of $\alpha-S_{N-1}$ eventually stabilize at a finite constant.  The
degree-$N$ coefficient alone then forces
$v_p(q_N)\ge v_p(N!)-O_{\alpha,p}(1)$.  For each finite set $T$ of
mismatching primes this gives
\[
 \liminf_{N\to\infty}\frac{\log q_N}{N}
 \ge\sum_{p\in T}\frac{\log p}{p-1}.
\]
The right side is unbounded as $T$ increases by
\eqref{rc:eq:mismatch-divergence}, proving the first assertion.

For irrational $\alpha$, choose an integer $A_0>0$ making $A_0\alpha$
integral.  Then $A_0N!$ clears every coefficient through $N$, proving
the upper bound.  In a fixed normal closure $L$ choose a conjugate
$\alpha^\sigma\ne\alpha$.  Since
$q_N(\alpha-S_{N-1})/N!$ and its conjugate are algebraic integers,
so is their nonzero difference
\[
 \frac{q_N}{N!}(\alpha-\alpha^\sigma).
\]
Its absolute field norm from $L$ to $\rcQ$ is at least one.  Writing
$h=[L:\rcQ]$ gives
$(q_N/N!)^h\lvert N_{L/\rcQ}(\alpha-\alpha^\sigma)\rvert\ge1$,
which proves the lower bound with a fixed positive constant.
\end{proof}

\subsection{Factorial truncations in characteristic \texorpdfstring{$p$}{p}}

\begin{proposition}\label{rc:prop:pcurvature}
On $\mathbb G_m\times\mathbb A^1$ over $\mathbb F_p$, let
\[
 V=\partial_z+(y-1/z)\partial_y,
 \qquad T_p(z)=\sum_{j=0}^{p-1}(-1)^jj!z^{-j-1}.
\]
As derivations of the rational function field,
\begin{equation}\label{rc:eq:pcurvature}
 V^p=(y-T_p(z))\partial_y,
 \qquad T_p'-T_p=-1/z.
\end{equation}
The rank-one foliation generated by $V$ is not stable under $p$th powers
at any prime.  At a good reduction of $(1,\alpha)$, its curvature class
vanishes at that point exactly when
\begin{equation}\label{rc:eq:pointwise-congruence}
 \alpha\equiv T_p(1)=S_{p-1}\equiv\delta_p\pmod{\mathfrak p}.
\end{equation}
\end{proposition}

\begin{proof}
For $m\ge1$, induction gives
$V^m(y)=y-\sum_{j=0}^{m-1}\partial_z^j(1/z)$, since $V$ acts as
$\partial_z$ on rational functions of $z$.  At $m=p$ this is
$y-T_p(z)$, while $V^p(z)=0$.  In characteristic $p$ the $p$fold
iterate of a derivation is a derivation, because the interior binomial
coefficients in the iterated Leibniz rule vanish.  Its values on $z,y$
therefore prove the first identity of \eqref{rc:eq:pcurvature}.
Termwise differentiation of $T_p$ telescopes against $T_p$; the last
derivative term has factor $p!$ and is zero, giving the second identity.

If $V^p=fV$ for a rational function $f$, comparison of the $z$ components
would give $f=0$, whereas $y-T_p(z)$ is not the zero rational function.
Thus the generic foliation is not $p$-closed.  The class of $V^p$
modulo the line generated by $V$ vanishes at $(1,\alpha)$ if and only
if its vertical coefficient is zero, since $V$ has horizontal component
one.  This is the first congruence in
\eqref{rc:eq:pointwise-congruence}.  The equality with $\delta_p$
modulo $\mathfrak p$ follows because every $j!$ with $j\ge p$ is
divisible by $p$.
\end{proof}

The full formal differential system also has an explicit $p$fold
iteration.  In characteristic $p$, with $E,F,D$ treated as differential
module coordinates satisfying \eqref{rc:eq:four-system}, one has
\begin{equation}\label{rc:eq:full-p-iteration}
 E^{(p)}=E,\qquad D^{(p)}=D-T_p,
 \qquad F^{(p)}=F-z^{-p}E-T_p.
\end{equation}
Here is a direct check of the less immediate last formula.  Iterating
$F'=F+(E-1)/z$ gives
$F^{(p)}-F=\sum_{j=0}^{p-1}\partial_z^j((E-1)/z)$.
The non-exponential part is $-T_p$, while its exponential coefficient is
\[
 \sum_{j=0}^{p-1}(\partial_z+1)^j(1/z)
 =\partial_z^{p-1}(1/z)=-z^{-p}.
\]
The polynomial identity
$\sum_{j=0}^{p-1}(X+1)^j=X^{p-1}$ holds in $\mathbb F_p[X]$;
the last equality uses $(p-1)!=-1$ in $\mathbb F_p$ and is also valid
for $p=2$.  Neither $\gamma$ nor $\delta$ enters these module identities.

\subsection{Checking the algebraization hypothesis rather than assuming it}

Bost's algebraicity criterion for leaves concerns a smooth variety over
a number field, an involutive subbundle of its tangent bundle, and a
rational base point.  Besides a complex Liouville condition on the leaf,
Theorem~2.1 requires the reduced foliation to be closed under $p$th
powers at almost every finite prime \cite[Theorem 2.1]{rc:bost}.
Proposition~\ref{rc:prop:pcurvature} shows that this required global
closure fails at every prime for the displayed Euler foliation.  The
pointwise congruence \eqref{rc:eq:pointwise-congruence} is not a
substitute for that hypothesis.  It is also weaker than the exact
equality $\alpha=\delta_p$ that produces radius one.

Likewise Theorem~\ref{rc:thm:radius-product} leaves no positive product
of these particular Taylor radii.  This is a coordinate-specific
calculation for the displayed parameter and local series.  It neither
identifies these radii with all intrinsic formal-germ sizes or capacities
nor rules out other arithmetic geometries, reparametrizations, or
arithmetic Gevrey estimates.  Those distinctions matter in the broader
formal-germ and capacity theory of \cite{rc:bost-chambert-loir}.
The concrete missing implication is that the real canonical boundary
condition might force compatible nonarchimedean conditions; no such
implication has been proved.  A general rank-three relation is still
further from a number-field initial-value problem, since it need not
make $\delta$ algebraic at all.

\section{Galois transport and exact joint-span bookkeeping}\label{rc:sec:galois}

\subsection{What an automorphism actually transports}

If $\gamma=a+b/e+c\eta$ with $a,b,c\in k$, then for an abstract
automorphism $\sigma$ of $\rcC$ the valid identity is
\begin{equation}\label{rc:eq:galois-transform}
 \sigma(\gamma)=\sigma(a)+\frac{\sigma(b)}{\sigma(e)}
                              +\sigma(c)\sigma(\eta).
\end{equation}
Conjugating only the algebraic coefficients while holding all three
transcendental values fixed would require a separately justified
automorphism with those fixed points.  The numerical relation itself
may obstruct that prescription.  Nor does an arbitrary field
automorphism commute with the limits defining $\exp$ or $\rcEin$:
there is no general identity
$\sigma(e^\alpha)=e^{\sigma(\alpha)}$ or its $\rcEin$ analogue.
The continuous automorphisms of $\rcC$ are the identity and complex
conjugation; the latter fixes the real constants here and gives only
the realness already proved in Proposition~\ref{rc:prop:bad-representation}.
Also, the argument $1$ itself is fixed by every field automorphism,
so its algebraic orbit cannot produce new evaluation points.

If $\gamma$ were algebraic, its algebraic conjugates would be roots of
its minimal polynomial, but their identification with the analytically
normalized Euler constant would still require proof.  Coefficient traces
and norms have the same issue: all constants in the relation must be
tracked under the embeddings.  Without additional field-intersection
information they do not turn an algebraic-coefficient relation on fixed
transcendental numbers into a rational-coefficient relation on those same
numbers.  Genuine E-function lifting can sometimes justify descent by
first producing functional relations; $D$ is not an E-function, and
extending that step to this mixed relation is precisely missing arithmetic
input.

\subsection{Circuits constrain the exceptional set but do not propagate it}

If $1\in B$, Theorem~\ref{rc:thm:badset} implies that every finite
subset of $B\setminus\{1\}$ is multiplicatively independent.  Indeed,
given an integer multiplicative relation on that subset, add at the
point $1$ the negative sum of its coefficients.  This produces the
forbidden affine relation.  It does not, however, assert that another
bad point exists.  The configuration $B=\{1\}$ satisfies every
circuit restriction and all finite-rank bounds in this part.  Changing
the geometric base leaves $1=r^0$ in every progression, so intersecting
such bounds cannot eliminate this configuration.

The exact scaling formula also shows the unremoved terms:
\[
 E_1(rx)=E_1(x)-\int_x^{rx}\frac{e^{-t}}t\,dt.
\]
By \eqref{rc:eq:connection} the difference introduces $\log r$ and
$\rcEin(rx)-\rcEin(x)$.  Without an additional relation eliminating
these new quantities it does not transport rank failure from $1$ to
$2$ or $4$.  Density-one full rank, or even sequences of full-rank
points approaching one, likewise gives no continuity principle for
arithmetic independence of values.

\begin{proposition}[Joint span on a geometric progression]\label{rc:prop:joint}
Fix positive algebraic $r\ne1$, let $L=\log r$, and let
$S\subset\rcZ$ contain $m\ge2$ distinct integers.  Put $x_n=r^n$ and
\begin{align*}
 V_S&=\rcSpan_k\{1,e^{-x_n},\rcEin(x_n):n\in S\},\\
 J_S&=\rcSpan_k\{1,e^{-x_n},\rcEin(x_n),E_1(x_n):n\in S\}.
\end{align*}
Then
\begin{equation}\label{rc:eq:joint-span}
 \dim_k V_S=2m+1,\qquad
 J_S=V_S+\rcSpan_k\{\gamma,L\},\qquad
 \dim_k J_S=2m+1+s_S,
\end{equation}
where
\[
 s_S=\dim_k\frac{V_S+\rcSpan_k\{\gamma,L\}}{V_S}\in\{0,1,2\}.
\]
For nested sets $S\subset S'$ the parameter $s_{S'}\le s_S$.
If $n\notin S$, then
\begin{equation}\label{rc:eq:joint-increment}
 \dim_kJ_{S\cup\{n\}}-\dim_kJ_S
 =2+s_{S\cup\{n\}}-s_S.
\end{equation}
Along any nested sequence the total loss relative to increments of two
is at most two, and the parameters eventually stabilize.
\end{proposition}

\begin{proof}
The dimension of $V_S$ is Lemma~\ref{rc:lem:evalues}.  For each
$n\in S$ the connection formula gives
\[
 \rcEin(x_n)-E_1(x_n)=\gamma+nL.
\]
This places every $E_1(x_n)$ in the right side of
\eqref{rc:eq:joint-span}.  Conversely, choose distinct $n,n'\in S$;
the difference of these two identities, divided by the nonzero integer
$n-n'$, expresses $L$ as an element of $J_S$, and either identity then
expresses $\gamma$.  The vector-space equality and its dimension formula
follow.  If $S$ is enlarged, $V_S$ increases while
$\rcSpan_k\{\gamma,L\}$ is fixed.  The dimension of its image in the
quotient therefore cannot increase.  Subtraction of the dimension
formulas gives \eqref{rc:eq:joint-increment}.  A nonincreasing sequence
in $\{0,1,2\}$ stabilizes, and the total decrease is at most two.
\end{proof}

This is an exact formula with an undetermined arithmetic parameter; the
finite E-value input alone does not determine $s_S$, its eventual value,
or a stabilization index.  The use of unscaled $E_1$ values is essential
to the stated formula.  Multiplication of one local block by the common
nonzero scalar $e^{-x}$ preserves its $k$-rank.  Multiplication of
different blocks by different transcendental scalars does not generally
preserve the rank of the combined vector, so
\eqref{rc:eq:joint-span} cannot simply be rewritten with all $E_1(x_n)$
replaced by $D(x_n)$.

For distinct points the local E-value spaces
$\rcSpan_k\{1,e^{-x},\rcEin(x)\}$ intersect only in $k$, as follows
by subtracting two representations and applying
Lemma~\ref{rc:lem:evalues}.  At a bad point $r^n$, the value
$\gamma+nL$ lies in that local three-dimensional space.  This explains
the cancellation geometry behind Theorem~\ref{rc:thm:badset}, while
leaving the actual position of $\gamma$ at $n=0$ undetermined.

\subsection{Scope of the continuation}

The unconditional conclusions are the exact rank-three alternatives,
the exceptional-set restrictions, the formal arithmetic nondivision
theorem, the spectral error identity and counterexamples, the local
radius and denominator calculations, the curvature obstruction, and
the joint-span identities.  Their proofs use identified classical or
published inputs.  They do not prove the designated rank-four assertion,
the irrationality of either constant, or a new transcendence conclusion
for $\gamma$ or $\delta$.  A successful completion would need an
unconditional arithmetic argument excluding an isolated rank defect at
one, respecting the canonical analytic normalization.  Neither a renamed
zero-removal gate nor a general spectral or differential property
supplies that argument.

% END INTEGRATED SOURCE

% BEGIN INTEGRATED SOURCE: ending.tex
\clearpage
\section{What the integrated argument decides}
\label{int:endpoint}

The construction and exclusion results in this manuscript have different
logical outputs. The binomial and prime results characterize arithmetic
behaviour under explicit hypotheses. The moment and Fourier sections
solve the auxiliary reconstruction problems that they state. The
functional results exclude differential-algebraic descriptions of
specified transforms. The rank-four investigation restricts where a
rank failure can occur, and the final spectral and $p$-adic calculations
test proposed transfer mechanisms. None of these deductions alone
identifies the arithmetic nature of the designated constant $\gamma$.

Three gaps deserve separate names only for clarity, not as claims of
progress from renaming the problem.
\begin{enumerate}[leftmargin=*,itemsep=5pt]
\item \textbf{Nonzero arithmetic output.} A short vector in a logarithm
elimination lattice may lie in the kernel of the rational output map.
A useful infinite family must retain nonzero output while satisfying
the needed height and error estimates.
\item \textbf{Arithmetic specialization.} Analytic division by $z-1$
after a zero at $1$ does not automatically preserve the arithmetic
Gevrey class. In the restricted four-dimensional family, the requested
lifting statement is equivalent to the designated rank-four target,
once the formal nondivisibility result is used.
\item \textbf{Compatibility between places.} A real asymptotic
normalization does not by itself identify the initial value with the
$p$-adic factorial sum. The exact Taylor-radius calculation measures this
missing equality; it does not establish it.
\end{enumerate}

The counterexamples locate the scope of such gaps. A deformation that
preserves a restricted list of identities may lose the canonical
normalization of the original analytic function. Conversely, the
positive integer Jacobi example preserves positivity, integer moments,
infinite spectrum and canonical resolvent normalization, and its Weyl
function is differentially transcendental, yet its value at $1$ is
rational. A proof for the original Euler system must therefore use
additional information distinguishing that system from the example.

The final status is
\[
 \boxed{\rho(1)=4\text{ remains unproved in this manuscript}.}
\]
Neither an explicit contradiction in the algebraic-gamma alternative nor
one in the transcendental-gamma alternative has been obtained. This
statement concerns the results established here; it is not an
impossibility theorem for future methods.

\section{Source record, corrections and reproducibility}
\label{int:sources}

The requested starting point is the evaluation question quoted in the
Korean reader's guide. The following source families cover the
mathematical developments from that evaluation through the last
cross-field investigation. They are internal research manuscripts and
proof archives. The earlier September 11 Two-World manuscript is the
object of the initial review, not a second copy inserted into this
edition. Manuscripts 001 and 002 are historical antecedents; the reused
foundations needed in this range are reproduced where invoked.

\begin{longtable}{@{}>{\raggedright\arraybackslash}p{0.08\textwidth}>{\raggedright\arraybackslash}p{0.42\textwidth}>{\raggedright\arraybackslash}p{0.43\textwidth}@{}}
\toprule
Stage & Recovered source & Content retained in the integrated text \\
\midrule
\endfirsthead
\toprule Stage & Recovered source & Content retained \\\midrule
\endhead
R & \nolinkurl{KEIO_TwoWorld_Review_Addendum_20260920.tex}
& Review corrections; sharp interval exponent; rational and algebraic
floors; collision-ring information; periodic deformation. \\
A1 & \nolinkurl{KEIO_Gamma_AB_Construction_20260920.tex}
& Binomial closure; Borel division; multiplication and Fourier
constraints; coefficient-rank criteria. \\
A2 & \nolinkurl{KEIO_Gamma_AB_Extension_II_20260920.tex}
& Minimal prime certificates; denominator and density refinements;
finite Fourier ranges; multiple poles; categorical comparison. \\
M & \nolinkurl{KEIO_Gamma_Moneyball_20260920.tex}
& Prime-based independence; denominator defects; Hankel construction;
countermodels; imported zero theorems; parity-moment reconstruction. \\
S & \nolinkurl{KEIO_Gamma_Snowball_20260920.tex}
& Functional differential transcendence; recurrence restrictions;
nonlinear descent; two-trace inverse data; local denominators;
residue-profile transfer and supporting foundations. \\
C & \nolinkurl{KEIO_Gamma_Core_Research_20260920.zip}
& Moving-degree arithmetic cost and its proof; fixed-transform,
positive-kernel and moment-elimination limits; arithmetic reassessment. \\
R4 & \nolinkurl{KEIO_Rank4_Proofs.zip}
& Differential system; rank-three classification; geometric-point
exception bounds; exact formal quotient obstruction. \\
X & \nolinkurl{KEIO_Cross_Field_Proofs.zip}
& Laguerre resolvent; integer Jacobi and Meixner counterexample;
difference-Galois proof; adelic radii and $p$-curvature;
Galois and multipoint elimination audit. \\
\bottomrule
\end{longtable}

\subsection{Corrections that affect interpretation}
The initial review corrects several claims in its predecessor: uniform
framing statements need a common interval; finite approximation ratios
are not asymptotic irrationality or transcendence exponents; and
restricted countermodels do not establish independence of the actual
analytically defined constant. The later independence-selection
statements are existential: a suitable density-zero set can be deleted,
not an arbitrary one. The functional results do not identify the
arithmetic nature of the value at zero; for the entire transform
$H(z)=\int_0^1 e^{zx}(1-\log x)^{-1}\,dx$, that value is exactly
$H(0)=\delta$.

The generic rank-three case must also be read correctly. Deriving
transcendence or even algebraic independence of $\gamma,\delta$ from
that case is not a contradiction: the rank condition is stronger than
those assertions. Finally, positivity or differential transcendence
of an auxiliary function cannot be treated as an unproved
value-transcendence principle.

\subsection{Rebuilding the document}
The delivered TeX is a single source, including all mathematical text,
macros and bibliography; it does not use external \verb|\input|,
\verb|\include|, graphics or a separate BibTeX database. Compile it twice
with XeLaTeX. Standard TeX packages, \texttt{kotex}, and a Korean font
(Noto Serif CJK KR, with a fallback to UnBatang) are required. These
are typesetting dependencies, not additional manuscript source files.

The historical computational archives remain separate research sources.
Exact finite checks supplement the proofs and are identified as such;
finite decimal observations have not been converted into claims about
unbounded sequences. The present compilation validates the TeX build,
cross-references and PDF rendering and incorporates the mathematical
audits described in the text. It does not claim formal verification or
an external assessment of originality.

% END INTEGRATED SOURCE

\clearpage
\addcontentsline{toc}{section}{References}
\begin{thebibliography}{99}
\bibitem{KEIO} KEIO, \emph{The Two-World Search for Euler's Constant}, supplied 52-page manuscript, 11 September 2026.

\bibitem{FRliouville} S.~Fischler and T.~Rivoal, Values of $E$-functions are not Liouville numbers, \emph{J. Ecole polytech. Math.} \textbf{11} (2024), 1--18. Theorem A and Theorem 1. \url{https://rivoal.perso.math.cnrs.fr/articles/liouville.pdf}.

\bibitem{FRrational} S. Fischler and T. Rivoal,
\emph{Rational approximations to values of E-functions}, manuscript
dated 3 July 2025 in the source audit,
\url{https://rivoal.perso.math.cnrs.fr/articles/Eexponent.pdf}.

\bibitem{Beukers} F.~Beukers, \emph{A refined version of the Siegel--Shidlovskii theorem},
Annals of Mathematics \textbf{163} (2006), 369--379.
\url{https://annals.math.princeton.edu/wp-content/uploads/annals-v163-n1-p08.pdf}.

\bibitem{FRmixed} S.~Fischler and T.~Rivoal, \emph{Relations between values of arithmetic
Gevrey series, and applications to values of the Gamma function},
Journal of Number Theory \textbf{261} (2024), 36--54.
\url{https://www.imo.universite-paris-saclay.fr/~stephane.fischler/ssmixte.pdf}.

\bibitem{DLMF} NIST Digital Library of Mathematical Functions. Sections 5.5, 5.11, 6.2, 6.9, 18.2 and 18.9. \url{https://dlmf.nist.gov/}.

\bibitem{FRoperators} S.~Fischler and T.~Rivoal, \emph{Arithmetic theory of E-operators},
arXiv:1406.5995 (2014), especially \S1 and the cited global
regular-singularity results for G-functions.
\url{https://arxiv.org/abs/1406.5995}.

\bibitem{Roth} K. F. Roth,
\emph{Rational approximations to algebraic numbers},
Mathematika 2 (1955), 1--20.
\url{https://doi.org/10.1112/S0025579300000644}.

\bibitem{Elliott} J. Elliott,
\emph{Birings and plethories of integer-valued polynomials},
arXiv:1109.3848.
\url{https://arxiv.org/abs/1109.3848}.

\bibitem{Chabert} J.-L. Chabert and G. Peruginelli,
\emph{Polynomial overrings of $\Int(\Z)$},
J. Commut. Algebra 8 (2016), 1--28.
\url{https://arxiv.org/abs/1503.06035}.

\bibitem{PerezMarco} R. P\'erez-Marco,
\emph{On the definition of Euler Gamma function},
arXiv:2001.04445v2, Section 2.11.
\url{https://arxiv.org/html/2001.04445v2}.

\bibitem{BadeaGrivaux} C. Badea and S. Grivaux,
\emph{Around Furstenberg's times $p$, times $q$ conjecture:
times $p$-invariant measures with some large Fourier coefficients},
Discrete Analysis 2024:10.
\url{https://arxiv.org/abs/2303.01089}.

\bibitem{GreenSanders} B. Green and T. Sanders,
\emph{A quantitative version of the idempotent theorem in harmonic analysis},
Ann. of Math. 168 (2008), 1025--1054, Theorem 1.1.
\url{https://arxiv.org/abs/math/0611286}.

\bibitem{KM} D. Kalinov and A. Mandelshtam,
\emph{Sums of binomial coefficients evaluated at $\alpha\in\Qbar$,
and applications}, J. Algebra 587 (2021), 171--214.
\url{https://arxiv.org/html/2010.14494v3}.

\bibitem{HK} N. Harman and D. Kalinov,
\emph{Classification of simple algebras in the Deligne category
$\operatorname{Rep}(S_t)$}, arXiv:1901.05080v3,
especially Corollary 3.2.4.
\url{https://arxiv.org/html/1901.05080v3}.

\bibitem{Milne} J. S. Milne, \emph{Algebraic Number Theory},
course notes, sections on Frobenius and Chebotarev.
\url{https://www.jmilne.org/math/CourseNotes/ANT.pdf}.

\bibitem{TZ} J. Thorner and A. Zaman,
\emph{A unified and improved Chebotarev density theorem},
Algebra Number Theory 13 (2019), 1039--1068.
\url{https://msp.org/ant/2019/13-5/ant-v13-n5-p02-s.pdf}.

\bibitem{FP} F. F\"urnsinn and L. Pannier,
\emph{An Effective Version of the $p$-Curvature Conjecture for Order One
Differential Equations}, arXiv:2510.00892v2 (12 March 2026), Lemma 4.4.
\url{https://arxiv.org/html/2510.00892v2}.

\bibitem{Previous} KEIO research program, internal research manuscripts and proof archives, 20--21 September 2026. Their proof-bearing content is integrated in the present edition; see Section~\ref{int:sources}. No external peer-review or priority claim is implied.

\bibitem{FRzeros} S. Fischler and T. Rivoal,
\emph{Zeros of $E$-functions and of exponential polynomials defined over
$\overline{\mathbb Q}$}, arXiv:2503.20345.
\url{https://arxiv.org/html/2503.20345v1}.

\bibitem{Simon} B. Simon, \emph{The Christoffel--Darboux Kernel},
\url{https://arxiv.org/abs/0806.1528}.

\bibitem{Kolchin} E. R. Kolchin, \emph{Differential Algebra and
Algebraic Groups}, Academic Press, 1973.

\bibitem{KR} C. Krattenthaler and T. Rivoal,
\emph{Arithmetic properties of the Taylor coefficients of
differentially algebraic power series}, 2025.
\url{https://arxiv.org/abs/2502.09259}.

\bibitem{KW} H. Kadiri and P.-J. Wong,
\emph{Primes in the Chebotarev density theorem for all number fields},
\url{https://arxiv.org/abs/2105.14181}.

\bibitem{BTV} T. C. Burness and H. P. Tong-Viet,
\emph{Derangements in primitive permutation groups, with an application
to character theory}, \url{https://arxiv.org/abs/1403.7666}.

\bibitem{sc:BRS} A. Bostan, T. Rivoal, and B. Salvy, \emph{Explicit degree bounds for right factors of linear differential operators}, arXiv:1906.05529v4 (2020). \url{https://arxiv.org/abs/1906.05529}.

\bibitem{sc:HP} T.~Hessami Pilehrood and K.~Hessami Pilehrood,
\emph{On a continued fraction expansion for Euler's constant},
arXiv:1010.1420 (2010).
\url{https://arxiv.org/abs/1010.1420}.

\bibitem{sc:VW} T. Wolfs and W. Van Assche, \emph{Rational approximation of Euler's constant using multiple orthogonal polynomials}, arXiv:2404.09799v3 (2025), Theorem 1.2 and Propositions 2.5--2.6. \url{https://arxiv.org/abs/2404.09799}.

\bibitem{sc:Sondow} J. Sondow, \emph{Criteria for irrationality of
Euler's constant}, \url{https://arxiv.org/abs/math/0209070}.

\bibitem{sc:Mestrovic} R. Me\v strovi\'c,
\emph{Variations of Kurepa's left factorial hypothesis},
\url{https://arxiv.org/abs/1312.7037}.

\bibitem{sc:WangYao} Wang and Yao,
\emph{Difference independence of Euler gamma function},
\url{https://arxiv.org/abs/2303.02767}.

\bibitem{sc:FRBessel} S. Fischler and T. Rivoal,
\emph{Transcendence measure for the values of the logarithmic derivative
of the Bessel function $J_0$ at algebraic arguments}, manuscript dated
8 September 2026 in the source audit,
\url{https://rivoal.perso.math.cnrs.fr/articles/mesureBessel.pdf}.

\bibitem{sc:RivoalSigma} T. Rivoal,
\emph{A class of arithmetic difference operators}, manuscript dated
17 September 2026 in the source audit,
\url{https://rivoal.perso.math.cnrs.fr/articles/sigmaop.pdf}.

\bibitem{sc:Powers} M. R. Powers, \emph{Transcendence Results for $\Gamma^{(n)}(1)$ and Related Sequences of Generalized Constants}, arXiv:2511.01849v7 (2026). \url{https://arxiv.org/abs/2511.01849}.

\bibitem{sc:Aptekarev2026} A. Aptekarev,
\emph{Multiple orthogonal polynomials and strong Szeg\H{o} asymptotics},
\emph{Russian Mathematical Surveys} 81:2 (2026), 209--273,
\url{https://www.mathnet.ru/eng/rm10281}.

\bibitem{sc:Burnol2026} J.-F. Burnol,
\emph{Some geometric series for Euler's constant},
\url{https://arxiv.org/abs/2603.29998}.

\bibitem{rc:rivoal} T.~Rivoal, \emph{On the arithmetic nature of the values of the gamma
function, Euler's constant, and Gompertz's constant},
Michigan Mathematical Journal \textbf{61} (2012), 239--254.
\url{https://rivoal.perso.math.cnrs.fr/articles/gammater.pdf}.

\bibitem{rc:teschl} G.~Teschl, \emph{Jacobi Operators and Completely Integrable Nonlinear
Lattices}, Mathematical Surveys and Monographs, vol.~72,
American Mathematical Society, 2000, Chapter~2.
Author's online edition:
\url{https://www.mat.univie.ac.at/~gerald/ftp/book-jac/jacop.pdf}.

\bibitem{rc:hardouin-singer} C.~Hardouin and M.~F.~Singer, \emph{Differential Galois theory of linear
difference equations}, arXiv:0801.1493 (2008), Proposition~3.9(2).
\url{https://arxiv.org/pdf/0801.1493}.

\bibitem{rc:ernvall-hytonen} A.-M.~Ernvall-Hyt\"onen, T.~Matala-aho, and L.~Sepp\"al\"a,
\emph{Euler's factorial series, Hardy integral, and continued fractions},
arXiv:2111.13649 (2021).
\url{https://arxiv.org/abs/2111.13649}.

\bibitem{rc:bost} J.-B.~Bost, \emph{Algebraic leaves of algebraic foliations over number
fields}, Publications Math\'ematiques de l'IH\'ES \textbf{93} (2001),
161--221, Theorem~2.1.
\url{https://www.numdam.org/item/PMIHES_2001__93__161_0.pdf}.

\bibitem{rc:bost-chambert-loir} J.-B.~Bost and A.~Chambert-Loir, \emph{Analytic curves in algebraic
varieties over number fields}, arXiv:math/0702593 (2007).
\url{https://arxiv.org/abs/math/0702593}.
\end{thebibliography}

\end{document}
